% Options for packages loaded elsewhere
\PassOptionsToPackage{unicode}{hyperref}
\PassOptionsToPackage{hyphens}{url}
%
\documentclass[
]{memoir}

\usepackage{lmodern}


\usepackage{amssymb,amsmath}
\usepackage{ifxetex,ifluatex}
\ifnum 0\ifxetex 1\fi\ifluatex 1\fi=0 % if pdftex
  \usepackage[T1]{fontenc}
  \usepackage[utf8]{inputenc}
  \usepackage{textcomp} % provide euro and other symbols
\else % if luatex or xetex
  \usepackage{unicode-math}

\defaultfontfeatures{Scale=MatchLowercase}
\defaultfontfeatures[\rmfamily]{Ligatures=TeX,Scale=1}

  \setmainfont[]{Tex Gyre Pagella}
\fi

% Use upquote if available, for straight quotes in verbatim environments
\IfFileExists{upquote.sty}{\usepackage{upquote}}{}
\IfFileExists{microtype.sty}{% use microtype if available
  \usepackage[]{microtype}
  \UseMicrotypeSet[protrusion]{basicmath} % disable protrusion for tt fonts
}{}

\makeatletter
\@ifundefined{KOMAClassName}{% if non-KOMA class
  \IfFileExists{parskip.sty}{%
    \usepackage{parskip}
  }{% else
    \setlength{\parindent}{0pt}
    \setlength{\parskip}{6pt plus 2pt minus 1pt}}
}{% if KOMA class
  \KOMAoptions{parskip=half}}
\makeatother

\usepackage{xcolor}
\IfFileExists{xurl.sty}{\usepackage{xurl}}{} % add URL line breaks if available
\IfFileExists{bookmark.sty}{\usepackage{bookmark}}{\usepackage{hyperref}}
\hypersetup{
  pdftitle={Narrativity Volume 1},
  pdfauthor={Mary Burger, Robert Glück, Camille Roy, and Gail Scott},
  hidelinks,
  pdfcreator={LaTeX via pandoc}}
\urlstyle{same} % disable monospaced font for URLs
\usepackage[a5paper]{geometry}

\setlength{\emergencystretch}{3em} % prevent overfull lines
\providecommand{\tightlist}{%
  \setlength{\itemsep}{0pt}\setlength{\parskip}{0pt}}
\setcounter{secnumdepth}{5}

%% %%% Bibliography Options %%%%
%% %% %% %% %%%%%%%%%%%%%%%%%%%%%%%%%%%%%%%%

\title{Narrativity Volume 1}
\author{Mary Burger, Robert Glück, Camille Roy, and Gail Scott}
\date{}

%%%%%%%%%%%%%%%%%%%%%%%%%%%%%%%%%%%%%%%%%%%%%%%
%% Custom packages imported here
%%%%%%%%%%%%%%%%%%%%%%%%%%%%%%%%%%%%%%%%%%%%%%%

\usepackage{adforn}
\usepackage{verse}
\usepackage{pdfpages}
\usepackage{stringstrings,ifthen}
\usepackage{caption}
\usepackage{lettrine}
% \usepackage[all]{nowidow}
%\usepackage{Typocaps, Acorn, Elzevier, Zallman, Rothdn}

\captionsetup[figure]{labelformat=empty}

\renewcommand{\baselinestretch}{1.24}

\renewcommand{\chapnamefont}{\normalfont\large\scshape\color{darkgray}}
\renewcommand{\chaptitlefont}{\normalfont\Huge\scshape\color{darkgray}}

\setsubsecheadstyle{\centering\normalfont\large\color{black}}
\setsubsubsecheadstyle{\raggedleft\normalfont\color{black}}

\usepackage{Zallman}
\newcommand*\capfont{\usefont{U}{Zallman}{xl}{n}}
\renewcommand{\LettrineFontHook}{\capfont \color{darkgray}}

%\baselineskip=1\baselineskip plus 12pt minus 2pt\relax

% Strongly discourage orphans and widows
% \widowpenalty10000
% \clubpenalty10000

%%%%%%%%%%%%%%%%%%%%%%%%%%%%%%%%%%%%%%%%%%%%%%%%
%\usepackage[T1]{fontenc}
\usepackage{url}
\ifpdf
\usepackage{pdfcolmk}
\fi
\usepackage{graphicx}
%%\usepackage{hyperref}
%% drawing package
\usepackage{tikz}
%% for dingbats
%\usepackage{pifont}
%\providecommand{\HUGE}{\Huge}% if not using memoir
\newlength{\drop}% for my convenience
%% specify the Webomints family
%\newcommand*{\wb}[2]{\fontsize{#1}{#2}\usefont{U}{webo}{xl}{n}}
%% select a (FontSite) font by its font family ID
%\newcommand*{\FSfont}[1]{\fontencoding{T1}\fontfamily{#1}\selectfont}
%% if you don’t have the FontSite fonts either \renewcommand*{\FSfont}[1]{}
%% or use your own choice of family.
%% select a (TeX Font) font by its font family ID
%\newcommand*{\TXfont}[1]{\fontencoding{T1}\fontfamily{#1}\selectfont}
%% Generic publisher’s logo
\newcommand*{\plogo}{\fbox{PL}}
%% Some shades
\definecolor{Dark}{gray}{0.2}
\definecolor{MedDark}{gray}{0.4}
\definecolor{Medium}{gray}{0.6}
\definecolor{Light}{gray}{0.8}
%%%% Additional font series macros
\makeatletter
%%%% light series
%% e.g., kernel doc, section s: line 12 or thereabouts
\DeclareRobustCommand\ltseries
{\not@math@alphabet\ltseries\relax
\fontseries\ltdefault\selectfont}
%% e.g., kernel doc, section t: line 32 or thereabouts
\newcommand{\ltdefault}{l}
%% e.g., kernel doc, section v: line 19 or thereabouts
\DeclareTextFontCommand{\textlt}{\ltseries}
% heavy(bold) series
\DeclareRobustCommand\hbseries
{\not@math@alphabet\hbseries\relax
\fontseries\hbdefault\selectfont}
\newcommand{\hbdefault}{hb}
\DeclareTextFontCommand{\texthb}{\hbseries}
\makeatother
%%%%%
\newcommand*{\titleWH}{\begingroup% Words in their Hands
\drop = 2\baselineskip
\centering
\vspace*{\drop}
\settowidth{\unitlength}{\Huge\uppercase{Narrativity Volume 1}}
\vspace*{\baselineskip}
{\Huge\color{red} \uppercase{Narrativity Volume 1}}\\[\drop]
%% \scalebox{8}[1]{{\wb{10}{12}4}}\\[\drop]
{\Large\scshape Edited by}\\[\drop]
{\Large\uppercase{Mary Burger, Robert Glück, Camille Roy, and Gail Scott}}\\[\baselineskip]
\vspace*{\drop}
\begin{center} \includegraphics{images/pandanus.png} \end{center}
\vspace*{\drop}
%% %{\Large\itshape with a Commentary by}\\[\baselineskip]
%% %{\Large A N OTHER}\par
%% \vfill
%% \scalebox{8}[1]{{\wb{10}{12}4}}\\[\drop]
{\Large ANTELOPE EDITIONS}\\
{\scshape \the\year{}}\par
\vspace*{\drop}
\endgroup}

\begin{document}

\pagenumbering{gobble}
\titleWH


\frontmatter


\setcounter{tocdepth}{0}
\tableofcontents



mainmatter

\hypertarget{kathy-acker-the-killers}{%
\chapter{Kathy Acker: The Killers}\label{kathy-acker-the-killers}}

\lettrine[lines=3, findent=0em, nindent=0.1em, lhang=0]{I}{ wrote a piece}
today and it's really just the beginning of something on narration. I've
got time trouble because of overwork. You know: you're broke, and you
take on too much work, and then you fuck up on everything. It's never
ending.

So understand this is simple and it's the beginning of something. You
know, it's totally open rather than anything else.

And I thought we were going to touch on two issues talking about
narration. One is identity. That's an obvious issue, the business about
who am I---you know, this is all postmodernist shit. What, therefore, is
real? Is there any real? And the other issue is raised by the quote from
this little card that advertises this panel about writing for century's
end. Constantly there's this trouble hovering in the air. The immediacy
of all the danger we're living in, of the possible end of both humanism
and of humanity, that we might all be going the way of the dinosaurs
very fast. And that culture is going the way of? of whatever---dinosaur
culture. And when we talk about writing these two things, these issues
hover there.

But let me start, since we're talking about narration, by telling you a
story. Very simple. My story begins with my friend Bob Glück who one
day, once upon a time, as story structure goes, told me that he has a
certain habit. A habit of circumventing his own habits by asking his
friends to give him a reading lesson. What should he, Bob Glück, read?
Perhaps in order to improve himself, or perhaps in order to do the
opposite. So Bob Glück goes to Kevin Killian, another friend, ``Kevin,
what do you want me to read?''

Kevin replies, ``\,`The Killers,' by Ernest Hemingway.''

Upon hearing this story, for this is a story, I replied, ``Oh Bob, how
weird, Hemingway?''

``Not weird'' replies Bob. ``All the English students at San Francisco
State love Hemingway.''

Two nights after this happened I had the following dream. First, I have
to tell you a few details about my childhood and about before my
childhood in order to elucidate this dream for you. I never met my
father. Though he was married to my mother, he left her when she was
three months pregnant with me.

When I was 26 years old, through an accident, I traced my father's
family. I wrote to them and they wrote back that they would accept me
into their family and we arranged to meet. I thought that I was going to
meet my real father, but I only met the first cousin. He told me (and I
think he was a little crazy---well not crazy but eccentric, because rich
people are never crazy), that perhaps I should not meet my father. Why?
Because my father had murdered someone who was trespassing on his yacht.
After remaining six months in a lunatic asylum, the state had excused
him of any murder charge. My father then disappeared. No one now knew
where he was, said the first cousin. And so I abandoned my search for my
father, for my life at that time was hard enough, and this new trouble
was simply not worth the trouble.

When I was 30 years old, my mother suicided.

Enough of my childhood.

I had the following dream: I dreamt that I was looking for my real
father. In my dream, I knew that it was dumb for me to do this because
my father was dead. Since I'm not dumb, I or the dreamer thought, I must
be trying to find my father so that I can escape from this house which
is being run by a woman. I go to a private detective. He calls me a
dame.

I say, ``I'm looking for my father.''

The private detective, who might be a friend, replies that my case is an
easy one. I like that I'm easy. We begin our search. According to his
instructions, I tell this private eye everything that I know about my
mystery. It takes me several days for me to recount all the details. It
was summertime in Dallas. Everything was yellow. I didn't remember
anything about this first period of my life, about my childhood. After
this not remembering, I remembered jewels. As soon as my mother passed
away, a jewel case was opened. The case, consisting of one tray, had
insides of red velvet. Perhaps I'm dreaming my mother's cunt. I'm given
a jewel which is green. I don't know where that jewel is now. I have no
idea what happened to it. This is the mystery of which I'm speaking.

The private eye pursued my matter. A couple of days later he came up
with my father's name: Olen. This name means nothing to me.

``Olen. Your father's name is Olen. Furthermore, your father killed your
mother.''

I thought, in order to dismiss thought, that's possible. The detective
continued to give me details about my father. He's from Iowa, and of
Danish blood. All this could be true, because how can I know anything?

Now when I woke up out of this dream, I remembered details about my
mother's suicide. She had killed herself eight days before Christmas. A
note in her handwriting, lying beside her dead body, said that her white
poodle was staying with such and such veterinarian. Nothing else. But
despite this note the cops were convinced that my mother was murdered by
a man whose name wasn't known. Nevertheless, it was Christmas and there
wasn't going to be any police investigation because the cops wanted to
return to their homes, Christmas warmth, and holiday.

For the first time ever in my life I had the following thought: My
father could have murdered my mother. After all, what if my real father
is crazy? At that moment I became very scared. If my father did murder
my mother, he could now be planning to murder me.

Now last week I was touring around and I found myself on a gig in
Roanoke, Virginia, where I met this absolutely wonderful writer, Richard
Dillard, and Richard told me that when he had been a boy he had
encountered his first proposition by a man during a showing of ``The
Killers'' at a local Roanoke theater.

That's the end of the story.

Also I should mention that in ``The Killers'' the name of the Swede
who's about to be slaughtered by those killers is Olen.

Now, is this a story, what I've just told you? Certainly it's narrative,
right? And each incident actually did happen. That is, it was real. But
a story? Stories are about something. And what I've just told you---it's
not about anything. It's not even about me. A story, you see, a story
has something to do with realism, and what I've just told you, though
each little bit was real or had happened, has nothing to do with
realism.

Listen to another voice as we proceed on this treasure hunt, or
narration. The voice of Julio Cortasar.

``Almost all the stories I have written belong to the genre called
fantastic, for lack of a better word, and are opposed to that false
realism that consists of believing that everything can be described and
explained, as was assumed by the optimism of 19th century philosophy and
science. That is, as part of a world governed more or less harmoniously
by a system of laws, principles, cause and effect relations,
well-defined psychologies, etc.'' That lovely world order---that
governing harmoniously by all these principles---we all know is now
over. Though Bill Clinton has a hype.

Okay, a little more on realism. Richard Dillard rightly calls that
narrative structuring named realism ``reductive.''

``But isn't that the central story of 19th century realistic fiction?''
he exclaims in his review of Alester Grey's marvelous novel \textsc{Poor
Things}. ``The central story of realistic fiction: we are born, we are
misshaped by biological inheritance, economic forces beyond our control,
and cultural biases beyond our recognition, and finally we die with our
failed dreams on our dry lips.''

Realism: reductive and dehumanizing.

``There's another order,'' says Cortesar, ``more secret and less
communicable. That the true study of reality lay not in laws but in the
exceptions to those laws has been the guiding principle in my personal
search for literature beyond all naive realism.''

Now return to the narration that I told you. I can't even say what it's
about: Bob Glück's conversation, my dream about my real father, Richard
Dillard's memory. Where in this narration lies the real? It lies in the
connections between the three sections; in the connections between the
``real'' events and the holes, the silences. In the slippages. Slippages
into what?

In Moby Dick, Melville speaks of reality as the interstices through
which all of us fall. There's another, a further way of putting this, of
proceeding on our treasure hunt. These interstices can be named chaos,
or places where language cannot be, or death. When storytelling, humans
attempt to cling to meaning. I think of narration, of narrative, or that
narrative which tries to encounter the real, as that which is
negotiating between two orders of time: clock time and chaos. The writer
is playing---when structuring narrative or when narrative is structuring
itself---with life and death. He or she is maneuvering between order and
disorder, between meaning and meaninglessness and so is making
literature.

However, these movements between clock time and chaos in written
narrative are more structurally complex. To begin, consider one aspect
of time in the novel: time. What? The time it takes to write a novel. A
novel's a big thing. I'm talking about novel narration at the moment
because that's what I deal with. It usually takes at least a year, often
many years. During that time the writer's life changes. So there's the
time of all the actual changes the writer is going through---the time it
takes to write the novel.

Two: the time it takes to read a novel. A novel is very rarely something
you read in one sitting. So that time it takes to read a novel
incorporates all the readers memories, all the interstices of time, the
time lapses between readings, all the returns to earlier parts of the
novel, etc. Three: the fictive time. The time within the story or the
narrations. Even Hopscotch---there's fictive time there. So in this
sense a novel, structurally, is a time triad.

Rather than continue with this---and you can start analyzing novelistic
structure---I want to just quickly skip to a final point. Consider again
realism. And I must say the more I think about it, I think writing is
more and more about time. Consider again this thing of realism (I'd like
to abandon the whole thing). Within the realm of realism lies the
assumption that language mirrors all that isn't language, right? A
table. That what a narrative is about: telling what is or should be
true. That a narrative mirrors reality.

Now do I need to say---how simplistic. But what I want to ask is, why
bother being so simplistic? Why bother with the lie of realism? Why
bother being so miserable, reductive, when one could play? If I'm going
to tell you what the real is by mirroring it, by telling you a story
which expresses reality, I'm attempting to tell you how things are. By
letting you see through my own eyes, I give you my viewpoints, moral and
political. In other words, realism is simply a control method. Realism
doesn't want to negotiate, open into, even know chaos or the body or
death. Because those who practice realism want to limit their readers'
perceptions, want to limit perceptions to a centric---which in this
society is always a phallocentric---reality. ``I am the one,'' says the
realistic writer, ``I'm telling you reality.'' And we can talk. I have
the same quarrel about narrowing anything to single identity.

In other words, behind every literary or cultural issue lies the
political, the realm of political power. And whenever we talk about
narration, narrative structure, we're talking about political power.
There are no ivory towers. The desire to play, to make literary
structures which play into and in unknown or unknowable realms, those of
chance and death and the lack of language, is the desire to live in a
world that is open and dangerous, that is limitless. To play, then, both
in structure and in content, is to desire to live in wonder.

\hypertarget{note}{%
\section*{Note}\label{note}}
\addcontentsline{toc}{section}{Note}

Kathy Acker gave this talk on a panel called ``In Extremis, Writing at
the Century's End.'' The panel was presented by Small Press Traffic in
the New College Theater in San Francisco on April 29, 1993. Fortunately,
the event was videotaped by S.S. Kush for his Cloud House Poetry
Archives (1557 Franklin St.~San Francisco CA 94109, 415-292-5554). We
are grateful to Kush and also to the Acker estate. The tape was
transcribed by Quintan Wikswo and edited by Robert Glück. The panel also
included Dodie Bellamy, Earl Jackson Jr., and Benjamin Weissman.

\hypertarget{lawrence-ytzhak-braithwaite-last-exit-to-victoria}{%
\chapter{Lawrence Ytzhak Braithwaite: Last Exit to
Victoria}\label{lawrence-ytzhak-braithwaite-last-exit-to-victoria}}

\begin{verse}
(5mins. of Murder Alone with Wigger and Bloodland)\\
\end{verse}

~

~

~

\hypertarget{section}{%
\section*{1}\label{section}}
\addcontentsline{toc}{section}{1}

\lettrine[lines=3, findent=0em, nindent=0.1em, lhang=0]{A}{s a child I was told}
that not knowing the alphabet will cause illiteracy. It'll send you into
a drugged out gangland life of white trash nightmares and corner boy
peddling to homosexuals, who are professional players, obsessed with age
and are willing to drag it and you into emptiness. That in knowing the
letters, Ill know that they assemble to construct various images that
become words. Words are the narrative transformation of the images.
Printing a page of unbroken words is like a fresh tattoo. It captures a
moment/place, sentiment and period. It orchestrates the body in motion
as it flexes to move a pen/strike at a key/form a fist/lift a drink or
move to a rhythm. The words become the unspoken intertextuality of
ethnic, racial and cultural metaphoric speech. The metre of casual
dialogue = a rhythm/noise/visual/bass, a soundtrack to a post-literate
train of thought.

I dedicated \textsc{Wigger} to the sound of Slayer and Sonic
Youth---``Bloodland'' to Cannibal Corpse (George ``Corpsegrinder''
Fisher) = I dedicated them to sound. Slayer is for the fury and speed
and violence that the book has. Deathmetal is the living desire of the
neo-redneck burnout. It's all going after the sport of brutality---the
art of hurting someone. The walking jokes, with targets on their backs,
placed there by lilbitches, taking them for Mn Ms.~The only violence is,
the way the words appear on the page, marked by the slashes that connote
rhythm of speech and interrupted thought. They are like semi-colons = /
the // are colons and so are the = signs. Sometimes the - move out to
separate speech---someone takes lead//does a solo. Sonic Youth's sounds
appear in the form of the book dissolving. Deathmetal is just dealing
with the situation. The book and the story, like the characters, are
trying to hold themselves together. Brian and Jerehmia symbolize this
the most. They talk. No one listens.

\ldots{}a printer ejects swatches of stolen tissue that collects the
sound and images of what is considered low brow art and skill//

Hardcore (H.C.)/Country/Rockabilly

Thrash and Speedmetal/Rap/Deathmetal/Blackmetal

Cinema/Television and;

Comics Books

Brian appears to be rambling but he makes the most sense because he
speaks through metaphor and his heart, like Black speech and song---the
negro spiritual, seeming to sing about a river,

when it is a code,

for an escape,

at a certain time,

from the plantation.

Gangsta rap seems to boast

about an evening

or event

but is just shouting

validity,

existence and

enterprise.

Brians grasp on reality becomes inconsequential.

Jerehmias story is dumbshow.

Brian surrenders his ideals and his soul, by killing Spook (the Black
hustler friend).

Jerehmia is just having his taken and he's telling you about it---but I
knew people in Atlanta that didn't want to hear about any of that or
them. I knew a boy from Tennessee hurting worst than what was going as
oppression under those 100 dollar Afro centric hats.

\ldots{}Wordcore, Jamaican dub poetry, Rap and Rock Steady are disposed
to the Homeric boasts and catalogues of post-modern thugs/hoodlums and
desperadoes w/ the hope of kindness and compassion.

In donning the Black persona, symbolized through the silver jacket,
Brian finally does what everyone has been attempting to do throughout
the book. Brian is killed---his soul is killed, through the burden of
the weight of the Black youth---the Black persona, the persona of
deglamoured oppression. He has achieved the goal of being Black but he
is unprepared to handle something that the Blacks are raised to deal
with through centuries of struggle---youd suppose.

I met a boy with crystal smack grey eyes, who was offering to suck the
fender (for \$5) off an urban assault vehicle for a painkiller.

If you listen and see the page, you'll see a tattoo and note well, that
some words are//

-Rhythm

-Lead and

-Bass

Jerehmia, he has no desire to be Black. He knows what the shit is and
he's got his own. He just watches the seemingly liberal Perry go catch
back the Black boy trying to bail on the whole scene. He just knows that
nothing that's happening, is about him or the niglet. It's about the
little men with big shirts and the chicks in bed or sitting off from the
Bar b q..

\ldots{}strange how all the youngmen holding up walls waiting on passers
by had deep drawls, no smiles and trailers in their hearts.

The book and the story are collapsing. The spaces and sparsness of the
narrative is there to get you to listen carefully and read again and
close read and not take things at face value. People say more,
especially, when they've had a lifetime of hurt, through very few words
and the words look in ways that they feel---speeching in accapella. The
words sound the way they feel. The book shows patterns of speech.

Mcluhan would argue the global village: technology is good---the book
people are lacking in the rich grammar of the TV = This is referred to
today as the post-literate generation---John Cages bastard children turn
to violence and the acoustic space of video and computer games and the
delicacy of words typed over television screens.

Mcluhan would argue\ldots{}

\ldots{}the acoustic

~

~

~

\hypertarget{section-1}{%
\section*{2}\label{section-1}}
\addcontentsline{toc}{section}{2}

~

The Comic book uses frames and juxtaposition to connect the
disconnectedness of thought and words, be it through the gruff noireness
of Frank Miller, Dave Sim or Martin Wales ``Kinder Nacht,'' or the
smooth lines, colours and ruff justice of Todd McFarlane. They collect
the street sounds and activities the way Rap culture manipulates and
emulates, via spin art and samples, the ready made terror of interurban
life. To frame is to make perfect the moment of the fingers striking the
keys, which can be the repetition of words and phrases = outlets ripping
and shredding, w/ the knowledge that framing is certainly an attempt to
make perfect, to make the words as enticing and elegant as the bovine
euphony of a fascist's goose step.

Brian is condemned to sit, numbed out on pills, in a room, filled with
drag queens, watching them transform. Just like Brian's desire, and the
other characters desires, to be Black, the drag queens attempt to be
women, without knowing, that beyond the glamour, is a lot of hardship
and struggle. Drag queens die, as well, when they finally achieve their
dream of being a woman. It's not as much fun or better than they thought
it would be. They are unprepared for the daily struggle and threats. You
can see this on the page. They are boxed in, allowing themselves to be,
finally, separated from the freeness and flow of the book. They must
stay within the parameters. The glamour is gone. They have to bear the
actual pain. Brian has got the auditory recollection of rape in his
head. Brian is numb. Sos Jerehmia, but he's numb from the beginning, any
way. Hes just recollecting, now---no need to figure. He cant hold it
together, nor does he want to. His flag is upside down, Oswalds
dead---nothings right. You separate the body from the mind and its all
good.

\ldots{}a word can be the hums or bops in the background as when a funk
musician beats the strings w/ his thumb---the buzz and scratch of a tat
gun on a steadied arm or back---the stilled breath and firm muscle of a
word or image engraved on a belly. It gets tied into the paragraph. Some
words take on the repetitive ecstatic riff of ska or reggae or the decay
of the crash/noise of H.C. or metal. They recline on the page or in the
air as if they were all going to amount to a junkies last sigh. Instead,
they collapse into the lost, disorientated and somewhat satisfied image
of Orsen Wells after thrashing a room at Xanadu.

\ldots{}and the book and story crumble into a final sustained note, done
through a note to Andrew from Jerry. ``Bloodland'' has a crack of gun
shot and shit and a heart beat. As it began on a note---it ends. His
(Jerrys) speech is proper and distant from the language of the book. He
has become a literary character of the Modernist sense.

Jerking the Modernist approach to distance, into ``groovy times,'' is
the dumbshow to the contemporary adoration of the absurdist tragic
comedy = anti-romantic. It presents itself, today, through the jaded
situationism of talkshow culture. The sort that ``Blast'' and ``Counter
Blast'' utilized through the headlines and print of newspapers. The
newspapers place horror and glory side by side in a folio that has a
calculated randomness. It will always assemble to state 1 thing = these
events happen. It is only through such genres of media
(music/comic/print/television), that we can see that each has a thing in
common/coupons and advertising = rhetoric and metaphor. If the passage
of a text takes on the sound and image of the disjointedneness of casual
speech and media, then the results are not only the coded tales being
woven, but the presentation of the brutality that leads to the violent
outcome.

Jerehmia, he just figures\ldots{}hes just feedback from a bass, right
now.

\textsc{Wigger} and ``Bloodland'' are what the printer jetted out---a
conceptual pastiche alphabet from the corner or gutter. It assembles the
images in sequence which persuades the narrative to be pulled out from
under a character. Its what got them blocked off---formatted, layed out,
then pile-drived with carefully selected skin grafts of onomatopoeias.
The words takes Audens diver, diving w/ his brilliant bow and takes it
beyond the enthralled spectator. They suppose that he has forgotten to
test the water and has come out a dumb struck paraplegic = what just
happened/where are our people/who are our people/do we have any people:
NO. However, they still keep tabs on the moment, place sound and period
that left the object devastated, hung over, and w/ an image permanently
scratched and coloured onto his flesh or hanging in a window.--

Oi Cheers,

Laz

\hypertarget{mary-burger-flange-or-premises}{%
\chapter{Mary Burger: Flange or
Premises}\label{mary-burger-flange-or-premises}}

\lettrine[lines=3, findent=0em, nindent=0.1em, lhang=0]{C}{an I say it went on}
for a week. Can I say I woke up each day, and went through each day,

``You were there, and you were there, and you, and you!''

Word went out, a consciousness with a consistency, looked at itself
looking at itself, halted the infinite regress, the mirror facing the
mirror, a language made possible a belief in location.

I wanted to speak as someone who had experiences. I wanted to speak as
someone who knew others. To have things happen, and to speak of them,

someone told a story to which I was unable to respond. The story was the
story itself.

~

I want the other thing. Hinges and holes and string. The impossible
appearing to be,

an engine, a large machine, an apparatus, a stage for a scene.

~

Can it be a narrator who doubts its premises?

The crumbling, if something will, whom it will happen to.

\begin{quote}
"Lyric, whose operation is display, and narrative, whose operation is
seduction.

Who will speak for the beside-the-point?" (Aaron Shurin, Narrativity)
\end{quote}

As what happens then, if what happens is always something more.

~

~

\hypertarget{david-buuck-despite}{%
\chapter{David Buuck: Despite}\label{david-buuck-despite}}

\hypertarget{section-2}{%
\section*{1}\label{section-2}}
\addcontentsline{toc}{section}{1}

\lettrine[lines=3, findent=0em, nindent=0.1em, lhang=0]{O}{ne narrative theory}
is that Agents are social beings. Equidistant from Los Alamos and
Roswell, the artist Walter de Maria places poles in the landscape in
order to construct a grid. The land is not the setting for the work but
a part of the work. The poles are angled at eighty-eight degrees else
lightening always strikes twice. The primary experience takes place
within the grid. Weather is total.

One narrative theory is that social beings are agents. The grid is the
imaginary frame by which it becomes possible to sense the spatial
landscape. The grid dwarfs Agents, the grid in turn being dwarfed by the
landscape, provoking Agents to realize their own relation to the spatial
landscape. In that the grid itself cannot be comprehended in its
totality, neither can the spatial landscape. Castles are sand that can't
be shelter to kings. In that the spatial landscape cannot be
comprehended in totality, neither can there be kings. Rabbits nibbling
on apple cores are nervously aware that there can't be shelter from
weather. Humility is nervous awareness of Agents not being rabbits.

There are three agents in the landscape who are not kings. The Agents'
relation to the grid becomes paramount. Agent A's scale becomes
paramount. Agent B's silence becomes paramount. Agent C's isolation
becomes paramount. Humility is nervous awareness of Agents not being
kings. The invisible is real.

The spatial relation of sky to the landscape is such that the clouds are
pressed flat against an imaginary plane of glass. The imaginary plane
rests upon the grid. One narrative theory is that Agents' minds function
as the imaginary plane. The action is the air and wind and rain.
So-called natural elements are grounding such that Agents' social being,
in isolation, in being dwarfed, is grounded. That this is not separate
from Agents' social being, but rather central to Agents' realization as
social beings, is one effect. Rabbits nibbling on apple cores are
nervously aware that lightening always strikes twice. Humility is
nervous awareness of Agents not being total. The primary experience
takes place within the social grid. Distance is total.

\hypertarget{section-3}{%
\section*{2}\label{section-3}}
\addcontentsline{toc}{section}{2}

One narrative theory is that one never dies in photographs. Equidistant
from Hiroshima and Nagasaki, the artist Kenji Yanobe builds nuclear-safe
capsules within which he is photographed while being in Chernobyl. That
the artist Kenji Yanobe would be in Chernobyl inhabiting sculptural
objects for the ostensible sake of safety is troubling to one. That the
photographs of the artist Kenji Yanobe inside nuclear-safe capsules
while being in Chernobyl are aesthetically pleasing is troubling to one.

One narrative theory is that in desiring the apocalypse it is
continually deferred. The artist Kenji Yanobe builds nuclear-safe
capsules for sexual acts. Such acts are expressed outward as exhalations
of perfumed music, compressed air modulated through anthropomorphic
trumpets, so-called safe sex ironized into cartoon farts. That
nuclear-safe capsules for sexual acts are aesthetically pleasing is
troubling to one.

One narrative theory is that one's apocalypse is inside. Such apocalypse
is not nuclear-safe, such sex is disharmonious, such art is a cartoon
that is actual. Agents inside nuclear-safe capsules are not safe from
themselves. Agents cannot be insulated in machines that are cartoon, in
art that farts its own radioactivity.

There are three Agents in the capsules who are not cartoons. Agent
Yellow breathes its own air. Agent Red eats its own waste. Agent Orange
is the machine that desires its own apocalypse. That the technology
necessary to provide safety for Agents is the same technology that
produces the effects from which Agents desire safety is troubling,
three, two, one. Castles are cartoon machines that can't be shelter to
kings. There are no kings, only agents.

\hypertarget{section-4}{%
\section*{3}\label{section-4}}
\addcontentsline{toc}{section}{3}

The sky is always in relation to Agents' landscape. Even underground,
though that's a dream. In dreams Agents' own minds are the relation of
the imaginary plane to Agents' own underground landscape. The grid's
only function is to provoke Agents' realization of Agents' own landscape
in actuality. The grid exists so as to become invisible. The grid exists
so that it doesn't have to. The invisible is real. Dreams are real.
Shelter is context.

One narrative theory is that the sun never sets. Agent Yellow, from a
distance, disappears into the landscape. Agent Red is still
disappearing, though Agent Orange is still here. Perspective is
constitutive of Agents' social being. Agents have different perspectives
that are relational. Total perspective is ideology that can't be shelter
to kings. Perspective is constitutive of Agents' social bearing. Context
is shelter.

Being dwarfed is perspective that reduces distances between social
beings. Relations between social beings provoke social bearings within
the grid. Roswell and Los Alamos are two landscapes that provoke such
social bearings. That the historical landscape dwarfs Agents, provoking
such social bearings, is troubling to one. Hiroshima and Nagasaki are
two landscapes that provoke such social bearings. Social bearings
realize social grids that produce the social landscape. That Agent A
might become grounded, in its isolation, is one effect. That Agent B
might become grounded, in its social being, is one effect. That Agent C
might become grounded, in its social bearings, is one effect. The social
grid is constructed by the many social beings. Its construction is
relational to the many social bearings. Perspective dwarfs the
historical landscape. The sun never sets.

\hypertarget{section-5}{%
\section*{4}\label{section-5}}
\addcontentsline{toc}{section}{4}

Dystopia is this topia, the cynically willed non-place of the
apocalyptic imaginary. Its filmy visuals are projected onto the
imaginary plane of glass that rests upon the social grid. The sky is the
image of the impending apocalypse. Its spectral traces appear
transparent when viewed from the ground up. The sky is always falling,
else the bottom falls out. Social beings labor at constructing the grid
to support the imaginary plane. Agent-sweat is radioactive. Agents drink
it as refreshment while awaiting the coming attractions.

One narrative theory is that one never dies in film. At the end of a
long day of grid-building it is nice to lie on the ground and watch the
bottom of the sky fall onto the imaginary plane. The ground is the
imaginary plane of consciousness upon which social beings construct the
social grid. Ideology is the nightmare made material by certain social
bearings. The ideological grid exists so that it might become invisible.
The ideological grid exists so that it doesn't have to. The invisible is
reel-to-reel.

One narrative theory is that Agents are never awake, except in history.
So-called nightmares are in actuality the history that Agents desire to
have been merely dreamt. Underground is the collective repressed. At
night it dreams itself through rabbit-holes into the social grid. One
narrative theory is that one never dies in dreams. Rabbits nibbling at
nuclear cores are nervously aware that history always strikes twice.
Humility is nervous awareness of Agents awaking from dreams and being in
history.

\hypertarget{section-6}{%
\section*{5}\label{section-6}}
\addcontentsline{toc}{section}{5}

The artist Kenji Yanobe's sculptural objects, in that they provide
habitation, in the form of objects of safety and art, while at the same
time being distancing, in that Agents are further distanced from the
actual world, while still being in it, are troubling to one. That the
artist Kenji Yanobe would be in Chernobyl inhabiting sculptural objects
for the ostensible sake of art is troubling to one. That the art
necessary to provide safety for Agents is the same art that produces the
effects from which Agents desire safety, is troubling, four, three,
agents, two, one. Agents are never safe from themselves.

The artist Kenji Yanobe's photographs, in that they document art
providing safety, while at the same time being distancing, in that they
are further distanced from Chernobyl, while still being in it, are
troubling to one. The photographs are viewed outside of Chernobyl, in
place of Chernobyl, as art and as safety. They are already happening,
though in actuality they are not Chernobyl. Distance is safety. There's
no distance in actuality. There's no safety.

The artist Kenji Yanobe has not yet built nuclear-safe capsules for
rabbits nibbling apple cores. One narrative theory is that nervous
awareness is not enough to prevent an apple core breach. That the wind
from Chernobyl might kill the rabbits is troubling, four, them, three,
to one. That the repeated vocalization of ``Chernobyl'' might become
aesthetically pleasing by virtue of its repetition is troubling to one.
There are sixty-six dead rabbits in the grid. Their bottoms have fallen
out.

\hypertarget{section-7}{%
\section*{6}\label{section-7}}
\addcontentsline{toc}{section}{6}

Sentry posts guard the perimeter. Their dogs are in nuclear-safe
capsules, making it difficult to chase the rabbits. Acid rain leaks into
the underground reservoir. Social beings labor into a future that is
still undecided. Agent-sweat is radioactive. Agents drink it as
refreshment while awaiting the sky to fall. One narrative theory is that
the sentry posts face in rather than out. Each pole waves a flag
advertising the coming attractions. One narrative theory is that there
is no perimeter.

Social beings that are also rabbits that live in underground bomb
shelters. Castles made of pamphlet bomb shreddings can't be shelter for
rabbit-kings. Rabbits have infrared night vision that always strikes
twice. They drink from the radioactive reservoir that runs from Roswell
to Los Alamos. That the wind from Los Alamos might kill the rabbits is
troubling to one. The sentries put the sixty-six dead rabbits into piles
to climb up and up. Else the bottom would fall out. That the repeated
vocalization of ``Los Alamos'' might become aesthetically pleasing by
virtue of its repetition is troubling five, four, three, to one.

One narrative theory is that dreams are underground nuclear tests. They
mushroom into consciousness through rabbit-holes and anthropomorphic
trumpets. Cartoon farts are air-raid sirens that go up and up. The music
is the bomb. It makes the bottom fall out.

\hypertarget{section-8}{%
\section*{7}\label{section-8}}
\addcontentsline{toc}{section}{7}

The underground compartment containing radioactive warheads is called
``The Pit''. The warheads are angled at eighty-eight degrees, else the
bottom falls out. Agent-sweat is radioactive. Agents drink it as
refreshment while awaiting the bottom to fall out. That the bottom would
fall out onto Agents' heads is troubling, six, five, four, three,
troubling to one.

One narrative theory is that there is an underground nightclub called
``The Pit,'' where rabbits push up and up. The player piano is
programmed to perform eighty-eight degrees of sound, else the bottom
falls out. Music is the distance of time within sound's detonation. The
speed of sound is measured by the distance between lightening and its
thunder. Agent Yellow never strikes twice. Agent Red is not a rocket
scientist. Agent Orange is the bomb.

The speed of sound is measured by the distance between the bomber and
the detonation of its payload. The warhead is programmed to perform
eighty-eight degrees of separation. Its coin slides into its slot. Its
music is the bomb, whose bottom falls out. The speed of light is
measured by the detonation of the bomb within Agents' own bottoms
falling out.

\hypertarget{section-9}{%
\section*{8}\label{section-9}}
\addcontentsline{toc}{section}{8}

One narrative theory is that one never dies in sculpture. Kenji Yanobe's
sculptural objects evoke a nostalgia for an apocalypse that never
happened. It is a nostalgia for the future of a past that is continually
deferred to a present that hasn't happened yet. It's already happening.
The so-called-and-ever-wished-for last instance never arrives, and yet
the present is-and-always-is this very last instance. The apocalypse is
always happening.

One narrative theory is that one consumes or is consumed. The artist
Kenji Yanobe has produced a never-ending simulation of a survival cart
race, which will protect Agents from radioactivity as long as Agents
keep inserting coins into their nuclear-safe capsules. Agents must exit
the capsule in order to insert the coins into its exterior
anthropomorphic coin slots. There is no future here, only the present
which enables Agents to survive as long as Agents keep consuming. The
rabbits chew through the string that is used for the coin-slot coin
trick. The only safety is to insert coins into another's coin slots in
hopes that the other will do the same. Distance is the opposite of
safety.

The artist Kenji Yanobe stands at the sentry post facing out, toward the
ruins of the future. Rabbits have infrared vision that can see through
the nuclear-safe capsule. The inside is lined with rabbit fur. It's a
never-ending simulation of eighty-eight degrees inside. Agent-sweat is
radioactive. Agents drink it as refreshment while awaiting the survival
cart race.

\hypertarget{section-10}{%
\section*{9}\label{section-10}}
\addcontentsline{toc}{section}{9}

The social grid exists so as to become invisible. The social grid exists
so that it doesn't have to. That the invisibility of the social grid
might provoke the realization of the ideological landscape is troubling
to one. Even underground, though that's a dream. The so-called social
elements are the imaginary plane of social beings. The imaginary plane
stretches out into the distance, held up by the many social beings.
History never disappears.

The imaginary plane reaches from one pole to the next, grounded by the
many social bearings. The present disappears into the historical
landscape, though it's still here. That such disappearance into the
historical landscape might be the forgetting of social bearings is
troubling to one. Even underground, though that's a dream. History is
the imaginary plane of social bearings. Here it is now, stretching out
into the distance, held up by the many social bearings.

The social grid extends indefinitely into the future. Social beings
labor at building it outward. The apocalypse of totality is continually
deferred. Humility is nervous awareness that Agents are grid-workers.
That the social bearings necessary to destroy the grid are the same
social bearings that build the grid is troubling, eight, seven, six,
five, four, three, agents, troubling, to, one.

\hypertarget{section-11}{%
\section*{10}\label{section-11}}
\addcontentsline{toc}{section}{10}

The Agents are in their survival carts. The carts fart perfumed
radioactivity through anthropomorphic trumpets. Agent Orange breathes
its own air. The music is air-raid sirens that go up and up. The
rabbits' infrared vision burns through the string that is used for the
coin-slot coin trick. Agent Red eats its own waste. It's a never-ending
simulation of eighty-eight degrees inside. Agent Yellow drinks its own
sweat. Agent-sweat is radioactive. Agents drink it as refreshment while
awaiting the coins to run out.

Agent A puts its coin into Agent B's coin slot. Agent B puts its coin
into Agent C's coin slot. Agent C puts its coin into Agent A's coin
slot. There are no kings to make more coins, so the grid-poles are
dismantled and melted into currency. Without the poles, the imaginary
plane's bottom falls out. The grid collapses into itself. The perimeters
collapse into their centers. Agents trade in their narrative theories
for the few remaining coins. The sun is setting. The social bearings are
detonated. In the last instance coins fall to the ground in piles that
go up and up. The sky is falling.

Agent Yellow's survival cart is red. Agent Red's survival cart is
orange. Agent Orange's survival cart is yellow. As they collapse into
one another, the detonation provokes the sunset. From a distance the
shutters click shut. That the sunset is a film that is aesthetically
pleasing is troubling, ten, nine, eighty-eight degrees, seven, sixty-six
rabbits, five, four, whom, three, agents, troubling, to, one. Action.

\hypertarget{postscript}{%
\section*{Postscript}\label{postscript}}
\addcontentsline{toc}{section}{Postscript}

One narrative theory is that there is always another narrative theory.
Action is choosing. The film reel is stuck in the machine. The so-called
sunset is merely the melting celluloid of one possible future. The
machine is the apocalyptic imaginary that desires its own foreclosure.
Agent Orange has killed off Agents A and B, Agents Red and Yellow. But
Agent C has managed to elude such foreclosure. New grids are being
constructed by social beings with different social bearings. Action is
choosing made material by certain social bearings. Humility is nervous
awareness that action is total. It's already happening.

One narrative theory is that one never dies in writing. The writing
cannot in actuality provide safety. It is a machine desiring its own
apocalypse. In actuality it is its own social being, desiring its own
narrative theory. Here it is now, stretching out into the distance, held
up by its social bearings. Agent C is writing itself onto the imaginary
plane. The script has yet to be finished. It's already happening.

~

\hypertarget{jacques-debrot-narrative-is-boring}{%
\chapter{Jacques Debrot: Narrative is
boring}\label{jacques-debrot-narrative-is-boring}}

\lettrine[lines=3, findent=0em, nindent=0.1em, lhang=0]{N}{arrative is boring}
because it precludes the direct actualization of the world through our
perception. But there is no alternative; narrative is both a disabling
and a necessary condition of perception. For this reason, boredom is the
affective condition that writing predominantly negotiates and exploits.
However, it has always struck me that boredom, despite its
stigmatization, is, in reality, a complexly ambivalent emotion.
Somewhere Roland Barthes--in \textsc{Roland Barthes} I
believe---describes boredom, suggestively, as a form hysteria (or is it
stupidity that he calls hysterical?). Whatever the case, I try to make
my writing both as boring and as stupid as I possibly can---in Barthes'
words, ``I agree to pluralize myself, to permit free cantons of
stupidity to live inside me.'' Error, play, bad taste, incomprehension,
artifice, and a lack of truth or reasonableness are, thus, the narrative
potentialities I attempt to fulfill.

Paradoxically, with the loss of narrative's legitimizing and explanatory
functions, in much of the most interesting contemporary writing there
has been an increasing recourse to, rather than a retreat from,
narrative forms. My own work is transversed by multiple, parataxic
narratives and entangled discourses. Ideally (and impossibly) I would
like to produce a narrative that would need to be read in the same way
one ``reads'' a painting; that is, by orienting oneself gradually, by
means of various intensities and expressive deformations of
surface---the textual equivalents, if they exist, of the mechanisms of
collaging, dribbling, scraping, scratching, scumbling, attaching,
interposing, reversing, permutting, and so on---to the entire
environment and atmosphere of the writing; a method aligned not so much
with an aesthetics of fragmentation, but rather, one that proceeds
through the \emph{accretion} and flux of many ephemeral ruptures and
stases.

This, obviously, entails a fundamental reconsideration of what is
traditionally classified as narrative form. However, for me, the central
question is not a taxanomical, or a typological one. More important than
the question of whether writing such as mine is narrative, is the
problem of how, in what sense, this writing can be read as narrative.

The situation of communication---the roles of author and reader, the
relations of subject positions within the narrative itself, as well as
the historical contexts in which the meanings of the narrative
circulate---cannot be excluded from the kind of reading my work attempts
to invite. But even more specifically, I am interested by the aesthetic
situation in which my writing occurs. By this I mean I am concerned most
of all with my feelings, and with my tastes and my preferences. Like
everyone's, my experience of art, as Thierry du Duve argues, ``begins
and ends with aesthetic judgement.'' It is impossible---at least on
aesthetic grounds---to be insincere to one's tastes. However our
preferences are not objective, and our feelings, although they are our
own, are also social constructions. Inasmuch as literature and narrative
are no longer given domains but generative practices, what I want to
experience in the act of writing, or reading is to feel, as du Duve puts
it, the conflict of values both within myself, and within and for the
cultural field as a conflict of feelings, ``to give in to the reflexive
feeling of dissent. . . . To give my consent to the felt absence of
concensus.''

It is this ongoing dissent--that is to say, my resistance to any single
unitary or consummating construction---that \emph{is} the self from
which, as much as the goal toward which, I write.

~

~

\hypertarget{jeff-derkson-text-and-the-site-of-writing}{%
\chapter{Jeff Derkson: ``Text'' and the Site of
Writing}\label{jeff-derkson-text-and-the-site-of-writing}}

\begin{quote}
``That the widespread textualization of the outside world in
contemporary thought (the body as a text, the state as a text,
consumption as a text) should itself be seen as a fundamental form of
postmodern spatialization\ldots{}'' (158)

--Jameson, \textsc{Postmodernism: Or the Cultural Logic Of Late
Capitalism}

``For its part the poststructuralist version of postmodernism worked to
exceed both formal aesthetic categories (the disciplinary order of
painting, sculpture, and so on) and traditional cultural distinctions
(high versus mass culture, autonomous versus utilitarian art) with a new
model of art as text.'' (72).

--Hal Foster, \textsc{The Return of the Real}

``A provisional conclusion might be that in advanced art practices of
the past thirty years the operative definition of the site has been
transformed from a physical location---grounded, fixed, actual---to a
discursive vector---ungrounded, fluid, virtual.''

Miwon Kwon, \textsc{One Place After Another: Notes on Site Specificity},
October 80 (Spring 97): 95.

~
\end{quote}

\lettrine[lines=3, findent=0em, nindent=0.1em, lhang=0]{I}{’ve started with this}
top-heavy loading of quotations to highlight a general movement from
textualization of (generally) the ``outside world'' and place/site
alongside a similar movement within (generally) art and site-specific
art in order to establish a frame to address how the ``place'' of
writing and cultural production could be rearticulated. Aside from a
poststructuralist tendency to read structures as texts, how does this
``textual turn'' implicate writing as a constructive act? Has the
modernist literary project of writing an imagined world been deflected
through poststructuralism so now the world is imagined as not only
written, but as text?

Parallel to this ``textual turn,'' which leads intriguingly to Mowing
Kwon's vague ``discursive vectors,'' there has been both a related
mapping impulse and a constructive intent directed at ``place.'' The
mapping impulse is both ontological and geographic (writing as a mapping
of a mind or of subjectivity, writing as part of the process of
realizing ``place''). In popular media, texts, particularly novels, are
given a primarily ontological role, of telling us something about the
places we live in, and by extension to tell us something about
``ourselves'' or to illuminate the author as subject. Geographically,
the constructive intent is perhaps clearest in a national literature's
assembling of images and icons to create the imagined place of a nation
and to the related levels of regional literatures and urbanist texts.

In contemporary art, a sort of sociological turn and mapping has
emerged. Hal Foster cites Dan Graham``s''Homes for America" which
``mapped'' typologies of American suburban houses as well as the taste
that they constructed and were based on, and ``Twenty-Six Gas Stations''
by Ed Ruscha as examples of the ``sociological mapping\ldots{}implicit
in some conceptual art'' (185). More recently, and more ironically and
more internationally, there is Komar \& Melamid's ``The People's
Choice'' ({http://www.diacenter.org/km/}) which uses official polling
agencies to survey a nation's preferences in visual art based on
approximately ten preferences (ranging from ``Favorite color'' to
``Prefer indoor or outdoor scenes'') and then realizes the ``most wanted
painting'' and ``most unwanted painting'' utilizing the information.
While this project is lightly politicized around issues of taste and in
its adaptation of the current political tendency to base policies on
poll results, it is linked to current site-specific art practices which,
as Kwon notes, ``routinely engage the collaborative participation of
audience groups for the conceptualization and production of the
work\ldots{}.'' From this collaboration, these site-oriented works ``are
seen as a means to strengthen art's capacity to penetrate the
sociopolitical organization of contemporary life with greater impact and
meaning. In this sense the possibilities to conceive of the site as
something more than place---as repressed ethnic history, a political
cause, a disenfranchised social group---is a crucial conceptual step in
redefining the public role of art and artists''(96). A recent
example---and there are many---of this laying bare of the historical
determinants of place is realized in Stan Douglas' Nu\emph{tka} which
presents a ``Canadian Gothic'' of late-nineteenth century Nootka Sound
on the Northwest Coast of B.C. through interlacing video images of the
area and disembodied voices of the Spanish and English colonizers.

This chronotopic imagining of place as the site of a repressed racial
(ethnic, class and gendered) history has been an obvious project of
literature. In Canadian literature, this project was determined both
discursively and historically. Discursively, the embarrassingly narrow
yet dominant critical trope assigned to the national literature the role
of providing a history to a country strategically defined as having
none. This necessitated a ``return of the repressed'' in literature to
counter the dominant literary (national) historical projects.
Historically, Canada has imagined itself as bicultural and this
framework worked to suppress the histories of groups other than the
French and English. Within this very generalized framework, official and
aestheticized responses to this historical repression have emerged.
Small-town history chronicles that celebrate a town such as
\textsc{Morden in Manitoba}, which can be bought at City Hall (here in
Austria these projects are very similar and are called Heimatbuch
{[}roughly, ``homeland book''{]}) or a book such as Andreas Schroeder's
\textsc{The Mennonites} which provides, in a coffeetable book format, a
history of the Mennonites in Canada. More well known, and with a larger
cultural impact, Joy Kogawa's novel \textsc{Obasan} brought forward the
history of the Japanese-Canadian internment. Yet the moment that
novelists are taken as historians is as problematic as when artists are
believed to be sociologists or social workers. For instance, Michael
Ondaatje's \textsc{In the Skin of a Lion} aestheticizes the history of
working-class people in Toronto, as if the workers appreciated the
modernist aesthetics of the worksite and the solidarity of shared labor
rather than complained about the relations of production and wages.

That these textual and visual models present place and site as being a
manifestation of history can lead to a deflecting of the present. For
instance, Roy Miki speculates that \textsc{Obasan} could, from one
vantage point, ``become an object of knowledge as a Canadianized text
that teaches us about racism in our past'' (\textsc{Broken Entries}
143). The implication is that racism is relegated, chronotopically, to
the past. Douglas' \textsc{Nu\emph{tka}}, in its return of repressed
history, could mimic the trope of First Nations' culture as a relic of
the past, as being only determined by European actions, as the passive
site where the history of dominant culture is acted out.

Here I want to speculate on (and politicize in a different way), Kwon's
idea of place as a ``discursive vector'' as a means to situate oneself
within a present site. I write ``politicize'' because it is possible to
propose a move from text to discourse as a movement from a static
structure open to analysis (whether virtuoso or standard) to a
constructing determinant of place and subjectivity. Here place/site
would not solely be determined by histories (dominant or repressed:
emergent or residual), but by competing constitutive discourses that
both affirm and erode the local/national and the everyday in the name of
the global. Without solely reducing globalization to the effect of a
discourse, it is possible to perhaps clarify the effects of
globalization by understanding them as being discursively enacted at one
level, as having a constitutive effect. Discursive is not synonymous
with ``fluid,'' ``ungrounded,'' as Kwon proposes, but is precisely
constitutive and grounding. In this formulation, the particularisms of a
place/site and its histories are not just the oppositional force to
globalization (or the corrective to dominant historical narratives), but
an aspect of place that can be utilized by globalism; that is, a
dialectic of local and global, or site and nonsite (if a place is
imagined as siteless in its loss of particularities due to
globalization).

My own turn toward the discursive effects of globalism in relation to
place/site arose because of the contradictions I saw in the city in
which I was living. The city of Calgary in Alberta imagines itself as a
regionally-based ``open-for-business'' kind of city, free of, but also
wary of, the ills of larger cities, yet a city of international
standards. This discourse of regionalism serves to cloak the existing
relations of production which Calgary---as an ``oil town''---is linked
into, and determined by. Global capital through the multinational oil
and gas industry whose corporate logos hover above Calgary's gridded
streets is a key determinant of both the social relations of the city
and how the city imagines itself. When the oil industry is profitable
(due to high crude prices or a manufactured crisis), the city thrives
and new homes are built, rents rise, corporate headquarters shuffle,
cigar bars and fusion-cooking restaurants open---in short, a lot of
money is made by a few people but the scramble for profit is on. This
imaginary masking of real relations, a kind of cognitive masking, is
what Arjun Appadurai cites as fetishistic: ``The locality (both in the
sense of the local factory or site of production and then extended sense
of the nation-state), becomes a fetish that disguises the dispersed
forces that actually drive the production process'' (\textsc{Modernity
At Large}, 42). In Calgary, the global forces then have a direct effect
on the planning and layout of the city, as well as the architecture; the
social space of the city (itself a constituting aspect), is determined
by the management of the effects of globalization and localization.
Seeing how the effective discourses of place/region and nation could
serve as a mystifying factor, blocking the real relations of the city, I
moved away from an investigation of place and toward the discursive
determination of a place/site.

As a writer, the emphasis on the constitutive discourses of site/place
can supplement site as repressed histories. These discourses must also
be seen, alongside repressed histories, as historical developments, as
constitutive elements of the repressed history and emergent history. The
site of writing then becomes imagined as a place of intersecting
discourses: not groundless and fluid but both determined and
determining.

\hypertarget{robert-gluxfcck-long-note-on-new-narrative}{%
\chapter{Robert Glück: Long Note on New
Narrative}\label{robert-gluxfcck-long-note-on-new-narrative}}

\lettrine[lines=3, findent=0em, nindent=0.1em, lhang=0]{T}{o talk about}
the beginnings of New Narrative, I have to talk about my friendship with
Bruce Boone. We met in the early seventies through the San Francisco Art
Institute's bulletin board: Ed and I wanted to move and Bruce and Burton
wanted to move---would we all be happy living together? For some reason
both couples dropped the idea and remained in our respective flats for
many years. But Bruce and I were poets and our obsession with Frank
O'Hara forged a bond.

I was twenty-three or twenty-four. Bruce was seven years older. He was a
wonderful teacher. He read to transform himself and to attain a correct
understanding. Such understanding was urgently political.

Bruce had his eye on the catastrophic future, an upheaval he predicted
with a certain grandeur, but it was my own present he helped me find. I
read and wrote to invoke what seemed impossible--relation itself--in
order to take part in a world that ceaselessly makes itself up, to
``wake up'' to the world, to recognize the world, to be convinced that
the world exists, to take revenge on the world for not existing.

To talk about New Narrative, I also have to talk about Language Poetry,
which was in its heroic period in the seventies. I treat diverse poets
as one unit, a sort of flying wedge, because that's how we experienced
them. It would be hard to overestimate the drama they brought to a Bay
Area scene that limped through the seventies, with the powerful
exception of feminist poets like Judy Grahn, and the excitement
generated by movement poetry. Language Poetry's Puritan rigor, delight
in technical vocabularies, and professionalism were new to a generation
of Bay Area poets whose influences included the Beats, Robert Duncan and
Jack Spicer, the New York School (Bolinas was its western outpost),
surrealism and psychedelic surrealism.

Suddenly people took sides, though at times these confrontations
resembled a pastiche of the embattled positions of earlier avant-guards.
Language Poetry seemed very ``straight male''---though what didn't?
Barrett Watten's \textsc{Total Syntax}, for example, brilliantly
established (as it dispatched) a lineage of fathers: Olson, Zukofsky,
Pound, etc.

If I could have become a Language poet I would have; I craved the
formalist fireworks, a purity that invented its own tenets. On the snowy
mountain-top of progressive formalism, from the highest high road of
modernist achievement, there was plenty of contempt heaped on less
rigorous endeavor. I had come to a dead end in the mid-seventies like
the poetry scene itself. The problem was not theoretical---or it was: I
could not go on until I figured out some way to understand where I was.
I also craved the community the Language Poets made for themselves.

The questions vexing Bruce and me and the kind of rigor we needed were
only partly addressed by Language Poetry which, in the most general
sense, we saw as an aesthetics built on an examination (by subtraction:
of voice, of continuity) of the ways language generates meaning. The
same could be said of other experimental work, especially the
minimalisms, but Language Poetry was our proximate example.

Warring camps and battle lines drawn between representation and
non-representation---retrospection makes the argument seem as arbitrary
as Fancy vs.~Imagination. But certainly the ``logic of history'' at that
moment supported the idea of this division, along with the struggle to
find a third position that would encompass the whole argument.

I experienced the poetry of disjunction as a luxurious idealism in which
the speaking subject rejects the confines of representation and
disappears in the largest freedom, that of language itself. My
attraction to this freedom and to the professionalism with which it was
purveyed made for a kind of class struggle within myself. Whole areas of
my experience, especially gay experience, were not admitted to this
utopia, partly because the mainstream reflected a resoundingly coherent
image of myself back to me---an image so unjust that it amounted to a
tyranny that I could not turn my back on. We had been disastrously
described by the mainstream---a naming whose most extreme (though not
uncommon) expression was physical violence. Political agency involved at
least a provisionally stable identity.

Meanwhile, gay identity was also in its heroic period---it had not yet
settled into just another nationalism and it was new enough to know its
own constructedness. In the urban mix, some great experiment was
actually taking place, a genuine community where strangers and different
classes and ethnicities rubbed more than shoulders. This community was
not destroyed by commodity culture, which was destroying so many other
communities; instead, it was founded in commodity culture. We had to
talk about it. Bruce and I turned to each other to see if we could come
up with a better representation---not in order to satisfy movement
pieties or to be political, but in order to be. We (eventually we were
gay, lesbian, and working class writers) could not let narration go.

Since I'm confined to hindsight, I write as though Bruce and I were
following a plan instead of stumbling and groping toward a writing that
could join other literatures of the present. We could have found
narrative models in, say, Clark Coolidge's prose, so perhaps narrative
practice relates outward to the actual community whose story is being
told. We could have located self-reference and awareness of artifice in,
say, the novels of Ronald Firbank, but we didn't. So again, our use of
language that knows itself relates outward to a community speaking to
itself dissonantly.

We were fellow travelers of Language Poetry and the innovative feminist
poetry of that time, but our lives and reading lead us toward a hybrid
aesthetic, something impure. We (say, Bruce Boone, Camille Roy, Kevin
Killian, Dodie Bellamy, Mike Amnasan, myself and, to include the dead,
Steve Abbott and Sam D'Allesandro) are still fellow travelers of the
poetries that evolved since the early eighties, when writers talked
about ``nonnarrative.'' One could untangle that knot forever, or build
an aesthetics on the ways language conveys silence, chaos,
undifferentiated existence, and erects countless horizons of meaning.

How to be a theory-based writer?---one question. How to represent my
experience as a gay man?---another question just as pressing. These
questions lead to readers and communities almost completely ignorant of
each other. Too fragmented for a gay audience? Too much sex and
``voice'' for a literary audience? I embodied these incommensurates so I
had to ask this question: How can I convey urgent social meanings while
opening or subverting the possibilities of meaning itself? That question
has deviled and vexed Bay Area writing for twenty-five years. What kind
of representation least deforms its subject? Can language be aware of
itself (as object, as system, as commodity, as abstraction) yet take
part in the forces that generate the present? Where in writing does
engagement become authentic? One response, the politics of form,
apparently does not answer the question completely.

One afternoon in 1976, Bruce remarked on the questions to the reader I'd
been throwing into poems and stories. They were self-consciously
theatrical and they seemed to him to pressure and even sometimes to
reverse the positions of reader and writer. Reader/writer dynamics
seemed like a way into the problems that preoccupied us, a toe in the
water!

From our poems and stories, Bruce abstracted text-metatext: a story
keeps a running commentary on itself from the present. The commentary,
taking the form of a meditation or a second story, supplies a succession
of frames. That is, the more you fragment a story, the more it becomes
an example of narration itself--narration displaying its devices--while
at the same time (as I wrote in 1981) the metatext ``asks questions,
asks for critical response, makes claims on the reader, elicits
comments. In any case, text-metatext takes its form from the dialectical
cleft between real life and life as it wants to be.''\footnote{Hayden
  White notes this etymology in \textsc{The Content of the Form}\emph{.}}

We did not want to break the back of representation or to ``punish'' it
for lying, but to elaborate narration on as many different planes as we
could, which seemed consistent with the lives we lead. Writing can't
will away power relations and commodity life; instead, writing must
accept its relation to power and recognize that at present group
practice resides inside the commodity. Bruce wrote, ``When evaluating
image in American culture, isn't it a commodity whether anyone likes it
or not? You make your additions and subtractions from that point
on.''\footnote{Frederich Engels writes in his letter to Joseph Block,
  September 21, 1890: ``We make our history ourselves, but, in the first
  place under very definite assumptions and conditions.'' Glossing this
  line in \textsc{Marxism and Literature}, Raymond Williams writes,
  ``What this restores, as against the alternative development, is the
  idea of direct agency: `we make our history ourselves.' The `definite'
  or `objective' assumptions and conditions are then the qualifying
  terms of this agency: in fact `determination' as `the setting of
  limits'\,''.}

In 1978, Bruce and I launched the Black Star Series and published my
\textsc{Family Poems} and his \textsc{My Walk With Bob}, a lovely
book.\footnote{Copjec, Joan, \textsc{Read My Desire}, p.~148.} In
``Remarks on Narrative''---the afterword of \textsc{Family
Poems}---Bruce wrote, ``As has now been apparent for some time, the
poetry of the '70s seems generally to have reached a point of
stagnation, increasing a kind of refinement of technique and available
forms, without yet being able to profit greatly from the vigor, energy
and accessibility that mark so much of the new Movement writing of gays,
women and Third World writers, among others. Ultimately this impasse of
poetry reflects conditions in society itself.''\footnote{Lacan, "The
  Mirror Stage.\emph{"}}

We appreciated the comedy of mounting an offensive (``A critique of the
new trends toward conceptualization, linguistic abstraction and process
poetry'') with those slenderest volumes. My poems and stories were set
``in the family,'' not so antipsychological as they might have been
given that we assumed any blow to interiority was a step forward for
mankind.

We contended with the Language Poets while seeking their attention in
the forums they erected for themselves. We published articles in
\textsc{Poetics Journal} and \textsc{L=A=N=G=U=A=G=E}, and spoke in talk
series and forums---a mere trickle in the torrent of their critical
work. If Language Poetry was a dead end, what a fertile dead end it was!

New Narrative was in place by the time \textsc{Elements of a Coffee
Service} was published by Donald Allen's Four Season's Foundation in
1982, and Hoddypole published Bruce's novel, \textsc{Century of Clouds}
in 1980. We were thinking about autobiography; by autobiography we meant
daydreams, nightdreams, the act of writing, the relationship to the
reader, the meeting of flesh and culture, the self as collaboration, the
self as disintegration, the gaps, inconsistences and distortions, the
enjambments of power, family, history and language.

Bruce and I brought high and low together between the covers of a book,
mingling essay, lyric, and story. Our publishing reflected those
different modes: stories from \textsc{Elements} appeared in gay
anthologies, porn magazines, \textsc{Social Text}, and \textsc{Soup}.
Bruce wrote about Georges Bataille for \textsc{The Advocate}\emph{. } I
wanted to write with a total continuity and total disjunction since I
experienced the world (and myself) as continuous and infinity divided.
That was my ambition for writing. Why should a work of literature be
organized by one pattern of engagement? Why should a ``position'' be
maintained regarding the size of the gaps between units of meaning? To
describe how the world is organized may be the same as organizing the
world. I wanted the pleasures and politics of the fragment and the
pleasures and politics of story, gossip, fable and case history; the
randomness of chance and a sense of inevitability; sincerity while using
appropriation and pastiche. When Barrett Watten said about \textsc{Jack
the Modernist}, ``You have your cake and eat it too,'' I took it as a
great compliment, as if my intention spoke through the book.

During the seventies, Bruce was working on his doctorate at UC Berkeley.
His dissertation was a structuralist and gay reading of O'Hara, that is,
O'Hara and community, a version of which was published in the first
issue of \textsc{Social Text} in 1979. He joined the Marxism and Theory
Group at St.~Cloud which gave birth to that journal. Bruce also wrote
critical articles (especially tracking the ``gay band'' of the Berkeley
Renaissance).(\footnote{Lacan, "The Mirror Stage.\emph{"}}) Bruce
introduced me to most of the critics who would make a foundation for New
Narrative writing.

Here are a few of them:

Georg Lukacs: In \textsc{The Theory of the Novel}, Lukacs maintains that
the novel contains---that is, holds together---incommensurates. The epic
and novel are the community telling itself its story, a story whose
integration becomes increasingly hard to achieve. \textsc{The Theory of
the Novel} leads to ideas of collaboration and community that are not
naive, that is, to narrative that questions itself. It redistributes
relations of power and springs the writer from the box of psychology,
since he becomes the community speaking to itself. I wrote
``Caricature,'' a talk given at 80 Langton in 1983, mostly using Lukacs'
book, locating instances of conservative and progressive communities
speaking to themselves: ``If the community is a given, so are its
types.''

Louis Althusser: His essay, ``Ideological State Apparatuses,'' refigures
the concept of base/superstructure. Terry Eagleton rang the following
change on Althusser's bulky formula: Ideology is the imaginary
resolution of real contradictions. By 1980, literary naturalism was
easily deprived of its transparency, but this formula also deprives all
fantasy of transparency, including the fantasy of personality. If making
a personality is not different from making a book, in both cases one
could favor the ``real contradictions'' side of the formula. If
personality is a fiction (a political fiction!) then it is a story in
common with other stories---it occurs on the same plane of experience.
This ``formula'' sets those opacities---a novel, a personality---as
equals on the stage of history, and supports a new version of
autobiography in which ``fact'' and ``fiction'' inter-penetrate.

Althusser comes with a lot of baggage. For example, he divided science
from ideology, and ideology from theory. Frankly, Bruce and I pillaged
critical theory for concepts that gave us access to our experience. In
retrospect, it might be better simply to ``go with'' cultural studies.
To the endless chain of equal cultural manifestations (a song by REM,
the Diet of Worms, Rousseau's \textsc{Confessions}), we add another
equal sign, attaching the self as yet another thing the culture
``dreamed up.''

Georges Bataille: Bataille was central to our project. He finds a
counter-economy of rupture and excess that includes art, sex, war,
religious sacrifice, sports events, ruptured subjectivity, the
dissolution of bodily integuments---``expenditure'' of all kinds.
Bataille showed us how a bath house and a church could fulfill the same
function in their respective communities.

In writing about sex, desire and the body, New Narrative approached
performance art, where self is put at risk by naming names, becoming
naked, making the irreversible happen---the book becomes social practice
that is lived. The theme of obsessive romance did double duty,
de-stabling the self and asserting gay experience. Steve Abbott wrote,
``Gay writers Bruce Boone and Robert Glück (like Acker, Dennis Cooper or
the subway graffitists again) up the ante on this factuality by weaving
their own names, and those of friends and lovers, into their work. The
writer/artist becomes exposed and vulnerable: you risk being foolish,
mean-spirited, wrong. But if the writer's life is more open to judgement
and speculation, so is the reader's.'' (\footnote{Gluck, ``Caricature,''
  19.})

Did we believe in the ``truth and freedom'' of sex? Certainly we were
attracted to scandal and shame, where there is so much information. I
wanted to write close to the body---the place language goes reluctantly.
We used porn, where information saturates narrative, to expose and
manipulate genre's formulas and dramatis personae, to arrive at ecstacy
and loss of narration as the self sheds its social identities. We wanted
to speak about subject/master and object/slave. Bataille showed us that
loss of self and attainment of nothingness is a group activity. He
supplied the essential negative, a zero planted in the midst of
community. His concept of transgression gave us lots of fuel, as did his
novels of philosophic pornography.

Now I'd add that transgressive writing is not necessarily about sex or
the body---or about anything one can predict. There's no manual;
transgressive writing shocks by articulating the present, the one thing
impossible to put into words, because a language does not yet exist to
describe the present. Bruce translated Bataille's \textsc{Guilty} for
Lapis Press when I worked as an editor there. We hammered out the
manuscript together, absorbing Bataille gesturally.

Five more critics. Walter Benjamin: for lyrical melancholy (which reads
as autobiography) and for permission to mix high and low. V.N.
Voloshinov: for discovering that meaning resides within its social
situation, and that contending powers struggle within language itself.
Rolland Barthes: for a style that goes back to autobiography, for the
fragment, and for displaying the constructed nature of story---``baring
the device.'' Michel Foucault: for the constructed nature of sexuality,
the self as collaboration, and the not-to-be-underestimated example of
an out gay critic. (Once at 18th and Castro, Michel pierced Bruce with
his eagle gaze and Bruce was overcome!---he says.) Julia Kristeva: for
elaborating the meaning of abjection in \textsc{Powers of Horror}.

Our interest in Dennis Cooper and Kathy Acker produced allegiances and
friendships with those writers. Kathy moved to San Francisco in the fall
of 1981; while getting settled she stayed with Denise Kastan, who lived
downstairs from me. Denise and I co-directed Small Press Traffic. Kathy
was at work on \textsc{Great Expectations}. In fact, Denise and I appear
in it; we are the whores Danella and Barbraella. Kathy's writing gave
Bruce, Steve Abbott and myself another model, evolved far beyond our own
efforts, for the interrogation of autobiography as ``text'' perpetually
subverted by another text. Appropriation puts in question the place of
the writer---in fact, it turns the writer into a reader.

Meanwhile, Bruce and I were thinking about the painters who were
rediscovering the figure, like Eric Fischl and Julian Schnabel. They
found a figuration that had passed through the flame of abstract
expressionism and the subsequent isms, operating through them. It made
us feel we were part of a crosscultural impulse rather than a local
subset. Bruce wrote, ``With much gay writing and some punk notoriously
(Acker the big example), the sexual roots of aggression come into
question. There's a scream of connection, the figure that emerges
ghostly: life attributed to those who have gone beyond. So in Dennis
Cooper's \textsc{Safe} there's a feeling-tone like a Schnabel
painting:the ground's these fragments of some past, the stag, the Roman
column, whatever---on them a figure that doesn't quite exist but would
maybe like to. The person/persona/thing the writer's trying to construct
from images---''(\footnote{Steve Abbott, ``Notes on Boundaries, New
  Narrative,'' \textsc{Soup: New Critical Perspectives} \#4) 81.})

In 1976, I started volunteering in the non-profit bookstore Small Press
Traffic and became co-director not long after. From 1977 to 1985, I ran
a reading series and held free walk-in writing workshops at the store.
The workshops became a kind of New Narrative laboratory attended by
Michael Amnasan, Steve Abbott, Sam D'Allesandro, Kevin Killian, Dodie
Bellamy, Camille Roy, and many other writers whose works extend my own
horizon. I would start by reading some piece of writing that interested
me: Chaucer, Robert Smithson, Lydia Davis, Ivan Bunin, Jim Thompson, a
book of London street games, Thomas Wyatt, Sei Shonagon. We were
aspiring to an ideal of learning derived as much from Spicer and Duncan
as from our contemporaries.

Most writers we knew were reading theory. Later, guided by Bruce, we
started a left reading group at Small Press Traffic, attended by Steve
Benson, Ron Silliman, Denise Kastan, Steve Abbott, Bruce, myself and
others. The personal demolished the political, and after a few months we
disbanded. From that era I recall Ron's epithet (which Bruce and I
thought delicious) \emph{The Small Press Traffic School of
Dissimulation}.

More successful was the Left/Write Conference we mounted in 1981 at the
Noe Valley Ministry. The idea for a conference was conceived in the
spring of 1998 by Bruce and Steve Abbott, who sent letters to thirty
writers of various ethnicities and aesthetic positions. Steve was a
tireless community builder, and Left/Write was an expression of New
Narrative's desire to bring communities together, a desire which
informed the reading series at Small Press Traffic, Steve Abbott's
\textsc{Soup} (where the term New Narrative first appeared), Michael
Amnasan's \textsc{Ottotole}, Camille Roy and Nayland Blake's
\textsc{Dear World}, Kevin Killian and Brian Monte's \textsc{No
Apologies}, and later Kevin and Dodie Bellamy's \textsc{Mirage}. We felt
urgent about it, perhaps because we each belonged to such disparate
groups. To our astonishment, three hundred people attended Left/Write,
so we accomplished on a civic stage what we were attempting in our
writing, editing and curating: to mix groups and modes of discourse.
Writers famous inside their own group and hardly known outside, like
Judy Grahn and Erica Hunt, spoke and read together for the first time.

Out of that conference the Left Writers Union emerged; soon it was
commandeered by its most unreconstructed faction which prioritized gay
and feminist issues out of existence. At one meeting, we were instructed
to hold readings in storefronts on ground level so the ``masses of San
Francisco'' could walk in!

During this decade--1975-1985--Bruce and I carried on what amounted to
one long gabby phone conversation. We brought gossip and anecdote to our
writing because they contain speaker and audience, establish the
parameters of community and trumpet their``unfair'' points of view. I
hardly ever ``made things up,'' a plot still seems exotic, but as a
collagist I had an infinite field. I could use the lives we endlessly
described to each other as ``found material'' which complicates
storytelling because the material also exists on the same plane as the
reader's life. Found materials have a kind of radiance, the truth of the
already-known.

In 1981 we published \textsc{La Fontaine} ** as a valentine to our
friendship. In one poem, Bruce (and Montaigne!) wrote, ``In the
friendship whereof I speak\ldots{}our souls mingle and blend in a fusion
so complete that the seam that joins them disappears and is found no
more. If pressed to say why I loved him I'd reply, because it was him,
because it was me.''(\footnote{La Fontaine, ``Perukes'' (Black Star
  Series, 1981) 63.})

In using the tag New Narrative, I concede there is such a thing. In the
past I was reluctant to promote a literary school that endured even ten
minutes, much less a few years. Bruce and I took the notion of a
``school'' half seriously, and once New Narrative began to resemble a
program, we abandoned it, declining to recognize ourselves in the
tyrants and functionaries that make a literary school. Or was it just a
failure of nerve? Now I am glad to see the term used by younger writers
in San Francisco, writers in other cities, like Gail Scott in Montreal,
and critics like Earl Jackson, Jr., Antony Easthope, Carolyn Dinshaw,
and Dianne Chisholm. Bruce and I may have been kidding about founding a
school, but we were serious about wanting to bring emotion and subject
matter into the field of innovative writing. I hope that these thoughts
on our project---call it what you will---are useful to those looking for
ways to extending the possibilities of poem and story without
backtracking into the mainstream, or into 19th-century transparency.

~

~

\hypertarget{rob-halpern-notes-toward-essaying-i-narratively}{%
\chapter{Rob Halpern: Notes Toward (Es)saying ``I'',
Narratively}\label{rob-halpern-notes-toward-essaying-i-narratively}}

\begin{quote}
I was already being prepared for {[}the world's{]} tournaments by a
training which taught me to have a horror of faulty grammar instead of
teaching me, when I committed these faults, not to envy others who
avoided them.

~ ~~~~~~~~~~~--Saint Augustine, Bishop of Hippo
\end{quote}

\lettrine[lines=3, findent=0em, nindent=0.1em, lhang=0]{O}{nce upon a time}
is always here. Crisis, my origin: this fault. Someone is called where I
am compelled to explain. To make narrative critically: this response to
crisis. Otherwise it's narrative by default simply because there is no
other way actively to be here.

~

\textsc{From the Sanskrit root} \emph{gna}, or know, our word narrative
derives bound to practice. In Latin, the root meaning appears to have
doubled, manifesting an otherwise latent critical difference. There is
at once \emph{(g)narus}: knowing, expert, skillful; and \emph{narro}: to
relate or tell. (\footnote{Hayden White notes this etymology in
  \textsc{The Content of the Form}\emph{.}}) A fault---mine, and
ours---runs thru the space between. \emph{Fault}: deficiency, lack,
scarcity; as well as slip, error, mistake. Also, more critically:
\emph{fault} as dislocation or a break in continuity of the strata or
vein. \emph{Here}, knowing is held in suspense and knowledge is always
potential. The process of telling/relating needn't connive with already
determined knowledge; narrative might, for example, radically deviate,
turning away from, or perverting, the very practice of authorization,
questioning rather than securing the position of the one who tells.
Here, between the overdetermined and the indeterminate, \emph{narrative
determines}, that is, it constrains and enables, just as it struggles
and responds, testing the tension between knowing and composing.
\emph{Here, we make our narratives ourselves, but, in the first place,
under very definite assumptions and conditions.} (\footnote{Frederich
  Engels writes in his letter to Joseph Block, September 21, 1890: ``We
  make our history ourselves, but, in the first place under very
  definite assumptions and conditions.'' Glossing this line in
  \textsc{Marxism and Literature}, Raymond Williams writes, ``What this
  restores, as against the alternative development, is the idea of
  direct agency: `we make our history ourselves.' The `definite' or
  `objective' assumptions and conditions are then the qualifying terms
  of this agency: in fact `determination' as `the setting of
  limits'\,''.})

\hypertarget{i.-poetically}{%
\section*{I. Poetically}\label{i.-poetically}}
\addcontentsline{toc}{section}{I. Poetically}

\textsc{At once wayless and wayward}, I positions a self \emph{here} at
the site of narrative beginnings: place of mutability, place of many
turns. The crossroads of legend where every way is lost and every way is
open. Site of impossible possibility and paradox, where extremes
coincide. At once \emph{here}: where aversion keeps one in
place---\emph{I: this paralyzed fixture}---disinclined to turn another
way; and \emph{here}: away from which every turn is perversion, an
aberrant way away from the proper. \emph{Here} is forever shifting, the
site of all potential, site of suspense: where \emph{aporia} (knowing no
way, having no word, possessing no resource) and \emph{euporia} (knowing
so many ways, having everything to say, possessing it all in abundance)
converge. Here is where I intervenes in its own story, the permeable
interface of averse states---for example: impoverishment and
abundance---as if with \emph{two faces averse, and conjoined}. In the
space between states I spreads, this narrative surface: at once this
promise of disruption, this promise of continuity.

\textsc{Crisis} is the point at which old orders no longer cohere and
\emph{here} looms suddenly bereft of foundation. Here, for example, is
where Socrates steers his interlocutor: toward \emph{aporia}-place of
lost ways. But what if this crisis could be experienced
\emph{euporically}, as a moment of possibility precisely here where a
way appears impossible. Here occasions the paradoxical coincidence of
antithetical terms. Not the paralyzing absence of a way but the
accessibility of always other ways. Not naively, not as if everything
and anything were possible, but realistically, as if the site of
narrative were the only place where the impossible might open onto
possibility.

\emph{Here}: wayless, I is exposed; and \emph{here}, wayward, I resists.
It seems so familiar: experiencing the simultaneity of political
possibility and impossibility. The fullness of the imagination
paradoxically constrained and enabled by a politically invested
environment which threatens to render that fullness innocuous if not
impotent. How to breathe this air that animates such cynicism, how to
live in this element that triumphantly promotes irony as the ``task of
the day''. Somewhere between the pessimistic and the utopian, where
these extremes coincide, where pessimism reveals its utopian bent: here,
one might ask how to think the moment when ``nothing is possible''
becomes ``only the impossible is possible''. And here another realism
takes the world for its measure and referent: here, where the limits of
history, the bleak horizon that pens us in, meet the pressures of the
imagination.

\textsc{Discrete identities and sovereign subjects}---like the kings and
saints of history---are what they are by virtue of narrative strength.
But I am more interested in narrative frailty, narrative faults: proxies
and prostheses, artificial extensions, these interventions in the world.
Rather than narratives saturated with ends known in advance of every
beginning, I am interested in narratives that destabilize the coherence
of the I which, like a promissory note or stand-in for some future
fulfillment, suffers this contractual obligation to self-identify over
time. The self-same I, untroubled by difference, might fill the space
between disruption and continuity, but it cannot fulfil it. The
self-same I: at once the coherent \emph{effect} of narrative and the
grammatical position to which one recursively returns to ensure that
coherence and whereby narrative possibility is filtered or censored to
preserve an ideal of intelligibility. Coherence, here, is but the
accomplishment of an optimal performance, a labored production of social
and linguistic grammars the faults of which are seamlessly concealed
\emph{so as no appearance of any wound or hurt outwardly might be at
once perceived}. But there are other possible effects, other ways,
veering toward away or versing waywardly. What if, for example, we were
able to avow I's fictitiousness without and affirm our necessary, if
only provisional, belief in it. Here, rather than persuading, narrative
might perform its agnostic potential: the promise of possible
impossibilities that holds sway in the ambiguous space between. This
refusal of certainty, this averse turn away from knowledge, this
perverse turn toward knowing: this is not a declaration of incoherence.
Between impoverishment and abundance (where impoverishment becomes
itself a figure for abundance) the I traverses narrative
orders---disrupting and maintaining---while it is simultaneously
traversed by them.

\textsc{I goes between: a love story.} In Plato's \emph{Symposium},
Diotima locates Love in the space between ignorance and wisdom.
Offspring of Poros (way, resource, expedient) and Penia (wayless,
poverty, lack.) Love paradoxically embodies this antinomy and cannot
give a reason for itself. According to Diotima, Love judges
\emph{without being able to give proper reasons}. ``It is judging things
correctly without being able to give a reason,'' says Diotima, ``surely,
you see that this is not the same as knowledge, for how can knowledge be
devoid of reason?'' Jacques Lacan derives from this---in however faulty
a way---one of his characteristic figures for love: giving that which
one does not have to give. I am thinking of narrative like Love: wayless
and wayward, determined to do precisely what it cannot do. Its reasons
are but ruses. Between \emph{aporia} and \emph{euporia} Love is risked
agnostically, like a narrative that risks its own foundation. Sentenced
and committed, here, I goes between, exposed to every possibility of
rejection and loss, abjection and shame.

\textsc{Augustine and Stein:} Someone is called where I is compelled to
explain. ``Well well is he. Explain my doubts, well well is he explain
my doubts.'': this is how Gertrude Stein begins ``Regularly Regularly In
Narrative'' in \textsc{How to Write}. But the work narrative is called
upon to perform (``Explain my doubts'') is alien to the work narrative
will ultimately do. Explain my doubts: at once the intended imperative
and the impossible end. My doubts will never be properly explained, my
uncertainty never properly reasoned away. Explanation fills the space
between ignorance and knowledge, but it cannot fulfil it. Recalling
another essay of Stein's, explanation is composition while composition
itself is figured as this ``beginning again and again and again
explaining composition and time.'' And taking the perverse turn toward
the threshold of early modernity, there is Saint Augustine determined to
found the authority necessary to begin a book whose very goal it is to
authorize its subject. He, too, is concerned with beginning, explanation
and time about which he writes in the \textsc{Confessions}: ``If no one
asks me, I know; if I want to explain it to someone who asks me, I don't
know.'' Between faith and doubt, past and future, knowing and
not-knowing, between `he' (the partial referent to which past experience
is attached; it interferes with and informs writing's present) and `you'
(a total referent---my audience, my God, the other I love---to which
``something more'' that escapes representation is attached and toward
which narrative tends) between all these faulty binarisms Augustine
navigates a way to forge an `I': the vehicle and the effect of an other
faith. Explain my doubts: the self-imposed imperative to narrate, to
compose a self. A crisis in knowing runs thru composing just as
narrative traces the fault along which I becomes to articulate or align
itself with the world \emph{in ways}. In the long afternoon shadows
thrown by the likes of Augustine and Stein, I am drawn to this, the
faultiness of narrative, its promise.

\hypertarget{to-i-this-plot}{%
\subsection*{To I this plot}\label{to-i-this-plot}}
\addcontentsline{toc}{subsection}{To I this plot}

\begin{quote}
Grant me Lord to know and understand whether a man is first to pray to
you for help or to praise you, and whether he must know you before he
can call you to his aid. If he does not know you, how can he pray to
you? For he may call for some other help, mistaking it for yours.

~~~~~~~~~~ --Saint Augustine
\end{quote}

The problem of beginning, the order of calling and knowing, is a
perennial preoccupation of both narrative and history. In our postmodern
era in which history has illusorily evanesced, narrative promises at
once to rescue the I from an indifferent individualism and to turn the
present away from indifference. Augustine's double bind is here: in
order to gain knowledge and understanding, one must call or praise; but
without knowledge and understanding one cannot know who or what it is
one calls. Mistaken audience is tantamount to mistaken identity. Who is
listening. Who is out there. Here, along the fault that runs between
audience and identity, narrative and history continue to produce each
other. Here, an I is ``emplotted.'' Here, Augustine's address, his ``how
am I able to speak to you?.'' the fragility of its beginning, awakens
him to the danger of being dispossessed of precisely that to which his I
professes to cling: communicative possibility. But who is the subject of
that possibility if the I who depends on past experience for recognition
is displaced, transfigured by the new experience of narrating a
conversion: the advent of an other I? I is this essay, this attempt,
this fragment, this experiment---like composition is explanation---this
reach transversing the pronominal distance between ``Explain my doubts,
well, well is he explain my doubts'' and ``how am I to speak to you?.''
In reaching out, I, discomfitted, doubles over the hinge of this plot
and risks a world. Between belief and doubt, this communication without
guarantees: the very condition of narrative possibility.

\hypertarget{a-crisis-of-audience}{%
\subsection*{A crisis of audience}\label{a-crisis-of-audience}}
\addcontentsline{toc}{subsection}{A crisis of audience}

\begin{quote}
How shall I call upon my God for aid, when the call I make is for my
Lord and my God to come into myself?\ldots{} Does this then mean, O Lord
my God, that there is in me something fit to contain you? Can even
heaven and earth, which you made and in which you made me, contain you?
Or since nothing that exists could exist without you, does this mean
that whatever exists does, in this sense, contain you? If this is so,
since I too exist, why do I ask you to come into me? For I should not be
there at all unless, in this way, you were already present within
me\ldots{} But if I exist in you, how can I call upon you to come to me?
And where would you come from? For you, my God, have said that you fill
heaven and earth, but I cannot go beyond heaven and earth so that you
may leave them to come to me.

~~~~~~~~~~ --Saint Augustine

So then beside as any one can come to be certain of then if it is as it
is that is an audience is what it is what is it if an audience is this,
pretty soon then can feel again that an audience is this, and then
introspection can go on but the habit of this thing makes it cease to be
this, because the audience and is it this keeps going on\ldots{}That is
to say can does any one separate themselves from the land so they can
see it and if they see it are they the audience of it or to it. If you
see anything are you its audience and if you tell anything are you its
audience, and is there any audience for it but the audience that sees or
hears it\ldots{}And all this has so much to do with writing a narrative
of anything that I can almost cry about it.

~~~~~~~~~~ ~--Gertrude Stein
\end{quote}

Bewildered by the question of containment, uncertain exactly which way
to turn, Augustine's move in relation to his audience is paradoxical.
Turning toward what ``there is in me'' he intimates, indeed creates, the
depth of a private self. Simultaneously he turns toward ``heaven and
earth'' revealing the self's dependence on exteriority. But is there an
essential difference between the two dimensions? Must one
\emph{aversion}---one turn away---achieve priority over and above the
another turn? Or might a mutable perversion---a turning aside from truth
or right, a diversion to improper use, a recognition of necessary
distortion, error, fault---intervene thus making coextensive interior
and exterior? Privative and excessive, I becomes this permeable surface,
this interface of in and out: this narrative tension between self and
other. Paradoxically, for Augustine, only a turn away from the
destabilizing question of audience, in other words, only silence or
not-writing, can ensure his possession of a self untroubled by spatial
location; only a turn away from either way will ensure comfort and allow
him to remain confident in the knowledge of such a self \emph{without
knowing it at all}. Self-knowledge as stasis annuls self-knowing as
narrative process. The very fact of Augustine's writing is testament to
his recognition that aversion can offer no way; it cannot ensure a
truthful I, only a proper self, that is, a property which must risk
itself in the act of narrating a story of origins, a return to the
fault. As soon as one begins to narrate oneself, one is adrift, exposed
to the hazard of contingency and doubt, relatively unmoored but bound to
the world nevertheless. In the face of political danger, I mobilizes
narratively: hazard becomes promise as I perverts a way away from
established norms. Looking back I still obeys, only paradoxically.
\emph{Here} crisis troubles the ground on which Augustine is authorized
to speak as ``himself.'' Every word of address threatens to expose the
subject to the vertigo that looms when the secure ground of
``self-knowledge'' gives way to groundlessness. And \emph{here} it is
precisely the move to explain my doubts, to provide for them a
foundation, that makes foundation impossible. Truthful I's are improper
selves, selves held in suspense, in potential and \emph{in situ} in the
space between. In the \textsc{Confessions}, this abyss, as it opens onto
narrative possibility, becomes paradoxically the very ground of love and
faith.

\hypertarget{intentional-fallacies}{%
\subsection*{Intentional fallacies}\label{intentional-fallacies}}
\addcontentsline{toc}{subsection}{Intentional fallacies}

\begin{quote}
They who write narrative and history do not do what they say they will
do when they start out to do what they are about to do.

~~~~~~~~~~ ~--Gertrude Stein
\end{quote}

\begin{quote}
What the storyteller narrates must necessarily be hidden from the actor
himself, at least as long as he is in the act or caught in its
consequences, because to him the meaningfulness of his act is not in the
story that follows.

~~~~~~~~~~ ~--Hannah Arendt
\end{quote}

Around every one of I's articulations, a story begins to congeal, a plot
thickens. Once upon a time, I stands in for this memory of the present,
a lapidary effect of all the stories I've ever told. Here is where I
unbecomes along this fault in the world, this cleavage or interruption
that opens onto narrative in the discomfiting space between where I
never fits. Here is where I encounters oblivion and possibility
simultaneously on the threshold of another sentence, another scene,
another world. ``It was at the threshold of a world such as this that I
stood in peril as a boy'' {[}Augustine{]}. But it is not, as in
Augustine, ``my past foulness and carnal corruption of my soul'' that
defines the narrative fault; rather the I itself---the risk and the
promise.

\hypertarget{something-more}{%
\subsection*{``Something more''}\label{something-more}}
\addcontentsline{toc}{subsection}{``Something more''}

\begin{quote}
Do heaven and earth, then, contain the whole of you, since you fill
them? Or, when once you have filled them, is some part of you left over
because they are too small to hold you? If this is so, when you have
filled heaven and earth, does that part of you which remains flow over
into some other place?

~~~~~~~~~~ ~--Augustine
\end{quote}

\emph{That part of you which remains} cannot be represented, cannot be
assimilated to any way. You---audience, condition and limit of my
self---an unrepresentable horizon coincidently inside and out of me.
``For without you, what am I to myself but the leader of my own
destruction?'' {[}Augustine{]}. Or, as Sandra Bernhardt titles her film:
\textsc{Without You I'm Nothing}. Her self composed thru an appeal to
the other. Turning toward `you.' toward the invitation to explain, I
becomes, thru narration, \emph{in relation to}. Relation alone saves one
from destruction while paradoxically establishing the bond between `you'
and `I' which erotically threatens I's preservation as a discrete and
coherent self. Narrative's prosthetic origin: this other, who both is
and is not I. \emph{You may} be I's determining ``last instance''---an
absolute referent modulated thru innumerable social circumstances that
determines `I' in the final analysis. You, an Archimedean point doubled
and reflecting two horizons, in and out, creating this field of apparent
depth in the space between. You is equidistant from and coincident with
I's every utterance. In relation to you, I's meanings are measured. Thus
is I realistic without being real. Risking nothing short of love, the
voice aims at ``a particular absolute\ldots{} it aims beyond particular
objects to that `something more' that exceeds them.'' (\footnote{Copjec,
  Joan, \textsc{Read My Desire}, p.~148.}) That part of you that always
remains.

\hypertarget{ii.-politically}{%
\section*{II. Politically}\label{ii.-politically}}
\addcontentsline{toc}{section}{II. Politically}

\textsc{Narrative as counter-statement}. What is the difference between
being called to state (called to call a thing by its proper name) and
being called to \emph{the} state (called to recognize one's place within
the established order of things)? I ask this with the belief that
narrative has everything to offer the effort that would resist these
imperatives, enabling responses to always other calls. Called to state
and called to \emph{the} state: names, like bodies, are assimilated to
established narrative orders. But narrative needn't accede to the status
of statement, or doxa: rather, it is narrative as counter-statement, as
paradox, that enables one to refuse stately calls.

\textsc{Edward II, a digression.} The story of the English king Edward
the Second. In the sixteenth century, Holinshed, the popular chronicler
or historical narrator to whom both Marlowe and Shakespeare turned for
material, tells us that King Edward ``began to hold the nobles in no
regard, to set nothing by their instructions, and to take small heed
unto the good government of the commonwealth, so that within a while, he
gave himself to wantonness, passing time in voluptuous pleasure, and
riotess excess: and to help them forward that kind of life, the foresaid
Peers who (as it may be thought, he had sworne to make the king \emph{to
forget himself, and the state to which he was called} {[}my italics{]})
furnished the court with jesters, ruffians, flattering parasites,
musicians and other vile and naughtie ribalds, that the king might spend
both days and nights in jesting, playing, blanketing and in such other
filthy and dishonourable excercises.'' {[}Chronicles of England,
Scotland and Ireland, 1587{]}. ``The foresaid Peers'' is none other than
Gaveston, the king's minion to whom Christopher Marlowe has the Edward
of his drama implore, ``Knowest thou not who I am? Thy friend, thy self,
another Gaveston.'' To forget oneself and to forget the state to which
one is called are virtually synonymous offenses that render Edward a
`bad subject' rather than a good sovereign. They are capital crimes
serious enough to land Edward in the tower, where he is sodomized to
death with a red hot iron spit. As Holinshed reports, the method was
chosen ``so as no appearance of any wound or hurt outwardlie might be at
once perceived.'' Death appears woundlessly so that Edward can appear in
state. Were this wound to surface it would assume the figural status of
a counter-statement. Here, between life and death, continuity and
discontinuity; here, around this sign of struggle, a different narrative
would take shape paradoxically, this intervention in the orders of
state. Thus would the wound's articulation make it all ``connectedly
different.''

\hypertarget{narrative-difference}{%
\section*{Narrative difference}\label{narrative-difference}}
\addcontentsline{toc}{section}{Narrative difference}

\begin{quote}
When you consider the very long history of how everyone ever acts or has
felt, it is interesting that nothing inside them in all of them makes it
connectedly different.

~~~~~~~~~~ ~--Gertrude Stein, ``Composition as Explanation''
\end{quote}

From Augustine to Stein, narration has provided the way for
re-articulating the world in the space between, a space inassimilable to
already established orders, a space where discontinuity and continuity
together define a fault ``outwardly.'' But what about politics? What
about poetry? And what about the division of intellectual labor that
traditionally separates the two? Marx's \textsc{Eighteenth Brumaire} is
often cited in this regard: ``The social revolution of the nineteenth
century cannot draw its poetry from the past, but only from the
future.'' But what if the inquiry into the relation between narrative
and ``social revolution'' were to focus on neither past nor future but
rather on this space between? Right now, I am reading Stein's
\textsc{Narration} beside Louis Althusser's essay on ``Ideology and
Ideological State Apparatuses''. ``The tenacious obviousness of the
point of view of production alone, or even that of mere productive
practice are so integrated into our everyday `consciousness' that it is
extremely hard, not to say almost impossible, to raise oneself to the
point of view of reproduction.'' {[}Althusser{]}. In Stein's text, the
question appears as follows: ``Can anyone separate themselves from the
land so they can see it?'' And, as if in response to both questions, she
writes ``I am so certain so more than certain that it ought to be done.
I know so well all the causes why it cannot be done and yet if it cannot
be done cannot it be done it would be so very much more interesting than
anything if it could be done even if it cannot be done.'' What is to be
done is precisely what cannot be done. Nevertheless, the imperative is
sustained: to recognize the unrecognizable, to do the undoable. Here,
the narrative imperative, like a Luddite wrench in the machinery of
state, would prevent the apparatus from reproducing itself. ``To make a
thing a thing that they recognize while they are writing make it
something that had no existing before that writing gave it that
recognition, they tried to do this by changing something'' {[}Stein{]}.
This must be done. The point of view of reproduction is the point of
view of an antinomy defined by the apparent impossibility of doing
mutually exclusive things, like continuing and discontinuing, localizing
and sublimating. ``Revolution'' is a narrative not at all unlike the
narrative of conversion and it too is defined by paradox according to
which continuity with past experience must be maintained in order to
compose a discontinuous way \emph{here}. Stein might be discoursing on
ideology and Althusser, at great pains to demonstrate how Marx succeeded
precisely ``by changing something.'' might be writing about narration.
And while it is dangerous to reduce politics to literary theory---and
this is not at all my intent---it is critical to ask whether any such
success at making it ``connectedly different'' is conceivable without
narrative? To make it all otherwise: a deceptively simple imperative. To
find one's way maplessly---aporically and euporically---along the fault.
To essay the unthinkable, to attempt to do it all by other means: that
is to practice narrative ``as a socially symbolic act.'' Fredric Jameson
articulates this as a dialectical imperative: ``To reckon one's own
position as an observer into the critical thinking in process\ldots{}one
{[}then{]} no longer has to posit an end to history in order for
historical thought to take place.'' {[}\textsc{Marxism and Form}{]}.

\textsc{Ideology \emph{is} a narrative practice}. Narrative situates us
historically in the world. Once upon a time opens onto this critical
scene of telling and distributes an historical \emph{here} equally thru
the space of narration. The scene is primal in the sense that it is
domestic; it frames a world within the world while promising a way to
encounter the frame. Narrative is also the vehicle of ideology: it
represents (in Althusser's formulation) one's imagined relations to real
conditions of existence; it is charged with the demands of history, the
production and reproduction of a world. And while narrative can also
intervene in the space between, narrators cannot separate themselves
from the land in order to see it. In other words, there is no outside in
relation to ideology, just as there is no outside of narrative. To deny
narrative, to resist its claim upon us, is not so much resisting
ideology as it is ideologically blind to the fact that without narrative
we are bereft of the means of counter-strategy in the face of dominant
and oppressive ideological orders. This is telling, this is our common
practice whereby the limits and truths of the world are negotiated. In
the making of history, every composition is the site of as many
decompositions. Rather than the solution to crisis, narrative is the
performance of it (this fault to which I return and from which I emerge
again and again); not a repetition of prefigured limitation, but a
renegotiation, an approach, a reaching out that defines the world. Here,
`I' appears as a provisional place-holder at the site where a subject is
called and assumed. The `call' is a critical trope in Althusser's theory
of ideology where `interpellation' is said to hail the subject into
social being. Accordingly, the ahistorical individual is called to
assume the historical position of a socially intelligible `I'. But I
like also to think that, appearing as a third term between the
determined and the indeterminate, `I' is equally the effect of
misrecognitions, incommensurable both with the name by which it is
called and the name I call myself. Mistaken audience, mistaken identity.
`I' holds the promise of some non-identical content: `you' my audience
and `he' my self. Narrative denaturalizes whatever `I' and naturalizes
only the need to narrate. The space of narrative, then, is one of
intervention, conversion, transfiguration: where one renders oneself
other than oneself in a paradoxical quest for a more truthful `I'.
Selves are called and recalled in turn. Augustine's ethic of conversion
and Edward II's ethic of friendship both articulate ways of potential
resistance to a prior call by way of a new response. Ethics are
ideological, but they allow one to narrate one's imagined relations to
real conditions differently, that is, to imagine it differently.

\hypertarget{critical-faith}{%
\section*{Critical faith}\label{critical-faith}}
\addcontentsline{toc}{section}{Critical faith}

\begin{quote}
Only, how are they to call upon the lord until they have learned to
believe in him?*

~~~~~~~~~~ ~--Saint Augustine
\end{quote}

But how do we come to believe in the I that is called? Stein writes of
``\ldots{}names being not existing because anybody can know what any
body else is talking about without any name being mentioning, without
any belief in the any name being existing.'' Without risking belief,
communication lies in state. And belief, situated on the other side of
doubt and uncertainty, must be born in the space between, the space of
the call itself. This recalls Althusser recalling Pascal writing (more
or less) that one gets down on one's knees and belief follows. But with
a critical difference: belief needn't refer to the uncritical
reproduction of a dominant call and response; it might rather refer to
one \emph{abandoning} oneself to a space between incommensurable calls.
Perhaps this would be the space of neither belief nor doubt but of a
more critical faith whose subject, whose `I', is attentive to the
paradox: for what (or who) is the subject of \emph{abandon} if one's
self is the object? While belief is necessarily ideological, ideology
isn't unequivocally ``bad'': in fact, it is necessary. Like narrative,
ideology enables. Here again, the crisis, the plot: narrative's origin,
this doubling along the fault which I essays. Augustine abandons himself
to narrative and in doing so the stakes of communication are raised
without offering the security of a predictable outcome. If coherent
selfhood is only an effect of narration, and if narration is always a
potentially destabilizing force, then security is always already
preempted when one risks a beginning. This is both the promise of
narrative and its betrayal. And this is what makes Augustine's
\textsc{Confessions} an exquisite act of love, as political as it is
poetic. ``If poetry is the calling upon a name until that name comes to
be anything if one goes on calling on that name more and more calling
upon that name as poetry does then poetry does make that calling upon a
name a narrative it is a narrative of calling upon that name.''
{[}Stein{]} The poetic tradition, however, itself ideologically
informed, has connived with the market of names---authorizing it in
ways---a devaluing trade in which names are used and exchanged
calculably as guaranteed referents, their faults concealed. \emph{So as
no appearance of any wound or hurt might outwardly be at once
perceived.} Thus has the very need to call been obviated and the space
between, from which another way may be reached, occluded. I believe,
however, that in relation to crisis, it is the response-ability of
narrative to restore our awareness of the need, and intensify our desire
for another way.

And only by way of narrative might we unriddle ``this knot of imaginary
servitude that love must always undo again.'' (\footnote{Lacan, "The
  Mirror Stage.\emph{"}})

\hypertarget{bibliography}{%
\section*{Bibliography}\label{bibliography}}
\addcontentsline{toc}{section}{Bibliography}

Althusser, Louis, \textsc{Lenin and Philosophy and Other Essays}.
London: New left Books, 1971.

Arendt, Hannah, \textsc{The Human Condition}. Chicago: University of
Chicago Press, 1958.

Augustine, \textsc{Confessions}. translated by R.S. Pine-Coffin, London:
Penguin Books, 1961.

Augustine, \textsc{The Confessions of Saint Augustine}. translated by
John K. Ryan. New York: Doubleday, 1960.

Copjec, Joan, \textsc{Read My Desire}. Boston: Massachusetts Institute
of Technology, 1995.

James, Fredric, \textsc{Marxism and Form}. Princeton, New Jersey:
Princeton University Press, 1971.

Stein, Gertrude, \textsc{Narration}. Chicago: The University of Chicago
Press, 1935.

White, Hayden, \textsc{The Content of the Form}. Baltimore: Johns
Hopkins University Press, 1987.

Williams, Raymond, \textsc{Marxism and Literature}, Oxford University
Press, 1977.

\hypertarget{laird-hunt-the-avenue}{%
\chapter{Laird Hunt: The Avenue}\label{laird-hunt-the-avenue}}

\begin{quote}
\emph{In my skull is an avenue I stroke}\footnote{Clark Coolidge,
  \textsc{The Crystal Text} (I think). It seems likely I might be alone
  in seeing in these eight words, words ripped right out of their
  context, a possible model for narrative.}

--- Clark Coolidge
\end{quote}

~

\lettrine[lines=3, findent=0em, nindent=0.1em, lhang=0]{T}{here is an avenue}
in my skull too---albeit one that is poorly maintained, cavernously
pitted, strewn with rubble, whole segments blasted away; one that is
curved, possibly circular, that, like the backgrounds in cartoons,
maddeningly repeats itself; one that is ill-marked, with many a false
turn-off and many a false vista; one that is skewed out of proportion,
that is frequently unsafe, almost always unsavory; one that is troubled
by converging lanes, of which there are hundreds, that even resembles a
parking lot in places and is probably haunted---\emph{hell of a
place}\footnote{For some years now Brenda Coultas, in her narratives,
  has been conjuring worlds to which multiple qualities of the preceding
  description might apply. Whether in her book \textsc{Early Films}
  (Rodent Press 1996) or in her recent chapbook \textsc{A Summer
  Newsreel} (Second Story Books, 1999) Coultas is engaged in
  constructing usefully disjunctive, lyric-enhanced investigations into
  landscapes that have come apart. For me, part of the great thing about
  Coultas' writing is that she brings an indomitable hunger for beauty
  to bear on her cracked and spilling subjects.}. And yet, {[}I,
narrator{]}\footnote{E.g., Pamela Lu's terrific \textsc{Pamela: A Novel}
  (Atelos, 1998). A bildungsroman of sorts, an intricate chronicle of
  the coming awake of a highly-syntaxed consciousness, in which the
  greater part of the narrative seems to occur in between commas.},
stroke it, speak through its mess, speak of its mess, multiply it. It's
an interesting dilemma---aspects of which are taken up in Adorno's
seminal essay on the place of the narrator in the contemporary
novel---the narrator, with no story to narrate, \emph{narrates
anyway}\footnote{``The expression that there is nothing to express,
  nothing with which to express, nothing from which to express, no power
  to express, no desire to express, together with the obligation to
  express.'' Samuel Beckett (in 1949 conversation with Georges Duthuit).},
a story that has been blown to bits. Of course not everyone sees it that
way. We live in {[}an age of errata, of misinformation, of
disinformation, of hoax{]}\footnote{The tendered premise in a recent
  episode of the X-Files: since World War Two Americans have had an
  unquenchable appetite for `bogus revelation'; having determined this,
  the United States Government has built, presumably as a cover for its
  more radical weapons of mass destruction programs, an elaborate web of
  alien abduction and sighting hoaxes. This hoax, we are led to believe,
  may or may not be one. At the end, intrigued to a ghastly degree, we
  are left both believing (the elements of documentary inherent in the
  X-Files aiding this process) that there are aliens around and that the
  Government has cooked up (and deployed) weapons even more diabolical
  than the ones it has copped to. This tv `fiction' seems to me only
  slightly less credible than what the evening news, reporting live from
  the White House! serves up, or than the gamut of what can be found
  gathered up daily in \textsc{The New York Times}. (The day after
  watching the X-Files episode I read allegations by Puerto Rican
  petitioners at the United Nations that the United States Government
  was testing uranium-capped bullets at its facility on Vieques Island;
  during the debate an Iraqi delegate stated that similar weapons had
  been used during the Gulf War.)}; perhaps it is little wonder that
there continues to be such a hunger for narratives that, as Adorno
describes them, largely by way of 19th century techniques \emph{mimic
the real}\footnote{Harry Mathews, who, in part of his work, has
  vertiginously built upon (see \textsc{The Conversions and Tlooth},
  recently reissued by Dalkey Archive), the wonderful, strange
  narratives of Raymond Roussel, had this to say in a recent interview
  in Rain Taxi: ``The great discovery I made through him {[}Roussel{]}
  was that you didn't have to write about what happened the day you fell
  in the swimming pool, but that you could invent your own swimming pool
  that had never existed before, full of quicksilver, lighted by
  imaginary lamps.'' Which, abstracted, might lead to a formula
  something like, The writer is in the world; the writing is in the
  writer. Or, looking at it from a reader-oriented perspective, we might
  arrive at: The writing is in the reader; the reader is in the world;
  Either way we are dealing with a considerable degree of mediation. One
  I am interested in.}; that say to us, with disarming earnestness,
this, my friends, is \emph{how it was}\footnote{``The more strictly the
  novel adheres to realism in external things, to the gesture that says
  `this is how it was,' the more every word becomes a mere `as if,' and
  the greater becomes the contradiction between this claim and the fact
  that it was not so.'' (Adorno, ``The Position of the Narrator in the
  Contemporary Novel'')}. In the domain of the fictive narrative, I tend
to have little patience for such works. I am much more interested in
(and seem only capable of constructing) narratives that are to some
degree aware of \emph{the provisional nature of their own
authority}\footnote{Among contemporary works, the narratives of W.G.
  Sebald (\textsc{The Emigrants, The Rings of Saturn}) are exemplary in
  this regard, pointing up, as they do, against the backdrop of this
  bloody century (and with an apercu of others), the enormous difficulty
  of fixing experience, both personal and collective. Keith Waldrop's
  harrowing memoir \textsc{Light While There Is Light} (which calls
  itself fiction) is another example. As is David Markson's intriguing
  \textsc{Wittgenstein's Ladder}.}, in which \emph{the fictive quality
of recollection}\footnote{Ryonusuke Akutagawa brilliantly treats aspects
  of this in a pair of stories which were later combined and made into
  the much more famous film \textsc{Rashomon} by Akira Kurosawa. Aaron
  Shurin writes, in passing, in his 1990 essay ``Narrativity,'' of the
  `irreversible solidity of the past tense.' I've lately begun to
  marvel, because of this generally ascribed quality, that History
  continues to be written in it. It was Herodotus, Father of History,
  who kicked things off---neatly blending directly observed
  incident/object with reported incident/object with rumored
  incident/object with imagined incident/object. At least a good part of
  what gets built into the past tense, then, is hardly irreversible,
  hardly solid. Pretty slippery, in fact.} is acknowledged, in which
\emph{forgetting}\footnote{Walter Benjamin writes beautifully on this in
  the opening pages of his essay, ``The Image of Proust''
  (\textsc{Illuminations}): ``For the important thing for the
  remembering author is not what he experienced, but the weaving of his
  memory, the Penelope work of recollection. Or should one call it,
  rather, a Penelope work of forgetting? Is not the involuntary
  recollection, Proust's memoire involuntaire, much closer to forgetting
  than what is usually called memory? And is not this work of
  spontaneous recollection, in which remembering is the woof and
  forgetting the warf, a counterpart to Penelope's work rather than its
  likeness? For here the day unravels what the night has woven. When we
  awake each morning, we hold in our hands, usually weakly and loosely,
  but a few fringes of the tapestry of lived life, as loomed for us by
  forgetting. However, with our purposeful activity and, even more, our
  purposive remembering each day unravels the web and the ornaments of
  forgetting. This is why Proust finally turned his days into nights,
  devoting all his hours to undisturbed work in his darkened room with
  artificial illumination, so that none of those intricate arabesques
  might escape him.''} is considered the key constituent of memory, and
in which, finally, getting it right shares center stage with getting it
wrong.

~

~

\hypertarget{trevor-joyce-why-i-write-narrative}{%
\chapter{Trevor Joyce: Why I Write
Narrative}\label{trevor-joyce-why-i-write-narrative}}

\lettrine[lines=3, findent=0em, nindent=0.1em, lhang=0]{B}{ut, first}:
why not? A chain of telling that requires discrete and distinct figures,
already in some sense familiar, acting across a background, largely
unexamined and abandoned merely to suggestion, towards the achievement
of, or the failure to achieve, certain goals, whose importance holds the
earlier periods in suspense until in some final resolution all that is
significant crystallizes in a perfection of plot and motivation, and all
the rest, wanting any real brush with language, retreats once more to
ground: that's narrative as I knew it thirty years ago, and it seemed
then, as it does now, inadequate to the world of my experience.

To avoid the fate of perpetrating such stuff, I instead wrote a poetry
which found itself increasingly characterized by argument, deploration,
pleading, threats, until I realized that I didn't like how that worked
either, and stopped writing altogether for about twenty years.

But I went on thinking and planning and testing and learning, though
largely just in imagination. I took a shine to the detective genre in
which, in theory, any object or event can be a clue and, as such, be
exalted into meaning as the everyday world sifts through the riddle of
observation, inference and deduction. But, even there, the world existed
only towards an end in which all might be revealed, and the head of
Holmes had a strictly limited inventory capacity.

Eventually, it was only in the parallel forensics of Gombrowicz's
``Cosmos'' that I found some satisfaction with detection, since nothing
is revealed there, and the patterns and clues are rubbish and trivial
chatter remains fraught, even on rereading.

I had by this time resumed writing poetry, though with conventional
narrative either expunged or twisted, and temporal change borne instead
by repetition, either incremental or with calculated variation. My
models for this were in the refrain structures of folk-song, often
mediated through the likes of Yeats or Lorca, and in the interplay of
stasis and movement in Chinese parallel verse. I was forced to
recognize, however, that the `lyric' mode which I practiced was quite as
prone to exclude the incoherent world as was the mannered narrative I so
distrusted. I had also encountered Cage for a second time, and with more
understanding of how the play of ambient noise across the receptivity of
his spaces might circumvent those exclusions and admit what might
otherwise not be acknowledged.

I wrote a longish `bicameral' piece called \textsc{Syzygy}, consisting
of two halves, ``The Drift'' and ``The Net'' respectively. ``The Drift''
consists of twelve compact, elliptical, but distinctly `lyric' poems.
``The Net'' is a single poem of 72 long lines, comprising 24 three-line
stanzas. Significantly these two sections are made up of exactly the
same phrases reordered, rigorously and exhaustively mapped through a
one-to-one matrix, the exact structuring of which is not directly
relevant here. There is a brief set of notes added, which ground some of
the detailed references of ``The Drift'' within the empirical world.

For me, much of the significance of the poem is in the way in which
sequences of phrases in ``The Net,'' arrived at through the blind
deploymen of predetermined procedures, carry a force both of lyricism
and of narrative, intense though severely fragmented in both cases,
which revealed a meaning different from and additional to anything I had
deliberately written into the work. Having lived with it for over two
years now, I have come to understand that ``with the first dream of fire
they hunt the cold'' and how ``devastation fell attending headbone the
high'' while ``outside the foundries the clumsy the deadlocked
disintegrates'' though ``not a tremor manifests the rare the quickening
across these settlements'' . So, I had confirmed for myself that a
densely overdetermined language, functioning in its most intensely
personal mode of the lyric, could survive radical disruption and return
from that alienation a yield which the reader might gather. It gave the
world an in.

I had meanwhile read all the Fu Manchu novels of Sax Rohmer, intrigued
the Chinese Doctor's ceaseless attempts to appropriate the world by
abducting experts in all fields, whose task was to analyse and replicate
in enhanced form all aspects of the workaday world. It matched my sense
as a child that what I saw in mirrors was a world diligently assembled
by unseen agents to match that which I inhabited, and I watched
carefully for small discrepancies to justify that sense, but never found
them, and the specialists of Fu Manchu's underworld empire attended only
to the great realms of science, technology, politics, always were
disbanded by resurgent law before their analysis could address the
classified advertisments in the newspapers, the torn betting slips
outside a bookie's shop, the inconsequential gabbling of drunks in a
pub. The workings of what I came to know, through Marx, Adorno and
Benjamin, as phantasmagoria fascinated and appalled me. And how is
responsibility to be assigned across mock-worlds if not through the
causal chains which are narrative's stock-in-trade?

Therefore:

\begin{verse}
while detailed depositions state\\
how further on\\
within the wood\\
. . .\\
\end{verse}

\begin{verse}
the bright axe\\
blossom suddenly\\
\end{verse}

\begin{verse}
the long bones lever\\
up from it like anthers\\
and beyond the startling\\
calyx of teeth\\
\end{verse}

\begin{verse}
an avid buzzing perishable\\
fruit set thicken\\
and disintegrate\\
to load with sweet\\
\end{verse}

\begin{verse}
secure deposits\\
of afflicting gold\\
their remote cells . . .\\
\end{verse}

Here is not just one narrative, but two. Firstly, the sequence of
blossoming, the detail of anther and calyx expanding, the fruiting
adumbrated in the gathering of bees about the flower, and their
dispersal to the hive where they load the cells with honey. This
apparently natural and value-free sequence is overlaid on another causal
chain, which starts with the felling of a tree by a logger with an axe
(commercial or strategic deforestation has a long and significant
history in Ireland as in much of the developing world today). From this
action, the causal sequence is run backwards, seeking earlier sources
where responsibility may be assigned: the long bones of the arm, the
gasping jaw of the labourer, give way to the investors depositing their
profits in banks. This single instance is simple, but the structure of
reversed causality running back from an act of violence against the
person or ecological ruin, masked meanwhile by a natural sequence of a
bird fledging or a mineral cystallizing, is repeated twice more in the
poem, broken by cases for and against the possibility of asylum amid
such wilderness. Here, I began again to write narrative because the
forensic process it allows seemed to me necessary to any possibility of
living ethically, of recognizing and fulfilling due responsibility.

Another narrative genre which interested me is that of scientific
experiment, where the researcher actively intervenes in the course of
nature, attempting to limit the causal influences at work, so that one
element may be manipulated, and the change in another, dependent
element, observed and measured. It is intended that the record of
correct prediction and accurate calculation of effect may then grant
understanding of what was previously obscure.

``A body thrown vertically down from the top of a tower moves through a
distance of 88 feet during the third second of its flight. Calculate,
then, the speed of projection, and determine the speed at which the
sleeve begins to move upwards.''

And yet, it is in a field complex with uncertainty, that we attempt to
understand, to categorize, to measure, and the experimenter must attempt
to exclude all forces not considered relevant to the investigation, and
accurately account for all that ensues.

``When he attempted to speak to her, the patient jumped off the bridge
falling some 30 feet into about 20 feet of water. There is always a
chance therefore that the critical act or change may take place when the
observer's eyes are withdrawn.''

Yet we must dispassionately observe, measure and record.

``Mild plethora of the face ensued, it being divided into three parts,
namely: the forehead, fair complected, one; the nose, another, sand
present in abundance there admixed with small crustacean shells; and
from the nose to chin, exhibiting extensive tooth loss though with roots
intact, another. Notice the blood tinged fluid coming from mouth. Red is
warm and radiates across the ground.''

In such a manner, sometimes we come to face what can scarcely be
countenanced.

``It may happen that we are not aware of all the conditions under which
our researches are made. Some substance may be present or some power may
be in action, which escapes the most vigilant examination. Not being
aware of its existence, we are unable to take proper measures to exclude
it, and thus determine the share which it has in the results of our
experiments.''

Though the outcome be uncertain, may we still presume to have advanced
knowledge in certain quarters? Lacking the machinery of suspense, to
what end may such a narrative aspire?

``A man has several bones in all, and beauty is lost when meaning and
form are split asunder. The handsome man must be swarthy, and the woman
fair, etc., the genitalia, both internal and external, without injury.
Provisional diagnosis: probable drowning. And had we exhausted all the
known phenomena of a mechanical problem, how can we tell that hidden
phenomena, as yet undetected, do not intervene in the commonest actions?
I will not tell you about the irrational animals, because you will never
discover any system of proportion in them.''

The plausibility of narrative is increasingly an issue for me, and not
without reason, perhaps, given the dominance of contending
master-narratives in the interpretation of the Irish past and,
consequently, in my present world. From such master-narratives it seems
worth trying to retrieve as much as may still have value. I have tried
this in a recent long work called ``Trem Neul''---part prose, part
`verse'---which I see as, in part, an attempt to recoup part of the
history of my world from what Beckett terms `the uniform memory of
intelligence'. ``Genealogies. The elementary tables. Dictionaries,
assembled in blind frosts. Grammar and chronology. Libraries. Index: the
Encyclopaedia, damascened with ice. So is the perfect body of knowledge
dislimned.''

How may one conduct a narrative of change, of loss and recovery, of
breakage and continuity, without presuming the existence of distinct
agents, freeing them incredibly from their ground, and committing a
plotwork of events, utterly plausible because familiar and foreknown?
Scholarship has shown us how the integrity and closure of the human
agent was arrived at in early Greek poetry, must we take that achieved
unity as more than provisional? Can we not tell without it?

``We build ourselves through the world and each through other, and this
proceeds to death as the world alters with experience.'' The bodies in
this plot are not distinct, either from one another or from their
ground, they emerge to make themselves, enjoy a transitory closure, and
then resolve again into a ground which offers further figures. ``When
the biology of your body breaks down, the skin has to be cut so as to
give access to the inside. Later it has to be sewn back like memory,
when it may house all knowledge. Memory is our comfort and our attire.
Fashioned with our hands it is the accomplishment of our dreams and
lapses; always a meaningful pattern though never an abiding one; a
shifting harmony of sub-patterns. Pretend I'm lost and try to find me.''

Because I wish to work comprehensively with the world which I inhabit,
however, and because that world is current with named identities, I have
admitted one such, and one specific narrating voice to speak of him. I
have permitted myself also two specific occasions in time, one exactly
situated starting at about 11 a.m. on February 15th. 1838), one not.
Between the two is a connection, and my essay is to account for that, to
recount it. ``Do not think it coincidental that memory should begin to
fail just as taxonomies become a prominent tool for thought.''

As, I believe Braudel and the Annales school refocused history on the
wholeness of ordinary lives, their habits, orientations and crises, by
attending patiently to things left, I see my present course, of
rendering the experience of connectedness, as being sustainable only
through exercising a similar patience with language. No longer expecting
to find there an exact mirror-image of the world I know, but rather to
have it deliver me one I don't.

\begin{verse}
We will read\\
every day\\
in the afternoon\\
When shall we learn\\
to write?\\
We shall soon\\
learn\\
I once went\\
to Europe\\
but I do not now\\
remember\\
what I saw\\
there\\
\end{verse}

It is in preferring to concentrate on the unpredictable ground rather
than to people it with puppets of my own making that I have elected in
these more recent pieces to work increasingly through collage. Each
fragment of language I adopt is already tale-bearing, a vector: carrier
from a prior host, director of action also across the space of my world.
And the point of such a narrative? Interim figures on an interim ground;
preservation of the complex weave of actions, not denouement; attachment
understood embraced abandoned; wanting executive or summary.

\hypertarget{notes}{%
\section*{Notes}\label{notes}}
\addcontentsline{toc}{section}{Notes}

All quotations are from \textsc{Syzygy} (Wild Honey Press, Dublin,
1998), available on the web at the Sound Eye site
http://indigo.ie/\textasciitilde{}tjac/Poets/Trevor\_Joyce/Syzygy/syzygy.htm

\hypertarget{kevin-killian-poison}{%
\chapter{Kevin Killian: Poison}\label{kevin-killian-poison}}

\lettrine[lines=3, findent=0em, nindent=0.1em, lhang=0]{I}{’m standing}
on a flat plain, and then, or so it seems, a little hole appears in the
sand ahead of me (like that movie \textsc{Tremors})? The hole grows
larger in diameter, this is my sanity, and all the little pieces of my
sanity are breaking up and slipping down into the hole. That's what it
feels like. I'm trying to write this piece, ``Poison,'' about the ways
in which the writer's personality dissolves as it weaves in and out of
the sentences he or she so painfully struggles to produce. While writing
it I notice a host of familiar symptoms. Nobody calls me on the phone. I
feel so isolated. I can't hear very well and wonder if I'm going deaf
like Beethoven, like Brian Wilson. When people do speak it's with loud,
ultra-charged voices, as though they're annoyed with me. I feel like I'm
losing my mind and with my mind, the meaning of life I once held onto. I
used to think that people are basically good at heart. I used to think I
would be happy someday, but now I feel differently, that there's no
chance for me, since that hole before me is opening up and soon
everything I ever clung to will be sucked down into it and I'll be
homeless and curled up on the gutter outside my former apartment on
Minna Street, a crack-ridden block South of Market in San Francisco. I
just bought a car and there's not really any place to park it.

How did I write a whole book? I'm trying to remember. In particular how
I wrote \textsc{Bedrooms have Windows}, which I loved so much too, I
wrote it in part as a shipwreck victim sends out a message in a bottle;
in particular to my dear friend Terry Black, with whom I'd lost touch a
few years before. The book has many appeals to him to get in touch with
me. (Before the Internet, through which, apparently, everyone is
available to or traceable by everyone else, this seemed my only
recourse.) I longed to see it in print, feeling that he would pick it up
and call me. But after it was published a mutual friend sent me Terry's
obituary, he had died in Richmond, Virginia, of AIDS, the same month the
memoir came out. This was not the answer I had hoped for. Part of me
felt that \textsc{Bedrooms have Windows} killed Terry Black---detailing
as it did our sex life and its creation, the way we had made up sex to
answer certain suburban needs for the authentic, the ``real,'' the
colorful. Naming names, his. Implicating others, him. What portion of
one's personality is a fiction? It wasn't going to do any good to
realize this was a sentimental fantasy, part of the mind's response to
the inexplicable horror of AIDS, part of my own need to find myself on
centerstage always. I went back to ``Poison'' as I delivered it
originally, as a talk in Bob Gluck's series, \emph{In Context}, a series
of talks delivered week after week at Intersection, once an important
writing venue in San Francisco. This was a happier time for me---April,
1987. When New Narrative writing seemed wide open, a place where
something entirely new under the sun could be created. And that we were
doing it, doing so. But ``Poison'' I found out was imperfect, it didn't
help me. ``Last week Dodie Bellamy's talk stressed the paradox of
writing as a two-way street-the importation of the world into the self,
and the generous export of the self back into the world. Bob's talk may
have been allied to my own idea of writing as an expression of the death
drive from \textsc{Beyond the Pleasure Principle}---what he calls
`Freudian pleasure based on an instinct to return to the inanimate.'\,''

Every writing act is an act of dying, or killing, or mortification.
Every time I write it's to expose to the air of the page a false part of
my personality. ---I guess this goes directly against Bob's theory of
writing, and links mine closer to Dodie's, though not in any way she'd
like or approve of. My talk, I thought, would be a patchwork of
quotations, writing about me written by others; some that I wrote about
other people, lots I wrote about myself, and also work by others that
didn't have me in it at all---to give a wide scale against which I could
test my propositions. And the first was from Alan Davies' book
\textsc{Name}, which I found useful in its treatment of the tie between
language and self; is language a function of the word or of the self? Or
vice versa?

\begin{verse}
The cryptic tongue.\\
We are getting ourselves\\
in the mood to have been\\
done with having been done\\
with this again.\\
It’s all very irreversible\\
which is what\\
makes its guts open up.\\
I wonder how long it will\\
be until this writes\\
itself, in\\
my direction. And we\\
haven’t proved\\
that it isn’t true yet.\\
When I think of you\\
the sentences come,\\
but I don’t.\\
\end{verse}

Then I read a passage from Dodie's book \textsc{The Letters of Mina
Harker}, in which Mina, a fictional character, reflects on her love
life. In this passage I felt myself inextricably named and described, my
human body a vessel for a flood of narrative concerns.

Flaccid, KK's penis is endearing, so velvety and shy---but the trouble
with babies (as my mother always said) is that they grow up---your bed
inflates to the breaking point with thirty-three-year-old male desire
panting and prodding \emph{the thing inside burst through her belly,
horror props, sausage links and ketchup} around my neck KK fastens a
locket filled with a snip of his hair \emph{to protect me from evil} I
cross the street with my eyes closed, cars screech then cease to exist,
the atom remains unsplit forever, cells multiply at a reasonable rate
every death is from a natural cause LOVE LOVE LOVE LOVE LOVE LOVE LOVE
LOVE LOVE remember when his nails were half an inch long, thick, hard,
yellowed---he clipped them off for me parting my capillary pink flesh
without a scratch \emph{all it took was one ``ouch''} claws retract,
breathing softens. He extends his palm from the bathtub and says, ``Sit
on it,'' human form follows function, in Cocteau's \textsc{Beauty and
the Beast} gloved arms poke out of the walls holding candles, their
flames trembling as Beauty recedes down the endless corridor \emph{I
wish I could walk through mirrors} our entwined bodies tighten into a
circle, a champagne bubble about to be swallowed by Marilyn Monroe
\emph{pushing the metaphor to the breaking point, in a word: orgasmic}
when we fuck we are two great hands shaking \emph{his cock a thumb} in
an explosion of light the bearded creator in Blake's watercolor points
from the heavens---mortal heads bow or stare up in awe and terror the
way I do whenever I'm naked \emph{a woman's hair is never thick enough
to hide her thoughts} KK reaches for a condom, fumbles with its little
blue capsule PRESS FIRMLY ON DOT AND PULL APART \emph{town fathers pack
data in time capsules burying them underground, schoolchildren dig them
up, crack them open a hundred years later. Things.}

I used to ask Dodie, ``Shouldn't it be, `When we fuck we are two great
hands shaking his cock a \emph{giant} thumb?' Wouldn't that make the
passage clearer?'' She said, ``No.'' That's all. Just ``no.''

At first I was embarrassed by this passage and many others like it in
her writing, for if at any time my penis is flaccid I don't want to know
about it, nor the world to suspect it. At the same time, I felt
flattered, singled out by her language as I felt singled out by her
love. ``Beauty recedes down the endless corridor \emph{I wish I could
walk through mirrors}.'' The sense of the syntax issues a seductive
invitation into a mystery world---and I am that mystery. Then the cold
water hit me; with a start I came to and asked the difficult question,
is it I who is being described? \emph{The Letters of Mina Harker} seem
to describe the sex lives, the love lives of two actual people, but of
course they don't, they don't even especially want to: their veiled and
mediating nature hints at this:

---extends his palm from the bathtub and says, ``Sit on it,'' human form
follows function, in Cocteau's \textsc{Beauty and the Beast} gloved arms
poke out of the walls holding candles, their flames trembling as Beauty
recedes down the endless corridor

---and so forth, then the phrase ``pushing the metaphor to the breaking
point'' nails it. Do you know the funny feeling you get when a stranger
waves on the street, you wave back, and then you realize the stranger's
waving not at you but at the little weirdo behind you in the brown
fedora? One's identification with the words that seem to conjure one up
like ``I Dream of Jeannie'' out of black letters on a white page is
like---just like-the way your heart then sinks and a blush colors your
face; except it's even more mortifying---that stranger on the street's
waving all right, and you wave back, then you turn around and it isn't
even a person that stranger is really waving at, but an atmosphere
perhaps, a draft of air. Only convention associates the wave (or the
sentence) to a corporeal body like one's own.

My next example I wrote. I quoted from \textsc{Bedrooms have Windows} to
give some idea of what goes through a writer's head when he decides to
name real people in potentially scandalous situations-sexual in this
case, but scandalous only insofar as they deal with the body; also to
suggest the peculiar vanity of the writer-and then I'll give it a fairly
close reading.

One evening several years ago I was lying in bed, after some
unsatisfactory fumbles towards ``safe sex'' with a writer I once
admired, Dennis A. He turned his head-exactly as he'd turned mine, an
hour earlier---and get this, he said, ``Why didn't you think to take
home that Tom Boll too and we could have had a threesome.'' ``Why didn't
I think,'' I replied, an echo of disbelief. ``Dennis, I did think; I
thought and thought. Had I thought any more, I wouldn't be having this
safe sex. Madness wouldn't have been safe from me.'' I felt attracted to
him when he spelled his name in plastic magnetized letters on Aaron's
refrigerator. Then he spelled mine with the same colorful letters.
Language fused. It was like William Carlos Williams. ``There were plums
in that icebox,'' I said to him. ``Forgive me. Forgive me. I couldn't
help it; they were so ripe and so purple and so cold.''

I doubt if ``memoir,'' the word ``memoir,'' really disguises from the
reader, or listener, that something very close to a real event is being
described. And naturally many people will correctly identify ``Dennis
A.,'' the writer ``I used to admire,'' with the Australian historian and
sociologist Dennis Altman. I figured out that I call him ``Dennis A.''
to make the narrator's personality kind of coy and obnoxious---and it
works, doesn't it? Maybe I was miffed because the sex we had wasn't very
perfect---or maybe I thought, ``He's forgotten about me, I'll employ
this sentence to haunt him.'' Writing as an act of revenge, and naming
names a superior way of taking it. Like a virus, the poison of this
passage will break down, over time, whatever goodwill and nice feeling
that I---my body---will have by then created---like some ghastly race
between life and death, immune system and viral infection. Bob's novel
Jack the Modernist begins with a paradigm of intention, deadly
intention, ``You're not a lover till you blab about it,'' where the
ugliness of the syllable ``blab'' is meant to suggest a whole medicine
cabinet's worth of emetic antidotes to the possibly-too-pretty word
``lover.''

The same watchful watchdogs will naturally conclude that this pick-up
scene takes place not just in anyone's kitchen, in the kitchen not just
of any old Aaron, but in Aaron Shurin's kitchen. Maybe they already know
of my admiration for his writing and his influence. When I first met him
I had just read his book \textsc{The Graces}, and I was so struck by it
I could hardly connect its author with a living human being; it was,
maybe, a gift from another planet like something out of Chariots of the
Gods. And since then I've gotten to know Aaron better ``as a person,''
but still blitzed by this admiration which---I see now!---is another
form of objectification, turned outward instead of inward. Hence there
was little psychic difficulty in turning his name, like a totem, into an
amulet to adorn my prose. Hence the repulsive, shy-making casualness of
using that name, ``Aaron,'' and the social-climbing note its use strikes
here, as though I were his intimate, or perhaps his boyfriend, and what
is it in actuality but an incursion on him, aimed squarely at him at his
most domestic and private, the kitchen setting, the letters on the
refrigerator, perhaps the implied sneer (leftover from the 60s) that the
letters are made of plastic. A plastic language. Doesn't sound like a
recommendation, does it? It's true that telling stories, ``narrative,''
does involve a local, in the sense that this quote of mine does have a
certain atmosphere, a sophistication, but it isn't really mine, it's
borrowed or stolen, the way you or I might borrow someone's boyfriend or
wife, return it to them and destroy a relationship like breaking a milk
bottle. (This isn't to suggest I haven't thought of the pleasure it
brings to writer and reader both, but---maybe because I'm a
Catholic---it's a guilty pleasure; there's pleasure in guilt too, and
even if there isn't there's the dying fall you get when you string one
word after another, after another, onto another, like bugle beads. The
pleasure of accretion.)

Next I'll read something I wrote for the ``Jack Spicer'' issue of
\textsc{ACTS}. My intention here was to get at what I saw as the malice
characteristic of Spicer's poetics, and I couldn't think of another
analogy except to describe an occasion from my own life, an occasion
when I felt malice directed at me.

My aim was to develop a visceral writing, a writing that would as
closely as possible parallel the effects of the anonymous letter:
insult, horror, shock and embarrassment. But did my face turn red when I
received one myself:

Dear Kevin Killian: Your piece ``Tom-Tom'' in the latest issue of
\textsc{No Apologies} aroused my suspicion that you must be an asshole.
First of all what is a gay man doing writing about a psycho-killer of
women and getting off on it? I'm suspicious that it's easier for your
fantasy screw to get off when the victim is a (dehumanized) ``Miss
Thing'' than if your homicidal maniac lusted after boys like---just
like-you . . . So what's your story? If it isn't a good one, this
dyke'll write you off entirely.

Threats. Intimidation. A whole map of misreading of my adorable piece.
Again and again I read this letter, each time with increasing unease and
paranoia. I howled into the open air, ``What have I done to deserve such
venom? Who wrote this tripe?'' Despite the writer's declaration, I
wasn't fooled into believing him a woman. In my heart I knew this letter
is from one of my so-called friends-of the male species. My only
questions were, which one---and why? I'd been put into the abject
position, one I was to maintain for a long time. Writing produces
deracination, I decided, by own best example. So much for theory. My
content upset ``her,'' I thought to myself. ``Take away the
psycho-killer storyline and there'd be nothing even a child would object
to.'' And ``her'' form upset me. Isn't ``dictation'' Jack Spicer's word
for ``receiving the letter''? You don't know where it comes from; so you
react badly. Martians are writing you, altering the furniture in your
room. I wanted a drink so bad I wound up in North Beach! Bats flying
from wall to wall, the whole schmear, and voices writing all my poems
and signing them with Dennis Cooper's name!

Here the dominant note of my writing becomes real plain: I'm talking
about hysteria. I mean that seriously enough, a line that runs a beat
too fast, that's capable of all sorts of unexpected connections (which
is good) but because it's hysterical is equally capable of making false
connections or ignoring the valid (which is bad). That is to say, in my
heart of hearts I don't really believe Dennis Cooper is signing all my
poems. Why, I don't even know what I meant by that. That's paranoid,
isn't it? But I had to find out, what are my rights as a narrator and
character?

When I put myself in this false position, I run the risk all divas take:
we might prove too much for even our greatest fans. And so I've noticed
the characteristic note when others write about me is comic, often
picking on my frailties as a real person and usually, by the way, noting
either that I love the stars, drink a lot of Tab, or also, you know,
drink a lot, of, you know, alcohol. The next passage is from a story by
Francesca Rosa called ``Canidae'' (the Latin word for ``dogs''). In it I
thought I recognized myself as the writer, ``K.'':

As K. speaks, a worried-looking dog pokes its head through the blue
gauze curtains separating the reading area from whatever is behind it. A
black dog, except that its fur is so sparse the skin shows through, a
gray and black dog then with running eyes; the wattles of its throat a
livid pink. It brings the rest of its body through the curtain and walks
upstage of K. who is telling us a story about Long Island. The dog
explores, listlessly sniffs at the floor, the podium base and K.'s
ankles. It walks back and forth a few times, turns a slow circle while
digging its teeth into the root of its tail, and then stops to stare at
us, the listeners, again. Not at all shy, confident, as if resigned in
its desolation, like one of those Kafka characters that have survived
their own death. K. does not, or decides not to, notice his center stage
companion, and takes us from Long Island to New York and then back
again.

When I first heard this story I remember she was sitting on stage and I
was sitting in the audience, our narratological positions reversed, and
my ears got red, and a voice rose up in me with a strangled scream,
``You can dish it out, Kevin, but you sure can't take it!'' In Bob's
talk,``Truth's Mirror is No Mirror,'' he says, ``I wonder if we are at
the point of reversing Flaubert . . . by accepting an artificial self,
with its own scale, depth and continuity. Eastern religion responds to a
`made-up' world with compassion---but with a fatalism that is the flip
side of Flaubert's scorn. To the degree we `see through' Flaubert's
scorn, we suffer from and enjoy a self-contempt that is close to
bragging . . . I wonder if it's possible to be aware of the artifacted
nature of the local and not be contemptuous of it?---to understand it as
a construct and be moved by its depth?''

I was telling my sister Maureen that I was giving a talk on the fiction
of personality, and she nodded and said she knew what that's about.
Sometimes, she said, when she's walking down the street she hears in her
head her theme music so she feels she's starring in her own TV show.
Elizabeth Bowen wrote, ``Nothing gets on the page that you started with,
and nothing you started gets on the page. To write,'' she said, ``is to
rave a little.'' To this formation I've added my own strategy, to rave a
little before I begin to write, to exploit the ``fiction of
personality'' and to see, if not blinded by the brushfire, what happens
then on the page. When I was seventeen and living in my parents' house
on Long Island I threw lots of parties, so did my friends Terry Black
and Lance Mallamo, we did so to enjoy ourselves and to write about them
afterwards. This is from a story one of them wrote about a party I gave,
a story that begins with me warning my parents to stay away from my
party and my fun and my personality and my fictions:

``All I know,'' pouted Kevin, ``is that I care about this house as much
as anyone else and I think you're a hideous couple even to think I'd let
anything happen to it. Why I just can't explain the chills that go up my
spine every time I see it. It's like finding a lemonade stand in the
middle of a desert. When I think about the beauty and splendor of this
house tears come into my eyes. Its white walls, its grapefruit shingles,
its collapsible mailbox, the one-of-a-kind sodded lawn, and, to top it
all off like the cherry on a banana split, I think our Spot, our
pedigreed dog, out in the backyard rolling in mud and barking at all
passersby as if to say, `I'm black and I'm proud!'\,''

Kevin's fit of love and tears was interrupted by the doorbell. ``I just
want to say one more thing,'' he said in a hurt voice. ``If I see either
of you downstairs from now until tomorrow morning at ten o'clock I'll
set both of you on fire. Don't think I don't mean it either because I'm
pretty sure that I do.''

With that he stormed out of the room, slamming the door behind him.

Later in the story Terry himself is in a car driving to my party.

``Terry,'' said Mary Phipps, a girl in the back seat. While drinking a
swig from a bottle of cheap apple wine she continued, ``I'm glad you're
acting like yourself again instead of like Kevin Killian.''

``Thank you,'' the courteous Terry replied through gritted teeth. Mary
Phipps smiled at him and he could not help smiling back. They were true
friends and he knew he should listen to what she had to say whether he
liked what she said or not. So he did.

``Not that we don't like Kevin,'' joined in another good friend of
Terry's named Sam Rye. ``It's just that you're not him and we liked you
better as you.''

Terry resigned himself to the friendly rainstorm of advice, opinions and
admonitions that followed. He could not help wondering though why it was
that Kevin always got almost all the attention and he, Terry, so little.

But I've argued myself into a corner if I insist that New Narrative
makes room for the stupid, the overblown, and the nasty---these are
three different veins of my writing, each with its own jet flow. The
bizarre thing was that, after I delivered this talk on April 27, 1987,
the woman who had written me the anonymous poison-pen letter lifted her
hand from the audience and said, ``Oh, I wrote that---before I knew you
better.'' Do we get to know people that much better that we can change
our minds so quickly about them? I was stunned. I remain stunned. I
think of life as a big empty desert place---cooler than Death
Valley---but just about as big---``I know! Let's call it Life Valley!''

Where, as I pause for a second, clearing my head, trying not to write,
dragging myself from that particular abyss of memory and missed
opportunity, there a small hole appears before me, about fifteen feet
before me. And all the things I brought with me to this valley are in my
trailer, fifteen yards behind. And one by one I lose them down the
ever-expanding hole---cans, jars, movie magazines, photos, food and
books. The hole keeps caving in on itself. My little home on wheels is
silver, rounded like the new Volkswagen models, and drives like a dream.
Its doors shear off with a sudden crunch of metal, bright in the noon
air, they slither across the desert floor into the hole. I'm next
perhaps, hold on to me.

\hypertarget{chris-kraus-hunger-technology-emotion}{%
\chapter{Chris Kraus: Hunger -- Technology --
Emotion}\label{chris-kraus-hunger-technology-emotion}}

From my LA diary ---

Los Angeles, sometime in the late 90s ---

\lettrine[lines=3, findent=0em, nindent=0.1em, lhang=0]{M}{y heart and stomach}
flip while waiting in the endless gourmet take-out line at Say Cheese on
Hyperion. This is the third full day not eating \ldots{}. I stare
through thick plate-glass at tureens of baby peas in mayonnaise. Ten
bucks a quarter pound, they're canned. Little bits of foreign cheese
displayed on the top shelf like so many sad specimens. English tilton,
camembert. From the bodies of imprisoned animals to the air-conditioned
case, it's obvious this food was never touched with love or
understanding. The chubby woman up ahead of me seems to think this food
is good. She is luxuriating in the moment when she speaks her choices to
the shop girl, even though the girl is bored and hardly even listening.
I'd hoped to trick myself to eat by ordering the most exquisite food but
now this place offends me. Say Cheese, Say Choose. She wraps the names
of foods around her tongue, pleased with her passable pronunciation. Why
do I hate everything? The food here is so vastly overpriced, it no
longer smells like food, it smells like bills and coins and plastic.

If I'm not touched it becomes impossible to eat. It's only after sex,
sometimes, that I can eat a little. When I'm not touched my skin feels
like the flip side of a magnet.

The Alien penetrated me very slowly as we sat together on the bed. (This
is Ulrike Meinhof speaking to the inhabitants of Earth \ldots{}.As the
rope was tightening around my neck I lost perception but regained all my
consciousness and discernment. An Alien made love with me \ldots{})
Uncovering his body takes my breath away. The paleness of it underneath
the soft dark hair. The Alien was naked. I had several of my clothes on.
We're very still. Fibrational quivers between our bodies in the dark.
``This's exactly how I imagined it would be. So smooth.'' It now becomes
possible to say anything. Low voice. ``Don't move.'' ``I like to hear
your breathing.''

Like me, the Alien is anorexic. Sometimes we talk about our
malabsorption problems. Everything turns to shit. Food's uncontrollable.
If only it were possible to circumvent the throat, the stomach and the
small intestine and digest food just by seeing. After several weeks the
Alien decides that he will no longer make love to me because I'm ``not
the One.'' Aliens spend their lifetimes on this planet testing,
searching. They get dewy-eyed, nostalgic about hometown virgins.

I'm in my kitchen making chicken noodle soup for the Alien. It's his
fifth day of withdrawal from valium and heroin. He can't walk, can't
sleep. I want so much for him to eat. Even though he says he doesn't
love me, I can't believe it's true. Therefore, I want to help him. ``How
about a nice piece of wholewheat toast?'' I ask, ladling out his soup.
``Don't take offense by this,'' he says, ``but there's something I have
to tell you. Your cunt smells bad. If you washed the way you should, I
would've done the things to you I do to all my other girlfriends.'' I
gasp. Soup spills. ``Sorry,'' he says. ``I guess I should've mentioned
it when we were dating.''

Food stripped of all its color, nutrients and smells and then
reconstituted, like my expensive hair (he loves it), Ravissant Salon,
\$300, like suburban small town cunts drenched in Massengil.

If I could only eat, a little ---

Although no one, to my knowledge, has analyzed the work of Frederich
Nietzsche through the occurrence of his blinding headaches, the poet
Kenneth Rexroth reads Simone Weil's philosophy through her anorexia.
Both are ``egregious nonsense \ldots{} unholy folly.'' Rexroth puts lays
the blame where it belongs; on the Catholic men, Gustave Thibon and
Father Perrin, who took her seriously. ``If only,'' Rexroth speculated
in \textsc{The Nation} (1957), ``she had sought out an unsophisticated
parish priest, who would have told her `Come, come, my child, what you
need is to get baptized, obey the Ten Commandments, forget about
religion, put some meat on your bones and get a husband' \ldots{}''

What you need is a good fuck, he said to me.

In \textsc{Holy Anorexia}, the scholar Rudolph Bell wants to take the
magnificence of the medieval female saints and drag them down to his own
level. He does this by conflating them with contemporary teenage girls,
who he finds pathetic and ridiculous. St.~Catherine, St. Theresa, and
Hildegaard van Bingham are all essentially the same; they're solipsistic
brats. The collective trans-historic She, the holy anorexic, ``emerges
from a frightened insecure psychic world to become a champion of
spiritual perfection \ldots{}. Her will is to do God's will, and she
alone claims to know God's will.'' The holy anorexic is a manipulative
vixen; she ``commands the war against her body and therefore suffers
deeply at every defeat, whether it is a plate of food she gobbles down
or a disturbing flagellation by nude devils and wild beasts. Then with
varying degrees of success, the holy radical''---like the newly slender
teenage girl---``begins to feel victorious \ldots{}''

Like witches, or female writers, thinkers, artists, who use the names of
others when chronicling their lived experience, holy anorexics are not
merely people to be differed with; they must be despised.

Shouldn't it be possible to leave the body? Is it wrong to even try?
Hungry yet repelled by food, Weil wrote: ``Our greatest affliction is
that looking and eating are two different operations. Eternal beatitude
is a state where to look is to eat.''

"The Alien is in my eyes. He's flooding my eyes. He's completely
penetrating me, every bit of me in my eyes.

He's in my eyes, he's spreading into my brain. Oh God,

he's in my mind. He's making me feel things in my body that I don't
feel. He's making me feel feelings, sexual

feelings. And he's there. He's everywhere. My body's changing."

David Jacobs, 1988 interview with an Alien Abductee

\begin{itemize}
\tightlist
\item
  from \textsc{Aliens \& Anorexia}, Semiotexte/Smart Art Press 1999
\end{itemize}

\begin{center}\includegraphics{images/"pandanus.png"}\end{center}

In \textsc{Aliens \& Anorexia} I am attempting to make contact with the
writer and philosopher Simone Weil. At the moment of her death in
August, 1943, Weil became an Alien, i.e., a legend who transformed
politics into tragic poetry. Not all Alien encounters are hostile and
dispassionate invasions. There are others who see the Aliens as their
friends. Aliens encounters, like narrative, happen essentially in
realtime. In order to make contact with the Aliens, it is necessary to
carve out little pieces of yourself to let the Aliens come in.

In Paris, the librarian Florence de Lussy is editing Weil's
\textsc{Collected Works} for Gallimard. Weil's been dead for more than
fifty years. The project's overdue. The edition will contain eleven
volumes, and yet hardly any of what Weil wrote was published in her own
lifetime. Because she was an amateur philosopher, teaching philosophy in
French girls' lycees until 1941 when Jews were banned from working for
the government, Weil's writing has a narrative quality mostly absent
from the philosophy of her own time. She wrote articles for the leftist
press, reports, position papers and communiques during her years as a
trade union activist. Concurrently, and in the years following her
disillusion with the trade union movement, she wrote notebooks and
voluminous letters to her friends and colleagues. In solidarity with the
dispossession of the workers who her labor colleagues claimed to
represent, Weil sought out the experience of dispossession in the person
that she knew the best, herself. ``If the `I' is the only thing we truly
own, we must destroy it,'' she wrote in \textsc{Gravity and Grace}.
``Use the `I' to break down `I'.'' She was despicable; according to
Bataille, ``odious, immoral \ldots{} a dirty hook-nosed Jew.'' Had
anyone taken her seriously enough to prompt her to write professionally
for publication, her work would not have happened. She was writing to
find out what she thought.

``The body is a lever for salvation,'' Weil wrote in her notebook in New
York. ``But in what way? What is the right way to use it?''

Simone Weil was a performative philosopher. Because her texts are really
notebook writings, there isn't ever any subject that's apart from her.
Which is not to say she's writing ``memoir'' or ``autobiography.''
Channeling her subjects through her person, Weil does what writers do.
She is constructing a narrative in realtime---arriving at a state of
openness, witnessed by her audience, the reader---in which thoughts fly
in and out according to who's listening. In Weil's philosophy, just like
in narrative or phone sex, it's not the story that we're really hearing,
it is the fact and act of telling it. Her thought approaches
narrative---an emotional transparency that occurs when someone else is
listening to you.

\hypertarget{rachel-levitsky-on-narrativity}{%
\chapter{Rachel Levitsky: On
Narrativity}\label{rachel-levitsky-on-narrativity}}

\hypertarget{a.-alphabets}{%
\section*{A. Alphabets}\label{a.-alphabets}}
\addcontentsline{toc}{section}{A. Alphabets}

\lettrine[lines=3, findent=0em, nindent=0.1em, lhang=0]{W}{hy?} Because
I love to quote myself. Because I have a fantasy of saying it all.
Because of the dense foundation. Because it is that difficult to imagine
it's a girl. Because, I want to write a beautiful sentence. (Violet
Leduc) Because, A fine rain was filling him with hope and despair.
(ibid) Because there is so much to explain when you are not, say,
Stanley Kubrick, like why the back of the naked torso (female) is
relevant to the subject matter. Because I should have written of the
degradation of an aging pussy before Richard Hell did. Because the world
needs to be created and stated over and over again if one is a woman.
Because some of my favorite women writers feel they must write as men
because it is impossible for them to write action as a woman. (See
Christina Peri Rossi; Iris Murdoch told a college class in which sat
Marcella Durand, that writing as a man was easier.) Because I am afraid
a penis will be placed in my absences when what I need is a dildo, in
other words I need plastic, rubber, elasticity, strength, reliability,
trust, safety. In other words to my dismay, my absences, my blanks, my
ambiguities are most comfortable to men, meaning, they will define them,
but it is women that I want.

\hypertarget{b.-bully}{%
\section*{B. Bully}\label{b.-bully}}
\addcontentsline{toc}{section}{B. Bully}

Because I'm a poet I will tell you this is about narrativity and I hope
you will believe me. Because it is narrative and I am a poet I expect
you to believe me. I got together with Judy the Pothead DumDum. She
didn't ask me about my writing or say that I looked good. Trashy but
good. She didn't say it, so I am saying it. I wasn't drawing a cat, dog,
house. I was scribbling. I didn't make sense to you or to myself. I was
being messy which was a clear expression. I want you to notice me but
don't expect you to. I didn't wear underpants and it wasn't a bad dream.
My shirt is both dirty and stained, just what you'd like to accuse me
of. My lips taste good and my pits stink. You like the smell of my hair
but its just Vidal Sassoon, bought cheap at the five and dime, no,
that'd be Rite Aid, everything that is anything these days is
everywhere. I am practically nowhere. I am trying to remember to stay
off the drugs most times I forget why. My ends are splitting and I can't
afford the expensive haircuts and don't trust the lady at the five and
dime. My hairs true color reforms from the root. You'd prefer to be
fooled. Call me cat, dog, house frau in a light blue frock. I can do it
for a moment and my thin lips rock and roll in almost any shade, though
the counter boy'd have me believe one works better than another. I
haven't got a date or room on my card for dinner. Martinis are
expensive. I make them at home like Laura Dern's mother in Wild at
Heart. I take other women's misfortune as a warning. She's a wizened cat
so waits, knowing Ill move first and fall on my reputedly pretty face. I
made the point that the pronouns don't matter but watch how they do, in
vitro, embryo, room with a view. All dogs are as faithful as me. I've a
new middle name, its hard to be ethical. Watch the one who emits so
much, so much love, s/he will betray you most stunningly or most poorly.
Jade Netanya looks up from the middle of the photo, grief stricken, the
dead guy wielding the eight-inch hammer divinely inscribed with the
ineffable name of god was her fiance, though she's a lesbian. The rabbi,
the unbalanced baal shuva, Jade, my sister, fellow poets. She was
waiting for Ms.~Perfection/Beatific Amor but No One likes to be alone.
Its like being diabetic which uninterceded leads to blindness and
amputations. Nowadays in more adults and poor fat children. He was a
menace to the children. The complaining neighbors came to feel guilty as
Jews will, they didn't invite him over enough to dinner. Instead, they
made a lovely funeral. About me. As a child I was often having others
take off their pants and squat. Expunge or include? A? or P?

\hypertarget{c.-cut-me-ups}{%
\section*{C. Cut me ups}\label{c.-cut-me-ups}}
\addcontentsline{toc}{section}{C. Cut me ups}

I think it's Nicole Brossard who said, surely among others, that in
poetry the character is always far away, whereas in prose there are
characters more alive than herself. Let me look for it, aha, here it is,
page 27, in ``She Would Be The First Sentence of My Novel.'' The
characters, they will talk and argue among themselves, they will
disintegrate and reappear, they will separate and unite. They will
explain, as Yaya and Grace do constantly, how it is that they both care
to be lovely and don't give a damn about happiness or what you think of
them, how they understand that the monstrous is the more potent state,
and embrace it, but politically and for survival, they must cover it up
constantly. Grace, who primps like a lady, is as disturbing as Yaya, who
emits the foulness of hell on earth, when she (Grace) walks into a room.
The lookers on are confused by their own discomfort. The obvious
grotesque (Yaya) is easier to ignore. Men claim ugliness for themselves,
it is free, without cost, abandonment. It is theirs, they own it, they
saddle us with the cost of beauty. I saw a woman who'd been freed from
her cancerous chin, now an absence covered by scarred and taut skin. Of
course class, capitalism, if she had been rich . . . The next day there
was another, a s/he, a distortion. I fell in love. And today, there was
another, my cheating heart. A beautiful creature in a convex mirror,
reading \textsc{The Adventures of Buffy the Vampire Slayer}. In my
dreams I am another, idealized, the one with whom I am in love. On the
train, I am a monster, I hold onto the hope that you will not turn away.
I wait for the novel in which I am not a ghost. The theory gave so much
hope---those French (what is it about the French?), the Lesbian Body
(bloody), De Beauvoir's the potent female finger on the button
(machine). Feminine plural, a sex that isn't one, alter egos, as drag,
Dred, Queer Confab, Bound. I began to believe I was beautiful on the
inside. Another woman became a mother. Another lesbian a wife. At the
Lesbian bar a screech, a piercing cry on the heals of a bitter
complaint. Forty pounds heavier, those horrifying saggy tits, from the
mouths of babes, the new lover of a good friend's ex-lover. Political
moment? Nowhere to be found. At the mall, at the Body Shop you may send
postcards to senators. I escape to Coney Island, as yet mall-less. Here
I can feel this---we are in the interim, the momentless, bearing the
interminable wait for the next moment. When the next time comes, I will
still be writing the Lesbian body, forty pounds heavier, with sagging
tits.

~

~

\hypertarget{pamela-lu-my-narrative}{%
\chapter{Pamela Lu: my narrative}\label{pamela-lu-my-narrative}}

~

\lettrine[lines=3, findent=0em, nindent=0.1em, lhang=0]{F}{irst there was this thing},
this experience, that I felt myself having at last. I had felt the
experience coming on for some time, without knowing quite what it was
about to be, what it might resemble, or when exactly it might arrive.
Perhaps I had read about it once, glimpsed a description of it catching
a character, a familiar protagonist, a stand-in, in its sure, capable
hands. There was the possibility that it might miss me entirely as I
loped along the wrong side of the alley, ordered the wrong label of
scotch from the wrong west-end saloon. A possibility that I might lift
my eyes just in time to watch it happen to my neighbor instead and, all
too late, jump to my feet to exclaim, ``Hey, that's my horse! Well,
maybe not my horse exactly, but absolutely its blood relation!'' To
which the experience would reply, tenderly and oh so Jamesian, ``Not
this time, not you. Maybe next year, to the soul you will have become if
only you had learned to abandon the regret of your current, expiring
ghost \ldots{}''

Then I was sure that it had happened, not merely to a surrogate, but to
me personally, that I had not only jumped into its hands but entered it
in total, inhabiting its limits with the whole of my being. The
unanswerable question of it flogged me with the rapture of body and
soul. I called myself into existence by my proper name, and watched my
former life flow past in parallel streams, burbling and breaking up into
little tearful books. I demanded unambiguous signs of my entry into the
actual world. And this all happened in the middle of time, while I or
some other writer was fast asleep, dreaming up masterpieces that were
bound to be forgotten. While I was too busy becoming the character who
would have become me in an ideal life, friends carried on in my absence.
I felt the need to track the movements of my friends, to check in
periodically and make sure no one had died or disappeared. Usually they
could all be accounted for, milling about in packs around the pool
table, as an unidentified narrator slipped through their midst like a
generous thief, strewing stolen, anonymous pears left and right.

In the slow crush of time, characters aged, my friends grew older, and
the unthinkable happened to people. I did not, however, in the moment in
which I was destined to live, understand any of this. I had been
impressed by the nature of fate and outcome in E.M. Forster novels --
how, in the absence of divine intervention, the hand of fate was given
over to mortals who schemed compulsively, recklessly, and blindly to
bring about often unfortunate results, perhaps even a tragedy. There
were ambitious do-gooders whose ambitions caused more harm than good,
prejudiced Victorians who lashed out politely and brutally at the weak,
colonial subjects who pleased their masters by losing out again and
again. The characters lived; they brought life down upon themselves;
they wrestled with causality and causality won. Is that why I'm always
confusing causality with casualty?

A friend from school once observed that fiction students as a group
seemed bigger and physically stronger than their poet counterparts. How
much would this statement suggest a three-ring bout in the writers'
circle, with scrappy but effete sonneteers grappling at the well-muscled
backs of market scriveners? What kind of over-dense poetic whimsy could
tickle the plodding narrative to unclog its drain to the real story?
This story was so real it was almost untellable; one would have had to
be already dead, installed in some paradise or afterlife, in order to
understand it completely. The neutron star of its as yet undiscovered
language weighed the earth down in orbit, while sending out regular
pulsations of superhot, ecstatic light which caused even fatalists to
see. And yet the parameters of the story demanded the most outrageous
courage of those who would experience it: biblionauts who would sail
paper spaceships deep into the fiery outskirts of existence and return,
defenseless, with a strange, material testimony. Was this, after all,
the true knowledge and danger that awaited our fumbling lives?

~

\hypertarget{nicole-markotic-narrotics-new-narrative-and-the-prose-poem}{%
\chapter{Nicole Markotic: Narrotics: New Narrative and the Prose
Poem}\label{nicole-markotic-narrotics-new-narrative-and-the-prose-poem}}

\lettrine[lines=3, findent=0em, nindent=0.1em, lhang=0]{I}{ am interested}
in a dialogue about ``new'' narrative, which is perhaps not so much
\emph{new}, as newly theorized. Many prose writers do not consider
themselves fiction writers, yet at the same time are not really part of
on-going poetics discussions which, for the most part, do not focus on
narrative. Although I also write prose fiction, I consider my prose
poetry and other alternative, interdisciplinary, and innovative
sentences to be a neoteric prose that both challenges and expands
language boundaries.

For me, the prose poem is a poetic strategy embedded within the
structure of narrative, and a feminist response to patriarchal language
and forms. By embracing both prose syntax and poetic disruptions, the
prose poem defies conventional linear grammar and refuses to satisfy my
desires for either poetry or story. My desire is for so much more than
causal, linear, rational and persuasive normative sentences. In my novel
and in my poetry, I try to live between the promise of narrative and the
fulfillment of the habit of fiction.

By expropriating two distinct genres while presenting a form of writing
that is seen to consist of both, prose poetry presents itself as hybrid
writing; a hybrid that I explore as crossing between the desire to
exceed formal considerations, and the narrational insistence surrounding
such considerations. In my prose poems, I try to embrace the problematic
of poetry that ``looks'' like ordinary prose narrative, yet invites
disjunctive readings which may extend beyond traditional poetic forms
and conventional conceptions of narrative. My pieces ``fail'' as
poetry--yet continually celebrate the erotic contradiction such failure
narrates.

The prose poem, for me, offers itself as a ``genre'' that permits and
encourages me to move freely from sentence to line-break to full-stop to
repetition. Without entirely settling on its own definition, the prose
poem rattles my airtight nerves, upsets my crabapple go karts,
insistently and lasciviously rubs my discourse the wrong
way---inventing, always reinventing itself. Its very possibility denies
its own existence. For me, the prose poem remains liminal: not merely
transgressive, but indeterminite, caught in the phoenix act of
disappearing, reappearing as its own \emph{possibility}.

My interest in this form begins at the level of the sentence. What makes
a sentence a line of poetry rather than an excerpt from an essay or
novel? The answer, for me, is the surrounding words and sentences, the
position of those words on the page, the complicated ache of the forced
return at the end of each ``line.'' Ear and eye games, the playing of
interruption and proceed. Not stop-and-go, but go-and-go-farther. As
soon as I take a step towards the horizon, the horizon
reconfigures---itself and me.

What interests me about the sentence as poetic form is that poets who
are \emph{least} interested in representational language are most often
the ones who turn to a grammatical structure which invests heavily in
sentential logic. The use and abuse of the sentence, then, seems to be
the focus of the poem. And because one sentence follows the next--which
offers a visual hypotasis but a semiotic parataxis---the poem insists on
a reading that is far more linear (for lack of a better word), than is a
line-break poem with disparate images stacked one above the other. The
one poetic device designated as prose or \emph{not} poetry, the
sentence, becomes the choice unit of composition for poets who intrigue
me, shock me, paratactically get under my skin.

Within each sentential piece can emerge a plasticity resistant to
notions of purity in either prose \emph{or} poetry. Gertrude Stein chose
the sentence as her basic unit of composition. Her sentences release
readers from the semantic baggage traditionally loaded into narrative
prose as they re-define, distort, conflate, skew, and otherwise render
the sentence pliable and visceral:

\begin{quote}
A kind in glass and a cousin, a spectacle and nothing strange in a
single hurt color and an arrangement in a system to pointing. All this
and not ordinary, not unordered in not resembling. The difference is
spreading. (\textsc{Tender Buttons} 9)
\end{quote}

Such prose rumbles inside the belly of poetry, or perhaps poetry shifts
awkwardly within the abdomen of prose. Readers, seduced by sentence
structure, discover themselves trespassing an erotics of prose,
transgressing away from familiar and known fiction offerings.

Writers engaging in what Stephen Fredman calls ``poets prose,'' struggle
to create new ideas from \emph{inside} the conventions that constitute
the either/or structure of the prose poem. This structure is the chance
for prose writers to both follow narrative and to deviate away from
story. For example, Robert Kroetsch, in his poem ``The Sad Phoenician,''
pursues an elusive and/but dialogue that generates a bifurcated and
pending narrative:

\begin{verse}
 and        the Phoenicians gave us the whole works\\
 but        what does that matter to a world that ignores\\
              them, the Greeks got all the credit of\\
              course, because they stole the alphabet\\
 and        the girl from Swift Current, she more or less\\
              took everything\\
 but        the kitchen sink, claiming all my books, my\\
              my records, my prints; she moved in with that\\
              photographer from Saskatoon, the one who\\
              takes those sterling pictures of the wind\\
 and        I should sue\\
 but        she follows large flocks of birds, I hear,\\
              calling my name\\
 and        pleading\\
 but        why she developed a thing for adverbs, that’s\\
              too rich for my blood, I want to tell you\\
 and\\
\end{verse}

Readers wander through an and/but tug-of-war that is also a narrative
wandering, from the Phoenicians and Greeks to Saskatchewan, to the very
grammar of the telling. It is this \emph{telling} that gets so many new
prose writers into trouble---isnt plot the very basis of novels? How can
one tell a story and avoid fiction? And why do we want to?

Feminist writers must resist and reinvent patriarchal language to make
it our own. We cross a border every time we pick up a pen or turn on the
computer or hum into a tape deck. The border is a visible line on the
page that the prose poem physically outlines. For many women, this is
not simply a gendered border, but one that has been criss-crossed by
social lines such as race or class. Harryette Mullen, in her prose
poetry book, **S*PeRM**K*T**, investigates the connective links between
the social and the linguistic:

\begin{verse}
Eat junk, don’t shoot. Fast food leaves hunger off the hook.\\
Employees must wash\\
hands. Bleach your needles, cook the works. Stick it to the frying\\
pan, hyped again.\\
Another teflon prez. Caught in the fire around midnight, quick and\\
dirty biz. Smoked\\
in the self-cleaning oven.\\
\end{verse}

Mullen's disconnected sentences can be read to produce meanings that
address subjects ranging from drugs to poverty in contemporary USA. But
the structure of the ``narrative'' does not create fictional characters
and circumstances to battle these issues; rather, Mullen's words entice
at the same time as they startle reading practices. Word association and
colloquial speech patterns reinvent the chance to re-enter the text from
another crack. This text does not close the book on the story, it
continues to remind that reading is a socially-coded act, and that a
reader's subjectivity is constructed through how she enters language.

Language enters me, and my sentences struggle against what I know to be
``just'' a sentence. I hesitate in doorways and on fences, I plan my
route along edges and cracks in the sidewalk, I pencil in margins made
up of words only words, and then I cross over into blank space. Not
because it is ``new'' or ``white'' or ``virginal'' or any other dominant
appropriative term, but because the blank space is where the prose poem
offers a \emph{visible} margin, a territory (known or not) that hangs
onto the edges of writing, clings to the edges of edges. I get edgy when
I try to read or write a prose poem, anxious, excitable, downright
perplexed. I hang about on the inside of this perplexing form because it
keeps me on my language toes: unsure and insecure, tip-toeing across
whatever lines a poem insists upon.

I ``plot'' at sentences and subterfuge, dream of a poetry that
withstands ragged-right fashions, yet continues to address narrative
within disjunctive non-lineated prose. I look to the prose poem as
mutation-- by definition unstable, punctuated by eternally reproducing
desires.

I want to carry words to where theyre already going, but not necessarily
to where I think they're carrying me. I strive to enact an and/but
\emph{narrative erotic} that is neither fiction nor poetry, that invites
readers to see beyond seeing and write beyond reading. I veer away from
the end of that prose line, from the punct signifying the closing
ceremony, the final page next to the cover, the certain and the
certified. I question the metaphorical imperative that insists toes are
what we tip on, and grasp instead the reassurance of instability,
embrace the inability to embrace a stable subjectivity, and echo
cheerfully what Kathleen Fraser describes as ``fragments of a wholeness
only guessed at'' (\textsc{Each Next}). My prose poems want/desire/lust
after such guessing contortions.

\hypertarget{ashok-mathur-the-narrative-site-narrative-as-preliminary-investigation}{%
\chapter{Ashok Mathur: The Narrative Site: narrative as preliminary
investigation}\label{ashok-mathur-the-narrative-site-narrative-as-preliminary-investigation}}

TEXT MISSING

\hypertarget{laura-moriarty-from-nude-memoir}{%
\chapter{\texorpdfstring{Laura Moriarty: \emph{from} Nude
Memoir}{Laura Moriarty: from Nude Memoir}}\label{laura-moriarty-from-nude-memoir}}

\hypertarget{section-12}{%
\section*{9}\label{section-12}}
\addcontentsline{toc}{section}{9}

Ten minus one

She counted but was

not believed. The syllables like a mantra. The answer. Counting in order
to rid herself of numbers, of time. The increments of the present. The
same count. A decimal minus one reminds one never to go over or to go
back. Not a sequence but a repetition like a song. A petition like a
question. Never again. Duchamp as a form of grief. A diary. A display of
the invisible. Of visible decay. The mode here is eliminative. It is the
only mouth she has left.

``Piece of Ass Lost.'' The necropolis pictured by a dead poet in
Ellroy's \textsc{Clandestine}. Poet cop. The mouthpiece of fate.
Exploitation as love. Women who don't get older. Women as men. It
doesn't matter. The U.S. number one in violent death just past Mexico
and Brazil. First in the First World and the Third. The guardian of the
infested spirit. Someone follows a woman. She is not a muse but a fate.
She talks too much.. They drink together.They eat. Is she one of many or
a singular masterpiece. Impossible to follow him down that street. The
action is complete to the degree that it is not legible. Or transferable
to another medium. It is not unfinished but undone. Not the crime but
the gun. The piece of string with which you hope to find yourself.
Missing. She was missing her head. What we have is a reconstruction.

Diana watches herself on TV. She wanders absently around the house. She
is not dead. She is an artist. The interview reveals everything she
hopes for. There is a queen and a corpse. She finds her cigarettes and
begins to smoke. She is seventy. Voluptuous as parchment. Thickly
written. She takes into herself a sense of death. She lets it out.

She wonders how to preserve access without giving into a deadened sense
of hierarchical exclusion. Now that authority has been shown to be the
shuck it is. To provide a method for reading, to alter the activity of
reading itself. In retrospect or in the sense that it had already taken
place. She read the scene before the crime. Or we wouldn't be having
this conversation. His ambition is naked, mechanical. He also wants to
read.

The vaginal scroll

She performs (Schneeman)

Memorably laid out

Source or origin

Like a physical note

Of itself sufficient

Opened like the book

She wrote

The reconstruction is sloppy. Bad. Not convincing. Not made to be so.
Black velvet lines the unseen back of the door. To soften the blow. Her
character pokes through her performance like bones through flesh. Judy.
Judith. She loves him but he loves her double. Her twin. Who is
convincing. A con, snare or fox. A pest. He likes the angularity. The
bending back. But he falls during the kiss. He is injured. ``I hate
this'' she says. He doesn't miss her then.

The scroll unwinding and the performer

Takes over removing

Her shroud her look

Distracted in the picture

She takes off

Her glasses her theatricality

Will be attacked in place of her person

He argues her down

Her desire is unknown

Impossible to predict its hold on him. ``\ldots{} a diagram of the cruel
geometries of desire.'' The reviews are ambivalent. The scriptedness of
their exchanges are painful to him. Her double jumps into the bay. The
actor after. Also his double. Later on a sound stage wet close-ups. The
real actors. His strained face dazed with obsession. The paleness of her
hair, dark dress spread out but clinging also. Her limpness in his arms.
He stares down at her, climbing imaginary steps out of the sea. Her high
heels in silhouette. The twist of her waist is displayed to him, to us.

\hypertarget{section-13}{%
\section*{1}\label{section-13}}
\addcontentsline{toc}{section}{1}

Madeleine is weak but alive. Unconscious. He takes her to his apartment
on Russian Hill. He kicks open the bedroom door, placing her carefully
on his bed. He first takes off her shoes. She stirs but doesn't wake. He
holds her in a sitting position, unzipping her dress in the back. Pale
wide shoulders, wet slip, brassier. He pulls her dress and slip over her
head. He becomes aware that her breasts are against him. Nipples like
buttons. He works quickly against her waking. Unfastens garters. Peels
off stockings. He reaches around her waist and has some trouble with
clips. He begins to breathe hard but quietly. Her head falls back. Her
neck is long and white. His hands are dark against her skin. He is
careful with his hands like a surgeon. He draws the blanket over her.
Even in this drowned sleep there is a defensiveness to the set of her
features. He stares only briefly. His audience, himself, is aware of the
contrivance of his composure. But is unaware that her unconsciousness is
an act. He knows only in retrospect. Perhaps she knows that he is both
protector and killer. She seems to be in control. Also looking back. But
she is out of her depth.

The Bride does not refuse

this stripping by the

bachelors, even

accepts it since she

furnishes the love gasoline

and goes so far as to help

towards complete nudity

by developing in a

sparkling fashion

Her intense desire for orgasm clings to her like his red silk robe. When
she runs she knows he will go where she goes. She has tricked him but at
too great a cost. She sees this as a job. He sees it as rescue and sex.
But she can't be saved from the danger of being false. She can't be had.
The hopelessness skews their perception. Heightens their senses.

The artist likes to see the woman go too far. The woman likes it too.
She manages the twists with some virtuosity. She puts herself beyond her
skill. The character changes in desperation but is unable, each time, to
survive. Her eyes in multiple shots. His eye in the credits. Jimmie. A
simple expedient but effective. Death to itself. And Kim. The actor is
left.

To be disassembled

Exit stage right

And another rescue

Marooned on a penal colony in space

She remembers his memory

His voice a slurred machine

``My nakedness creates you,'' he says (dead)

She brings up the interface

Turns it on

A golem. Like herself. Skin like flesh only not in her mind. Inside the
memory, the wetware an obscene cream as if robots ran on semen. She
turns to him. In pieces. In a kind of mechanical pain. He carries in him
the mothers who didn't survive. It is evident in his soft address. He is
her confidant. Her history. She confides. They confer in a jacked-in
version of love. His love is wired in. They know things together that
otherwise only she knows. ``It's better than sex,'' she says. He says "
No, it's not."

Kienholz by definition

The bionic man

``We can rebuild him''

Not a faceless door

But organ donor

Plugged in

What happened to him

If she is Death in Orphee

He is her driver

An assisted death. They would be lovers in another life. She prefers his
later more vicious fiction but he hasn't grown into it yet. Their times
are not synced.

\begin{verse}
He in fact She assisted\\
Was scattered It was speaking\\
                    Not belief\\
She longed for The expert\\
She was Not love\\
\end{verse}

~

~

\hypertarget{lisa-robertson-for-soft-architects}{%
\chapter{Lisa Robertson: For Soft
Architects}\label{lisa-robertson-for-soft-architects}}

~

Lucite

(a didactic)

(because the present is not articulate)

Sit us on Lucite gently and we will tell you how knowledge came to us.

First the dull mud softened resulting in putrefaction, lust and

intelligence, pearl globs, jeweled stuff like ferrets, little theatres
of

mica, a purse containing all the evil smells of daily life. Then just
the

one vowel, iterate and buttressed and expiring; leaning, embracing,
gazing.

It devised with our claw identity for the sake of food. Selves, it says,

feeding us, I adore you, you know. Like a boy blowing from a tree, we

decided, we were paid, we were free. We incessantly prepared for the

future. On the title page, two angels blowing on the trumpets of fame
held

up a globe decorated with three fleur-de-lys and topped with a crown. We

learned habits and tricks; we faked happiness and relief. We were a
single

grin with lips pasted back. We said we saw Europes of hallucination,
fatty

broths sprinkled with deer, stenciled eagles, serpents and lurid rags.
That

was a format of saying, a frayed ligature. We were fading into the
presence

or absence of food, of sleep, of doubt.

Enough of the least. Doubt wants no space. Sincerity takes too long in
an

aggressive emergency. The present is not articulate. Also we feel a
sense

of duality. We wear out the art. We start to modify our vocables--
flick,

pour, dribble estrangement's sex. Since it is we who are one, and we who

are scattered. We're this pair or more that can't absorb one another in
a

meaning effect. We feel palpitated by daylight and its deliberate
plants. A

description makes us. Among all these ejaculations and judgements we
feel

its elsewhere sculpt our body.

We would be walking down the street in the poetrycity. Gauze would be

everywhere. The day would be big, halting, gracious, revocable, cheap.
We'd

be the she-dandies in incredibly voluptuous jackets ribboning back from
our

waists, totally lined in pink pure silk, also in pure humming, and we'd
be

heading through the buildings with knowledge---that is, ephemeral

knowledge, like leafage or sleeves or pigment. The streets are salons
that

receive abundantly our description. The soft buildings are charming. And

our manners are software. We feel sartorial joy. We'd be at the river

watching the fat water lap the blonde built part, loving temporal

improprieties, the bright trash floating in slow liberation. We'd be

applying our make-up at noon by the river, leaning on the balustrade,

thinking about a little shun, a little fight, a little sofa. We'd be

thinking about hinges. We'd feel for our pen. Something might seduce us.
A

likeness. Samesame pouring through it.

Knowledge comes to us.

~

~

\hypertarget{camille-roy-monstrous}{%
\chapter{Camille Roy: Monstrous}\label{camille-roy-monstrous}}

~

\lettrine[lines=3, findent=0em, nindent=0.1em, lhang=0]{O}{ne of the forms}
of narrative I write is software. Its lucrative. About four years ago I
used stock options to buy a house right around the corner plus one block
from one of the worst housing projects in San Francisco. A couple
thousand people live there. It gives my neighborhood the highest child
hunger rate in the city. Our first night in the house someone got
murdered, just before midnight . It was a block a way but the shot
sounded like it was in our back yard. One shot, a pause, then another.
Purposeful. Somehow I knew it was intended to kill, and not just a
couple of kids shooting at the moon. Plus the neighbor told us hed had
his car stolen 3 times.

Impenetrable poverty plus dumb fuck rules, class and race segregation:
Id moved into the only San Francisco neighborhood that duplicated on a
smaller scale what I grew up with. It annoyed me.

Locality, forever. \emph{Skewed}. Something huge gets mutilated as it
slides through a stuffy tube. Were on the beach very far to the west,
watching what pops out. It contains all of American culture. I came here
so tightly wound. Born on 43rd street south side Chicago \& havent been
back since I left the hospital.

From my dining room window, at the rear of the house, the project looks
strangely vacant. There just never seem to be many people around. The
buildings proceed down the hill towards the old industrial port like
giant shabby steps, but there is never anyone on the racks of balconies.
Ive rarely driven through it. Structurally, its sort of a dead-end
place, the way its laid out, like a suburban subdivision: streets point
into it, then twist up like spaghetti. The few drive-through streets are
dotted with dealers scoping out the passing cars. Im just talking about
the roads.

When I first moved in I often found myself dreamily staring out the
dining room window. I wanted to check out one of those balconies. The
view would be amazing, they practically hang over the bay. Developers
have been salivating over that piece of land for years. Nowadaze they
are nibbling at the edges of the project, building expensive live work
lofts for software designers on adjoining vacant industrial land. Its
weird. Different economic classes get spliced together via crimes. The
mode of interaction being criminal. So one day I mentioned to a friend
of mine that I didnt get it, how did dealers get kids to work for
them---playing courier, or delivery boy. What would a dealer have on a
kid? Why get involved with some jacked-up, scary asshole? I felt like an
idiot as soon as the words left my mouth. Patiently, step by step, my
friend explained how it was done, until I could have done it myself, as
obviously he had. All I had to do was ask. Knowledge. The getting \&
taking and the tearing up. Did I want to go there?

Of course I did. One day I walked in, took a place on the balcony next
to all of my friends \& drank their salty water. I listened to the
radio. I watched as a crack lady ran down the street behind a white dog.
Then the dog was scratching at the door. When I woke up that sound was
the shade, bumping against the window frame. And I was thinking, as I am
always doing and my thinking told me this: \emph{This is what I want.
Its inside my system of attractions. Im penetrated by the present and
its always the same: chronic anger. Awful but refreshing}.

From up here, it is all visible. From down below, also. Radiant
contradiction. Eyeballs: the severely vivid mechanism. Finally what is
seen is not a target but just circumference, expanding. Highlights
scatter across the field.

I walked into the projects a couple of weeks ago. Its right around the
corner, why not just walk? It was a friends birthday. She told me where
she lived, but it wasnt easy to find. The apartments didnt have numbers
on them, you had to just know. I asked a bunch of people. Kids were
running everywhere. How come I hadnt seen them from my back window? I
look whiter than usual, I thought, looking at my hand. Up here and not
even shopping, that made me odd. People looked at me skeptically. I felt
skeptical about myself, but slick, as in greased. I wanted to fall off
my little ledge. Bored with what had gotten dished up as myself. The
backwash of swallowing it. That nausea.

The balcony was great. I hung with my friends and listened to the radio.
They played that song I like, the one about money. Later we went out to
eat birthday steaks.

California is shallow. Thats true. Though it thrills me that I can walk
across the city without getting beat up for crossing some invisible
dividing line of racial turf. Of course I could get beat up for
something else. \emph{Im so easy to please.} Its the instability at the
heart which is to say the heartlessness of just washing away faster.

Im supposed to write about narrativity but these problems of locality
are where I get started. For me writing grinds itself into whats
familiar yet unbearable. Add mobility to that and, voila, narrative.
Disjunction is the formal consequence of this ripping and tearing, and
its packed with information, almost to the point of being insensible.

The streets I walk measure me. They measure you too, through mechanisms
both criminal and friendly. Including that knowledge is a kind of
spectacular innocence---the moment of saturation feels dazzling, but
there is probably no point. Still I love it, formally and erotically and
intimately. Its all about nested structures. I entrust my twisted little
pieces to the warm nest of the sick social body, and I feel our bond. It
nourishes me.

To theorize my point of view, to pursue critical formalism as a ritual
and as a grasp for power, let me put it this way. Narrative provides
context so that the rupturing of identity is recognizable. I think we
are impossible beings. We ruthlessly evade scrutiny, yet recognition is
the beginning of transformative emotion. Its a feeding process. You dont
know if youre creating a monster.

As a narrative writer I improvise recognition. Its like a location from
which mutant beings emerge. This feels true, in life they never stop
emerging. Look---they even swarm through this text. I allow it because
Im terrified and seduced. To encounter them via narrative is to
formalize a moment of surrender.

(An earlier version of this piece was presented at the Portable Talk
Series, SUNY Buffalo.)

\hypertarget{leslie-scalapino-narrating}{%
\chapter{Leslie Scalapino: Narrating}\label{leslie-scalapino-narrating}}

\lettrine[lines=3, findent=0em, nindent=0.1em, lhang=0]{I}{ just write,}
without advance plan as to a plot or form the writing driven by a wind
of my being in some urgency at the time. Each work I've done has been
written in that way. Each is different. The next time I'll write
dependent on what's occurring then. Here I'll say something about Defoe
(Sun \& Moon, 1995) since of my prose it was the hardest on me (in a
state of extremity) and required the greatest extension. `Narrating' is
for me:

Inextricably related not simply because events `are related', but
because the relation intrinsically is oneself. The phenomenal world of
nature is inseparable from one's `conception' (visual and tactile) and
also does not come from oneself. One is one's relation to the rose
disc-cloud on the desert as to social occurrence, such as people's
minute nuances or war:

"disc floating on the desert, with cattle that come to the edge of the
blue water and the white desert not coming from it.

One has to agree with them or is not there which in society can not push
out from themselves the red soft hanging disc. that is separate.

misery in having to be in agreement with them giving up being there
though one can't do that."

One can not be as social definition, yet that exists. That social
relation is inherent conflict as oneself in which one cannot exist or be
`causes' the syntax to be a `place'. Surfaces collapse on each other,
existing together as plates in paragraphs, clarity only possible by
inclusion of, or being, the conflict. Conflict is the writing's shape
configured and is attention in space.

Sleeping dreams are `compared' to the literal ``day''-as-surface as if
both day and dreams were surfaces as thin discs and placed side-by-side
the ego not in either of these but there also:

``One sees oneself as simply a shallow behavior that is thin. and so one
is free. And then turned inside so that one's dreams are sent into one
from the day, as they are it. So there is no REM. People rest in the
thin area. One laughs frisky. love in behavior is there, in reverse.
That is the street, the way it is seen.''

Also fictional events and real events (the Gulf War) are placed
side-by-side regardless of their size:

"An object connotes by a prior connection or culture.

Where the object loses, or begins, its connotative existence is to be
just then non-connotative. (To see its edge\ldots{}) would place the
piece at the point of dissolving. forming as an object and socially.
which is nothing. An object's form isn't anything as that. So it's the
point of the form and the conceptual meaning separating. (or rather,
being the same)\ldots{}Trying to get to or to see where actions separate
from their social existence which is what they are that is their
connotative being, so seeing an action without that."

Action is seen `without' itself in the sense that action (that which is
enacted, as physical movement) is `other than' writing, the one can not
be the other they (both) are also phenomena as conflict, thin discs that
impinge on each other:

"Action collapses on itself and is compressed, in that it is description
of itself. It is thus in the present-time, still and calm.

It does not diverge from itself, supposedly. That's a conception which
people already know. Therefore speaking enables one to see the diverging
or separation that is this present. These actions occur as if to make
the pupil of an eye (some other's) dilate, and be held open. The present
takes place as `some other.'

Life is dilated.

It does not diverge because description cannot be separate. as it is of
itself."

Defoe, the writer, created fictional characters who were supposedly real
giving documentary accounts. They were to be imagined as entities, which
had no reality. Paragraphs are plates of actions which are not single
and do not begin or end there in order for there to be reflection (as
the text) throughout everywhere that is not mirroring of authority as
authority is not (is not in) any single one of those actions:

"To name it will be merely caught again in their authority as they will
recognize that as themselves and will be in their trap again.

If there is no area between his dreams and waking life, there is no
identifying with anyone.

Walk through stream of hot metal bumpers of cars that then move like
plates shifting. Mass of plates and she's veering in it. On tin tail of
motorbike of yellow teeth thrashing in neck. He's ahead then. On it
again thrashing neck. The bike skidding and swaying in narrow channel
amidst bumpers. It bucks forward. Neck bike veering and bucking out on
the vast tin sheet. the sunlight is reflecting off. He wriggles free,
lunging. The motorbike lunges forward. Crashes into car's side in the
mass. Slippery blood on his head neck slashing. She's slashed hard
slamming into car's side. Rider on foot then, veering and is way ahead
then out on stalled mass.

Seeing him out on the vast reflected expanse."

Prior sight we see things before they occur: whatever dream occurs,
that's waking life. Perception does occur in people (or `in' something
happening) what is it?

"Perception occurs before the context it is not in a setting.

Without seeing it as it was going on while seeing it, really. The man in
being authority is the process of down fawning on the feet of the other
man while this is merely an image a negative that is imposed on it. Like
an angel.

Is it that seeing into the future as what's going to happen to people it
then just is that. It occurs. Isn't deflected. That isn't continual
change which is aware as we are in it.

We can see it in ourselves. that's why it's aware.

People see that later. Not when they're young. it's separate from what's
actually there which is to be seen directly. has to be.

There is no `inner reality' to the destruction that this person is
creating for themselves from within, which is seen from the outside by
others.

So suffering is in oneself, and has no `inner reality'."

The relation of creating a spatial shape as the writing to changing what
occurs (a friend dying from leukemia was having a bone marrow transplant
at the same time as the bombing of Iraq the illusion that one could
avert and change this reality while it's occurring in the mind relation
that's the writing yet the friend died and the bombing destroyed) is an
impossible terrain/conflict. The illusion itself took one apart. It's an
illusion that is bathos, like jump-starting a car. There is nothing
commensurate. Fictional jump-starts:

Real boys (and others) starving in the Sudanese desert. One is in or has
the dreams of one's fictional characters to make a particular surface.
"There is no resemblance between the boys and the boys in the desert,
that's forward in time supposedly. really, as it has reality in itself
and is separate from this. One non-reality against others in oneself, it
is seen between them. as oneself is simply that.

It isn't changed when it gets to this. Or by it later.

Dreams are excluded from this, which occur. That's convenient because I
can't remember them. dreamed one which met this red (actual) sunset that
just occurred.

This isn't it. Flattened to my inherent nature. I just woke up. The man
who was the driver deserted by the boys who don't resemble them on the
desert, has a dream:

Venus was hanging above the red smear, which I didn't see in it. The
neighbors with whom had formerly been at war for many years agreed to
smuggle oil to them.

It doesn't resemble them and can be reproduced and only seen in reading
as lovely isn't the market. Seen as lovely it therefore isn't the
market."

In writing: ``The seeing is the event, rather than arising from its
occurrence.'' And: ``People's experience lags behind the time they're
in.'' In other words, Defoe is trying to find the relation between
seeing (apprehension) and occurrence (events occur `outside'). There is
a relation one makes:

"Make writing that forms before motions.

As when I was in a rage, it formed first I'd forgotten it so the anger
occurred, I was in misery. Walking was at base and therefore observation
occurs.

Not observation occurring about it after.

Fictions are so simple that no observation occurs after.

So the fiction occurs first.

Defoe occurs first and so has no validity. It's not bathos and is only
distorted to be first or as being that.

So there's no Garden of Eden. There's not authority at all, that's to be
at base.

Walking occurs first and observation is not derived from it.

I want to still observation first. I started to walk one time with no
writhing and didn't remember beginning walking. I had nothing to writhe
about. There isn't a start of that, that's separated from walking.

One can eliminate power in oneself by not eliciting placating. Memory of
what hasn't been seen before. So fiction (I mean all conceptualizing)
occurs where the real event isn't isn't seen, it's eliminated, but is
occurring or has. Openly doesn't resemble it.

There's present-time first. That's what fiction is, when it is entirely
stilled. It doesn't resemble the event."

Fictionalizing drives the mind-terrain into a state of separation. There
the dream abuts itself. A real event (such as here an anti-war
demonstration) also abuts itself only:

"The dream I had is closer to being because it abuts it then. before
it's occurred. This is in the present. The blossoms on the plum trees
opened. On some. Birds sit, and fly amongst branches of bare trees. It
has no reverse side, and is thin.

Birds skim low by me.

Wall of police swinging their clubs sweep forward. young punks with
tattoos emerge from the crowd, turning the cop car over and burning it.
They receive beatings.

\ldots{}A memory as just a thin disc, as it is seen that way. There is
no event.

Warm, the birds stirring and flying. I had a memory of being on an empty
street in San Diego near-by the ocean; just that. that as a disc because
it became so stagnant it seemed it was going to collapse, and be still.
it flapped. my mind had become so tired its resistance was going. The
mind being weary as clarifying in itself.

\ldots{} `false joy' because they're lying to us. but inside it is not
their `false joy.' It could be one's own, that is a thought. which in
itself is joy which bangs against it. The dark cattle stood earlier with
the disc of moon on the other side of them why they were dark; when the
sun was on the side (of them) on which they were dark.

That is in memory as the actual event.

The chest in one is struggling in constriction of joke. Trash. The air
cannot get into the rib cage. The thought of a bicyclist comes up which
is a memory and therefore separated from one, who is small. That was him
and so flattened shooting ahead. His back, that makes one laugh
constricted inside so that it's him.

So the flapping swans here so they're the same as the sky simply flying
is that (in time), before.

The entire day completely dilated is flapping in one's recesses so it is
behind in the back of one. Stagnant in that sense, as it's flapping and
then narrowing it in front though it's still the day.

It's the day in a sense, narrowed to a thin disc in front in which the
entire dilated day behind in the recesses of oneself bangs shuttering on
it. The day collapses in that sense when that's slowed and interminable.

One rests.

One just sees from one's social group perspective say and then the huge
bright day is dilated flapping with that; so it's involuntarily let go
as it's large and while hitting the recesses in front of one in the
present, remains there shuttering."

Comic book style, the apprehension occurs via ``the other'' who has a
hole (or eye) in her side, everything seen alongside and flattened. The
visual slots or frames of the comic book (only conceptually implied) can
include all layers and perspectives at once, horizontal and vertical.
The Other is in love with James Dean who in death is turned into a deer.
She pursues the deer to rescue James Dean as the black rain falls. She's
traveling through refugees dying from blood blisters on the roads. Her
own blood blister is gently lanced by a soldier who's driving the jeep.

The incommensurable relation between outside-action (in the world) and
interior-action mind phenomena is overtly the subject. The shape the
writing takes is related to use of fictional and
autobiographical/historical event (such as the bombing of Iraq) at the
time it's occurring. The time it's occurring the paragraph-plates are
real-time moment of exterior event brought into compressed relation of
interior occurrence/horizon together and always separated. The intention
is that fictional illusion be at the moment of real-time to: implode it
as real-time action.

These are illusions in the practical sense of being `only' writing
(writing has no relation to present or historical reality it has no
reality, is it as well, being mind phenomena. So the `ordinary' small
action is {[}to be{]} as much `reality' as events that are devastating).
I am trying to divest hierarchy-of-actions. `Hierarchy-of-actions' voids
people's occurrences (that is, individuals' actions are relegated to
inconsequential or invisible). Such hierarchy substitutes `overview' of
`history'/interpretation/doctrine therefore to divest
`hierarchy-of-actions' is certainly a political act. (In one's/reader's
/viewer's conceptualization then {[}is the intention{]}). What I'm
referring to as `divesting hierarchy-of-actions' by definition has to be
in oneself.

Making illusions and noticing them makes them so overt that they
`become' bathos. (I don't regard that as ``alienation,'' however.)
Fundamentally anarchism (viewing that as being observation itself) is
necessitated. Besides \textsc{Defoe}, my other prose works (crossed with
being poetry) are: \textsc{The Return of Painting}, \textsc{The Pearl,
and Orion/A Trilogy} (reprinted by Talisman, 1997) and \textsc{The Front
Matter, Dead Souls} (Wesleyan, 1996). Also \textsc{Orchid Jetsam} (in
manuscript) and \textsc{R-hu} (forthcoming Atelos Press).

~

\hypertarget{gail-scott-bottoms-up}{%
\chapter{Gail Scott: Bottoms Up}\label{gail-scott-bottoms-up}}

\lettrine[lines=3, findent=0em, nindent=0.1em, lhang=0]{M}{ontreal, Quebec}:
I return to the site of an old novel. This will not always be a street
of failure, even if still playing the former glamour of the seedy.
French-language covers of country music are playing on the radio. Men
(sit) alone at tables. One with haggard bearded face. His lower jaw (is)
toothless. Teeth in dreams = sex. The low façades opposite parading
peeling tin cornices. An air of stagnation possibly verging on cliché. I
take some verbs out; then put them in again.

It is the hottest summer ever. Meteorologists are, with divine
hindsight, admitting global warming. To escape the heat, I go to a dark
afternoon theatre. Where a sassy high school troupes performing Büchners
\textsc{Wozeck}. The adolescents sweet songs juxtaposed on the gloom and
violence under, somehow dissolving earlier-decade interpretations
foregrounding worn-out existential despair. Crossing a parking lot
behind a smoked meat shop, I am thinking a story, to avoid eternally
returning (to keep it moving), might be structured thus: an ``older''
woman writer circulates at a ``fringe'' theatre festival of young
artists. She has lost her love. Feels discombobulated. Nothing. She
skirts the crowd asking several 20-year-olds: anything good? She will
then string a humiliated (but raunchy) love tale together, interlaying
bits of youthful sweetness and anger, representing to her in her present
state of emptiness, a more contemporaneously rapid conflation of
beginnings and endings. A word about narrative:

At another table, a man shakes his purple-shirted shoulders in the hot
bar air. To ``Caught With Your Pants Down.'' Drinking straight tequila.
The same table where X and I sat. Discussing why our favorite writers of
new and experimental fiction are not rich and famous. ``Too abject,'' we
say. In a bag at our feet, some cheaply extravagant dollhouse boudoir
furniture purchased for her little son. From the store with the saggy
awning opposite. Leaving, we forget it. Outside another man passing.
With several teeth missing. Teeth = action. But why the narrator (the
woman) endlessly remarking. Such toothless ghosts of subjects. Among
dozens of better-furbished faces. Is her ``local'' constructed out of
some kind of distancing or bias?* Or ``mere'' projection of loves
losses? I am thinking \textsc{not} abject the subject. Unless negatively
framed. At point of crossing space of (non)-reception between the work
and the world. When site (interior)s signed minor (inadequate). While
exterior=dominant. Yet, intrinsic and extrinsic syntax are not separate;
rapidly they mate. Should I be happy the lover having sex with no one.
If having only virtual sex with me? I take some verbs out. And put them
in again. A word about the waitress:

Dominatrix in clinging leopard-spotted polyester. Small breasts.
Ponytail. Fortyish. Saying (in québécois): UN verre de vin-maison ne te
ferait pas de mal, one glass of house wine will not ruin the stomach.
I.e. its terrible. For 75 cents more you can have a quarter liter.
Adding: you look en forme, terrific. WE tell each other--in guise of
complimenting--we haven't changed in years. I opting for the smallest
quantity of red vinegar. Now young woman passing by window. Huge breasts
encased in tight halter top. Some little cap over short hair.
And---upper incisors gone. Like the meteorologist, I admitting only to
coincidence of circumstances. Projecting accumulated lack on
unsuspecting bodies. (Waitress bringing unasked for quarter-liter).
Which bodies, in being thus sentenced, already frozen in past gesture or
alias. While environment drifting elsewhere. As in that restless white
river the lover (much later) wanting to attack with kayak. Drifting
eternally toward full-color horizons. The silver flashes indicating
minor notes of salmon struggling upstream. Sitting close beside her, I
encouragingly declaring such projection, battered by tergiversation of
context, likely best position. For contemporaneous narrator. Just as
lover putting hand under skirt cut on the bias to flare out attractively
at the bottom. And saying she still wanting something with ``us.'' A
story is something---you can put your teeth into:

But what might happen here? With (soon) a closed sign across the window.
Two blocks down. A polished glass and marble multi-story, multi-venue,
multi-media repertory theatre. Having replaced the crushed blue façade
of the smoky cozy café-cinéma where we used to see old or experimental
European French movies. Built by Québécois Soft Image inventor Daniel
Langlois. Were like that here: always oscillating between the
ultra-modern and the crumbling. The sound systems so efficient, one
feels, listening to the Buena Vista Social Club, as if in some club in
crumbling Havana. Though the octogenarian musicians having mostly ceased
performing on their blockaded island. Because ``nowhere'' to go. Until
American Ry Cooder choreographing their (global) return. I am thinking
of when Québec liberationists claiming solidarity with Cuba. Which
solidarity currently updated to meaning ``corporate opportunity.'' At
last, the Cubans playing Carnegie Hall. I am thinking how any narrative,
if fronted, marketed from economically preferential geographical site,
avoiding being troped as exotically remote, a remoteness that must
(paradoxically) scale global walls or languish. I am thinking of Cuban
sunsets, suffused towards evening with a little blood. Of the womans eye
growing humid, sentimental. When, in the movie, the elderly guitarist
exclaiming: ``There is nothing like a night of love.'' To keep it
moving, I take some verbs out and put them in again. A word about
identity:

Now the waitress in her leopardskin is leaning over the purple shirt.
Whos reading a sexy mens magazine. Pausing, he pulls a huge fluorescent
flashlight still wrapped in hard plastic cover from a bag. She points
out (still speaking {[}naturellement{]}, French): its no good without
batteries. Laughing. I can smell him. The conversation segueing to
European-French actress Catherine Deneuve. The waitress declaring in
drawn-out Québécois dipthongs the Europeans looking down on, Deneuve
``plate,'' boring: ``Always the same tone.'' Outside a friend passing by
window. Looking for a date. Very strong shoulders. Her long red hair
streaming out behind. I know she has latex gloves in her bag. Walking
naturally, she ``vogues.'' Causing her to get thrown out of womens
washrooms. For being ``wrong'' gender. I am thinking how certain minor
or unusual demeanors, if flaunted, getting read extrinsically as
unappealing or threatening. (In Canada, ``Québécois'' often = flamboyant
or risqué; coincidentally, same thing for ``queer''). This can be
devastating. Or interesting. Friend passing again, very feminine in
wrap-around skirt, curly red hair. Yet walking as if faking it. I am
thinking how narrator (other). IF not knowing how friend
herself-self-perceiving, able to create a relationship with reader
permitting permutation of both obvious and unspeakable. The body being
capable of gestures contrary to understanding. I am thinking about
narrator as double. About evening out subject/object weight at each end
of sentence. A word about affect:

Many empty tables. Two beefs in dark corner. Chairs fading outward.
Towards sunwashed waitress, purple shirt and harsh light of street.
Waitress, flashlight and me. A sign above our heads: ``Les crevettes de
MATANE arrivées!!! {Fraîches.} MATANE shrimp in!!!! Fresh. The drawn-out
Québécois fraîîîches savorously feminine and plural. The waitress and
purple shirt still complaining Deneuve lacking edge. Offering as
replacement, Québécoise country music star. Recently coming out
with''fabulous CD." Though ``star'' still having to work nightly in
bars. ``Very tough. If youre pregnant.'' Outside, young nomad in
artfully torn outfit and \$50-dollar-haircut. Asking for change in
anglicized French. They come here from the suburbs. Not to feel guilty.
In this huge dark room with stuffed animals on walls, I am thinking of
simulation in a diminshing French-speaking city, where everyone
considering themselves minority; yet in some way also faking it. I am
thinking of consequence for narrative, so often straining toward
``natural'' on site where continental popular culture in perpetual
transmigration from Germanic to Latinate. Of Mickey Mouse, Patsy Klein,
Sylvester Stone, moving lips in English. While saying something else
entirely. I am thinking of living on a Renaissance Ponte Vecchio: less
bridge between than Babel of echoes.** Whose interpretation requiring,
precisely, resistance to meaning. I am thinking of narration as
introjection. Many friends have dentures. A word about hybridity.

Two blocks up. Little stages. Blue-and-white fleur-de-lys flags for the
fête nationale. The Argentinean tango I missed. Somehow having erased
Québec ``patron'' Saint Jean Baptistes bleeding and miserable little
lamb. Featured in earlier decades national feast-day celebrations.* Fête
progressing north. Group doing samba. Québécois folk songs. Block
production of Molière farce. Free good food. Everyone speaking French
``with accent.'' Whole family dragging mahogany dining table out on
sidewalk. Now eating supper. Later, I, sleepless, listening to
thunderclaps on John Cage piece. To compensate for dryness. Leopardskin
waitress offering another glass of vinegar. Muscled guy in tight
sleeveless t-shirt and gleaming chain with cross. I am thinking of
narrative as opaque barrier. There is no limit between life and deathly
fascination. Viva Che! on wall. Now I must go before another cigarette.
Buy food. Take messages. If any. The disease is not under control.
Police car outside.

~

\hypertarget{bottom-notes}{%
\section*{Bottom Notes}\label{bottom-notes}}
\addcontentsline{toc}{section}{Bottom Notes}

This piece is dedicated to Carla Harryman.

It is in dialogue with other writers and other texts, including the
writing of, and conversations with, Dianne Chisholm, Robert Glück
(notably re: the narrator as bottom), Carla Harryman, Camille Roy, and
Sarah Schulman.

** Textual references are to Barrett Wattens \textsc{Total Syntax}, a
cover citation from \textsc{Sight} by Lyn Hejinian and Leslie Scalapino,
and Sherry Simons essay ``The Ponte Vecchio and the Comma of
Translation.''

Saint Jean Baptiste Day, June 24, is a national feast day in Québec.

~

\hypertarget{juliana-spahr-narrative}{%
\chapter{Juliana Spahr: Narrative}\label{juliana-spahr-narrative}}

\begin{verse}
The water is blackish, green, and dark.\\
It gathers from its\\
separated state, gathers from rain, gathers into stream.\\
It gathers in the mountain.\\
It gathers then travels, collects to become brackish.\\
Water travels, falls over a cliff, is falling.\\
Water is falling.\\
Falling onto rocks.\\
Because water falls it is.\\
Because water streams it is.\\
Because water collects into a pool it is narrative.\\
The pool is cold.\\
The pool is sheltered from sun by a cliff.\\
The pool is filled with rocks.\\
Water gathers over rocks, on rocks.\\
Moss gathers over rocks, on rocks.\\
Three arrive at this scene.\\
This scene where water is cold.\\
Where rocks are mossy.\\
When the three enter the pool their feet slip on rocks that are\\
mossy.\\
As they slip, they slip deeper into water.\\
Slip on mossy coldness.\\
Slip into water so cold it makes their chest close.\\
Slip deeply in.\\
Water gathers them, gathers them by their slipping.\\
Water covers them.\\
Because water is brackish it is narrative.\\
Because water is brackish they don’t open their mouths.\\
Because water is brackish they are closed in their immersion.\\
They do immerse.\\
But they don’t open as they swim to falling water.\\
They stand beneath falling water.\\
Water beats down on them.\\
Beats on their shoulders.\\
Beats on their heads. Beats on them.\\
They stand beneath water that is a roar of a falling.\\
They stand in roar, in narrative.\\
But water is brackish.\\
So they do not kiss.\\
Do not open their mouth beneath the waterfall.\\
They let brackish water fall over them.\\
Over their heads.\\
Over their lips.\\
Over their eyes.\\
Over their ears.\\
Over their hands which they hold up to feel the roar landing into\\
their palms.\\
Down their body.\\
Brackish water.\\
Closed bodies.\\
Water gathers them to this place, this narrative place.\\
Water covers them and they are covered with brackish water.\\
They do immerse.\\
They don’t open.\\
Immersion seals them off.\\
There is no open mouth.\\
No opening.\\
No exchanging.\\
Because it is brackish it is narrative falling over them.\\
It falls over.\\
It falls over them.\\
Sealed yet together.\\
They have come together with brackish.\\
They have let brackish wash over then even if they don’t let it\\
inside.\\
Because water falls it is narrative.\\
Because they are immersed it is narrative.\\
Because they love each other while separated it is narrative.\\
Because the rocks are slippery with green moss it is narrative.\\
Because they slip on the rocks it is narrative.\\
Because they slip deeper into water.\\
Because they allow the slip, are prepared for the slip and love its\\
immersion it is narrative. Because it involves love it is narrative.\\
Because it leaves them alone it is narrative.\\
\end{verse}

\hypertarget{anne-stone-why-write-narrative-when-you-could-be-the-next-guest-on-the-jerry-springer-show}{%
\chapter{Anne Stone: Why write narrative when you could be the next
guest on the Jerry Springer
show!}\label{anne-stone-why-write-narrative-when-you-could-be-the-next-guest-on-the-jerry-springer-show}}

\lettrine[lines=3, findent=0em, nindent=0.1em, lhang=0]{I}{ met someone}
whose tattoo had been sliced out of his arm, leaving a jagged scar in
place of the name of city he was born. I understood the urge, to inject
a place under your skin, only to have it removed like a second stomach
or a wart. Lately, I've been thinking a lot about place. If you ask me
where Im from, I can answer in a lot of different ways, none of which
would tell the whole story, or even manufacture the kind of sense I
would want to live with. I could show you the scar on my leg and point
to an intersection, or trace the scar angling up from the first phalange
on my index finger, pointing out the parking lot where the Streetsville
liquor store used to be. I could tell you about strip malls, or stock
yards, or franchised donut shops that remind me of Dresden the morning
after, draw you a map of the territory in the fine sprinkling of ash
collecting on a Formica counter. I could tell you about the derelict
cornfields that used to be where that donut shop is now, or the first
place I lived in St.-Henri, show you the abandoned granary that burned
to the ground seven years ago and has since been replaced by industrial
trailers. Or maybe the building I lived in itself, which is sealed off
now. I might point out the remains of a much older structure, a series
of barn-board sheds. I would go to one of those sheds each afternoon to
write. The pick-up-sticks anatomy of that place inspired local kids into
acts of arson. I dont know if it was my own narrative tendencies or the
steady supply of thirteen year olds with zippos, but it was the first
part of the building they chose to tear down.

I should tell you, before we go any further, that the scar is really on
my pointer finger. You should know this story is not particularly
reliable and, as a friend told me once while I was staring absently at a
menu, the map is not the territory.

Somehow, during the months I spent in that shed, writing what would
become \textsc{jacks: a gothic gospel}, the smell of the packing plant
my father took me to as a girl infiltrated the barn-boards. The plant
was on his delivery route, and I used to stop off at the stock yards to
pet the shit-stained flanks of cows as they were herded into slaughter.
St-Henri reeks of my childhood and the Streetsville liquor store is as
organized as a months supply of birth control pills, the bins of
national and international wines formatted on the principle of a 28-day
dial, the kind that comes complete with seven sugary pills to remind you
of the color red. When I think of all the places that have come to
dereliction, I understand myself in terms of the word tenacious. I know
it is only my thinking that makes it so. I dont know if they hand out
birth control pills to fourteen year olds anymore, the kind that remind
me of a flattened representation of a monthly agenda or a rolodex with a
daily regimen of business cards. I do know that the smell of those
stockyard cows do not belong in this essay. Nor does the way they kill
cows in the Eastern Townships, which is much more and much less
personal.

I also think its important to tell you, before I go any further, that I
am not a vegetarian.

On my friends farm, the breeders have names like Sophia or, out of
earshot, are called after her daughters-in-law. The cows she will
slaughter and wrap herself are named after beige-toned bureaucratic
dociers. When I first met the woman who owns this farm, I asked about
the cows and was given the names of the breeders and the bull, the
others implied by a slurred ellipsis or existing as the space-off to her
story. I persisted, asking about one of the cows shed passed over.
``T-2306? Or do you mean that one?'' she asked, pointing to Sophia,
``theyre sisters.'' A little later that day, she promised to save that
cows thigh-bone for my dog.

The cows exist in terms of place. The cows of the stockyard and the cows
of the townships farm are the same breed, indistinguishable, except for
their deformation in terms of place. Even their smells are different.
Sophia and T-2306 are twin sisters, born on the same day, from the same
uterus.

If I were to write about T-2306, I would want something of that tension
to come through. The format of the place infiltrating something which
might be appalled identity. Or perhaps by identity I mean body-as-place.
If I were a photographer, I might take a long-exposure photograph, the
kind that effaces her body if it implies movement, if her body-as-place
shifts in relation to geographical setting. The immobile objects and
landscape would exist in the picture as perfect solids. It would be an
unremarkable photo, except for the way you could see right through the
cow to everything behind it. Except for the way the picture would be
taken as daylight gave way to dark, and something of that lumen would
gather the hide together into the appearance of a bruise.

The length of exposure is an interesting way for me to begin to think
about length in narrative. As I am given to understand the technique, if
you lengthen the time that the shutter is open, the camera slowly tucks
whatever light there is into an objects surfaces. Ive seen night photos
taken with a long exposure, the light isnt reflected off the surfaces of
an object so much as collected there, staining the fabric. The light
appears to be a dirty liquid, seething up from under. The implications
fascinate me, the way that pulling sight long can deform what is there,
or ways of seeing it.

Thankfully, I am not bound to a single snapshot, because I want to
examine (and re-examine), lit (et relire), bed (and embed), these
surfaces, under-pinning the impoverished look of St-Henri, for instance,
with the reek of a west-end Toronto slaughter house, twenty years
distant in time, and 600 kilometers away. I want the space-off of the
shed I started writing \textsc{jacks} in to be that terrifying plant.
God: What is that? What is it? The dingy white towers next to the
stockyards are organized vertically, even though, as everybody knows,
cows are organized horizontally---they are grazers, right? So, between
the slaughterhouse and the packing plant some terrible deformation of
what it means to be a cow occurs, the axis of cowness undergoes a
grotesque skew. How do you even fit the word ``cow'' in a building that
tall? What could a building with elevators possibly have to do with a
cow? When I was a girl, I wasnt afraid of the men with electrical prods,
or the smell of shit and death. But I was terrified by those vertical
towers. I knew they embodied a cold and impersonal force. If I called
out to the men prodding the cows, for instance, I could get them to talk
to me, stop what they were doing for a while. I doubted my presence
could effect a pause when it came to anything in or about those
buildings. It wasnt so much that I didnt exist, as it was that in their
terms, I could not possibly exist. I lacked the relevance to insist my
own axis likewise undergo some miraculous Holstein shift from x to y.
The plant excluded all terms but those present in its formula. If I had
Goëdal for a companion on those walks through the slaughterhouse, they
might have made another kind of sense. As it is, the memory refuses my
attempts at sense-making, and continues to fascinate me twenty years
later. The shed I came to narrative in has been torn down, and the
vertical granary -- the one that existed at the place where Beaudoin
stutters to a stop at rue St-Ambroise, the one that once implied the
axial-shift undergone by a field of wheat -- burned to the ground one
winter night as I watched. These places have a grammar that I am only
now beginning to grapple with (or make-up).

I should probably confess that Ive apprehended a similar syntax in
unfamiliar restaurants. Once, a maitre-D unfolded my body like a table
cloth, the surface tension flattening and lengthening my body to the
skin of water, it delineated a lumen that precisely. Sensing some
terrifying grammar, I found that the only punctuation legible to me were
the round bottoms of oven-warmed plates pressing into my back like a
series of medieval cups. The kind 17th century doctors heat with candles
before placing in rows on either side of the spine, creating a series of
tiny vacuums to draw the suspected impurity up from under. I find
restaurants inscrutable. I mean, who would think to spell the word
``spent'' on a plate in a diagonal line of cutlery? But god, the relief
to discover there was a silent means to signal an impending departure.

What happens in the meantime? As your liminal-linen body hosts a series
of plates? And yes, physically, the body is a place that can, with some
effort, be occupied by one's self, or by others.

I am not at a restaurant. If I close my eyes, I lose my boundaries, the
air is tepid against my skin, no rough frisson strokes this surface. I
am not sensitive enough to maintain a secure sense of where I end and
something other begins. That requires a lepers perpetual gaze to the
extremities. If I were to close my eyes and begin typing by rote, you
would have to remind me: Where are my arms just now? my fingers\ldots{}

I am only sure of the tips. As they pad against the keys. One of my
fingers describes an unstable arch. Ive pulled a tendon with this
incessant tapping. It must be a very tiny tendon, the kind of thing you
might extract from the haunch of a bisected frog. Eyes closed, I am the
stoop of shoulders, the curve of a sore finger, and a series of
insensitive pads Ive developed as a result of my obsessive predilection
for narrative -- long narrative.

This is where I am today, and not at a restaurant packing plant. I am at
the place the scar on my index finger describes, an intersection between
place and the body. Writing narrative affords me a lateral way in which
to think about the body, about place, about the body in place, about the
body as place. In \textsc{Hush}, written very much like those
long-exposures I talked of earlier, the women exist in terms of
place---a narrative focus that, on the surface, would deny movement. But
writing against the vertical axis, implying it, even as I choose to
write on the horizontal plane, I can let a particular image accrete to
itself all sorts of other resonances---not just in the way a particular
image occurs in relation to her-body, Roses body, Loralies body, but how
it recurs, in relation to other events or objects or the place itself.

Plot is a term that I hear twice, I hear the English plot and the
français plotte, or cunt---so it is a gendering of narrative. I am not
about plot. But am about this other form of plotting in narrative. Of
telling a story that unfolds at the level of language, where the words I
use recur over the length of the narrative. The repetition becomes
skewed, and what it touches, likewise becomes communicable. It isnt
mood, or even sensibility Im aiming at, though thats what descriptive
words are employed for a lot of the time. No, what is here is as
convoluted as the vaginal tract of a snail, and though yes, snails
plottes are bundled relatively small, they do possess immense surfaces
to play on.

I should admit Ive forgotten the word for a word that is descriptive.

Where am I going next? I want you to imagine a picture taken with a very
long exposure. No.~I want you to imagine you think of the phrase
``glory-hole'' every time you hear the word exposure, and then I want
you to imagine a picture taken with a very long exposure. A picture
taken on a night at the circus. The temporal trace of the antics of
acrobats and Dahlmer-clowns rendered in fly-wheels of color. Are you
sufficiently distracted? Good. Because until a moment ago I was holding
a Polaroid in my hand. It was a very old picture, or else, someone has
gone to a lot of trouble to make it look that way. The edges are black
from fingering. If I was going to show you the picture, evidently, I
changed my mind, because while you were looking at the circus-shot, I
slipped the picture into my pocket and lit a cigarette.

I think you should know that I am familiar with the practical
applications of luting agents.

\hypertarget{michelle-tea-explain}{%
\chapter{Michelle Tea: Explain}\label{michelle-tea-explain}}

~

\lettrine[lines=3, findent=0em, nindent=0.1em, lhang=0]{I}{ get asked}
sometimes to defend the way I write. To explain how it's um writing I
guess. Because I'm not making stuff up so why do I think that anyone
wants to hear about my crappy life. Well then don't read it, punk! It's
all I can think about, my own life and your life too. I mean, what's
everybody doing? I really want to know.

What if fiction writers had to answer dumb questions like that. Like,
What makes you think we believe this crap? How the hell do you know what
that geisha was thinking, white boy? Or, What are you afraid of? Are you
trying to hide something? Who are you, where can I find you skulking
about in these pages?

Oh I don't want to divide and conquer the writers. I like fiction, I
just read \textsc{Jesus Saves} by Darcey Steinke and it nearly put me in
a coma. My stomach hurt so bad I had to put it down, rest it on the dash
of a van that was rumbling through woodsy mountain landscape, manly
land, hunting land, in a state famous for serial killers, and in the
pages of this made-up story a little girl is tied up captured and doomed
and I know it's true, I mean, these things happen. I read Jenn Banbury,
\textsc{Like A Hole in the Head}, and it made me want to fuck up on some
spectacular way, as opposed to my normal, low-key fuck ups. Those
stories were really good.

But I just think it is so exciting that life can be so meaningful, that
you can smear it on a page and it's really important. It's as much of a
story as any fiction, but it's a document too. It happened, so it's
history. I'm reading \textsc{The Unsinkable Bambi Lake}, it's before me
on the table right now beside a rack of blueberry muffins and a pile of
bills that can't seem to pay themselves. When I read Bambi's book I get
to hear all about how cool my city, San Francisco, was in the 70s. All
the stuff I missed, victorians stuffed with glamorous and dramatic G.A.
queens, glitter and dustbunnies in the corner, Bambi backstage in Berlin
giving Bowie shit for selling out way before that Dancing in the Streets
fiasco. Dangerous boys and hustlers, getting on estrogen, really just
living a life, her own. This question about `memoirs', I think: are they
asking about form or structure, or really, am I being asked if my life,
if Bambi's life and Eileen Myles' life, and Cookie Mueller's life, are
they asking if our lives matter? I think they are.

I'm lucky because I get to be friends with some of my favorite writers.
I sat with Eileen Myles at a coffee shop on Valencia and she told me
that Charles Bukowski gave her permission to write. He just claimed it,
his terrain. The horseraces and the bars, the night, the lamp that lit
up his lousy hotel. Why not him. A poet. Why not Eileen. I read her
book, \textsc{Chelsea Girls}, and it all clicked. ``Bread and Water,''
what is she doing in that story but being poor. Scrounging for money,
tallying debts to friends, having her period, drinking and not drinking.
It's a genius story and the whole book is full of them, electric. I was
zapped. Why not me. My poverty and the girls that don't love me and how
drunk I got the other night. How I was a prostitute. It seems to be
literature when guys write about it, its practically become a genre, men
writing about their transcendental trips to the cathouse, their orgasms
and revelations. Or men writing about women's lives in general. Straight
people writing about queers and white people writing about every other
race on the planet. The writing that I love, it's the Other telling the
part that got left out, the truth. Not only a writer and a historian,
but a spy. There's this awful copy shop near my house, I go there all
the time because I'm too lazy to walk up to Kinko's. The guys at this
place are such jerks. I had a bunch of my books and he said, Are Those
Your Books? Yeah. You wrote them? Yeah. He makes this suspicious little
scrunched-up face. Are You Sure? he asks. He means it. Looking at my
dirty fucked up hair and tattoos scrawled up my arms and whatever else
he saw. You Just Don't Look Like You Would Be A Writer. Yeah well keep
an eye out for yourself in my next novel, asshole.

I'm at work and I get a phone call and it's this girl. It's 2:30 in the
afternoon and she just woke up. I can hear birds chirping in the
background like she's in some fabulous jungle and I imagine light
slanting through fat glossy leaves and her all rumpled from sleep in a
slippery slip that sticks to her bones. Do You Have A Bird, I asked. I
Do. What Kind? A Cockatiel. I wonder what its name is but don't ask
because I want to keep it wild the way she is. She got arrested by the
cops for tagging Freedom Is Dying in a phone booth, they hauled her off
to jail. I feel this teeny obsession with her forming, and it feels a
lot like needing to write. A sort of zinging that grips you from the
inside. There's a place for her in my life right now and so I need to
write her down. The other night we were both high in that club that's so
red, a deep low red that made everything red, my dress and torn
fishnets, her shiny legs she kept twisting around, unfurling like a
mermaid. The red beer label curling on the bottle that she begged off
the bartender for free. I want to tell you all about her fantastic
strangeness, the sexy awkward way she walks like she's going to wipe
out, how she's paranoid, the excellent purity she sort of radiates out
from her eyes all black and smudgy, her dirty hair. I could tell you
about her jagged poems and drawings and the tattoos on her fingers and
the sad Johnny Thunders song on her answering machine. I could tell you
her whole story as I know it, and I want to, but it's not mine to tell.

And then she calls me the next day kind of wiggy, says I Got Tripped Out
Thinking I Was Going To End Up In One Of Your Stories. All the time
someone gets mad at me or else I worry they will. Where does my life
stop, where does it slide into someone else's story and then it's their
job to tell it. My ex-girlfriend is going to kill me when she finds out
I wrote about that time she flipped out and was banging her head against
my bedroom wall and I had to kick her out and she told me to fuck off. I
mean, that's personal. And can I mention that thing about her parents?
It seems so crucial to understanding who she is as a character. But
she's not a character, she's a person. And she had this affect on me.
Can I talk about it? When I went to Ali's house and she was cutting her
arms and there was blood on her yellow sheets and there was the razor,
kind of crusty on her dresser. I said Ali, I Feel Like That's Part Of
The Story, Can I Write That? She said, Tell The Truth, Michelle. But
then, Ali's a writer too.

Sometimes, maybe even right now, I don't want to write. It seems so
boring. There are so many rules and I made them all myself: not here in
the house, in a cafe or a bar. You have to take a walk, and think about
it while your walking, try to go back there, feel all hurt or fucked up
and kind of lost, feel really in love, you were just with the person,
you smell like her, remember how you ate a burrito afterwards, what it
felt like to crave the steak. You have to drink something while you
write, either coffee or beer, or sangria. If you're getting a late start
and the sun is down already then it should be a bar. You don't want the
place closing up around you when you're just hitting your stride, you're
going to stay up late immersed in all this life. Also you have to smoke
continuously, so pick a place that will let you do that. Maybe you
should get on a Greyhound and go somewhere else a place with romantic
terrain where you can be with your cigarettes. I think, maybe I should
just be in a band. Play some drums. Maybe I should take pictures, that's
as true as any story I've got to tell. But it's so expensive. Once I
didn't have any money, and I hardly knew anybody and I was so scared I'd
have to live out on the street. And I thought, it's ok I can still
write. Paper and pens are easy, they're cheap or easily stolen, you
could probably get people to give them to you. And if something really
bad like that ever happened, and whenever something does happen that
really sucks, it's hard, I think well ok I'll write about it. It's so
consoling and so redeeming. When I have a friend who doesn't write and I
see her having a hard time, like she's lonely or on too many drugs or
something fucked has happened I think oh I wish she could write it.
Write it all out. She could be an enormous angel. Probably she is
anyway, but then she could see it all lit up on the page, swimming
there, shining back up at her, good lighting for the movie that is her
life. She could write the soundtrack. She could bask in herself, it's
the greatest. Writing is the greatest, and writing the truth in
particular. Your little slice of it. You've got that, and you've got
your body, and I think that's it.

~

\hypertarget{quotes-from-the-editors}{%
\chapter{Quotes from the Editors}\label{quotes-from-the-editors}}

\hypertarget{quotes-from-bob-gluxfcck}{%
\section*{Quotes from Bob Glück}\label{quotes-from-bob-gluxfcck}}
\addcontentsline{toc}{section}{Quotes from Bob Glück}

In guise of an ``editorial,'' we gave ourselves the task of contributing
quotes that seemed exemplary or important to us as writers at this time.
Informed by science (Burke's notion of coarse-graining, implying the
production of prose as pattern), by references Biblical, Benjaminian,
post structuralist, post-Marxist, Oulipo-ian, etc., these cut-outs,
compounded, comprise our narrative about Narrativity. For now. Signed
Mary. Camille. Bob. Gail. ~

\hypertarget{from-antonin-artaud-anthology-ed.-jack-hirschman}{%
\section*{\texorpdfstring{From \emph{Antonin Artaud Anthology}, ed.~Jack
Hirschman}{From Antonin Artaud Anthology, ed.~Jack Hirschman}}\label{from-antonin-artaud-anthology-ed.-jack-hirschman}}
\addcontentsline{toc}{section}{From \emph{Antonin Artaud Anthology},
ed.~Jack Hirschman}

From this pain rooted in me like a wedge, at the center of my purest
reality, at the point of my sensibility where the two worlds of body and
mind are joined, I have learned to distract myself by the effect of a
false suggestion.

For in the space of that minute the illumination of a lie can last, I
manufacture a notion of escape; I rush off in any wrong direction my
blood takes. I close the eyes of my intelligence and open my mouth to
the speech of the unspoken; I give myself the illusion of a system whose
vocabulary escapes me. But from this minute of error there remains the
feeling that I have snatched something real from the unknown. I believe
in spontaneous bewitchments. It is impossible that I shall not some day
discover a truth somewhere on the routes my blood carries me.

\hypertarget{from-the-gospel-according-to-thomas}{%
\section*{\texorpdfstring{From \emph{The Gospel According to
Thomas}}{From The Gospel According to Thomas}}\label{from-the-gospel-according-to-thomas}}
\addcontentsline{toc}{section}{From \emph{The Gospel According to
Thomas}}

Jesus saw some little ones nursing. He said to his disciples, What these
little ones who are nursing resemble is those who enter the kingdom.
They said to him, So shall we enter the kingdom by being little ones?
Jesus said to them, When you make the two one and make the inside like
the outside and the outside like the inside and the above like the
below, and that you might make the male and the female be one and the
same, so that the male might not be male nor the female be female, when
you make eyes in place of an eye and a hand in the place of a hand and a
foot in place of a foot, an image in place of an image then you will
enter the kingdom.

\hypertarget{from-foucaultblanchot-maurice-blanchot-the-thought-from-the-outside-by-michel-foucault}{%
\section*{\texorpdfstring{From \emph{Foucault/Blanchot, Maurice
Blanchot}: "The Thought from the Outside" by Michel
Foucault}{From Foucault/Blanchot, Maurice Blanchot: "The Thought from the Outside" by Michel Foucault}}\label{from-foucaultblanchot-maurice-blanchot-the-thought-from-the-outside-by-michel-foucault}}
\addcontentsline{toc}{section}{From \emph{Foucault/Blanchot, Maurice
Blanchot}: "The Thought from the Outside" by Michel Foucault}

If the only site for language is indeed the solitary sovereignty of "I
speak" then in principle nothing can limit it---not the one to whom it
is addressed, not the truth of what it says, not the values or systems
of representation it utilizes. In short, it is no longer discourse and
the communication of meaning, but a spreading forth of language in its
raw state, an unfolding of pure exteriority. And the subject that speaks
is less the responsible agent of a discourse (what holds it, what uses
it to assert and judge, what sometimes represents itself by means of a
grammatical form designed to have that effect) than a non-existence in
whose emptiness the unending outpouring of language uninterruptedly
continues.

\hypertarget{from-the-madness-of-the-day-by-maurice-blanchot-trans.-lydia-davis}{%
\section*{\texorpdfstring{From \textsc{The Madness of the Day}, by
Maurice Blanchot, trans. Lydia
Davis}{From The Madness of the Day, by Maurice Blanchot, trans. Lydia Davis}}\label{from-the-madness-of-the-day-by-maurice-blanchot-trans.-lydia-davis}}
\addcontentsline{toc}{section}{From \textsc{The Madness of the Day}, by
Maurice Blanchot, trans. Lydia Davis}

I am not learned; I am not ignorant. I have known joys. That is saying
too little: I am alive, and this life gives me the greatest pleasure.
And what about death? When I die (perhaps any minute now), I will feel
immense pleasure. I am not talking about the foretaste of death, which
is stale and often disagreeable. Suffering dulls the senses. But this is
the remarkable truth, and I am sure of it: I experience boundless
pleasure in living, and I will take boundless satisfaction in dying.

I have wandered: I have gone from place to place. I have stayed in one
place, lived in a single room. I have been poor, then richer, then
poorer than many people. As a child I had great passions, and everything
I wanted was given to me. My childhood has disappeared, my youth his
behind me. It doesn't matter. I am happy about what has been. I am
pleased by what is, and what is to come suits me well enough.

Is my life better than other peoples lives? Perhaps. I have a roof over
my head and many do not. I do not have leprosy, I am not blind, I see
the world---what extraordinary happiness! I see this day, and outside it
there is nothing. Who could take that away from me? And when this day
fades, I will fade along with it---a thought, a certainty, that
enraptures me.

I have loved people. I have lost them. I went mad when that blow struck
me, because it is hell. But there was no witness to my madness, my
frenzy was not evident: only my innermost being was mad. Sometimes I
became enraged. People would say to me, Why are you so calm? But I was
scorched from head to foot; at night I would run through the streets and
howl; during the day I would work calmly.

\hypertarget{from-crossing-into-poland-by-isaac-babel-trans.-walter-morison-the-collected-stories.}{%
\section*{\texorpdfstring{From "Crossing into Poland" by Isaac Babel
trans. Walter Morison, \textsc{The Collected
Stories}.}{From "Crossing into Poland" by Isaac Babel trans. Walter Morison, The Collected Stories.}}\label{from-crossing-into-poland-by-isaac-babel-trans.-walter-morison-the-collected-stories.}}
\addcontentsline{toc}{section}{From "Crossing into Poland" by Isaac
Babel trans. Walter Morison, \textsc{The Collected Stories}.}

Fields flowered around us, crimson with poppies; a noontide breeze
played in the yellowing rye; on the horizon virginal buckwheat rose like
the wall of a distant monastery. The Volyns peaceful stream moved away
from us in sinuous curves and I was lost in the pearly haze of the birch
groves; crawling between flowery slopes, it wound weary arms through a
wilderness of hops. The orange sun rolled down the sky like a lopped-off
head, and mild light glowed from the cloud gorges. The standards of the
sunset flew above our heads. Into the cool of evening dripped the smell
of yesterdays blood, of slaughtered horses. The blackened Zbruch roared,
twisting itself into foamy knots at the falls. The bridges were down,
and we waded across the river. On the waves rested a majestic moon. The
horses were into the cruppers, and the noisy torrent gurgled among
hundreds of horses legs. Somebody sank, loudly defaming the Mother of
God. The river was dotted with the square black patches of the wagons,
and was full of confused sounds, of whistling and singing, that rose
above the gleaming hollows, the serpentine trails of the moon.

\hypertarget{quotes-from-camille-roy}{%
\section*{Quotes from Camille Roy}\label{quotes-from-camille-roy}}
\addcontentsline{toc}{section}{Quotes from Camille Roy}

\hypertarget{excerpts-from-the-storyteller-by-walter-benjamin}{%
\chapter{Excerpts from "The Storyteller," by Walter
Benjamin}\label{excerpts-from-the-storyteller-by-walter-benjamin}}

The art of storytelling is coming to an end. Less and less frequently do
we encounter people with the ability to tell a tale properly. More and
more often there is embarrassment all around when the wish to hear a
story is expressed. It is as if something that seemed inalienable to us,
the securest among our possessions, were taken from us: the ability to
exchange experiences.\footnote{Benjamin, Walter. \textsc{Illuminations}.
  (Shocken Books, © 1968).}

One reason for this phenomenon is obvious: experience has fallen in
value\ldots{} {[}The nature of every real story is that{]} it contains,
openly or covertly, something useful. The usefulness may, in one case,
consist in a moral; in another, in some practical advice, in a third, in
a proverb or maxim. In every case the storyteller is a man who has
counsel for his readers. But if today "having counsel" is beginning to
have an old-fashioned ring, this is because the communicability of
experience is decreasing. In consequence we have no counsel either for
ourselves or for others.

Every morning brings us the news of the globe, and yet we are poor in
noteworthy stories. This is because no event any longer comes to us
without already being shot through with explanation. In other words, by
now almost nothing that happens benefits storytelling; almost everything
benefits information. Actually, it is half the art of storytelling to
keep a story from explanation as one reproduces it.

The value of information does not survive the moment in which it was
new. It lives only at that moment; it has to surrender to it completely
and explain itself to it without losing any time. A story is different.
It does not expend itself. It preserves and concentrates its strength
and is capable of releasing it even after a long time.

\emph{Boredom is the dream bird that hatches the egg of experience.}
Today people live in rooms that have never been touched by death, dry
dwellers of eternity, and when their end approaches they are stowed away
in sanatoria or hospitals by their heirs. It is, however, characteristic
that not only a man's knowledge or wisdom, but above all his real
life---and this is the stuff that stories are made of---first assumes
transmissible form at the moment of his death. Just as a sequence of
images is set in motion inside a man as his life comes to an
end---unfolding the views of himself under which he has encountered
himself without being aware of it---suddenly in his expressions and
looks the unforgettable emerges and imparts to everything that concerned
him that authority which even the poorest wretch in dying possesses for
the living around him. This authority is at the very source of the
story.

Death is the sanction of everything that the storyteller can tell. He
has borrowed his authority from death.

{[}{[}I love this essay, although it doesn't address what seems to me to
be a distortion in the general view of the narrative process:
story\textsc{telling}. In addition to the story which is told, there is
the story withheld. Attraction, narrative suspense, emanates from the
latter. The story withheld has its own (buried) structure and tension.
The reader marries the hidden story, not the visible one. That is why we
continue to read a story that inspires dread. An attachment has been
formed---surreptitiously. It operates like a craving, which is how I
recognize Benjamin's reference to death. The craving for the unknown, as
it is held by the act of storytelling, points to death (or its
substitute, the body). (Roy){]}{]}

\hypertarget{excerpts-from-fame-by-robert-gluck}{%
\section*{Excerpts from "Fame," by Robert
Gluck}\label{excerpts-from-fame-by-robert-gluck}}
\addcontentsline{toc}{section}{Excerpts from "Fame," by Robert Gluck}

There are contradictory reasons when I use "real" people in my work, and
the contradictions comfort me. Any literary practice should, I think,
derive from contradictory sources and motives\ldots{} I name names to
evoke the already-known, to make writing co-extend with the world and
history, and to examine the fiction of personality, as well as the
fiction of the word\ldots{}. I propose self-community-story as a tonic
for the loss of human scale; by naming self-community-story I
participate in their disintegration, their progress from invisibility to
something to be named and manipulated---to be reintegrated later in a
new context, in a third term that history must provide.\footnote{Glück,
  Robert. "Fame." \textsc{Poetics Journal}, Number 10, June 1998.}

Prolonged scrutiny can become an expenditure of self, a potlatch of
self. I've come to experience the unreeling of interiority and sexual
disclosure as such a loss, and also part of a historical trajectory.
It's a writing activity that privileges the aggression of naming---an
ongoing colonization of self into one's own language. Once something is
named, you are in relation to it. Name the disease to cure it.

We want to see a story as we see other representations: being hiding
behind appearance---that is, hiding and revealing the body. But the use
of real names {[}in my work{]} reorders connections and disjunctions. I
do mean fragmentation. I don't want to make the predictable distinction
between story and fragmented writing. Naming names creates an open form
that co-extends with the world. In a postmodern switch, it applies the
open form of modernism to content by putting quote marks around the
entire story, turning the story into a fragment, an example of a story.
The story floats---as gossip does---between the lives of the people who
are its characters, and the lives of its readers (in that thorny field
of reader/writer dynamics). The problem of figure and ground becomes a
social one, and some of what is existential in the content is subtracted
and reintegrated in the relation between reader and writer.

\ldots{} Using real names provides a relation between the writer and
myself that carries some risk, like performance art. What I witness is
always the same: any story hides and then reveals the body.

{[}{[}Two aspects of narrative writing have been subjected to
devastating critiques: the linearity of narrative, and the credibility
of the narrator. (The popularity of the memoir is in part due to the
diminished credibility of anything other than an autobiographical
voice.) The critical and political forces which produced these critiques
have pushed experimental narrative writers to the boundary where form
collapses into experience. The searching re-examination of narrative
premises becomes part of the practice. (Roy){]}{]}

\hypertarget{excerpts-from-paragraphs-blowing-on-a-line-by-gail-scott}{%
\section*{Excerpts from "Paragraphs Blowing On A Line," by Gail
Scott}\label{excerpts-from-paragraphs-blowing-on-a-line-by-gail-scott}}
\addcontentsline{toc}{section}{Excerpts from "Paragraphs Blowing On A
Line," by Gail Scott}

\hypertarget{entry-11}{%
\subsection*{Entry 11:}\label{entry-11}}
\addcontentsline{toc}{subsection}{Entry 11:}

She sees the present as a great tear, a great rip in the surface of
things: the gap of which is at first impossible for her to move across
(*\ldots{}does not thinking seek forever to clamp a dressing over the
gaping and violent wound of the impossibility of
thought\ldots{}\footnote{Gayatri Chakravorty Spivak, Translator's
  Preface to Jacques Derrida, On Grammatology (Johns Hopkins University
  Press, 1976).} It is awareness of this that makes linear narrative
impossible. A virtually tragic awareness, for the writing ``I," which
modernists began to "deconstruct" as they recognized that gap (accepted
to give up power) has, in her case, been''deconstructed" a priori by
social conditions. She is not only split between the self and the
"real", but also within the self\ldots{}\footnote{Scott, Gail.
  \textsc{Spaces Like Stairs: Essays by Gail Scott}. (Toronto, Ontario:
  The Women's Press, ©1989). Gail adds: ``This diary was written over
  the years I worked on my novel, \emph{Heroine}.''}

So the question becomes how to write across the almost \ldots{}
hysterical \ldots{} overdetermination of her gaps...

Maybe my resistance to the narrative conventions of the novel have to do
with what I think of as its Protestant qualities: its earnest
representation of the "real," its greed for action, its preference for
the concrete over the philosophical.

A compromise might be possible: structuring the story by means of the
fluctuations in her \ldots{} (dare I?) hysterical voice. The use of the
voice invoking a poetic meaning in excess of the sentence\ldots{}.

Blows: violent and unusual interruptions in her narrative texture-which
texture the writing is hopefully starting to discover. More and more I
intuit that it has to do with starting from a negative point: a crushed
ego that doesn't see its boundaries.

\hypertarget{entry-25}{%
\subsection*{Entry 25}\label{entry-25}}
\addcontentsline{toc}{subsection}{Entry 25}

During this whole process I have been caught between the need to
accomplish this thing, which is both to posit some new kind of subject
and to have the sense of being a subject myself; and yet to resist
repeating, by creating a ``feminist narrative," what Barthes calls
\ldots{} the staging of a new "father," a new hierarchy of acceptable
concepts. (\emph{Every narrative (every unveiling of the truth) is a
staging of the \ldots{} father.}) Maybe what my heroine discovers in
trying to write her novel is only that the novel doesn't suit the
(diffuse?) women's way of seeing things. Where there is closure (firm
conclusions) in''straight writing" there are spaces, questions in hers.
Even her anecdotes point to other possible representations, leave
themselves open for reader intervention.

Still, given that my firm and conscious intention has been to counter
(patriarchal) ideology in this process, the poststructuralist recipe for
taking apart everything from the sentence to the author won't entirely
do. Presumably I have to at least propose some other direction: language
slips all around us. One's response to that is a question of the
relationship between writing (what it comprises of consciousness) and
body. Again to quote Barthes: \emph{The text needs its shadow. This
shadow is a} bit \emph{of ideology, a} bit \emph{of representation, a}
bit \emph{of subject: ghosts, pockets, traces, necessary clouds.
Subversion must produce its own chiaroscuro.}

{[}{[}It is generally true that reviews of experimental narrative do not
have a critical vocabulary for what the writer is doing. The goal seems
to be attaching a label to the content without ``giving it away''.
Narrative practice, as a field of critical inquiry, is in a practically
fetal state. Theory borrowed from poetics is not adequate for describing
or investigating how narratives work, what is at stake in narrative, the
social field narratives both operate within and set up. The
interventions into the social field which experimental narratives
perform are usually in practice imperceptible, indistinguishable from
background noise, unless the critical piece is written by the writer
herself. If narrative writers want the deepest issues in their work to
be visible, they must supply the critical ideas that drive their
practice. This journal will hopefully be a place that supports the
development of that discussion. (Roy){]}{]}

\hypertarget{excerpts-from-interview-with-carla-harryman-by-chris-tysh}{%
\section*{Excerpts from "Interview with Carla Harryman," by Chris
Tysh}\label{excerpts-from-interview-with-carla-harryman-by-chris-tysh}}
\addcontentsline{toc}{section}{Excerpts from "Interview with Carla
Harryman," by Chris Tysh}

Tysh: ``The I expands," you wrote in "Realism." In this age of
decentered subjectivities, how do you deploy the narrative machinery of
the''I" and the various sensibilities it is associated with? Could you
discuss the techniques with which you defamiliarize your use of the
first-person pronoun and in general address the issue of the speaking
subject in your prose. Who speaks?\footnote{Carla Harryman. "Interview
  by Chris Tysh." Number 10, June 1998.}

Harryman: My answer to this question is complex and leads to
difficulties because I associate the I with power, control, destruction,
gender, and whatever might transpire or exhibit itself between the poles
of animation and passivity. Mine is an unwieldy set of interests, less
easy to organize into the kinds of discussion that the problems of
``confession" or "authenticity" lend themselves to; even so, as you will
see, I am concerned with the relation of the''self" to the limits of
knowledge.

When I think about the I, what first comes to mind is omniscience in
novels. (My understanding of literary device and the problem of
authorship comes as much from reading novels as poetry. ) I will use
Balzac as an example, because he has come up in my writing from time to
time. Balzac wanted to account for all the history of France as it
related to the ruptures of the bourgeois revolution. The expanding I
here would be Balzac-as-author-of-France, i.e., Balzac as authorial
agent/persona, rather than a deployment within the narrative text, which
in Balzac is driven by third-person narration. I am interested in
uncovering the hidden, expansive I. What happens when one concretizes
it, making it evident in the writing, is carnivalesque. My writing
parades various self-conscious positions in regard to constructs such as
authorship, omniscience, and agency. It chooses to reflect upon the
positions it takes and the relations it draws\ldots{}

The deployment of the I through the autobiographical, the fictive,
abstraction, argumentation is involved in the breaking up of unities.
The I, as I think you are suggesting, does not expand but disperses. If
the text is the social body, the I is mobile within it. If I is that
which is not the body, it attempts to identify with the body's
fleshiness. The I can be a code for seeking pleasure and examining
freedom. Then it becomes troubled by its demands, for the relation of
the self to the world is always troubling; it is never adequate and the
I in its freedom is too willing to give up civility, a cherished element
of comedy.

\hypertarget{quotes-from-gail-scott}{%
\section*{Quotes from Gail Scott}\label{quotes-from-gail-scott}}
\addcontentsline{toc}{section}{Quotes from Gail Scott}

\hypertarget{from-eros-the-bittersweet-by-anne-carson}{%
\subsection*{\texorpdfstring{From \emph{Eros The Bittersweet}, by Anne
Carson}{From Eros The Bittersweet, by Anne Carson}}\label{from-eros-the-bittersweet-by-anne-carson}}
\addcontentsline{toc}{subsection}{From \emph{Eros The Bittersweet}, by
Anne Carson}

``\ldots{}Whenever passion seems within reach, \emph{aidos} falls like a
veil\ldots{}. This \emph{aidos} is the archaic ethic of 'shamefastness'.
. . . Ahprodite is the divinity in charge of the perversities of
\emph{aidos} within the novel. She is the chief designer and chief
subverter of the story's changing triangles, both patron and enemy,
inspiring with a passion strong enough to resist all the temptations
that she herself proceeds to hurl against it\ldots{}. Aphrodite's role
in novels is an ambivalent, not to say paradoxical, one like the role of
Eros in archaic poetry.'' Eros, p.~79-80

{[}Carson's Greeks lined up Eros versus 'shamefastness,' which I
contemporaneously interpret not as 'chastity' but as 'cover-up'; class
has entered discourse: shame is private and social (Scott){]}

\hypertarget{from-the-shell-and-the-kernel-vol.-1-by-nicholas-abraham-and-maria-torok}{%
\subsection*{\texorpdfstring{From \emph{The Shell and the Kernel} (Vol.
1) by Nicholas Abraham and Maria
Torok}{From The Shell and the Kernel (Vol. 1) by Nicholas Abraham and Maria Torok}}\label{from-the-shell-and-the-kernel-vol.-1-by-nicholas-abraham-and-maria-torok}}
\addcontentsline{toc}{subsection}{From \emph{The Shell and the Kernel}
(Vol. 1) by Nicholas Abraham and Maria Torok}

\ldots{}What haunts are not the dead, but the gaps left within us by the
secrets of others\ldots{} \ldots{}It is reasonable to maintain that the
'phantom effect' progressively fades during its transmission from one
generation to the next and that, finally, it disappears. Yet, this is
not at all the case when shared or complementary phantoms find a way of
being established as social practices along the lines of \emph{staged
words}\ldots{} To stage a word, whether metaphorically, as an alloseme,
or as a cryptonym-constitutes an attempt at exorcism, an attempt, that
is, to relieve the unconscious by placing the effects of the phantom in
the social realm. \ldots{}The phantom's periodic and compulsive return
lies beyond the scope of symptom formation in the sense of a return of
the repressed; it works like a ventriloquist, like a stranger within the
subject's own mental topography. The imaginings issuing from the
presence of a stranger have nothing to do with fantasy strictly
speaking. They neither preserve a topgraphical status quo nor announce a
shift in it. Instead, by their gratuitousness in relation to the
subject, they create the impression of surrealistic flights of fancy or
of oulipo-like verbal feats.

{[}I like to think of the phantom dressed up in a cloak studded with the
latest consumer items, on his way to the zoo. I like the play between
contiguity and disjunction implied by attempting to converse with the
unspeakable. This makes strange sentences that swing "both" ways,
dialoguing with both the uncanny and the social: a rhetorical edge is
added to poetic thinking. (Scott){]}

\hypertarget{from-zoo-or-letters-not-about-love-by-viktor-shklovsky}{%
\subsection*{\texorpdfstring{From \emph{Zoo or Letters Not About Love}
by Viktor
Shklovsky}{From Zoo or Letters Not About Love by Viktor Shklovsky}}\label{from-zoo-or-letters-not-about-love-by-viktor-shklovsky}}
\addcontentsline{toc}{subsection}{From \emph{Zoo or Letters Not About
Love} by Viktor Shklovsky}

\ldots{}She loved me like her own daughter. When I was two months old,
she fed me cabbage soups and once she managed to poison me by gorging
herself on the pits from the cherry preserves that were being made at
our summer cottage. When I grew up, she came to see me, always with
presents; she remained standing and spoke to me in the formal way; then,
when everyone left, she sat down to drink tea with me and used the
familiar form . When I became an adult, I began to understand her
cheerful disposition: `My mistress lives with another woman; it's beyond
me-just like nuns!' And she'd roar with laughter-\/-\ldots{}She had a
special smell, like her wooden trunk when she lifted the lid: calico and
apples. A tlted nose, knowing eyes\ldots{} "One time she was working in
a very rich household. The house was robbed\ldots{} Papa asked her: 'And
where were you when the house was robbed?' 'At the Novo-Devichy Convent
visiting a nun.' 'Well, then, tell them and they'll let you go.' 'What?
And get a nun mixed up in such business, master!' So she didn't say a
word, but sat in prison until the thieves were found and she was
released. Then, after the revolution, Mama tried to persuade her to go
vote, but she said that after that business with the silver spoons, she
wouldn't go near the police station for love or money.

{[}This passage, though "authored" by Sklovsky, is actually by Elsa
Triolet, Shklovsky's beloved at the time. He integrated her letters into
his own fiction. I had intended to quote Shklovsky; whose manner of
stringing together seemingly cut-up prose lines to effect narrative
'continuity' is exemplary. Not only shame, but also authorship, get
partly buried here. (Scott){]}

\hypertarget{from-lyn-hejinians-introduction-to-jena-osmans-character}{%
\subsection*{\texorpdfstring{From Lyn Hejinian's Introduction to Jena
Osman's
\emph{Character}}{From Lyn Hejinian's Introduction to Jena Osman's Character}}\label{from-lyn-hejinians-introduction-to-jena-osmans-character}}
\addcontentsline{toc}{subsection}{From Lyn Hejinian's Introduction to
Jena Osman's \emph{Character}}

Performance requires the person who is the actor (i.e.~already a
character) to be in character, and this, in turn, cannot occur without
performance. This produces not a tautology (``performance requires
performance'') but a bifurcation-character occurring as a performance in
and of itself; or, as Osman says (in ``The Figural Cabinet''), "the
taking place is double." Such double-mindedness, which is to say this
'self'conscious, is central to Brechtian theater; it forms the basis on
which alterities can come into (the) play. It is characteristic, too, of
everyone who is forced into identity-consciousness, as W.E.B. Du Bois
pointed out long ago in The Souls of Black Folk, a work which itself has
been described (by Houston Baker) as a ``cultural performance." More
recently, the notion that''identity" is itself performative has been
greatly elaborated by theorists such as Judith Butler. Identity-the
sense of difference that belongs to one's self-sameness-is itself
``character," something that "takes place.'' This means both that it is
an occurrence or event-it happens-and that it involves taking up a
position among others (as might be indicated by stage directions:
\emph{Enter Woman, wearing a hat,} etc.)

\hypertarget{from-not-all-halfbreed-mothers-in-i-knew-two-muxe9tis-women-by-greg-scofield}{%
\subsection*{\texorpdfstring{From "Not All Halfbreed Mothers" in \emph{I
Knew Two Métis Women}, by Greg
Scofield}{From "Not All Halfbreed Mothers" in I Knew Two Métis Women, by Greg Scofield}}\label{from-not-all-halfbreed-mothers-in-i-knew-two-muxe9tis-women-by-greg-scofield}}
\addcontentsline{toc}{subsection}{From "Not All Halfbreed Mothers" in
\emph{I Knew Two Métis Women}, by Greg Scofield}

\begin{verse}
Not all halfbreed mothers\\
wear cowboy shirts or hats,\\
flowers behind their ears or\\
moccasins sent from up north.\\
\end{verse}

\begin{verse}
Not all halfbreed mothers\\
read \emph{The Star}, \emph{The Enquirer},\\
\emph{The Tibetan Book of the Dead}\\
or Edgar Cayce,\\
know the Lady of Shalott\\
like she was a best friend\\
or sister.\\
\end{verse}

\ldots{}

\begin{verse}
Not all halfbreed mothers\\
speak like a dictionary\\
or Cree hymn book,\\
\end{verse}

\ldots{}

\begin{verse}
Not all halfbreed mothers\\
\end{verse}

\begin{verse}
drink\\
\end{verse}

\begin{verse}
red rose, blue ribbon,\\
Kelowna Red, Labatt’s Blue.\\
\end{verse}

\hypertarget{quotes-from-mary-burger}{%
\section*{Quotes from Mary Burger}\label{quotes-from-mary-burger}}
\addcontentsline{toc}{section}{Quotes from Mary Burger}

\hypertarget{from-tristram-shandy-by-laurence-sterne-first-published-1760}{%
\subsection*{\texorpdfstring{From \emph{Tristram Shandy}, by Laurence
Sterne, first published
1760}{From Tristram Shandy, by Laurence Sterne, first published 1760}}\label{from-tristram-shandy-by-laurence-sterne-first-published-1760}}
\addcontentsline{toc}{subsection}{From \emph{Tristram Shandy}, by
Laurence Sterne, first published 1760}

\hypertarget{chapter-v}{%
\subsubsection*{Chapter V}\label{chapter-v}}
\addcontentsline{toc}{subsubsection}{Chapter V}

On the fifth day of November, 1718, which to the era fixed on was as
near nine calendar months as any husband could in reason have expected,
--- was I, Tristram Shandy, Gentleman, brought forth into this scurvy
and disastrous world of ours. --- I wish I had been born on the Moon, or
in any of the planets (except Jupiter or Saturn, because I never could
bear cold weather), for it could not well have fared worse with me in
any of them (though I will not answer for Venus) than it has in this
vile, dirty planet of ours, --- which o' my conscience, with reverence
be it spoken, I take to be made up of the shreds and clippings of the
rest; --- not but the planet is well enough, provided a man could be
born in it to a great title or to a great estate; or could anyhow
contrive to be called up to public charges, and employments of dignity
or power; --- but that is not my case; --- and therefore every man will
speak of the fair as his own market has gone in it; --- for which cause
I affirm it over again to be one of the vilest worlds that ever was
made; --- for I can truly say that from the first hour I drew my breath
in it, to this, that I can now scarce draw it at all, for an asthma I
got in skating against the wind in Flanders, --- I have been the
continual sport of what the world calls fortune; and though I will not
wrong her by saying, She has ever made me feel the weight of any great
or signal evil; --- yet with all the good temper in the world, I affirm
it of her that in every stage of my life, and at every turn and corner
where she could get fairly at me, the ungracious Duchess has pelted me
with a set of as pitiful misadventures and cross accidents as ever small
Hero sustained.

\hypertarget{chapter-vi}{%
\subsubsection*{Chapter VI}\label{chapter-vi}}
\addcontentsline{toc}{subsubsection}{Chapter VI}

In the beginning of the last chapter, I informed you exactly \emph{when}
I was born; --- but I did not inform you \emph{how}. \emph{No}; that
particular was reserved entirely for a chapter by itself; --- besides,
Sir, as you and I are in a manner perfect strangers to each other, it
would not have been proper to have let you into to many circumstances
relating to myself all at once. --- You must have a little patience. I
have undertaken, you see, to write not only my life, but my opinions
also; hoping and expecting that your knowledge of my character, and of
what kind of a mortal I am, by the one, would give you a better relish
for the other: As you proceed further with me, the slight acquaintance
which is now beginning betwixt us will grow into familiarity; and that,
unless one of us is in fault, will terminate in friendship. --- \emph{O
diem praeclarum!} --- then nothing which has touched me will be thought
trifling in its nature, or tedious in its telling. Therefore, my dear
friend and companion, if you should think me somewhat sparing of my
narrative on my first setting out, --- bear with me, --- and let me go
on, and tell my story my own way: --- or if I should seem now and then
to trifle upon the road, --- or should sometimes put on a fool's cap
with a bell to it, for a moment or two as we pass along, --- don't fly
off, --- but rather courteously give me credit for a little more wisdom
than appears upon my outside; --- and as we jog on, either laugh with
me, or at me, or in short, do anything, --- only keep your temper.

\hypertarget{from-the-confessions-of-saint-augustine-trans-edward-b.-pusey-d.d.-first-published-c.-400-a.d.}{%
\subsection*{\texorpdfstring{From \emph{The Confessions of Saint
Augustine}, Trans Edward B. Pusey, D.D. First published c.~400
A.D.}{From The Confessions of Saint Augustine, Trans Edward B. Pusey, D.D. First published c.~400 A.D.}}\label{from-the-confessions-of-saint-augustine-trans-edward-b.-pusey-d.d.-first-published-c.-400-a.d.}}
\addcontentsline{toc}{subsection}{From \emph{The Confessions of Saint
Augustine}, Trans Edward B. Pusey, D.D. First published c.~400 A.D.}

Afterwards I began to smile; first in sleep, then waking: for so it was
told me of myself, and I believed it; for we see the like in other
infants, though of myself I remember it not. Thus, little by little, I
became conscious where I was; and to have a wish to express my wishes to
those who could content them, and I could not; for the wishes were
within me, and they without; nor could they be any sense of theirs enter
within my spirit. So I flung about at random limbs and voice, making the
few signs I could, and such as I could, like, though in truth very
little like, what I wished. And when I was not presently obeyed (my
wishes being hurtful or unintelligible), then I was indignant with my
elders for not submitting to me, with those owing me no service, for not
serving me; and avenged myself on them by tears. Such have I learnt
infants to be from observing them; and that I was myself such, they, all
unconscious, have shown me better than my nurses who knew it.

Passing hence from infancy, I came to boyhood, or rather it came to me,
displacing infancy. Now did that depart, --- (for wither went it?) ---
and yet it was no more. For I was no longer a speechless infant, but a
speaking boy. This I remember; and have since observed how I learned to
speak. It was not that my elders taught me words (as, soon after, other
learning) in any set method; but I, longing by cries and broken accents
and various motions of my limbs to express my thoughts, that so I might
have my will, and yet unable to express all I willed, or to whom I
willed, did myself, by the understanding which Thou, my God, gavest me,
practice the sounds in my memory. When they named any thing, and as they
spoke turned towards it, I saw and remembered that they called what they
would point out by the name they uttered. And that they meant this thing
and no other was plain from the motion of their body, the natural
language, as it were, of all nations, expressed by the countenance,
glances of the eye, gestures of the limbs, and tones of the voice,
indicating the affections of the mind, as it pursues, possesses,
rejects, or shuns. And thus by constantly hearing words, as they
occurred in various sentences, I collected gradually for what they
stood; and having broken in my mouth to these signs, I thereby gave
utterance to my will. Thus I exchanged with those about me these current
signs of our wills, and so launched deeper into the stormy intercourse
of human life, yet depending on parental authority and the beck of
elders.

\hypertarget{from-our-lady-of-the-flowers-by-jean-genet-trans.-bernard-frechtman}{%
\subsection*{\texorpdfstring{From \emph{Our Lady of the Flowers}, by
Jean Genet, Trans. Bernard
Frechtman,}{From Our Lady of the Flowers, by Jean Genet, Trans. Bernard Frechtman,}}\label{from-our-lady-of-the-flowers-by-jean-genet-trans.-bernard-frechtman}}
\addcontentsline{toc}{subsection}{From \emph{Our Lady of the Flowers},
by Jean Genet, Trans. Bernard Frechtman,}

First published 1943

Our Lady of the Flowers, still and already wearing the light, baggy,
youthful, preposterously tin and, in a word, ghostlike gray-flannel suit
that he was wearing the day of the crime and that he will be wearing the
day of his death, came there to buy a ticket for Le Havre. Just as he
got to the platform, he dropped his wallet which was stiff with the
twenty thousand franc notes. He felt it slip from his pocket and turned
around just in time to see it being picked up by Darling. Calmly and
fatally, Darling examined it, for though he was a genuine crook,
nevertheless he did not know how to be at ease in original postures and
imitated the gangsters of Chicago and Marseilles. This simple
observation also enables us to indicate the importance of dreaming in
the life of the hoodlum, but what I want particularly to show you by
means of it is that I shall surround myself only with roughnecks of
undistinguished personality, with none of the nobility that comes from
heroism. My loved ones will be those whom you would call "hoodlums of
the worst sort."

Darling counted the bills. He took ten for himself, put them into his
pocket and handed the rest to Our Lady, who stood there dumbfounded.
They became friends.

I leave you free to imagine any dialogue you please. Choose whatever may
charm you. Have it, if you like, that they hear the voice of the blood,
or that they fall in love at first sight, or that Darling, by
indisputable signs invisible to the vulgar eye, betrays the fact that he
is a thief\ldots{}Conceive the wildest improbabilities. Have it that the
depths of their being are thrilled at accosting each other in slang.
Tangle them suddenly in a swift embrace or a brotherly kiss. Do whatever
you like.

\hypertarget{from-spacetime-geometry-cosmology-william-l.-burke-university-science-books-mill-valley-1980.}{%
\subsection*{\texorpdfstring{From, \emph{Spacetime, Geometry,
Cosmology}, William L. Burke, University Science Books, Mill Valley,
1980.}{From, Spacetime, Geometry, Cosmology, William L. Burke, University Science Books, Mill Valley, 1980.}}\label{from-spacetime-geometry-cosmology-william-l.-burke-university-science-books-mill-valley-1980.}}
\addcontentsline{toc}{subsection}{From, \emph{Spacetime, Geometry,
Cosmology}, William L. Burke, University Science Books, Mill Valley,
1980.}

What does it mean to make a model of the universe? We must ask such
questions, even though they seem to belong in the realm of philosophy,
because of the fundamental nature of the subject itself. A good analogy
for this model making is the construction of a map of a city. The city
and the map of the city are two distinct things that no one would
confuse. This confusion is harder to avoid in cosmology. To be useful,
the map must be faithful to some of the city's interesting properties.
Furthermore, the map is constructed according to some specific
conventions: major roads are in red; north is at the top; and so on.
Also, a map does not try to define what a road is. That you learn by
demonstration. People point out roads and non-roads to you, and you
acquire the idea of a road by induction. Such basic ideas, which are not
defined but are learned by demonstration, are called \emph{primitive
notions}.

The same things will be true of the models of physical reality that we
make. We will construct a map of the world according to some specific
but arbitrary conventions. The image is constructed not on paper, but in
terms of some mathematical structure, such as a linear vector space.
Because we can talk about physical reality only in terms of some model,
it is easy for us to forget the distinction between the object itself
and our model of the object. As in our map of the city, there are some
undefined elements in our model. For spacetime, these primitive notions
are \emph{event, free particle, clock,} and \emph{light signal.} We will
describe in words what these correspond to, much as a roadmap might
discuss what it does in fact consider to be an acceptable road.

An important feature of any model is \emph{coarse-graining.} The model
ignores a lot of things. The exact width, surface, and color of a road
cannot be deduced from the ordinary roadmap. In fact, the roadmap is
more useful than a photograph precisely because it omits such irrelevant
detail. Similarly, there is a coarse-graining in our primitive notions
for spacetime. For example, we will not describe a clock by making a
detailed model of its internal construction. Such details will not be
mentioned. The model presumes that they are irrelevant, and it is the
task of experiment to verify this.

\hypertarget{from-the-passion-jeanette-winterson-vintage-books-1989}{%
\subsection*{\texorpdfstring{From \emph{The Passion}, Jeanette
Winterson, Vintage Books,
1989}{From The Passion, Jeanette Winterson, Vintage Books, 1989}}\label{from-the-passion-jeanette-winterson-vintage-books-1989}}
\addcontentsline{toc}{subsection}{From \emph{The Passion}, Jeanette
Winterson, Vintage Books, 1989}

Lovers are not at their best when it matters. Mouths dry up, palms
sweat, conversation flags and all the time the heart is threatening to
fly from the body once and for all. Lovers have been known to have heart
attacks. Lovers drink too much from nervousness and cannot perform. They
eat too little and faint during their fervently wished consummation.
They do not stroke the favoured cat and their face-paint comes loose.
This is not all. Whatever you have set store by, your dress, your
dinner, your poetry, will go wrong.

How is it that one day life is orderly and you are content, a little
cynical perhaps but on the whole just so, and then without warning you
find the solid floor is a trapdoor and you are now in another place
whose geography is uncertain and whose customs are strange.

Travellers at least have a choice. Those who set sail know that things
will not be the same as at home. Explorers are prepared. But for us, who
travel along the blood vessels, who come to the cities of the interior
by chance, there is no preparation. We who were fluent find life is a
foreign language. Somewhere between the swamp and the mountains.
Somewhere between fear and sex. Somewhere between God and the Devil
passion is and the way there is sudden and the way back is worse.

I'm surprised at myself for talking in this way. I'm young, the world is
before me, there will be others. I feel my first streak of defiance
since I met her. My first upsurge of self. I won't see her again. I can
go home, throw aside these clothes and move on. I can move out if I
like. I'm sure the meat man can be persuaded to take me to Paris for a
favour or two.

\begin{quote}
Passion, I spit on it.

I spat into the canal.
\end{quote}

\hypertarget{mike-amnasan-the-universes-on-semantics-and-learning}{%
\chapter{Mike Amnasan: The Universes: On Semantics and
Learning}\label{mike-amnasan-the-universes-on-semantics-and-learning}}

~

\hypertarget{section-14}{%
\section*{1}\label{section-14}}
\addcontentsline{toc}{section}{1}

\lettrine[lines=3, findent=0em, nindent=0.1em, lhang=0]{I}{ can’t tell you}
what I was just thinking. As in Augustine's view of intuition, the
associations I registered were too free of any repeatable limitations
for me to verbalize the experience. Perhaps these associations were of
diverging thoughts that have not departed my mind. The most handy
example of something similar is the simultaneity of sense perception.
Each sense perception is specialized and in that respect removed from
the whole and yet also registered in the same moment. In the thought I'm
speaking of there were different concerns, we might even say a universe
of concerns none of which I can fully express. This complex event might
be considered incidental in regard to what I have learned to value. I am
now attempting to acquire a greater appreciation of something I cannot
verbalize, meaningful associations I can't excite a recurrence of
through keying words into a computer.

Semantic sensation is never original. It must always be familiar.
Language does excite original sensations, as in the sound of a speaker's
voice, but the semantic experience itself is never sound or vision, or
any other sense perception of the material universe. Original experience
of semantics would be like immediately comprehending a language we never
heard before. This kind of appreciation is possible with music. Music
we've never heard before can be immediately appreciated as music, but
semantics, like memory, must always be a response to what is already
familiar.

My concern is how to proceed. If I can only register verbally what has
already become familiar through cognitive means, my work with language
is \emph{not} directed toward spurring meaning for the first time. Has
there ever been a first time in regard to comprehending language? Is
anything we read utterly strange, or is it rather strangely familiar? We
may read something and make no sense of it, and later return to it and
find familiarity as if we always should have been able to comprehend
this particular passage. This parallels how we initially acquire
language through a growing familiarity with the effects of verbal
expression. We learn to fortuitously repeat limited effects. We grow to
appreciate what we had already experienced albeit as incidental and free
of the constraints of communication. Infants can distinguish between
phonemes their parents, having learned a particular language, can no
longer tell apart (Pinker 264), and meaning is similar in this respect.
To understand how this can work we must put aside the notion that
language makes meaning. If this were the case we would begin with a
rather meager supply of semantics that we would not want to reduce by
limiting what we come to appreciate. The verbal cannot originate
anything we come to value as verbal expression. Perhaps this is true of
all our senses; the original experience is what we take in from sensory
organs that does not result in sense perception, sense perception being
that which we perceive to the extent it is already familiar to us,
already among the effects we are prepared to make. In this sense the
original is not an object, but rather that which we are unprepared to
value: particular and unique effects we do not register.

Perhaps I simply want to learn, having reaped benefits from what I have
learned in the past by reducing the number of hard choices I must make
or that are made for me. Does the construction work, through which I
make a living, decide what sort of interests I can successfully pursue?
If cognition is verbal then it must follow that we can only appreciate
what language excites an interest in. If meaning originated in language
semantics could not have the benefits I hope for. The use of language
would be the limits of its effect, and all the incidental experiences
that I accrue on a construction site, when I seldom talk, and much more
rarely about what I'm now writing, would go for nought---or worse. My
time on a construction site would suspend the cognition I long to learn
from. What I then felt to be my most prized asset, the verbal, and all
my work with the verbal, would be suspended as I endured a painfully
irrelevant expenditure of time and energy. Furthermore my understanding
of the world might be contaminated by prevalent views that re-inscribe
my thoughts with unwanted expressions.

We can see that my disposition toward construction work is strongly
impacted by my regard for language, and this involves study that is, in
part, motivated by a desire to develop an understanding through which I
can regard different aspects of my life as of mutual value. I want this
view to be revised within different kinds of experience that I am never
fully "conscious" of---even if some rehearsal of what I cannot register
is, at this very moment, being performed as simultaneous events in my
mind. Language is not only directed toward fortuitous limitations, but
also toward realizing a universality that is always changing in its
simultaneous instances of social expression, expression that we cannot
oversee, or monitor in any way, but that our minds do, to some extent,
incidentally receive. Mental activity may be full of simultaneous events
some of which act like separate subjects we anticipate. It may be that
these subjects, if we are removed from interactions that disrupt our
expectations, must form into a consensus of what others will expect from
us, and limit what we can expect of ourselves, selves that come to
behave as sovereign beings constituted as a uniformity. Others become
mere foils, since they can only be modeled after this consensus that
determines what we are disposed to find.

Our only particular and unique experience is that purveyed through sense
organs that introduce stimuli from the material universe---and yet
language is often external, as in the example of lines of words on a
monitor. How is it so very different? How can it be lacking in original
experience if there is an external and physical text that we respond to
in the material world? The answer is that our experience of physical
language, whether as sound, or text, or raised dots in the case of
braille, is not semantic, but a different kind of sense: vision, sound,
touch. Even if vision, sound, touch, also respond only to what is
already familiar they still form our sense perceptions of the material
universe, while the semantic effect excited by language is never a
direct expression of material reality.

As you read these words they arouse an immaterial universe simultaneous
with your perceptions of material reality. A simplified version of what
happens is this: the stimuli from your eyes travels along the optic
pathway until it enters the region of the cerebral cortex associated
with language.\footnote{Steven Pinker states that ``{[}Gazzaniga's
  coworkers{]} found that all the areas that have been implicated in
  language are adjacent in one continuous territory. This region of the
  cortex, the left perisylvan region, can be considered to be the
  language organ'' (307).} This linguistic region responds only to the
familiar stimuli, the marks, that spur semantic sensations. Stimuli from
the eyes passes on to the visual area of the occipital lobe where it
excites the visual perception of the physical world. The visual follows
hard upon the semantic so that you register both without any discernable
interval and you cannot tell, from what you experience, that more than
one sense perception is responding to stimuli from the eyes.

The preceding description is far from an exhaustive account of what
happens when we are engaged by written text. Language could only have
developed because our response to the stimuli from our ears, eyes, or
the surface of our skin was never exclusively devoted to creating sense
perceptions of the material universe. Our perceptions may be more of an
afterthought. It is possible for a creature to have an eye, and to
respond consistently to stimuli from that eye without the luxury of any
visual impressions.\footnote{In the documentary Inside the Animal Mind a
  man, and also a monkey, both suffering from blind sight, are shown to
  respond correctly to visual stimuli on their right side, visual
  stimuli that they do not register having seen.} Such behavior would
not be derived from cognition minus sense perception, but rather a
reaction lacking in a consistent impression that could become either
reliable or doubtful.

When reading, we do not fix our attention on the series of characters in
order to receive visual impressions that we then interpret, transforming
the visual into semantic sensations. Sense perception does not work
sequentially in that way. We do not see something, and then, after we
have concluded seeing, hear a sound, and then take a turn at smelling.

Stephen Pinker, in \textsc{The Language Instinct}, refers to the
linguistic as a sixth sense. When I read his book, I began to feel less
trepidation in regard to proposing that semantics is the sense
perception of, and creation of, an immaterial universe. Here was a
respected linguist repeatedly returning to the idea that linguistics is
a form of sense perception as I was. But this mutual appreciation was
one version of myself, while I was also aware that I had a social
identity, within a community of writers, as a construction worker
writing about writing from a position of disadvantage. This identity
presently determines what is expected of me more than anything I could
say to the majority of those people who are familiar with what I write.
Being able to cite Pinker may up my chances of convincing others that I
am not merely avoiding, through intellectual exercises, the emotions,
that fueled earlier writing. The proposal that semantics is a form of
sense perception does not allow for creating new meaning as the direct
result of innovative language, and what I am saying will not be
understandable if approached with the expectations that we can create
new meaning solely through unconventional grammatical constructions.

If semantics must be consistent like any other sense perception, and we
can only read what we are already disposed to read, this is not to
suggest that there are no surprises in reading, or that we all read the
same, or that we should ideally read the same, or that we should always
use words in a similar manner. It does not even necessitate that we must
change our current expectations to learn, but rather that we must form
new expectations. This is an important distinction since we tend to
think of values as a set of values so that to change we must replace one
set of values with another. Perhaps this is the path of least
resistence, as in the example of fully adjusting to the demands of some
vocation we want to succeed within---but, what if we want to gain the
advantage of studying more than one discipline? The material for
interdisciplinary study would be available even if there were no
programs that encourage it. Perhaps the main concern, in this regard, is
whether or not we can become proficient without limiting the focus of
our interests. Can we learn, not only from close study with expedient
results, but also through a delayed cognitive response that includes
more material than we can consider within any close reading? In ``The
Love of Reading" Virginia Woolf speaks of "the process of
after-reading.'' After a book is read it takes on a shape "held in the
mind" (416). After-reading may occur days after the book has been set
aside, and this later experience may become related to other experience;
simultaneous mental events that were incidental at the time of reading
may become appreciated as meaningful associations. We can see that the
space-time of simultaneous events is non-sequential; intervening
experience does not conclude previous events. Our appreciation of
simultaneous events effects our conception of space-time, and even what
is continuous through intervening events that need not determine the end
of prior events. We have the example of sense perception in which case
specialized processes occur simultaneously; we register very different
senses in the same moment, and this simultaneity of perception may very
well have developed from simultaneous mental processes we never
register, processes as incidental to particular occasions as continuing
to breathe.

In the case of my own life as a construction worker there is no way I
could make useful contributions to philosophy or any other academic
discipline if close reading, or some similar attention, was exclusively
how we learn, and perhaps my confidence in being able to offer useful
models is deluded. Within a fairly isolated effort I cannot tell the
difference between adding to current understandings and my own ignorance
of other people's efforts that I may simply repeat. I may expend time
and energy in a direction others have long ago recognized the folly of.
To some extent this is rhetorical in that I have found confirmation for
my efforts in books I've recently read---but I do depend upon intuition
to an extent that may seem sloppy to anyone who expects strong textual
support for what I might propose.

My own social identity, and also my ability to think of learning as a
greater appreciation of what is already cognitively familiar, are at
stake in, to some extent, replacing sequential models with those of
combined influences, the influence of simultaneous expectations I can
never be fully aware of. Often when I write, I key words into a computer
as a one-time event that does not work well for later verbal
communication---even with myself! Though writing did excite an
experience that I believed I was fully expressing, the words I laid down
failed to spur a strong recurrence of that experience the following day.
Perhaps there were events occurring within cognition that were
registered as successful writing---but I was learning a new emphasis on
the familiar, at the same time as I was making language, and I was not
\emph{only} reading the effect of what I was keying into my computer. I
don't mean to suggest that this considerable cognitive delay is the only
way we learn in regard to language. Certainly there are occasions in
which learning is derived from a close focus on a passage that is
difficult to comprehend without repeated readings in which significant
aspects of language begin to emerge, and it may seem that we are
bringing meaning into "consciousness." But even this instance of close
attention requires a non-linguistic, cognitive response to make the
already familiar appreciable.

If we cannot learn without others, and if, as I believe, reading is not
sufficient in itself for changing verbal expectations, how can this
learning even occur sitting alone before my computer monitor? I do not
know the kind of physical incidents and the frequency of those occasions
required to learn. We may even learn from anticipating upcoming social
events in regard to which our participation will be expected, events
that we have not yet entered. One important element is that others we
learn from do not have to be proven. We do not unlearn if we find that
we gained a new appreciation of a subject through misreading what
someone said---or at least misreading what we later came to believe they
felt they were saying. We may conclude that we need to anticipate others
more in line with what we come to believe they meant, or we may continue
to look for the unintended meaning beneficial to our own study. The
social stimui that infuse cognition with original experience may not be
traceable, especially more enduring influences that occur over a greater
lapse of time. Since language is only acquired through cognition, new
linguistic expectations, that have a broad influence on our thinking,
cannot be the \emph{immediate} result of verbal interactions.

We must be disposed to anticipate what we read before we are engaged by
language to receive semantic sensations. If this is true then we cannot
become self-taught through isolated reading. Having said that, I should
point out that I am not talking about having rallied all our senses to
participate in a social occasion like a classroom discussion, and then
later on, after we arrive home, reading similar texts. Even the
expectation that we may benefit by repeatedly reading a difficult
passage must be learned as a change in what we are able to appreciate.
If rereading is sufficient for comprehension (that is, there are no
unfamiliar terms or an unconventional use of a word that requires us to
turn to other sources like a dictionary or the instruction of a
professor familiar with the author) we must still somehow realize that
the components of what we are reading are already familiar, and that we
must adjust our expectations when what we were initially looking for is
not discovered.

When I speak of isolated reading, reading in regard to which I have not
become well disposed toward a text by any prior social occasions, what I
then need to read effectively is the expectations of others to disrupt
my satisfaction with the values I was trying to read with. An important
consideration here, is whether new expectations, that allow me to
comprehend more, displace previous values or become simultaneous with
previous values, adding to the mix. Both possibilities occur. The latter
case does not involve displacing one social logic with another felt to
be superior, but simply developing the expectations through which to
read a text that my life has not prepared me for. I may then discover
what was already, incidentally, familiar to me.

I am trying to anticipate sufficient social involvement to lead to
expectations of further participation with what could be regarded as
incidental aspects of my own mind. My social position, is, and must be,
one of inferiority in relation to material I have already assimilated.
Despite whatever mutual effort I am working toward, social inferiority
is still an important aspect of my efforts. I may feel \emph{that} as my
physical body, alone in an unsympathetic universe, the vulnerability,
the loss of control over the sensations of my skin, my face, my
mannerisms and my speech if I speak out in a classroom or within some
other public event. This physical experience disrupts my sense of being
a coherent whole "conscious" of who I am, and these occasions have
always had this effect on me, but earlier in my life I felt this
disruption as revealing weakness, and perhaps it does, but it is the
weakness of risking a social life.

Poor beginnings left me too confident in my own mental resources in
comparison to the culturally-deprived interactions I was limited to.
Like Nietzsche's mosquito "feeling within itself the flying center of
the world" (42), isolated efforts tend to return me to an exaggerated
sense of ``self-importance." There is nothing within cognition and
semantics to prevent me from being highly successful at meeting my own
expectations, and imaginary others as well---though I will, at the same
time, feel a nagging suspicion that something is missing. I have become
deceived in regard to the expectations of those others who I anticipate
impressing. Augustine states''If I am deceived, I am" (33), and it seems
that he is talking about being deceived in a social context. If I
realize that I have been deceived in anticipating the expectations of
others I am a social being learning from the effect of rallying my
senses to deal with the vulnerability of interacting.

Paradoxically it would seem that we exist, as social beings, because we
are not in full control of our body, our appearance, voice, mannerisms,
what anything we express means. Judith Butler states, ``If one always
risks meaning something other than what one thinks one utters, then one
is, as it were, vulnerable in a specifically linguistic sense to a
social life of language that exceeds the purview of the subject who
speaks'' (87). If we have no disruptive memory of interacting with
others, we not only do not ``exceed the purview of the subject who
speaks,'' we slip into a false consensus of what the language we use
means.

Learning requires an infusion of original experience from the material
universe and, however much we may discount other aspects of learning,
while paying close attention to a text, reading, alone, will not likely
allow us to comprehend aspects of a text we were not already disposed to
understand any more than it will allow us to exceed our own consensus of
what the language we use will mean to others. The semantic universe of
each individual is private, and yet the disposition to appreciate
incidents of cognition differently than we have previously is not. In
that respect the potential of our minds is not entirely our own, in that
relying exclusively on our own values is a form of deprivation, a
reliance on a false consensus. Isolated reading and writing can locate
us within this deprivation; vision is relegated to locating words, as
the verbal appropriates the eye, and other forms of sense perception
become distraction or incidental to reading.

I exaggerate. Our other senses do, and must, remain significant while we
read. The historical examples that Martin Jay provides in
\textsc{Downcast Eyes} suggest that a greater attention to text can
result in the denigration of vision as the purveyor of "truth," but, as
Jay points out, this need not result in the denigration of vision as a
rich experience. I don't know that reading and writing as a form of
sensory deprivation has ever been studied, or how it could be researched
since anyone isolated in their efforts, in the manner I'm suggesting,
would not succeed in drawing attention to himself, and so, for me, this
must remain autobiographical, a concern I will return to.

\begin{center}\includegraphics{images/"pandanus.png"}\end{center}

When walking down a street we can register both the sight and sound of
people passing by and the traffic in the street. We can the feel a
breeze on the skin, the pressure of a backpack over one shoulder, and
the contraction of muscles, all at the same time, and this simultaneity
of sense perception is essential to our being---even though we do not
require equal focus on all possible forms of sense perception. Speech is
very different from written language in that, though the semantic
experience is no different in kind from that excited by written
language, the sound of a voice and even background music may be
considered enhancing. Taking a walk and enjoying a view may be
considered an occasion for conversation, as also may be the excitement
of participating in a classroom discussion. When reading, we can begin
to realize that the sensations excited by language are not vocal, but
rather sensations unique unto themselves.

When studying some subject at home we do not want to be distracted by
noises coming from a neighboring apartment---but we still hear. In this
case we may be preparing for experience within a classroom, and take
advantage of the different space-time of semantic experience to prepare
for expectations that may not be reflected in our immediate
surroundings. It will not likely help our solitary preparations to
anticipate greater involvement with the people next door---and yet in a
sense we do want this simultaneous experience to tell us of other lives
as we study, other lives that are not a diminishment of what we are
presently involved in. We should not be lulled, while studying, into
believing that life loses significance at the periphery of our immediate
experience of a particular text.

I cannot monitor all my influences---but this is not to suggest that I
do not appreciate what others have written. I would quote more often if
it worked out for me, but I edit so much that cited passages tend to
become progressively arbitrary over the course of writing---and I am
also effected by undo tension associated with citing. Citing can be
promotional in giving credit where credit is due. It can call attention
to other authors and works that may be of interest to the reader, and
perhaps more importantly it keeps others in the foreground as a constant
reminder that critical writing is not an individual effort. And yet
despite these important reasons to quote others, in my own experience in
community colleges, citing was introduced as having two major concerns.

The first was plagiarism: stealing another's ideas. The second was the
need to provide textual evidence to confirm the worth of anything we
propose. I particularly remember fliers handed out in each class
defining plagiarism and warning against it---though I never knew who
orchestrated the distribution of this cautionary note. My introduction
to citing told me that we were being watched to ensure that we did not
get away with committing serious infractions, and this emphatically
warned of our full responsibility in regard to recognizing when we were
borrowing ideas, and to make sure that we did not misrepresent any
viewpoints we paraphrased. These may be valid concerns, but not ones
that I feel should be forefronted in a way that turns the novice,
unfamiliar with the terrain, into a solitary being watched for possible
infractions, since citing should call attention to the first person
plural of critical writing, the others we cannot fully separate from our
own experience of learning.

I remember, as a teenager, lying drunk on my back, in the back yard of
my parents' tract home. I was too inebriated to do much more than move
my head from side to side, through which I could see the stucco back of
the house, and the garage which, among other things, sheltered my
father's automotive tools (the trade he taught in a local high school).
Letting my head roll in the other direction, I could see the redwood
fence at the back of the yard. Looking straight ahead I saw the sky,
representing my freedom to seek out a being identical to myself. I felt
my drunken stupor to be of little disadvantage for this impossible task;
I felt myself to be different to the bone from all others. Even if my
difference felt undeniable within a tract home development in the San
Gabriel Valley, it was an impression formed in relation to a very small
number of people, a tiny perspective within the sprawling suburbs of Los
Angeles. I didn't have sufficient evidence to confirm or deny what I
firmly believed---and I never would even if I expanded the world I took
an interest in. My disposition was not due to a lack of similarity to
others, but rather that I saw them much as I regarded myself as an
isolated being. I could not see others as having simultaneous, and
sometimes incongruous qualities. I was poorly disposed toward critical
abilities in that they were not expected of me or anyone else. Semantics
remained a mysterious by-product of physical language.

I felt an inexpressible internal conflict that formed me into a
mysteriously complex center, of little value to others; they were mere
foils for my troubled outlook. I felt that only I had a complex inner
world incongruous with my behavior and I regarded this lack of the
ability to act as a coherent whole as a terrible fault. I felt such
disquiet, such a need to rail against everything that made me feel
small. I was beleaguered by the anxious suspicion, barely registered,
that in richer circumstances it was possible to perform a more varied
life, a life in which projects in common with others could displace the
need to continually recreate an isolated self, reacting to oppression in
my mind.

Should I include autobiography? It may reflect the limits of my social
experience, and in that respect I often feel torn between first person
and first person plural. Finding it difficult to settle on a pronoun is
probably similar to a discomfort with citing, a lingering weakness, a
fear of leaping confidently into a position of equal standing with those
who I might quote. I want to describe my simultaneous experience as a
construction worker, even if it is an inferior position from which to
participate in a mutual critical effort. Should I tell you of running to
the underground MUNI in the dark of morning, while pulling a dirty sweat
shirt over my head, because I lingered too long reading over this I'm
working on? We don't really know how important this personal information
might be if thrown into the mix of critical writing. Perhaps the
singular, where a plural is wanted, carries with it an overblown sense
of the subject writing, which should be dissipated within a mutual
effort in which having been influenced by others is strongly highlighted
and rightfully so, and I'll return to that soon.

Sheet Metal work, my only income, has been recently altered by the
growing aesthetic preference for exposed duct, and this has created a
tension between the traditional building trades emphasis on hanging duct
as fast as possible (duct that will be hidden above a ceiling) and
producing the sculpture of visible duct work important to the appearance
of the finished building. The use of the eye for hanging duct is
restricted by the need to get the duct up quickly, the need to quickly
handle the duct into aesthetic lines. We are unlikely to see it painted
within a finished building---which is just as well in that the finished
vision will not include the heavy, galvanized air-tunnel, physically
hoisted in a dusty, gutted building, a richer physical experience than
keying words into a computer as I'm doing now.

One benefit of this work, an unexpected one considering that it would
seem to prevent me from pursuing the interests that I find more
compelling, is that it involves using multiple sense perceptions in an
exhaustive manner. Construction work can be a personal measure of the
value of what I have been writing in that it involves the potential
threat of viewing aspects of my life as mutually exclusive and harmful
to the realization of some hypothetical life more suited to my
abilities. This is something I am not likely to forget while writing,
and perhaps exhaustion has come to displace the tragic view of looking
at myself unsympathetically as if from outside the construction site,
from a campus, or anywhere else from which the view of myself as a
social being will be an expression of limitations, a small circumscribed
place in society that my desires should perhaps be shaped to fit. Where
the tragic view prepared me to pursue the seemingly impossible, as an
embrace of the unnatural (my solitary life passed before me as is said
occurs before death), exhaustion has taken the place of my previous
tragic view, becoming a universal exhaustion. Even though I may be a
poor scholar, \emph{we} are exhausting \emph{our} abilities.

It may be that I am inclined to make pronouncements without doing the
research required to substantiate my claims---of doing the wrong work!
This is the modus operandi I am most likely to be accused of, and that I
feel a need to justify in that I find that if I always need textual
evidence to back up whatever I propose, I cannot maintain momentum---at
least at this time in my life. This doesn't mean that anything I say
comes out of the blue, but rather that I don't keep track of influences,
especially since there tends to be considerable delay, and a great deal
of intervening experience before I arrive at what I may eventually
express with some confidence. I might then speak of this universal
exhaustion as a universal delay, a disjunctive interval, in which my
efforts do not come to a sudden a halt when I must turn my attention to
construction work. Construction work itself has more meaning as part of
simultaneous projects like that which I'm now working on.

\hypertarget{works-cited}{%
\section*{Works Cited}\label{works-cited}}
\addcontentsline{toc}{section}{Works Cited}

Augustine. \emph{The Essential Augustine}. Ed. Vernon J. Bourke.
Indianapolis: Hackett, 1974.

Butler, Judith. \emph{Excitable Speech; A Politics of the Performative}.
New York: Routledge, 1997.

\emph{Inside the Animal Mind}. Episode Three: "Animal Consciousness."
Nature Video Library. Produced by Green Umbrella, Thirteen/ WNET New
York and the BBC, 2000.

Jay, Martin. \emph{Downcast Eyes; The Denigration of Vision In
Twentieth-Century French Thought}. Berkeley: University of California
Press, 1994.

Nietzsche, Friedrich. \emph{On Truth and Lie in an Extra-Moral Sense}.
\emph{The Portable Nietzsche}. Trans. Walter Kaufmann. Ed. Walter
Kaufmann. New York: Viking, 1965.

Pinker, Steven. \textsc{The Language Instinct; How the Mind Creates
Language}. New York: Harper, 1995.

Woolf, Virginia. "Virginia Woolf: The Love of Reading." \emph{A
Bloomsbury Group Reader}, Ed. S.P. Rosenbaum. Oxford: Basil Blackwell,
1993. 415-418.

\hypertarget{betsy-andrews-note-from-the-underground}{%
\chapter{Betsy Andrews: Note from the
Underground}\label{betsy-andrews-note-from-the-underground}}

~

\hypertarget{working-notes}{%
\section*{Working Notes}\label{working-notes}}
\addcontentsline{toc}{section}{Working Notes}

\emph{As a poet, I'm so intensely enamored of craft that the line can
become my taskmaster, both beloved and despised. When I began the piece
herein, I intended it as a poem, endstopped and enjambed and all. But
there were so many stories converging---an autobiographical story, an
ancient mythological story, the stories of numerous industries (the
industry of patriarchy, the industry of meat, the industry of religion,
the industry of death), the story of the archetypal presence in the
day-to-day---that the messy horde of them resisted the line's enclosure.
The mob had a mind of its own, came to quorum instead in the alchemical
vessel of the very tall tale. So, it was that a choice made itself.}

\begin{center}\includegraphics{images/"pandanus.png"}\end{center}

\lettrine[lines=3, findent=0em, nindent=0.1em, lhang=0]{M}{y father was a lying drunk}.
My father's name was Sin. He chased my mother with a tire iron about the
yawning frontyard. Then he invented escape. As helpless in his absence
as presence, my mother stood in the same yard and wept when the
floodwaters rose. The neighborhood gossip, Tiresias, unzipped her
cornstalk ears. She tossed her lit cigarette at the garden snakes
fucking and pronounced herself a bard. She sat on her porch across the
street and composed the epic of my parents' frontyard which she sang in
verses just brief enough for the neighbors, in passing, in learn how to
hum. It was the din of the rain falling in squalls; sprung rhythm of
maelstroms unravelled. This is how I became acquainted with some of my
pastoral cycle. I learned to chant these slaughtering idylls from the
bard across the way. I was there in the midst of foul weather. But I
remember nothing I was not retold. These recollections are potsherds,
fragments of sunken relief dredged up and presented with someone else's
commentary, as if behind glass.

My memory is a museum of pieces of dubious provenience. The story I'll
tell is 5000 years old and is nothing if not debatable. How am I to tell
it, then? How about:

She ate me?

She ate me all the way back. She ate me all the way back into childhood,
all the way back through birth. She ate me back through conception. Back
through my parents' wedding before The Vietnam War. Back to my father
falling on a bloodied Korean slope (shrapnel fucking his thigh where
metal will seep from his skin in the shower forty-eight years hence).
She ate me before my mother's birth between The World Wars. She ate me
before The French-Indian War, The Spanish-American War, The
Sino-Japanese War, The Franco-Prussian War. She ate me before The Ten
Years' War, The Boer War, The Crimean War, The War of 1812. She ate me
before The Seminole Wars, The Opium War, The Peninsular War. She ate me
before The Seven Years' War, The First and Second Silesian Wars, The
Great Northern War, Queen Anne's War, King William's War, King Philip's
War, and The War of Devolution. She ate me before The Thirty Years' War,
The Wars of Religion, The Wars of the Roses, The Hussite Wars, and The
Hundred Years' War. She ate me before The Punic Wars, The Sacred Wars,
The Samnite Wars, The Messanian Wars, The Macedonian Wars, The
Mithridatic Wars. She ate me before The Peloponnesian Wars, The
Corinthian War, and The Trojan War. She ate me before The War Against
Hyksos, The War Against Palestine, The War Against Arcadia, The War
Against Thessalia, and The War Against Babylon. She ate me back into Ur.

Er. Utterance before argument. Er. The advent of war.

It was not entirely erotic. Entirely. It was grave.

She ate me, but I buried myself. I had buried myself before. Locking the
bathroom door as a girl and crouching down in the shower, my hands
cusping my face and my face urged at the tiles, I'd dream my dream of
the grotto. From Latin for ``crypt," from Greek for''vault," and in
feminine form, "to hide." Grotto, speleo-aqueous myth of safety as in
drip hole, gravity spring, lily pool, lost river.

``The possible use of limestone caves as shelters from the effects of
nuclear weapons has received considerable attention in the past few
years,''\footnote{George W. Moore and Brother G. Nicholas.
  \textsc{Speleology: The Study of Caves}. Boston: D.C. Heath and
  Company. 1964.} in the juvenile water of the grotto whispers the blind
cave fish.

Sin drew the living room blinds. Tiresias pressed her ear to the ground,
discerning the clatter of hooves. Sin sat on the couch in his shorts and
chained my mother, for his children's digestion, to the dissembling line
of his industrial tongue. He stunned and stuck and scalded and scraped
her. He split her along the backbone. While her head and her hocks and
her skin were removed, we milled in the lairage wild-eyed, waiting for
our turns. Picture the herd slunk from the mother, pushing and bumping
about:

Suddenly, at one end of the herd, a huge thrust forward. The first
maddened layer crushes the next which crushes the next, and so on. A
wave of compression washes the herd. But the herd is still milling in
random directions. The wave that started crisply, decomposes, becomes
disorganized. By the time it reaches the opposite end of the pen, the
pattern's a mess.

The same thing happens to a wave of sound. Close to the source of sound,
the compressions are patterns of tightly packed molecules. But the
molecules have no direction in mind. They simply mill about. The
compressions untighten as the wave of sound travels so that meaning
near-grazes our ears. When the pattern is totally thronged by the
milling, the herd hears nothing at all.\footnote{Sara Stein. \textsc{The
  Science Book}. New York: Workman Publishing. 1979.}

I woke to my own furious bleating. ``Kill her," I said. I thought this
would preserve us.''Kill her." And, "Kill her." And, "Kill her." My
siblings fidgeted deafened and dumb. My mother was processed and canned.
That night my father, laid drunk up in couvade, gave birth to much of
me.

We were feeding on hard-shelled, trap-doored things: augers and miters
and lightning whelks, spindles and turrids and helmets; horse conchs and
dogwinkles; mouthdrills and neptunes and clives.

I was going him one for one with the cocktail; a choking wash of an
audible note poured from a giant, irregular jug. Inky-mouthed, my father
dumped his boxcar eyes on me and spoke a hollow meander. The body of his
speech was dominion, the limbs of which were gifts to me. Then, atavism
buckling him under, he belly crawled into darkness. He had drunk his
alluvial fill.

Thus he lay till the morning of the day appointed for his inhumation.
The oven was being heated for the baking of the funeral meal when he
awoke, climbed out of his coffin and walked downstairs, where he was
found by the horrified cook, standing before the kitchen fire,
complaining that he `felt the cold.'\footnote{Kenneth V. Iserson,
  M.D.~\textsc{Death to Dust: What Happens to Dead Bodies?} Tucson:
  Galen Press, Ltd.~1994.}

Deboned ham, he lived for years after out of spite for my skilled
accretion of matter from out of his pinwheeling arms. To me, he had made
these gifts: the me of power and wickedness, of plundering cities, of
rebel lands; the me of strife and judgement, of quiver and standard and
sword and club and silence and running and fear. He made me the me of
consternation, of deceit and dismay and lament. He made me the me of
kindling fire and the piling of burning coals. He made me the me of
righteousness, of triumph and travel and beer. He made me the me of the
art of speed. He made me the me of the making of plans. He made me the
me of godship and forthright speech and deceitful speech and
grandiloquent speech and divinity. He made me the feeding pen.

Once I was so much of me, and he was a mummy of displaced vigor, he
gathered himself in his goatskins and prophesied the epoch of his own
replenishment. Sin developed an invasion theory. He called it
Revelations. He sat me on a scarlet beast and poured me a cup of
abominations. He attempted to swallow me whole.

By then I was known for Clouds Of Which Were Called My Breasts and, of
course, for my Wondrous Vulva.

I was called Lady of Myriad Offices: Queen of Heaven, Lady of Victory,
Light of the World, Leader of Hosts, Creator of People, Guide of
Humanity, Mother of Fruitful Breast. I was called Goddess of Goddesses,
River of Life, Shepherdess of Lands, Bestower of Strength, Framer of All
Decrees. I was called Opener of the Womb. I was called Lady of Date
Palms.

I was called Forgiver of Sin.

He had called me the Whore of Babylon. He had called me the Mother of
Harlots.

I called myself Bereft of Forgiveness. I called myself Indigestible.

That was before my chthonian striptease when flutes and facets and
flowstones were cleft. When squeezing, scrambling, chimneying, jamming,
rappelling myself into caveat, she ate me, which was as natural as a
logarithm. It was her due.

Tiresias' slumlord husband, having made of his property a catafalque and
leased his tenants to incandescence, was arrested amidst the cremains.
Tiresias packed her plucked feathers and eloped with eight billion
flightless confessors. I was left, big toes tied together, legging the
walls of the impermeable gist. What was I to make of all this bruised
and cutting matter? of the evisceration bequeathed me by pall-bearing
parents? (My mother had selected my casket from an overnight catalogue
while sipping a cup of coffee and sitting in a comfortable
chair.\footnote{Robert T. Hatch. \textsc{What Happens When You Die From
  Your Last Breath to the First Spadeful}. New York: Carol Publishing
  Group. 1995.} My father had won first prize from the Académie de
France for inventing a clawed forceps for pinching my corpse's nipples
to assure that I was dead.\footnote{Iserson.})

I refused to play their thanatomime. I left them intoning Vade in Pacem.
I took to speech.

I spoke blank verse. I spoke canzone. I spoke tanka and terza rima. I
spoke bouts rimés and rondeau and renga. I spoke eclogue, epithalamium,
madrigal, limerick, lune. I spoke senryu, sestina and satire. I spoke
calligram, cento and ode. Ottava rima, projective verse, quatrain,
ballade, epistle, the blues. I spoke ghazal. I spoke her to sonnets.

But she ate me; she furrowed my heart. Inside the asylum of her ear, my
language was rendered offal.

I had played awhile at cultic prostitute, at sheepfold, at castrate, at
the loosening of hair. The suspect was an effigy who moonlighted mowing
the yawning frontyard. He laid me in the lilies unmoored from Tiresias'
abandoned bed. He promised me strings and pulleys and bells. But when I
touched him, the webbing between my fingers grew opaque with loss of
blood. The mirror he held to my mouth clouded a phrase of silver sulfide
of which the text was ennui.

So I buried myself, and she ate me.

The causal meaning of my descent has puzzled somatic scholars. Scholars
have been inclined to dismiss as mere excuse my stated reasons:
bull-like passion, raw desire, libido, ego, gestalt.\footnote{Sylvia
  Brinton Perera. \textsc{Descent to the Goddess: A Way of Initiation
  for Women}. Toronto: Inner City Books. 1981.} The god gets
antediluvean kinged. The god gets hold of her leaking hair, and the god
gets to wager a bet with the Huns: if he lops off her head and commands
her to run, in how many steps will she crumble? The god gets the power
of conjugation. The god gets declamation. He gets flashes of lightning,
fits and wind. The god gets to steep in the rift.

I get the query fever. I get the conjectural choke.

The god gets oolites and porthole and spongework.

I get to peter out.

In the factory of piracy, god becomes god exponentially. I become
boil-in-bag. The fatal hemmorhage of me out of me onto an areal map of
the god: I am displaced into ambience; the god oxidizes and dilates. The
god is blisteringly explicable. I go wrinkled and clot like jelly while
the conduits of the graduate crawlspaces hook and barb their own lips.
Once every 20,000 stitches, they publish a rattle elucidating civil war
and embalment; I am painted as pale as putty and petrified like marble.

The bookkeepers take to my carcass like germs. They digest me into the
emergent naturalness of that improving stalagmite, the god. With the use
of this technology, the bookkeepers patent the concept of 'freshness.'
My rind for the fruitful media.

The morticians of the media feign post mortem inspection. My survival,
assumed to be a chronic rash, is expurgated from the thin skin.

But I live. I live to bury myself. I buried myself, and she ate me.

It may seem paradoxical, but at the mouth of the dissolving solution was
the fact that, in my depotentiation, I had begun to be decorated by
features that did the most toward making me interesting to the average
consumer in the age of space exploration.\footnote{Moore and Nicholas.}
I had been fashioned, for instance, with wingspan and tiny mythic hands
for roosting, armed with light explosives, beneath the eaves of the
uncooperative. Through the benevolence of chemical manufacturing
concerns, I'd been given an eerie dripping sound, the seepage of fluids,
and manifest changes in odor, flavor and appearance. I was
steam-jacketed, coagulated, and encased. I wore like condensation the
impervious rock of the crown, the anklets, the epaulets, the girdle, the
napkin, the national brassiere.

How did I coil clean in the patriotic wake? Like an ant with a slipknot
sucking its midriff after the hoopla of threading a shell. No one was
paying attention. First I dreamt I was an old bull-cow with horns,
knocked on the head, head rent like a pod releasing, who nevertheless,
as the shackle embraced her, leapt and went running down Rte 17, who was
shot in the back and dragged by a tractor back to be quartered and
plucked. When I awoke, I went to the oak trees and began to root at
their faults. If I'm to be a stuck pig, I thought, liken my final supper
to a sumptuous mould.

It was the me my father neglected to feign to give. It was the me of my
own escape, the wild boar at the bars. The capybara, the kangaroo, the
camel, alpaca, the moose. The antelope, the reindeer, guanaco, vicuña,
the yak. I rumbled amidst the exotic meats, a pilgrim on a bender for
genesis. Numerous eminences and depressions later, swirling like tea in
a cup, I tripped off my stirrups and creepthed upon the earth, the me of
the lump, so that when the goats of the underground breathed through my
ears at a trillionth a watt per square meter, I heard them as trumpets
wielding a knife, my ears bloated like eggs.

I scraped the many facets of my body off the pulpit, sorted and filed my
serrated edges, and took to the beckoning maw. The goats of the
underground waggled their beards. I straddled the glamorous threshold. I
was sandwiched between slices of contaminated hide; lamp, hard hat, chin
strap, sling, coveralls, knee pads, grapple. My wholesomeness was
immediately sniffed at, the nostrils of judgement flying. I reeked of
formaldehyde, hubris and want. Barred as I was at the gates of immersion
by fossils who chewed at my sutures, there I swelled, an acquisitive
cloud, resoundingly willing to bargain because:

My ear had become a compass; on the guillotine thump of the goat song, I
could navigate the azimuth of a note. It was the cavel of the muscle
that cinched and gave about her mouth. It was the augury of the slap and
the swallow of her heart. It was a fracture screwing its mantical smoke
screen up her throat. It was Cassandra in her stomach, the
thousand-petal sibyl of her scalp; it was the soothsaying foci of the
curve of her palms, her fingertips the Chaldeans, her elbows twisting
weathercocks, her breasts portentous gyres, her legs long scrolls of
omen, and her cunt prophetic ziggurat, repeating and repeating a clause
of joint and parting. It was the structure of the tragic and the
anti-tragic fusing. It was the future, and I craved it like a robber
craves a crypt.

Sarcophagus, a tin vault containing two essential pieces, sacro meaning
"flesh" in Greek, and phagus meaning "eater." But from flesh to eater we
travel thousands of miles through the dross. At seven hungry borders
bribe the seven crooked cops.

I offered up bravura, and the greedy bones bit in. I left them nibbling
the image I'd assembled in the mirror. The goats and I grew rubber
souled and sloshed into the vibrant air. Beyond the vestibule, the
atmosphere was profligate with battle. Draperies and canopies and
columns fell apart. Across the continents of spore and bug, an oracle
came warning. I heard prediction sweeping with her tongue across her
teeth.

Disaster is a sexed up promise, and I would leach the lodestone from my
psyche for a hit. As it happened, there were sledgehammers, forceps,
aspirators, clamps. She vandalled questions from inside herself and wore
them like a lab coat. In the autopsy of accusation, stones of gall and
truth. I hemmorhaged quietude, certitude, poise, balance, clarity,
safety, circumspection, humor, confidence, ken. The floor was slick with
my giving way. The goats of the underground donned mourning suit and
veil. There was to be a feast. Bacterium shuffled in.

I'd come in making choice and cut, a bundle of nerve and muscle. I was
sawed and reapportioned, matter rendered moot. Leg, loin, blade, rump,
topside, sirloin, brisket, chuck, ribs and clod and sticking. We fucked
alot, so fucking was a thing that I had left, but in the gut room zero
gravity, the gesture lost its heft. I took on strangle, lockjaw,
blackleg, taint. She was back-flood and breakdown; she was lethal. I was
hooked.

Bacon: a feature in a limestone cave that bleeds in colored streaks. The
name suggests both kill and cure, the double helix tourniquet that binds
the human palate. She was starving, and she loathed and longed for being
in my meat as when the slaughterhouse employee says,

Stick that bitch. I grabbed her and flipped her over. She looked up at
me. It was like she was saying, ``Yeah, I know it's your job, do it.''
That was the first time I ever looked into a live hog's eyes. And I
stuck her. And she bled,\footnote{Gail A. Eisnitz.
  \textsc{Slaughterhouse: The Shocking Story of Greed, Neglect and
  Inhumane Treatment Inside the U.S. Meat Industry}. Amherst, NY:
  Prometheus Books. 1997.}

and he means that it is intimate, more intimate than family, more
intimate than war. Despite the husbandry of scholars, we were not
sisters, were not enemies, not aspects of one self. We were lovers which
means bloodthirst, which means eaten and to eat. It is desire carves the
wild body into savor and into spoil. It is desire and disgust, that is.
But the dissection of the lover into delicious and grotesque is plastic
surgery through metaphor after the metonymic accident of lust. It is a
reconstruction full of mercilessness and mercy all at once.

Hung up, I turned the green of badness and of newness. But when you're
eaten raw and rotten, you sicken those who eat. Rumor has it Sin became
my out; but family matters less to him than dirt beneath his thumbnails.
The truth is that she spit me up. The siren she's remembered for was
bellyache, not birth. She ate me; I was hard to eat. I bolted from her
stomach, her feral burden sprung.

Above us, kingships moved upstream, regions weakened under struggle.
Surpluses were stored and spent, craftsmen fashioned temples,
independent cities rose; mosaics, altars, frescoes; cosmogony appeared
and faltered, priests were fed and housed and crucified; sheep were bred
and shorn and slaughtered, society collapsed. Below you, she was doubled
with the loss of too rich a meal. I came up flopping wet and naked and,
in some sense I was healed, of the rigor mortis Sin had put me in, the
hunger of the daughter.

So, who is eaten, then? Who is eater? Bad blood fills the stomach with a
host of famished stowaways. And my memories are ships adrift and whirled
in the eddy. The log books streak and bleed and run, and hard-shelled,
trap-doored things---the murex and the triton and the gemtop and the
cone---make meals upon the pages, growing plump, in turn, to make a meal
for a scoundrel faking at god. I said nothing is not debatable. Tiresias
long sailed from here, her cruiseship sunk and done, I tell stories that
she dictates from the hulls of former living. But the bowline of her
narrative in my mouth's a fraying twist. And Tiresias a suburbanite of
the artifice of up-and-up, this turn of gossip from the underground is
mine to rip and knot.

I know one more yarn, and it is this:

The warrior pin inks his face a feast of luscious coil. At death, he
meets a hungry hag who feeds upon his facial labyrinth. Sated, she
strokes his eyes and gives him the gift of the vision of the spirits. In
this way, he can navigate the landscape of the dead. If his face is
smooth and left untouched, if he has no mark of carving, like a vulture
she eats his eyes and picks his sockets empty. He's left to wander in
forever, blinded like Tiresias in the myth.

I said she ate me, though she was no hag, and I am not a warrior. These
winding sheets I give you as a gift from past the border between self
and dissolution, between semblance and disorder, between want and what's
leftover, between killing and the cure is the tale of simply sitting
through the pains of being dressed by life, for life and life some more.
For the lifting of the skin from bone is cursive to the examined body's
page where we write scripts for useful ruptures fed on ardor and on
rage.

\begin{center}\includegraphics{images/"pandanus.png"}\end{center}

\hypertarget{taylor-brady-the-onset-exits-and-the-damage-done}{%
\chapter{Taylor Brady: The Onset Exits and the Damage
Done}\label{taylor-brady-the-onset-exits-and-the-damage-done}}

~

Once upon a time, you could say anything nice, but you didn't say
anything at all. Diacritical ultimatums interrupted dinner with
definitions withheld from view so as to grow huge and ravenous in the
dark. I am making this up, of course, as others made it up then, and
thus the accusation leveled by a doctor and anatomist of fact named
Duplex, had it been voiced at that time, would resonate still and call
me to account, but since after all I \emph{am} the one who makes this
up, and since I did, and since what is done is done in that image, it
turned out that the good doctor had mortified his voice\footnote{Once
  upon a time a cloud of sobbing strings pretended to protect the figure
  of a woman in what was either a song or a landscape, hovering above
  her and around her like moonlight and water, but only in the sense
  that moonlight serves as the medium for speculation on nuclear fission
  ninety million miles away in the opposite direction, and water acts as
  a lens, focusing what it covers into a disclosure adding the
  picturesque to the particular. There had been an episode of some kind,
  starvation had loomed, she had been forced, not without some ironic
  pleasure, it's true, to eat the still-squirming flesh of an
  Intellectual in Red Bermuda Shorts, among other things. Some of these
  things were clear, while others remained titular. Under the titles
  were texts, lyrics she was engaged to write so that a man named Birdy,
  no longer young at all, but who saw himself now as he had been ten or
  fifteen years ago, solitary and intransigent and correct in a concrete
  box that housed a single room, plus toilet and kitchenette, might
  within those texts produce disengagements by which the letters would
  emerge as black designs against white paper, and behind them, whether
  reading or hearing them sung accompanied by guitar, the people of this
  landscape might then intuit the presence of himself and his Apostrophe
  - for thus, tropically, did he continue to think of her - she
  beautiful and he evil, or vice versa, or all one way or the other, so
  long as there was some sense of distribution implicit in their roles.
  He identified the luxury sedan in the commercial whose jingle made him
  think of her, and bought it with the last of his royalties from a
  previously pilfered tune, and named it As Such, driving it around
  everywhere, even on rainy indoor days, so that the city and indeed the
  nation might fondle the exciting and aerodynamic contour of everything
  that was surging within him precisely as it occurred. \emph{And the
  pleasure of being a lover, of living by love alone, was enhanced in
  his eyes, as a dilettante of intangible sensations, by the price he
  was paying for it}{[}\^{}brady\_endnote{]}.}1{]}(brady\_note1.html)\^{}
as an act of contrition for the grossness of a speech impediment as of
that evening still entirely prospective, and so no accusation was made,
though a list of sorts was made up. Have I proven my spontaneity to you,
and do you feel free of the memory of that night, reading over the
transcription? Have you swallowed your tongue, or mine, and can I make
it up to you?

\hypertarget{the-ensemble-exists}{%
\section*{The Ensemble Exists}\label{the-ensemble-exists}}
\addcontentsline{toc}{section}{The Ensemble Exists}

\begin{verse}
The Ensemble Exists\\
beginning microcosm\\
— Cecil Taylor,  liner notes to \emph{Air Above Mountains}\\
\end{verse}

\begin{verse}
it would have to be lost so thoroughly as to breathe\\
— Leslie Scalapino, \emph{New Time}\\
\end{verse}

\begin{verse}
You say you have this\\
talent for the said.\\
To rest in that is\\
true—as in default\\
of sense beneath whose\\
crust run faults, cracks\\
full of serious,\\
obvious wet.\\
The habit of that subject is the plan\\
of a beyond which values a succession.\\
However all and each may speak of chance\\
employment, every character is one\\
and only one. Two or more at most. Here\\
we have a logic of the subject very\\
frequently violated. Our one man\\
(or two or more) for the occasion,\\
Dr. Lyon Playfair, for example,\\
takes gravity to specify each toe\\
on which he stands alert, a point \emph{en pointe.}\\
Hold that representation. Admit only\\
the material as in reasoning only\\
to the real, as it really is. And there it is,\\
belief in metaphysical fictions,\\
        which  means there must be\\
        beneath this thing\\
        incited contour.\\
Or too late, if semiosis\\
will not be which sign is used.\\
Now such hypothesis seems\\
any number of reasons to be\\
recommended to the foldwise\\
irritation of the fissile ones\\
and twos who would pursue my\\
fondness for a measured stroll.\\
        I study the varieties\\
        of humanity that move me\\
        and still retain the element.\\
        Even a least, peripheral\\
        sensation of pleasure, of\\
        pressure, tongue sliding over\\
        fissure, has its excitation\\
        and molecular gridlock.\\
The principle with which one sets\\
this weight upon a true continuous\\
plumb, a line of chalked descent,\\
    is named without appeal,\\
    trafficking in outlines\\
    on the bare, black asphalt,\\
    just inside the \emph{simply there.}\\
\end{verse}

~

\hypertarget{the-onset-exits-and-the-damage-done}{%
\section*{The Onset Exits and the Damage
Done}\label{the-onset-exits-and-the-damage-done}}
\addcontentsline{toc}{section}{The Onset Exits and the Damage Done}

\begin{verse}
…to treat philosophy itself as \emph{impossible}.\\
— Louis Althusser, \emph{For Marx}\\
\end{verse}

\begin{verse}
One of these good old days, I’m gonna call your name.\\
— Al Green\\
\end{verse}

\begin{verse}
Though made however visible, the signal\\
success becomes an end for putting up\\
one’s evidence. The result is what sagged\\
thick beneath the contrary pole, sort of handed\\
down from the dishonor attaching to\\
productive early consequences of\\
your ownership. And so it is to be\\
"community" no longer but a stain,\\
dependent for its livelihood on yes.\\
\end{verse}

\begin{verse}
    Form of wasted effort\\
\end{verse}

\begin{verse}
    Apparatus evident\\
\end{verse}

\begin{verse}
    In evidence rendered\\
\end{verse}

\begin{verse}
    Leisure in "housewifing"\\
\end{verse}

\begin{verse}
    Rendered as a verb\\
\end{verse}

\begin{verse}
  to some few cheap false goods conspicuously\\
  outset in philosophy. The floorplan,\\
  the voyage to the backing of the mirror\\
  and back to front the openness of this\\
  enclosure. Its cross-eyed\\
  binocular\\
  windowing of options from the desktop\\
  horizoning the line of sight behind\\
\end{verse}

\begin{verse}
    the head.\\
    The drudged\\
    caption floats\\
    to title\\
\end{verse}

\begin{verse}
      justified columns that\\
      direct each direct\\
      expression of classed\\
      schemata, bent space.\\
      Past life replanted\\
      in the yard that\\
      hasn’t time for second\\
      planted feet to take\\
      the straddled upright measure of ruled, unruly stance.\\
\end{verse}

\begin{verse}
A certain number bend their energies\\
to shape, living up the walls. *Time enough\\
for severance—*give that repetitive\\
stress injury a final grunt by way\\
of compensation. When you think\\
to leave, your key drops through the slot, like mail,\\
stamps you with the ring of self-address.\\
\end{verse}

\begin{center}\includegraphics{images/"pandanus.png"}\end{center}

Once upon a time, there was a storyteller who introduced into a tale
\emph{one of those fragments of literal truth which liars, caught off
guard, console themselves by introducing into the composition of the
falsehood which they have to invent, thinking that it can be safely
incorporated and will lend the whole story an air of verisimilitude.}
Apostrophe runs away, as geological layers of forwarding addresses turn
the multiply-enveloped scene of the narrative into a palimpsest, with
reference to which the only certainty is that someone will receive a
paper cut rather than useful knowledge. However, in another tense her
face pressed out from behind the paper as from an interior, all
velvet-papered walls and lamplight, which the storyteller thought to
foreclose every time the story drew the curtain on a scene, the contour
of nose and cheekbone solidifying for a second in the soft billows of
that gauzy fabric which served only to insure that an audience's
interest remained with what the tale hoped to leave behind, in its jerky
gait stumbling from episode to episode.

Again, though, I am lying, as she draws the curtain back and lets the
light fall upon you there, a bit below her in the street outside, where
all presuppositions stand together in first and second persons, and
says, ``I am not a bad habit, not on the order of the forces 'driving
the urge to differentiate, to identify one and then two,' only for you
to make the subsequent half-turn into saying that the latter follows the
former, as misdirected speech follows curiosity about a lighted window
in a cinderblock rental duplex on a dark night with all the sodium vapor
lamps shattered, duping an accomplice into accompanying you to that
theatrical''bad neighborhood" in which the passing audience method-acts
its way into belief in the constitutive powers of its own direct address
to the emptiness of its relation to the stage properties, when in fact I
have been here, and have been telling the story of this circle of light,
all along."

\begin{center}\includegraphics{images/"pandanus.png"}\end{center}

Once upon a time, there was a man named Birdy, older then than now,
whose job was to remind a text on memory of where it had been and why it
had missed itself there, as if he were a third-person argument between a
first and a second person who tried to keep appointments, but whose
lives, in the event, were compromised by one emergency after another, so
that one or the other always got stuck waiting, picking up the check,
taking the wrong bus home.

``Otherwise," he said, coming to his senses already in an anaphoric mood
on the curled linoleum floor, decorated in a grease-beclouded floral
pattern, of a dimly familiar kitchenette, talking to no one in
particular, or else an offstage presence stuck in an abandoned
episode,''you'd simply live the singer-songwriter myth, and have nothing
to say for yourself but a lot of florid weeping and moaning. I knocked a
little hole in your wall because I care. Pornography drives the
photograph forward into full relation to its time, and becomes a method
to produce clear prospects on what was once obscure. I continue to love
your body, the taste of your tongue and sweaty ass, only because I
simply can't imagine you, and see not you, but a broad expanse of sky,
on the far side of every window into which I peep."

\begin{center}\includegraphics{images/"pandanus.png"}\end{center}

Once upon a time there was a man named Birdy, somewhat younger than he
had been in an episode just a moment previous, who gave the lie to
epistemology and remembered the exact contour of every pain his body had
ever suffered through, beaten into shape, even, perhaps especially, when
his mind had been on something else. Sitting in the dark and hearing,
from a door or two down the block, soft blowing on soprano sax, he
experienced someone else's clarity out of phase.

\emph{"I was afraid he was going to eat me,"} a particularly
disingenuous character in the family drama was just then reciting, in
rehearsal with the cop with whom, arm in arm, he or she was coming up
the walkway of half-buried railroad ties that kept the path from parking
slab to painted steel front door relatively dry and passable during the
summer floods. Of course, since there was no air conditioner, Birdy had
propped the window open earlier in the evening with a stack of crank
pamphlet literature, and so heard every word, but refused on principle
to believe in the reality of any of what was said under such
circumstances, and was maintained in his condition of exemplary,
self-consuming skepticism. By the same token, though his skull was
perforated as proof of the most certain kind, he could not grant the
veracity of reports that told of a mutant, cannibal freak lying just
beneath the surface of the duplex and chewing on the sleeper's head
during the sticky night, but maintained that this Mr. Article was, if
anything, a fiction of grammar granted temporary substance by the
growing thickness of ideological descriptions of landscape. Meanwhile
the sax was licking lascivious promises into his ear from distant,
microscopic horizons, so that, when his one hand found his other hand
scrabbling around his lap in the dark, his first impulse was to turn on
the overhead bulb and find a window envelope, write a check, and mail
away his first adjusted payment.

\begin{center}\includegraphics{images/"pandanus.png"}\end{center}

Once upon a time, my little bird, you were saying. Then the\\
therapeutic architecture, its medical practice sprayed thick over\\
walls and fixtures alike in anticipation of each new tenant, opened\\
wide in a tiny, repeated gasp of cancellation, swallowing your\\
evidence in this incredulous and barely articulate displacement of\\
pitch and timbre by pure, iterated rhythm, as if it were a quiet\\
conflict approaching silence, becoming a branded product, and we\\
promptly forgot it all, and had to wait to recall you as an\\
after-effect of our own belated recognition of a housing crisis.

\emph{Once upon a time, you could drive through the city stoned and
create a makeshift eclogue out of that fuzz of vague participation, much
as earlier drivers had created their pastoral investments with property
in rolling parkland. In the windows which were not plywood painted to
resemble glass, but behind which were walls painted to resemble
all-encompassing skies, people would perform the parlor drama of their
social obligation to your passing. Going through someone's trash, or the
stack of papers on someone's desk, or the pages inserted by a printer's
error into someone's typescript, was a hot and dirty intimacy. The
substance of a whole affair might take root in a similar sense of
euphemism.}

\emph{Ours,}

\emph{The Editors}

Apostrophe, watching it all in slo-mo, said to the unidentified others
sitting next to her on the dilapidated couch in the television's
blue-tinged glare while the cheap VCR hummed softly, ``Here's the part
where he begins to think he's outsmarting - who? not me, even he must
have given up anything but the broadest of analogies in that direction -
but someone, nonetheless, who turns her head from side to side, allowing
his right hand to drip something oily and bitter into the balancing
mechanisms of her inner ear, while his left buys out her lease. \emph{It
becomes possible to cut off the supply of light from a house.} Look out
for the landlord who tells you he sees through the vanity of possession,
for he's about to knock holes in your walls through which later to knock
a hole or two in you.''

``Tell me everything you saw, and what you think it means,"
stage-whispers Grace Kelly in the fascinated hush that cuts off the
supply of light, her absent presence in this vignette responding to your
recitation, which takes up each sentence in the discarded papers blowing
around against the all-encompassing sky and makes of it an opening whose
dimensions are given in the aperture of the vocative''O." Framed thus in
the round zoom lens of surveillance, do you now see the dashing figure
you cut from afar brought into the tightest close-up, made ugly and
atomic?

\hypertarget{nicole-brossard-soft-links}{%
\chapter{Nicole Brossard: SOFT LINKS}\label{nicole-brossard-soft-links}}

\lettrine[lines=3, findent=0em, nindent=0.1em, lhang=0]{T}{hese are the names}
of places, of cities, of climates that make characters. The clear
mornings, a fine rain fallen for twenty-four hours, rare images
originating in America and elsewhere, two natural disasters that oblige
us to stick together in the midst of cadavers, these are the gestures
quiet or purple, the shells, the icicles in the happy hour glasses,
crockery noises or a light stuttering that torments an instant, a slap,
a kiss these are the names of cities like Venice or Reading, Tongue and
Pueblo, the names of characters Fabrice, Laure or Emma. The words
sharpened over the years and the novels, words one spoke while breathing
badly, while laughing, while spitting, while sucking an olive, the verbs
we add to the lips' pleasure, to success, to a certain death. These are
words like knee or cheek or still others stretching as far as the eye
can see making us lean over the void, stretch like cats in the morning
these are the words that make one stay awake till dawn or take a taxi on
week nights when the city falls asleep before midnight and solitude
sticks in our jaws like an abscess. These are words spoken from memory,
from want or from pride, very often words pronounced with love while
placing hands on nape of neck or filling a glass of port. These are
words whose etymology must be sought, that must then be plastered to
what's called a wall of sound, in a manner like those who cry out in
pain and sigh with pleasure that wander in dreams and documents assault
the heart's mysterious obscurity. These are words like bay, hill, wadi,
via, street\footnote{These words appear in English in the original
  text.*} stasse dispersed through the dictionary between flame-trees
and neons, cemeteries, dismal and forests. These are words sound the
body of meanings that are claws or soft\footnote{These words appear in
  English in the original text.*} on our chests, cold, shivers, furrows
and fear in the back without waiting while we try to split the sleek
future tense with trenchant quotations. These are words swallowers of
fire and life, one no longer knows whether they're Latin, French,
Italian, Sanskrit, Mandarin, Andalusian, Arabic or English, whether they
hide a number, an animal or old anguishes eager to gush before our eyes
like cloned shadows filled with light and great myths.

Translated by Peter Dubé

\hypertarget{mary-burger-all-new-yorker-stories}{%
\chapter{Mary Burger: All New Yorker
Stories}\label{mary-burger-all-new-yorker-stories}}

~

\lettrine[lines=3, findent=0em, nindent=0.1em, lhang=0]{A}{ll New Yorker fiction}
pieces stop at the point where the person makes a bad discovery about
himself or herself or the world. That he is or she is a failure
personally---in love, usually, romantic love or familial love---or that
the world is a failure toward his personal or her personal sensitive
nature---that the world is violent, that unequal distribution of power
causes pain and unhappiness, usually to the less powerful, but sometimes
to the powerful as well.

Regardless of its narrator, its characters, its particular conceits or
conflicts, what anchors each story is the sick feeling at the end. The
same feeling that comes after a radiation treatment for cancer. The
queasy realization that all this, the technological sophistication, the
aggressive preservation of human life, is merely its own reward, not a
means to anything.

This is how the melancholic condition of privileged passivity confronts
itself. When a grown man behaves mercilessly towards his humble,
bewildered parents; when a wealthy young Latin American woman is
threatened with kidnapping and even death, right in front of her own
gated home, right inside her private limousine; when the mostly likable
gay neighbor, in the midst of his theatrical, flamboyant prime, withers
suddenly and agonizingly, and dies. When an urbane, not quite young
writer, supplementing her New York lifestyle with a teaching job in the
heartland, falls in love with a simple, decent man, neither a New Yorker
nor a writer nor even in any way ironic---and simply leaves him, his
marriage proposal, his simple, open face, to return to her lair in the
crowded city, her archly urban self in the witheringly sophisticated
intellectual world.

In "withering," a possibility that what is socially belittling might
also be personally devastating. That the intricate architecture of
social standing might collapse occasionally, like a poorly braced studio
set, to reveal something more like an experience arrived at through
contemplation.

When the modestly privileged, moderately young childless couple who
attend the adoption picnic focus on the same large-eyed, attractive kid
(that is, the querulous wife focuses. The ironic husband, whose
disavowal of the thing expected of him mimics Kevin Spacey's in
"American Beauty," instead alights for awhile with an aloof, jaded
teenager, whose unflinching lack of delusion contributes to the reader's
own worldliness), that is, the moderately young, uncomfortably nervous
wife alights on the same doe-eyed little boy attended to by an imposing
Texas judge and his Barbara Bush-edition wife. It is an orphan fair.
Like all bazaars, the most attractive goods are fought for, the least
attractive are left behind to grow even dingier.

In this version, the judge and his wife are the cunning players. They
are powerful, and they are malevolent. The judge and his wife invite the
more modest couple to their lair, a large house on a large piece of land
by a lake.

With "lair," variously, the idea that the seeming impotence of the
self-conscious intellectual (the New York writer snagged in the
loneliness of her cleverness, the cleverness of her loneliness), the
rueful split between the life of the heart and the life of the mind, is
a familiar condition, the condition that seems to be most easily
recognized as a sign of distress, of all not being right with the world.
The dissonance between what you know and what you need---this is the
price charged for entry to the club. And who complains? Does a young
gymnast object that her muscle-strapped torso and thighs set her apart
from other girls? The world-weariness might be counted as a privilege;
anyway it's displayed often enough to be interpreted that way.

Through the battlefield drills of drinks, dinner, after-dinner
segregation of husbands and wives, the older couple's battle plan
emerges. The Texas wife begins an oblique attack, a narrative of a
foster child they'd had, a black toddler taken from an addicted mother.
His mounting health problems and learning disabilities, together with
the difficulty in keeping a nanny---when enough details had been laid
out, the concluding sentence wasn't needed.

The off-register comments from the Texas wife about her foster child,
about finding the boy a nanny ``of his own kind" and a school with''all
different races, the Mexicans and the Chinese and the Indians and all
that," created another clear space in the story, a way to tell what is
right.

In the kitchen, away from the women, the judge's assault was direct: he
confronted the husband with ammunition gathered from his labyrinthine
legal connections: evidence of the husband's expulsion from college 15
years ago for selling pot, evidence that the judge, with a final
flourish of legal muscles, had already inserted into the husband's case
file at the adoption agency. This in the few days since the couples met.

But since the cunning of great power isn't really the story, or enough
of a story, it was necessary for the judge to punch the husband in the
nose and knock him down before the couple, and the story, could leave
the judge's house.

The women's exchanges, happening at the same time, were far less direct
than the men's but resulted in arguably more damage, like a barbed hook
that's slipped smoothly into the skin and then twisted.

It's true too that the young wife and husband wanted different things
(as different as a doe-eyed youngster from a disaffected teen), a
further cause of dissatisfaction.

The obliquity of the wives and the passive irony of the young husband
all form a rueful cognizance that weighs down to the point of paralysis.
And yet it's that ruefulness that cradles its hands around the
narrative's delicate crown jewel.

Like a lead musician the fiction narrator repeats and elaborates a
single sensitive melody line, or crescendos on a particular note of
knowing regret, a middle-aged loss of innocence, though evacuated of any
possibility that the afflicted might be rescued from the cause of their
desolation or even, astonishingly, change their own lives.

This somber performance seems to be the provenance, or the penance, of
the fictional narrator only. For the non-fiction narrator, the feat is
more acrobatic, if not therefore also more precarious.

The non-fiction narrator is permeated by a sense of ambiguity impossible
in her counterpart. While the fiction narrator lives with the certainty
that moral rectitude is real, that extreme power is obscene, the
non-fiction narrator is instead subject to the wild cards, false leads,
and missed signals of a lived life.

The non-fiction narrator tends to have an intimate relationship with
squalor, or danger, or pleasure, an intimacy marked by unguarded
revelations of vulnerability. She exchanges the reassurance of an
intricately self-referencing universe, in which each element contributes
clearly and inevitably to the whole, for rapport with the volatile, the
off-balance, the unpredictable. Her familiarity with squalor becomes its
own sophistication, her fluency in the idiom of disorder, with its
free-jazz syncopation and arhythmia, its own seductiveness. The effect
holds our attention like a sudden, beckoning whistle coming in through
the window of a quiet room.

In this version, the non-fiction narrator is free to both love and
despise her volatile, negligent parents, her alcoholic mother who
disappears on binges and returns home to demand attention and threaten
suicide; she's free to both admire and resent her slightly older, vastly
more sexually experienced sister; she's free to fear and adore the
coltish, sanguine boy she's had a crush on since age six. The
non-fiction narrator elides between bewilderment and passion for her own
still strange sexual feelings and for the obvious, pronounced desire
that persists between her parents even in the midst of their pitched
battles.

While the fictional narrator breathlessly cradles a fragile egg, the
non-fiction narrator darts in and out of the room bouncing a rubber
ball.

In the photo, a snapshot salvaged from one of those summers that were
both vacant and teeming, she is 15. Her long hair blown a little into
her face, head tilted a little to the side, lashes lightly mascarad or
just naturally thick and dark---she is the early 70's icon of innocent
sexiness, the womanly girl, what Ali McGraw, "Love Story," and John
Denver's "Sunshine" all immortalized, but what she seems to be without
trying, maybe without even really knowing.

It's that not knowing, or anyway not always needing to know, that the
non-fiction narrator settles for, but that the fiction narrator battles
like a riptide.

The one bobs in a tight ship just off shore, pulled between the current
and the undertow, while the rest shout and lunge in beach volleyball,
play with children, play with dogs, play radios, nap. She imagines her
voice as faint, or inaudible, from shore. The other shipwrecks in the
shoals and scrambles from her splintered raft. Her escape is always in
progress.

\hypertarget{works-cited-1}{%
\section*{Works Cited}\label{works-cited-1}}
\addcontentsline{toc}{section}{Works Cited}

J. Robert Lennon, "No Life." The New Yorker, Sept.~4, 2000. Pp. 74-81.

Mary Karr, "The Hot Dark." The New Yorker, Sept.~4, 2000. Pp. 42-49.

\hypertarget{margaret-christakos-often-she-finds-butchers-flirtatious}{%
\chapter{Margaret Christakos: Often She Finds Butchers
Flirtatious}\label{margaret-christakos-often-she-finds-butchers-flirtatious}}

\hypertarget{from-charisma}{%
\section*{\texorpdfstring{From
\emph{Charisma}}{From Charisma}}\label{from-charisma}}
\addcontentsline{toc}{section}{From \emph{Charisma}}

\hypertarget{working-notes-1}{%
\subsection*{Working Notes}\label{working-notes-1}}
\addcontentsline{toc}{subsection}{Working Notes}

\emph{Charisma\footnote{\emph{Cool for You}, New York: Soft Skull Press,
  2000} is ``for the chorus of us." I wrote Charisma over a period of
nine years, moving through and enervated by the political prescriptions
of middle-class feminist academia; craning my gaze on women through the
hourglass of my wounded relationship with''mother, sister,
grandmother``; diving through several portals of baby-shit and bliss to
wear the secret subcultural drape of motherhood; its repeated
re-shapings of Normal greeting my sense of''what happens" as a waking
frequency of turgid extremes. I realized that straight novels had always
bored me, why I masturbated to make them more interesting. I would need
fiction to become a foodlike procession, an irrefutable love-surrogate,
imbibed to stimulate the senses. (Going down on me.) I liked a
multi-channel environment, not for its staccato disavowal of the
constant but because I had been raised with an appetite to introject
several stories simultaneously, each contracting and expanding according
to its own pulse.}

\emph{Mothering a boy and then girl-boy twins often involves being
unable to regulate how much experience I confront; a surfeit of
sentience. An overflow of interaction, so that the negotiated space of
consensual relationship capsizes. As when one is in love with more than
one lover, or dreaming, or in hate, an oxymoronic prostrate agent (who
is) subject to the re-calibration of self-limits and society's clichés
about who/what is "hot." I wanted to write the storiedness of bisexual
subjectivity (even tho maybe rejected by queer command-central) as a
kaleidoscopic inner identity script of positive contradictions and
overflow, and to acclaim cunt-agency for mother subjects. A story that
managed to open gender itself like a cocktail umbrella, with the notion
of the gendered character as a laquered overarch of spokes. As persona
spoken and collapsed again, efficient and portable to other stories,
ducking in and out of brief theory showers.}

\emph{When you winked your way through the glass we knew it was a cover.
We wondered about ourselves then, about the social organization of
gender categories, about barriers to the heteroerotic--we mean hetero as
in plural, different and manifold--desires we could feel, we wanted you
for our tutors in studies of the self, ours and yours. To tell the story
because it's a-hankering to be told. The novel was more cyclonic in form
before the final edit; I attributed extra value to narrative
contiguities once I had hung out with non-writers and babies for a
while. I wanted the book to fly.}

\hypertarget{cameo}{%
\subsection*{CAMEO}\label{cameo}}
\addcontentsline{toc}{subsection}{CAMEO}

Often she finds butchers flirtatious; out on Danforth Avenue where the
anemic lamb carcasses suspend for three or four seconds her trust in
bodily safety, before the mind rescues her from moral crisis. Just a
piece of meat like any. It's the hook, the hook that hooks, that
perforates the eye, the brain, the guts of sympathy which must be let
down, rinsed and laid in the cedar drawer until autumn. A time when
ghouls are acceptable again and the skeletal dance of baby sheep is
seasonally admired. Dressed in his white blood-smeared coat, he is like
any irresistible and prodigious surgeon. He wears the coat, the lamb's
clotted juices and pristine plate glass. He sees the world through
hooked sides of things. Why be afraid?

Flirtation has to do with how the lambs are led along the corridor of
doom. How they go and go. How they continue up the ramp while their
cousins shriek. It has nothing to do with Easter. With men and beasts,
rather, and strangers who may be beastly though their flirting may
become on occasion beautiful. The natural frame of the plate glass
composes him in savoury labour, up to the wristwatch in precision. He
peeks through the space where the organs were to a woman's tenuous
shock. Hooked by the blood of someone else. Thing else. Some other
thing‚s public juice. And the fixed sway of the dead meat muscle---Cameo
is a sitting duck.

Cameo coaxes the lamb chops from their wax paper sleeve and lays them on
sizzling garlic. She scoops mint jelly into one of her mom's old
teacups, and amuses herself while the meat goes bloody by holding the
porcelain overhead---there they are, as delightful as when she was a
girl, green shadows flickering through the cup's delicate filigree.
Impatient to eat, she flips the chops, sniffing in their high corduroy
bloom, flexing her knees to an old disco song. Gradually the element's
blue flame singes her dreamy gaze and she fidgets with her too-tight bra
strap, then jets out her arms at right angles and croons into the
kitchen window, hopping from one foot to the other and pulsing her hips.
Soon, the smell is so velvety she is near euphoric. She unloads the
shining chops onto a plate, spoons coleslaw on top of the steaming jus
puddle that runs off them, selects her favourite bone-handled cutlery,
and sets these, a napkin and a beer on a silver tray. As she passes the
hall mirror, she transfigures the carefree, dimpled smile on her face to
a farcically enamelled grin, and, childlike, the round-jowled reflection
catches itself awkwardly, averting its eyes, a dinner guest of inferior
status. Well, fine. She just wants to fade in front of the McNeil-Lehrer
Report while reading Interview. Clarence Thomas is deplorable and River
Phoenix has siblings named Rain, Liberty and Summer, and Keanu Reeves is
part Hawaiian---so that's it---explaining his preternatural gaze. She
thinks about licking his eyelids. She doesn't mind such overlays,
allowing that the real Keen and his big-screen simulacra have little in
common. The pastoral fogs up and loses its colour when she remembers
Keanu went to the same high school as an ex-lover's sister, Korona, the
mesmerically sirenic just-past-teenagehood dyke who moves like a rock
stud at the Marzipan Room. Same eyes. She considers herself curled up in
a wet tongue licking those eyelids. Then the newscaster misses a beat
describing a murder-suicide in which the ex-boyfriend used two
industrial meat hooks to set up the nooses. She shuffles into the
kitchen, unloads pink bones into the garbage, rinses her plate and
balances it in the dish rack. Humming an unforgettable TV ad promoting
eternal brand loyalty, she flicks the one-cup switch on the espresso
machine and reaches for pungent cinnamon and powdery, bittersweet cocoa.

\begin{center}\includegraphics{images/"pandanus.png"}\end{center}

Cameo goes to the local butcher, a garrulous older Polish woman, because
the baby needs lots of iron and protein. Hardcore nutrition. This meant
meat and potatoes, meat and rice, meat and pasta when Cameo was growing
up. Lamb stew, pot roast, liver and onions, spareribs with ketchup
applied during the last ten minutes to get a braised country look. She
imagines the placental tissue forming, a red sirloin in the epicentre of
her psyche. Like drawings she happened upon at a gallery, of galactic
uteri with trees and dreams growing inside. Already her own dreams make
her anxious. She takes a baby boy home from a supermarket, and somehow
he can already talk. She writhes about rubbing herself while holding him
tucked under one arm, then says, Do you want to lick me? He wags his
tongue, starts slurping at her vulva. After a moment he stops and looks
up, cream all over his soft little chin. She runs the tapwater and says
maternally, Do you want a drink? She tells people who ask that she's
keeping him if no one claims him back. The dream's meaning escapes her
over the cup of coffee she leans into the next morning, before she goes
to the deli where the Polish butcher will wink at her, Ah, look at the
glow on that one, making a big deal, and Cameo will like the invasive
boastfulness of her, the European excess; her motherishness.

All week, in fact, she thinks a wordless blur of foodstuffs. The bakery
woos, Fresh Valentine Cakes Available. Signs like this keep Cameo on
track from day to day, remind her that English is the official language
even for her meandering body's new requirements. When the other butcher,
the young one out on the Danforth, slicks his sweet right eyebrow to the
ceiling, he is flirting in the queen's own tongue, though his first
language is probably Greek. In a fit of hysteria libidinosa he might
mutter in his mother tongue if she lapped his incurvate dimples. Or
would he wish this? Cameo silently paraphrases in pig Latin how he
undresses her in quick swipes, opchay opchay, opchay, why waste such a
precise and professional syntax! Try to make sense, baby-mother, she
chides herself, tenting her shoulders forward to contain the
cross-continental sound loop. A sign propped up in the cooler says,
Fresh Lamb Whole Or Pieces Suit Your Fancy. Speechless again, she sees
that the flirtatious butcher's open freezer is layered with the slight
upside-down leftovers of lambs dreaming the disembodied envelopes of fur
back onto their rib cages. She starts shivering. He says, What'll it be,
miss? With those eyes, he could be her cousin. When Cameo comes here she
dislikes the imbalance, like she's a plate in a dish rack about to
topple, but she does come. And each time: the frightening image of a man
with a knife imposes itself into the hem of her ahem-hawing as she
stands at the pink-smeared counter. For what order to be uttered, or
redress to arise. Sawdust in her nostrils, metal blade at her view's
rim. Always at the edges his hand the wrist the blood the appetite.
Sudden. Her eyes close.

\begin{center}\includegraphics{images/"pandanus.png"}\end{center}

Once Cameo had thought about mother's cream in a strange way, a way she
could not after that first time continue to speak of. She began to live
it in code, through images of food, since cream was a food to begin
with. There was no other sort of cream she could allow out of a mother's
body. When she thought of cream or milk her thoughts often turned quite
pleasurably to snow, and then her arms would be swimming in it, folding
it back onto itself like slow motion film images of milk spilling on a
tabletop. Liking this picture, she had always wanted to go swinging in a
snowstorm, but was thwarted by Parks \& Rec's policy of packing the
playground away in late November. Children were known to get overheated
and absent-minded in their snowsuits, and the city could not be
responsible for them burrowing sleepily into snowbanks and freezing to
death. She had almost frozen once, surrendering too long on the walk
home from school. The cold felt like milk being swallowed as it slowly
entered her, and she thought of her mother, paying the babysitter for
the hour of lunchtime she was wasting flatbacked on the bank, thinking
of milk and a creaminess she couldn't place but missed somehow, like
snow angels in summer.

\begin{center}\includegraphics{images/"pandanus.png"}\end{center}

Sometimes Cameo feels like the carcass slung on the silver crook, the
freshly stripped and gutted one mounted in the showcase to attract
discerning customers. When she stops to look through the plate glass she
finds her own face grafted on the plate glass looking back at herself,
soft into the desiring hook of her eye. Perhaps she is the woman she
likes to flirt with, midday, as much as she flirts with her threshold
for abjection. Yes, part of going out into the world each afternoon is
to catch herself in it, raw and available. She slowly learns how to
protect the butcher, who winks and waves from inside the shop, from her
gut hankering for the woman between them.

\begin{center}\includegraphics{images/"pandanus.png"}\end{center}

\hypertarget{marilyn}{%
\subsection*{MARILYN}\label{marilyn}}
\addcontentsline{toc}{subsection}{MARILYN}

Happiness has to do with how the sleep goes and goes away, flood water
evaporating into mist. Cameo's eyes brighten, she pushes upright and
steps into the word Hi. Speaking first in her own person and then in the
echo of her aunt's spritely eyes. Wanna come, sweetums, and meet
Marilyn? Her aunt deftly snaps closed the hand pouches on Cameo's
jumpsuit, wedges each cluster of toes into its rubber foot envelope.
Let's go fast, auntie says. Fast, Cameo doubles, the way she doubles
everything in her aunt's face, her sparkle, the discourse of her desire,
Cameo would say now. Discourse and desire.

That night, the mother and the aunt sit on chairs pulled tight to the
kitchen table. Their elbows grind into the wood surface connecting them.
A glass of milk slides around in its own white spillage. Don't you ever
take Cameo there again. How can you be so pigheaded? Don't you speak
about me, we're talking about you. You have no right to talk about me
like this, Chari, you have no clue. All I know---Is nothing, is pig
shit. Is what is moral. Your morality is a castoff, from Dad, from the
system, from the goddamned church. And what happened to you, how did you
get to be so selfish? Being true to myself isn't selfish. No, Seal, what
you are is confused, and I want you to keep my daughter out of your
selfish confusion. This isn't Cameo's problem, it's yours, can't you
admit it? If anyone here's got a problem, Seal, it's you, end of
conversation.

Cameo backs down the hallway chewing the tough thumbnail she has
detached absent-mindedly. Her mother's curt voice flip-flops like bad
perfume in her senses. But Marilyn gave her two balloons. Aunt Chloris
bought them all ice cream cones, and together they had swung her in
giant leaping steps up the grassy hill. She had been dangling like a
fish, giggling heartily, and the women both guffawed, a fish? You‚re a
whale of a girl, and don't you forget it!

Womens' creamy cream. Their cream. Mama's creaminess. Ooosshhh.

\footnote{\emph{Cool for You}, New York: Soft Skull Press, 2000} Pedlar
Press, Toronto: 2000

\hypertarget{brenda-coultas-failure}{%
\chapter{Brenda Coultas: Failure}\label{brenda-coultas-failure}}

~

\lettrine[lines=3, findent=0em, nindent=0.1em, lhang=0]{I}{’m a failed short story writer}
in the traditional sense. I write the way I write because I have no
choice. I wish I could write in a traditional narrative shape (plot,
characters, conflict), I mean that I don't do characters that begin to
talk and speak as independent entities with free will. I've always been
attracted to language more than plot and character. And I hate most
fiction. I hate the whole artificial structure of popular fiction yet
love artificial elevated language. My last attempt at straight fiction
has left me stuck with 6 pages of notes about Southern Indiana carnival
life. But every once in a while I fall in love enough to keep going.

I focus on sentences and images. I like to describe. I'm most influenced
by documentary film and photographic essays at the moment, and taking a
cue from visual artists and piling up a lot of shit (dumping memories,
images, found objects into a journal), then sculpting it for a shape. I
use narrative to connect, also I'm a sucker for a narrative riff and for
beauty. I'm called a poet and prose writer and I'm at home with both
titles: however, my main company consists of poets and poetry is a large
chunk of my literary diet. Of course, I follow the proponents of
narrativity, whose writers are working out the same issues.

My latest writing project focuses on the Bowery, one block away from my
home, specifically the brief section between Cooper Union and Houston,
an area that contains the remnants of SRO hotels and the remains of the
1890s Bowery that are slated to be demolished by The Bowery Development
Plan in the next decade. The artist coop that used to be McGurk's
Suicide Hall, named so because prostitutes used to fling themselves out
the windows, the Sunshine Hotel and various soup kitchens will be
extinct. The project is centered on documenting and reacting to the
layers of debris, including the human kind that layer the streets. The
intention is not to romanticize the suffering or demonize the Bowery
residents but rather to comment on poverty, class, suffering and my own
dilemma and identifications as a teacher and poet one paycheck away from
the street. It's the transparent medium I walk through with my own
poverty. Needless to say that for me, the Bowery has taken on a
metaphysical weight as a passage and frame of mind as well as a power
spot where ghosts and the nearly living compete for space with the cell
phonies who have come to replace them.

I have failed at traditional forms but I suspect that in the long run
those forms have failed me. Maybe, maybe not. Each writer must make and
break her own rules. As for now, I'm at work observing, discovering
obsessions, and devising a container to carry it all in.

\hypertarget{aja-couchois-duncan-however}{%
\chapter{Aja Couchois Duncan:
However}\label{aja-couchois-duncan-however}}

\hypertarget{notes-on-narrativity---writing-in-prose}{%
\section*{Notes on Narrativity - Writing in
Prose}\label{notes-on-narrativity---writing-in-prose}}
\addcontentsline{toc}{section}{Notes on Narrativity - Writing in Prose}

\lettrine[lines=3, findent=0em, nindent=0.1em, lhang=0]{P}{rior to my life as a writer,}
I was a student of anthropology. I still do not know exactly what this
means, other than it is printed on my college diploma. I do know that I
am obsessed with what makes us human and how this human configuration
struggles to differentiate itself from animal. This study of human is of
course a study of self and other, a study of self in family, a study of
family as self. However is about family, about the process of aging and
the injustices this culture (as if a singular referent could actually
take on so much meaning) inflicts on the old. However is an ethnography
of the last few weeks of my grandmother's life and her shouting right
back at me - you got it wrong and right, a little right. However is in
prose because this process of description is too messy for a pretty
mobile of string and glass. However is in prose because prose is the
viscous tissue of body and there is a lot of body, its excrement, in
this piece. However is in prose because, as James Clifford wrote in the
Predicament of Culture, ``participant observation obliges its
practitioners to experience, at a bodily as well as intellectual level,
the vicissitudes of translation.'' However is in prose because there is
too much, rather than too little, to say.

\begin{center}\includegraphics{images/"pandanus.png"}\end{center}

Not accounting for everything---who could---not just the alcohol and the
iron deficiency, but other failures as well. What happened exactly. The
itchy warmth of the car seat. The light from the house across the way.
My own need to itch and scratch and be done with it. Of course this body
no longer responds, I am earlier and inside taut skin, two breasts. Even
the car is not liquid of the body which fed them, taught them how to
walk into a nursing home and sign me up for a cruise around the parking
lot, my feet above the wheels of the chair. As if a wheelchair could
contain all that is inside of me. Such structures remember nothing and I
remember everything that transpired. It is not the recollecting but how
to frame it in what they know.

I know they know nothing.

Real-estate is not a sure thing, I told them after the fires, but they
forgot. So now almost paralyzed and past speech I can only watch them
arrange things, moving the couch, replacing the tattered cloth. Home is
after all only a settling of dust. I cannot move, so dust settles.
Settling. They put me near the picture, underneath it, an awkward
likeness. Pointy nose, unfriendly hair. My hair, like many things, is no
longer mine. A purple haze about my head. How they frame me.

The towel is rough and the shower too hot. Each time it is a different
nurse. Mrs.~Gagnon, they scold when my legs buckle or a rebellious ankle
twists. They set me on the toilet to rest, each of us exhausted by the
eighty-three pounds of yield. The bathroom so unlike my own and still I
get lost between them. The girls think the water angry and the grass
along the yard just weeds. Everything was planted with his own hands,
large scaly trout even though he was chicken legs and a crows fat beak.
Mom, they say and try to bend my elbow outside its range. They are
making up, or evening out. Those years I yelled and complained and
wished them unhappy marriages and thankless children. No need to wish
such things, every one of them cries themselves to sleep. Perhaps only
now. Or after.

Clutching the pillow, my name across the breast. They know me as my
order in their birth, first generation, narrow fingers scaring an uneven
yellow embroidery, grandma it would say if it were legible, granbra or
jomda. The girls hold my arms, leather handles on the suitcases they
carry, satisfied with my containment. A body folded into an origami
crane and stuffed into the front seat of a car, my car, now theirs. See
mom, just a few blocks, not that different. They think my silence not
thought or post-thought as if counting was beyond a brain once stroked.
Each distance the length of my tongue and I know they are lying, driving
me outside my own frequency so even the hairdresser will be lost and
cannot mix the exact shade of dusk that is my hair. What injury can
touch these few blocks, my home since before he died and I roamed them
like a doll house, my hooked body fitting perfectly through each door.

Whose life so easily told as words. They think I am past me or they hold
me in their earliest versions so I am no longer part of description, an
obstacle even to the sight of family. Perhaps I am these trees along the
freeway dusty with carbon monoxide or the eucalyptus behind the nursing
home shedding. The body takes much longer to break free of limbs,
forgetting my other daughter's name

~

\hypertarget{cory-frost-its-bits-world}{%
\chapter{Cory Frost: It's Bits World}\label{cory-frost-its-bits-world}}

~

\hypertarget{working-notes-2}{%
\section*{Working Notes}\label{working-notes-2}}
\addcontentsline{toc}{section}{Working Notes}

\lettrine[lines=3, findent=0em, nindent=0.1em, lhang=0]{T}{he primary idea here}
is that we have to forget where ideas come from in order to find them
interesting again. Non-English-speaking poets of English will take a
text-could be a phone book, could be a love letter, could be the
instruction manual for a pocket fisherman-and create astonishing
juxtapositions which in turn create narrative fireworks in the mind of
an ordinary, pre-programmed English-speaking reader like me. ``A
comfortable twig for the heroine in the town. I feel more refreshed than
anyone. My daughter has arrived at puberty.'' Their ability to do this
rests on their inability to remember or even understand the original
context.

When the pervasive, mass-produced poetry of (culturally and economically
Imperialist) jingles, slogans, and catchphrases is read through a filter
of foreign-ness-either cultural or linguistic-what results is something
infinitely more imaginative and disruptive. These mutant words and
phrases are the visible residues of desire, in that they are produced by
the conjunction of the foreign-language longing for the power of
English, and the longing of the English-speaking reader for the novelty
of otherness. The name ``Bits World'' is derived from the name of a
Japanese stationery brand, ``Bits Goods": notebooks bearing slogans such
as "I'm an useful and enjoyable stationery. Let's get along with me!"
The name also describes the writing/performance techniques I've been
experimenting with. (I sometimes use the more academic-sounding
moniker''discourse collage.")

In performance, this piece is accompanied by travel slides collected
from second-hand shops: pictures of famous landmarks (Saint Basil's
Cathedral in Moscow), not so famous landmarks (anonymous obelisks on the
side of the highway), people at the beach, road signs in various
languages. There's a shot of a wooden Santa Claus cut-out surrounded by
palm trees, shots taken from buses and trains, shots of totem poles and
castles and flamingoes: typical scenes from the great
20th-Century-Enlightenment project of Tourism. The slides emphasize with
irony that this is not a travelogue or a visit to a picturesque
alterity: the piece attempts rather to remove the distance from travel,
thereby making the exotic obsolete. But the slides are also meant to
distract the audience long enough to make them think they are listening
to a story. Which, I admit, despite the disparate origins of and
inspirations for these sentences, they are.

\begin{center}\includegraphics{images/"pandanus.png"}\end{center}

They were looking for a total transformation of mundane experience into
bliss. It had become necessary. The winter came fast and hard, and it
stayed a long time. The potatoes froze in the ground, there was not
enough wood for the stove, and the modem was too slow. There was a war,
and people were being asked to recycle aluminum. It was the off-season.
I am so sleepy right now, but I'll tell you something: you should regard
every anomaly as an opportunity to be awestruck.

When she was a small child she found a dead pigeon beside the railroad
tracks, and later in the same place she found a dead calculator. Her
friend went to the hospital. His cornea was wrinkled and it was making
his eyelids stick together. It was the off-season. He closed up the shop
and went out in search of a sandwich and a malted milk. "Oh God damn it
for Christ's sake," he said. Love. What was this thing that so many
poets made such a fuss about? When they first met, they stayed indoors
all the time. He was agoraphobic and she was afraid to be alone.

In those days everyone had a story to tell. Losing a toe to frostbite or
getting one's head stuck in a milk can. Outside, the streets were dirty
and cold. They needed distraction. They wanted to see everything, so
they would be distracted from everything else. He is a man in fine
fettle. He possesses a generous spirit, and she has nice manners. She
knew when she got up that morning it would not be an ordinary day. They
took a lot of things with them: it was not the things they were trying
to escape so much as the origins of things. They had to forget where
ideas came from in order to find them interesting again. She justified
the trip as ``radical chemotherapy for my cancerous soul.''

All the hotels were filled with celebrities they didn't know. People
were eating their dinners alone in restaurants. The mountains were
dream-like and seemed smaller than she had imagined. On the side of the
highway there were farmers selling Kool-aid. The amusement park
employees were not amused. Of all the rides, the bus back to the city
was the one they liked most. There was a restaurant that served nothing
but insects, and another one that served only tofu. ``Can you drive a
standard?'' he asked. They drove a Chevy to the levee and drove a Ford
to the fjord. They tried to make themselves more interesting by being
interested in more things. Some of the maps had to be stretched to make
them fit. ``You can't get there from here," he said.''It seems like I
can't even get here from here," she replied.

Everywhere they went, the souvenirs had been made somewhere else.
``There are so many beautiful places in the world," she said. "But they
are so disorganized.'' They couldn't picture themselves in any tragic
news reports, so when the boat was overcrowded and barely out of the
water, they didn't complain. The hotel room was exactly six feet long by
three feet wide by three feet high, and it had a TV. Oh, I'm so drowsy.
What should we do? Let's go for a stroll by the river, for a change of
pace. Do you want to see a movie? It seems silly to go to a movie
theatre. They called home and left messages for themselves. When they
found an internet caf they digitized postcards and emailed them to
friends. Tourism is the greatest gift of the twentieth century, they
thought. ``The reason I'm always having fun,'' she said, "is that I'm
doing what I like. I'm the happiest loser alive," she said. "My
specialty is nonchalance," he said.

Ten miles out of the city they saw the most amazing sight they had ever
seen. A lifetime of unspoken, unconscious supplication was finally
rewarded. They were late getting back to the train station and the train
was already beginning to pull away. They called out for it to stop so
they could get on. The conductor said, do you have tickets? Who are you?
New man, new woman, feeling lucky to be somehow caught up in this
tedious, cruel, and inexplicable existence. New woman, new man, carrying
an old suitcase with 30 lbs of brandy-filled chocolates and no spare
underwear. Reading books upside-down and backwards. Losing our wallets
with a sense of purpose. Sleeping with our eyes open. Haggling over
hagiography, canonizing cartoon characters. Boxing with St.Goofy. Paying
for dinner with underarm deodorant. Leaving a trail of discarded
punctuation marks. Taking out insurance and releasing it into the wild.
Fucking on the carpet, between the seats. I am a gentle philosopher of
the road, he proclaimed. I am the zeitgeist multiplied by a weltschmerz,
she declared. New man and new woman, picking up the pieces and moving
on, except sideways this time.

They played Scrabble in the observation car, by moonlight, with blankets
around them and strangers sleeping in the corners. They passed a frozen
waterfall. They found a newspaper and discovered that there had been an
election while they were gone. The war had moved on to a different part
of the world. The stars fizzled and went out. The leaves were vivid in a
way that made them feel light-headed. Sunrise above the clouds doesn't
compare with this, he admitted. "It's so\ldots{} ineffable!" she
exclaimed. "That exclamation mark," he said, ``that exclamation mark is
the last exclamation mark you'll ever need.''

Two people are travelling on a train. Watching tiny villages and fields
go by, having sex in the washroom. "It's bits world," she says. It's
bits world. Let's get along with me. I am an edifying and enjoyable
text-based performance artist, and this is most comfortable performance
you have ever run into. Wherever you may be, whoever you are, when you
think of usefulness, think of Bits.

\hypertarget{susan-gevirtz-from-thrall}{%
\chapter{\texorpdfstring{Susan Gevirtz: From
\emph{THRALL}}{Susan Gevirtz: From THRALL}}\label{susan-gevirtz-from-thrall}}

~

\begin{quote}
I am attractive but you are captivating." --- "Life has its attractions,
but death leaves one spellbound.

--- Baudrillard, \emph{Seduction}
\end{quote}

~

\begin{quote}
Even children connect the thought of death with ``a strong unconscious
sense of pleasure associated with the return to the mother's
womb\ldots{}''

--- Sprengnether, \textsc{The Spectral Mother}
\end{quote}

~ Is person place or person

\hypertarget{etymologys-nemesis}{%
\section*{Etymology's Nemesis}\label{etymologys-nemesis}}
\addcontentsline{toc}{section}{Etymology's Nemesis}

Soaring and diving with the aquarium sharks and casting another shadow
half imitation and half prebirth. Drinking of this aquarium water
tincture and thereby knowing something happened. Thus reaching for the
dictionary. Inadequacy's tablet, one nervous twitch before story's wish
to subtract by including too many particulars. Such soft wax of
impression. A highly sophisticated selection of pattern for which, in
the past, they had only the word ``memory.'' In Intelexa there is a
saying, ``Obfuscate so that the windshield's film will clearly wash over
you.'' This is not religious maxim but the slay of time's desire to
organize you into a bonsai beyond recall.

~

~

~ \#\# The Anchored Lost

Then she appeared---the woman of one hundred arms. First as still as
tree branches, then windmilling, but that is a word from the land next
door. Then the tale of the man who switched clothes with a berber woman
and went looking for the lost city in the sand. The desire to paint
without making a picture. This was the first story for the first day and
one was aware of that. ``Sux - intravoxa,'' they muttered attempting to
get on with things. But all that was construed as thing had reverted to
sky at the slightest glance, embarrassed sky.

~

~

\hypertarget{in-this-place-which-legend-posits}{%
\section*{In this place which Legend
Posits}\label{in-this-place-which-legend-posits}}
\addcontentsline{toc}{section}{In this place which Legend Posits}

A silence fell upon the land. More quiet than a map. The sounds of place
having to do with the irreducibility and seamlessness of story. Airlines
want us to forget this. One body feeds sleep to the one next to it. And
the next imagines it as her own. For example.

~

~

~

~

\hypertarget{paper-doll-day}{%
\section*{Paper Doll Day}\label{paper-doll-day}}
\addcontentsline{toc}{section}{Paper Doll Day}

\begin{verse}
The egg from the other chapter has followed us here.  is\\
the only word for it. A study of continuity necessary. Of meals eaten\\
as opposed to meals made, read about, bought or imagined.\\
Flowering thyme takes too long to say.\\
Handprints on cars.\\
Colander hours.\\
Archaeology you can’t touch.\\
Ruffled hour, gravel hour\\
\end{verse}

~

\hypertarget{renee-gladman-from-studies}{%
\section*{\texorpdfstring{Renee Gladman: From
\emph{Studies}}{Renee Gladman: From Studies}}\label{renee-gladman-from-studies}}
\addcontentsline{toc}{section}{Renee Gladman: From \emph{Studies}}

\hypertarget{working-notes-3}{%
\section*{Working Notes}\label{working-notes-3}}
\addcontentsline{toc}{section}{Working Notes}

\lettrine[lines=3, findent=0em, nindent=0.1em, lhang=0]{I}{ had been holding a vase}
in my hands that now lay shattered on the floor. The vase was ``the
novel.'' I'm not sure of the force that knocked the whole away from me
but I imagine that it was a slow accumulation. As I looked down at the
pieces, I could not see how they ever fit together. In one place lay
prose, in another narrative, and off to the side, in several, lay the
components of fiction. I never felt that I was writing ``fiction,'' so I
was surprised to see it there. What I found more difficult, however, was
seeing ``prose'' extricated from ``narrative.''

This confused me. Before this moment, I had thought of narrative as the
trajectory of language, as the direction in which it unfolds, and of
prose as narrative's rawest tone. But as I observed them lying there, I
saw each independent of the other. I knew that no matter what I did, I
would be writing narrative, but I could not fathom how I wrote prose
separate from that. The process is as natural as breathing, however this
crisis I experienced was emblematic of that occasion where we catch
ourselves in the mechanics of respiration and get caught up in the
mystery of that.

~

\hypertarget{section-15}{%
\section*{7}\label{section-15}}
\addcontentsline{toc}{section}{7}

\begin{verse}
\emph{for Tisa Bryant}\\
\end{verse}

A friend has pointed out a paralysis this evening, which she has been
ignoring for some time. The friend has had to ask why certain parts of
her body don't move, why there is failure in her language, why she
carries this failure in her body, making her body unable to move. It is
evening when the friend can no longer contain her observation and is
forced to point to this stillness that makes her a stranger in the room.

The friend points to the paralysis.

Our immobile one is sitting quietly in the corner of a room staring at a
glass of milk, seemingly enmeshed in daydream. There is condensation on
the glass, an indication that the milk has sat unattended for too long.

Earlier, the friend brought the immobile one this milk.

Then paralysis took its form.

---Now she has been made to see her paralysis. It impresses her and
makes her mad. She and the friend sit in the room mad at the paralysis.
As the milk grows warm.

The friend looks at the milk and asks to sip it. Our one says, please.

Now that the paralysis has been exposed both women stare ravenously at
the milk.

Both reach for the milk.

``Now that I've pointed this out to you, can you explain how it got
there?''

She answers.

``Does the measurement, though compelling you into silence, serve the
companion in anyway? Is a co-habitation possible?''

Again, she answers, this time with more emphasis on each word.

``I have known you for ten years,'' the friend says placing her chin in
her hand, amazed at what she's seeing.

\hypertarget{gad-hollander-from-serialscribbler}{%
\chapter{\texorpdfstring{Gad Hollander: From
\emph{Serialscribbler}}{Gad Hollander: From Serialscribbler}}\label{gad-hollander-from-serialscribbler}}

~

\hypertarget{working-notes-4}{%
\section*{Working Notes}\label{working-notes-4}}
\addcontentsline{toc}{section}{Working Notes}

\lettrine[lines=3, findent=0em, nindent=0.1em, lhang=0]{S}{erialscribbler}
(as the title suggests) is an ongoing work, so far only in prose, though
there's no rule about that. It started a couple of years ago, when I
sent out a piece of scribble on a postcard every day to someone I knew
around the world. The texts were continuous (a lot of drivel thrown in)
but each recipient only received a fragment. The idea was to
re-integrate other people's (not necessarily writers or artists of any
description) responses in my own writing. So far, that hasn't happened,
and after some 120 postcards I stopped sending them out, especially as
the internet seems to throw a slightly anachronistic twist on the
project. For the time being, then, it's still an open-ended project and
anything might happen with it.

From some working notes on Serialscribbler found in my notebooks from
the last few years:

The language eludes me just as I step into it. I do not have the power
to seize it, let alone transform it.

Serialscribbler---what is more than language in it--- now sit in
Waterloo station now hear \& be tempted to bark but not bark,
quarantine---look up "bark"--- then \& now reconciled (thru fog)---may
as well not pay \& suffer the consequences---what if someone gets up \&
leaves in the middle of a word---how do we find the middle of a
word---if poetry is an approximation, the world is a wild guess---guests
come \& go, trains depart \& arrive, film rolls.

Snatch scribbler---between picking up dog shit \& chasing after the
puppy, the memorable phrase or fragment might fall into place. More
likely it will not. Take in the decent weather while it lasts---a day or
two---then brace yourself for the onslaught of dark winter days.

Write without closure, as if each block of text were a prelude to the
soundwork or video. Take it one sense at a time and then return--- not
quite, but in a spiralling form rather than a circle--- serialscribbler.

\hypertarget{the-natural-suspicion-of-words}{%
\subsection*{``the natural suspicion of
words''}\label{the-natural-suspicion-of-words}}
\addcontentsline{toc}{subsection}{``the natural suspicion of words''}

Sometimes it happens that even Serialscribbler, who by definition is
busy at all times if not totally self-absorbed to the point of oblivion,
sometimes (I say) it happens that Serialscribbler will devote some
attention to another experience, eg, reading a stranger's book or
looking at a sculpture, even feeling it to gain some tactile experience
of it.

I write in prose by default . I don't know where to break the line in
free verse and I don't like having the line-lengths predetermined by
fixed verse forms. It's an interim solution. I write in prose but think
in poetry.

\hypertarget{serialscribbler}{%
\section*{Serialscribbler}\label{serialscribbler}}
\addcontentsline{toc}{section}{Serialscribbler}

I took editor to one side and had a quiet word with him. This is for my
own good, I whispered, and proceeded to disable his editorial tools,
lopping off his hands, extracting his tongue, etc. He seemed to
understand and, as I now read it, thanked me profusely with his lolling
head and rolling eyes. Then I began to practise my fiddling, keeping an
eye on his writhing body while folding, turning one word into another
the way dough is turned before being given time to rise. I dredge up the
analogy (more fiddling) with a tinge of regret, even shame, for bread is
my favourite food in all the world---I suppose I would kill for it---and
to mention editor in the same breath as bread seems almost blasphemous;
ah, but I remember there's a limit to words and it's just been reached.
So let's step back a step, retract that near-blasphemous breath, inhale
deeply, say: I fiddle is to say I scribble within the confines of my own
space, despite the neighbours' persistent complaints about the
screeching and scratching, their calling in the police who, in any
event, are powerless to stop me fiddling when I explain that the noise
is just a scribbling, a turning of one word into another in the manner
of kneading dough, careful to point out that I'm not inciting any
revolutionary acts by invoking the analogy and, at the same time,
carefully keeping editor out of view lest he give the wrong impression
or, perhaps, because I'm ashamed of him. Upon leaving the police shake
my hand---I've always had a way with the police, as one has with
children or animals---and I resume my fiddling. Editor is hurt,
professionally and physically, and since this has just been established
as his necessary condition, we are now both faced with a steep learning
curve which we need to get used to, each in his own way.

I am no virtuoso, it must be said, and fiddle only because I'm curious;
my fiddling is spurred on by a vague desire to discover the new sounds
rather than improve or polish the old ones. That's not to say I don't
appreciate the old ones, much in the way editor does, but it's not
enough to fulfil my desire. Editor, in contrast, has a way of expunging
the new sounds with a stroke of his pen, somewhat like an abortionist,
fossilizing them before they ever take flight. He believes my desire is
misguided, that there is a better way of channeling and controlling it,
that it should be put to better use. While he was fit and able he'd
always wanted to have a virtuoso at hand (forgive the expression) and
had dreamed of making me one of his own. Ah, but there are too many
virtuosos, cried Robert Schumann, spitting on the polished parquet
floor. Editor is confounded by the riposte. I can see his mangled body,
thinking: surely you mean virtuosi. He is certainly not humming a melody
inside his head, because as I snatched the pen and disabled him I
discovered, to my shock, indeed my shame (after all, he is/was my
editor), nothing in his internal makeup remotely connected with music.
Instead there were millions of micro-blueprints stacked on top of each
other, slightly curled at the edges like wafers, each bearing the
imprint of some formula for a melody, past or present or future: all
nuts and bolts and intricate scaffolding, which in its way presented a
dazzling construction that any virtuoso would have given his right arm
to play (forgive the expression), but without the slightest indication
of what voice or tone or timbre such melodies might be borne upon. In
fact, a moment before disabling him I informed editor that there are no
more virtuosos in the world, only students, and to my surprise (and
eased conscience) he seemed more relieved than disturbed by the news. I
was not telling him the whole truth, of course, but our relationship is
such that it forgoes the elaboration of truth (on my part because its
structural complexities would be overbearing, while on his part the
aesthetic and moral dimensions would defy his comprehension); rather the
smallest grain of truth is often enough to maintain the fragile rapport
between us. So for the moment editor and I live side by side, reconciled
to each other's exclusionary habits; I fiddling, he nodding his wobbling
head, keeping time for the sake of appearances. Occasionally, when our
respective tasks come full circle and we turn a mutual glance of
recognition towards each other, we will take a break together over a cup
of coffee or a glass of wine, our dreams oozing incoherently from the
corners of our mouths, dripping like bird shit on the cafe table, as
life flows to and fro on the pavement before our eyes---his, like mine,
darting back and forth with unabashed lust as one young woman or another
comes into view---only to have our little reveries promptly interrupted
by a waitress asking, ``Are you all right, sir?" So that I'm forced to
look up and around me, awkwardly, desperately. Editor has suddenly
absconded and left me to answer her simple query. In a futile attempt to
grasp a pen with my mouth and scribble''yes" on any scrap of paper, I
have, not for the first time in my life, accidentally overturned the
table and chairs. Remorse for my reprehensible behaviour towards editor
momentarily seizes me, and a shameful tear is about to fill my eye as
the waitress, brimming with kindness and efficiency, comforts me: "Don't
worry, sir, it's all right." Dutifully she bends down and shuffles on
all fours as she cleans up the mess around my truncated legs; her
breasts, liberally exposed to my parched and crippled gaze, sway
slightly with her scrubbing movements, as if to say: drink, drink my
compassion, o scribbler.

As a student of fiddling I am prone to daydreaming, especially when I
practise my bowing and fingering, which I do religiously yet by all
accounts hopelessly, for whether I practise once a week or every day of
the week, whether I fiddle on street corners or in cafés or in the
privacy of my room, the result is always, acoustically speaking, a
scribble. It has to do with distraction, of course, and the enticement
offered up by an image of loveliness passing by on the street, in a
café, or across the walls of my room. Oh, but I like that, remarks a
progressive tutor, a specialist on expressionism, his ear ever alert to
the permutations of the new sounds. Annoyingly, I must watch my posture
at all times---me, a cripple, watch my posture!---make sure my eye-line
does not stray too far from the strings, and always keep my feet firmly
planted on the ground. Such, at least, were the rules given at the
outset, when I was a mere sub-scribe, long before I became a student,
and although I'm now at post-graduate level the same rules still apply,
the cardinal one among them being "never drift into mindless reverie."
The trick, as another tutor advised us, is to dream without affecting a
dreamer's countenance, to cast one's mind off in some direction while
the (very same) mind stays focused, enrapt by the intricacies of the
hands' manoeuvres. As you can see, I haven't given up yet, perhaps
because I've noticed that over the years tutors come and go, experts
become fallible, if not completely discredited, histories get revised,
wars start up and arrive at their conclusions, stores change ownership,
buildings get torn down and re-built, and things generally move
on---whereas scribbling just evaporates into the atmosphere like so much
vacuous sighing and, despite its noxious fusion with the air we breathe,
is accepted as just another waste product with a tolerable, low-level
health risk, no more dangerous than the toxins given off by other human
wastes. The fact that nobody pays any attention to it, however, could
prove tragically fatal under certain circumstances---for example, if a
scribbler were to be retroactively regarded as talented, even a genius,
and we, human kind, belatedly berated ourselves for having been so deaf,
so blind, and beat our breasts with futile lamentations, and tore each
others' hearts out with proprietary fingernails, and waged suicidal wars
in the name of said dead scribbler, now canonised, his or her pen
emblazoned on the national pennant billowing proudly over our burning
libraries, day and night, etc. But in the case of students, even at
post-graduate level, such potential catastrophes are highly improbable,
as unlikely as a cataclysmic end of days, so there's really no need to
fret, no need to get anxious over our performance, no need (or indeed
desire) to seek encouragement from our peers or advice from experts,
because scribbling is merely a process, nothing more, a perpetual
practice session with or without a fiddle, and in any case, as
otologists have shown, the human ear's threshold for noise rises by the
hour, edging nearer to a level of infinity with every exclamation, every
grunt and moan, every scratch or tap of the finger, with every etched
apostrophe or comma on a blank page. As one of our most distinguished
tutors put it: \emph{Our potential capacity to listen to everything will
shake the foundations of our belief systems to their core.} I was
thinking, as it happened, about the word as the foundation of faith (not
my original thought, just pondering an old one) when he added:
\emph{It's something we should consider with the utmost gravity and
tackle with a measured rationality.} I think he was soft peddling doom,
imploring us to reinstate a more humane threshold for the ear's intake
of noise before it's too late. Too late for what, I wanted to ask, eager
to join the debate. But being a cripple, being doomed to my studies,
playing uncontrollably with my fiddle, I would have been unable to get
my point across to him, over there, on stage, behind the podium, without
upsetting the furniture and provoking a riot. All the same, upon leaving
I reflected that despite my restraint it was an interesting outing for a
bleak winter's afternoon, and I went away believing that my knowledge
had increased by two- or threefold.

Although I'm sometimes mistaken for a cripple, sometimes I'm not. It is
therefore just as true to say that I am as that I am not a cripple, and
it's a matter of political interpretation whether you believe one or the
other of these truths. My identity as a scribbler has not been affected
one way or the other in this respect, but this is probably due to the
perverse social conventions of our time. In the old days they would have
certified me, put me in a safe place behind walls for my own good, and
said no more about the matter; invisible, I would have been perceived as
either happy or unhappy. Today our techniques are more subtle and
sophisticated, especially in the free world, where we take pride in the
fact that we can air these issues openly and hold interesting debates at
our local cafés or on television or occasionally even on a bus, on the
way home from work, though that is still somewhat rare owing to our
collective fatigue after work. But those of us who don't work, who are
less weary, will not hesitate to launch into such democratic debate if
we should come upon a fellow non-worker on the bus, even heading home at
5 a.m. after an all-night party, just when our fellow workers are
setting off to work on that bus, though they may not be sufficiently
awake to participate in the debate and are possibly too preoccupied with
their night's dreams, reconstructing the details, thinking of plausible
explanations for implausible events, half-consciously tying together
random associations in their minds, their hands struggling to support
their drooping yawning heads, their bleary eyes peering into the
pre-dawn darkness as they silently ask themselves what does this or that
object mean in that incongruous setting, among such unlikely people, and
is the dream object connected with the real thing they stumbled upon
only a few days earlier, after not having seen the likes of it for
years, their brains persistently searching through this miasma of vague
objects and people for some kind of solution, a key, as if these freshly
painted dreams were riddles that only needed the correct answer to
facilitate a stress-free day at work and allow them to confront their
superiors with the confidence of people who know themselves, whose
selves have been revealed to them in dreams, regardless of how menial
their jobs might be and irrespective of their professional status, their
pecuniary position, their grooming or their physical appearance in
general---it being their democratic prerogative, theoretically, to be
themselves in any guise they choose---because they saw the light through
the dusk and, having opted out of a now raging democratic debate, heard
a voice in the wilderness and were born again on the upper deck of the
bus at 5a.m., en route to work. I know all this for a fact because,
having participated in that democratic debate, I witnessed their
revelations with my own eyes and quickly scribbled everything down while
it was all fresh in my mind, when their dreams were far less obscure
than they are today.

Since then reverie has been my precise goal, and my fiddling studies are
only a means of dealing with the pragmatic side of other dreamers'
dreams. One fine day, as we return from our afternoon ritual of ogling
nymphs on the banks of the boulevard, editor and I will come upon a list
of words posted on our door, with a set of instructions to consider
their place in my scribbling. Editor is keen to proceed, sensing an
opportunity to regain the high ground in our relationship, but I remind
him of our pact and avert his deployment for the time being. He is all
but dead, a mere copyist, but I too am suddenly paralysed by a bout of
amnesia, my lifelong years of study failing me miserably.
Conscientiously I copy the list into my computer and it comes out like
this:

Shame

The private

The secret in narrative writing

Representations of time

Event

Expectation (the future) in narrative

My computer (or rather its editor) automatically capitalizes the first
word in each line---not what I expected, but I leave it as it is. I
imagine the software writer will claim responsibility, perhaps a former
poet eager to accommodate traditional verse writers. But, capitalized or
not, the words are easily comprehensible, so I feel at least partially
redeemed by my scribbling studies, imagining that any reader would
understand them: the typography and size of a word has no bearing on its
meaning, I remind editor, though it may be argued there's a difference
in degree. Though I would have preferred the initial words without
capitals, it seems pointless (and too difficult) to change them now.
(Perhaps in future the software writer, a being of necessity in the
future, will include an instruction that would allow the computer's
editor to distinguish between a list of items and lines of verse;
perhaps another instruction might be added that would allow its editor
to discern a writer's intended nuance for a given word, or line, and
automatically capitalise such words---for example, "life" as typed in by
an immortal's fingers.) But this is beside the point. The point is I
don't understand why or how language and writing coexist, though I
understand and to some extent sympathize with the historical
expediencies that have conjoined them. But how, after so many
thought-bearing hours in such claustrophobic proximity, have they failed
to find their separate autonomous domains? Have they just grown used to
each other and learned to be mutually tolerant? Have they gone a step
further (backwards, it seems) and learned to coexist, perhaps even to
love each other in our thought processes? If that were so---though I
don't for a moment believe it---why? what for? We scribble, day in, day
out. If we advance, writing lies in wait, ready to devour us; if we
retreat, language stabs us in the back. Would it not be better if we had
neither---no language, no writing---if there were no scribbling, no
anti-scribbling, no pauses and deliberations over their duration, if
everything were simply a constant torrent of some alien construct
pouring into the well of truth, that shameless metaphor born from the
union of language and writing. I'm not alone in this my bafflement. My
incomprehension is borne out by the writers I read, writers who have
nothing to do with language, who don't know the meaning of the word,
writers as removed from representations of time, event, expectation and
futures as Greek shards, who have as little idea about shame, the
private, the secret in writing as that extrovert Yaweh, who have never
even heard of a narrative concept. Those writers I read, those are the
writers I write, that is, those writers I re-write; that's why they
write and why I write, why we, they and I, write, tuning our scribbling
fiddles to a sorrowful pitch. Editor reads the very same writers I
write, and yet invariably incinerates their pages, marking every word
for approval before burning them, and sometimes adding a last-minute
grammatical correction in his head. But editor reads only what has
already been written, texts, what may never be re-written, that's
impossible, what may only be read over and over and over. Whereas I and
the writers I read write only to re-write each other and will not allow
ourselves, for our verbs are many, to be re-written otherwise,
distracted by a shopping list of contingencies with moral and historical
imperatives, since we're ambitious to succeed our scribbling with our
failure. That shopping list belongs in my pocket. It is revealed (like a
furtive cigarette in the street) when I go shopping and use language as
a counter, when I'm talking face to face, seeking out your immediate
reply, my lips smacking each other, my jaws moving like pincers, my
tongue slithering between movements of jaws and lips, my vocal chords
dilating or contracting like synchronized swimmers, my lungs and every
other pertinent organ acting in time and with deceptive ease to allow me
to articulate something or other, this or that. Thus we talk and act,
sometimes sequentially, sometimes synchronously. The articulation of our
bones, as well as what we concoct with language, says something about
our life, in death, after the flesh has rotted away and the worms have
moved to another body, yet language itself can never exhume the writing
that went on (perhaps on a daily basis) before the final heartbeat.
While the writer, o scribbler, is steeped in shame, in some perverse
privacy or in the self's profoundest secret, the body of writing is
perpetuated in its own skin, marking out its own as-ifness the way
silence delineates the borders of language. Let us admit, let's say,
that I have murdered you and am now stricken with remorse, dumb with
shame and crippled by guilt; whatever I've done, whatever else I'm
doing, I am not writing. I am outside the writing, doomed to carry in my
language as author some pitiful mental inflammation throbbing with a
succession of phonemes, each in turn having been orphaned by its
murderous language, the same that suckled it, that gave it comfort in
the face of silence. Which is all well and good, is as it should be in
the realm of as-if, but as the writing proceeds what if, scribbling, it
revealed itself in some sacred alphabet, a revelation earmarked as a
future holy text, as though writing were shame and privacy and secrecy
itself, to be divulged in our death? And what if that as-if were an
angel, classical, orthodox, with the power to unleash language's
ultimate terror---though I and the writers I re-write don't know what
that might be---would editor allow it to fulfil its divine mission, or
would he intervene and, effectively, overthrow the lord in our
household? Such questions seem infinitely difficult when, of an
afternoon, editor and I return from viewing the maidens by the river,
slightly tipsy from their beauty.

I like, I say to him before reaching our door, a poetry that is
definitive in its ambiguities, but cannot be as passionate about it as I
am about a poetry of ambiguous definitions. Reaching for the list of
words on the door, he calls me a hypocrite.

I was so ashamed that, if writing were language, "so" would never
survive under the burden it claims to support, for shame draws us back
through time, as even the earliest scribblers evoking Paradise were at
pains to point out, and during that interregnum every act of love and
war perpetrated by the species has accrued to its surface. Since writing
is not language, however, it is easy to articulate the so of it, to say
it or spell it or memorize it for future use. And I hid, editor
confessed to me once, for I was so naked.

The nonsense we tolerate, I confessed back to him in a moment of
compassion, is nothing compared to the nonsense we produce, especially
when the latter is predicated by a collective pronoun and dressed, how
can I put it---I paused, I scratched---aphoristically? he offered.

(work in progress)

\hypertarget{eileen-myles-long-and-social}{%
\chapter{Eileen Myles: Long and
Social}\label{eileen-myles-long-and-social}}

~

\lettrine[lines=3, findent=0em, nindent=0.1em, lhang=0]{M}{y book}
(*Cool for You\footnote{\emph{Cool for You}, New York: Soft Skull Press,
  2000}) is absolutely an extension of my poetry practice, I've even
come to refer to it and maybe the class of novels it's a member of as
the poet's novel. Having written a novel I don't have to protect myself
from the disparaging term poet--yet I'm more clearly a poet than ever. I
mean a poem is an extravagant grandiose and trembling form, for better
or worse always alive, I think, and I've brought those weaknesses and
virtues into novel writing and I'm dying to do it again. As a younger
poet I was urged (in order to be important) to think large, to write the
long poem, but I think this is it. It's epic poetry in the sense that
the epic poem is a communal form, and long and social. I wrote it for
you. That's what the title means in a way. I know there's a more
technical definition for epic poetry but that seems the best way to get
it wrong. I'm more interested in what Joyce Carol Oates did in
\textsc{Blonde} than Seamus Heaney's prize cow. Poets should write
novels en masse and reinvent the form and really muck up the landscape.
I often think of Leslie Scalapino telling me that her long book the
return of painting, etc. was called a novel because a novel means you
start at the beginning and you read to the end. It was a way to
guarantee she could be read differently. I don't care about that
particular aspect, the front to back thing but being read differently of
course is nice. Time is my real subject and calling a book a novel
immediately adjusts the reader's perception of time.

My book is shuffled if you take the long view. When I started to write
the first chapter which was a story, the whole thing is a story for me,
it occurred to me that a novel could just be all these stories shuffled,
and I thought of books I loved long ago, Julio Cortazar's Hopscotch. I
read that book straight through, but I read it differently knowing that
he had that chart at the front suggesting you could alternately read it
this way, and that the book could permute all these different ways. I'm
in love with the form of the novel being molten, not putting the power
of reordering in the hands of the reader but to extend to the reader a
sense that the form of the book is as accidental as life. It's studied,
damn studied, the cool casual life as is the cool casual novel, but as
accidental as say the zabruder footage. Some guy happened to be there
with his camera, and as it turns out history was made. As it turns out
he gave us our only window onto that day in Dallas. When one realizes
that John Clare inadvertently gave us the first No Trespassing signs in
literature you realize bumpkins must write first, not last.

One more note on that---everything's visual and in even or especially in
advertising the oddly cropped shot is what's used. Literature as it's
sold today is so backwards. Hand held literature is of course what I'm
selling. It's not a memoir, it's a recording. And even purporting to be
a bad recording at times. As a female writer the pose of awkwardness is
very dangerous though because at this post feminist moment one should be
a top, one should win etc. But female in history is ground down,
anonymous, untold. That's the story that's interesting.

And another thing: more poet's novel. In \textsc{Chelsea Girl's} "My
Father's Alcoholism" was the story where I stretched myself, and tried
to write out of pools and see how they would build. I just literally go
into a room (one labeled my father's alcoholism) and begin inventorying
the memories, the substance of the memory, the materiality of it. The
fur of childhood, not the feelings. When I was done with one I would
wait and let another one grow. It really is like walking in the rain.
Can that generate a narrative,or a narrative feel? It sure works in
movies. But movies have the dark on their side. My hope is that by being
as solipsistically in my mind, downloading erratic shaped drops and
letting some momentum build on that order, rather than an order based on
action, you can kind of reverse our assumptions about activity. Or maybe
I mean self. There's a beautiful math to it, like music if you just let
yourself go and describe the experience of your life like it's some
public spectacle you are privy to and you‚ re waiting for the music of
the information along with everyone else. I feel less alone when I tell
my most private stories than any other time. This does link up to
performance. When you are acting or even reciting a text you wrote and
you're on a darkened stage under a pool of light there's a terror as
you're singing along with the text, its like driving, that moment at
night when you think shit I could just pile into the opposite lane and
that would be it. The self could just smash into the light and dark and
be gone forever. So you feel contained within the details by that sense
of danger and that provides the tension while you proceed lavishly
downloading "me" in every ribbony way I can imagine.

The limitations of "New York School" have been what Frank did, what John
does, are you like Jimmy. Mustn't forget Kenneth or Barbara. It's like
Mt. Rushmore. You can say some quick thing about them all---it's chatty
abstraction, it's American speech, and I've learned different things
from all of them but since they are just people I think it's tone that
people wind up miming, or concerns rather than electricity, the weight
and pause and incredible expanse of the exploring mind in speech. I
discovered what I sound like long ago and of course I go where they
wouldn't and couldn't so it's sort of like explaining why you prefer not
being a corpse. New York school mimed is worse than academic, it's like
my dad being a mailman coming home with the Ivy League clothes from the
Harvard dorms where he had his route. You have to blow that up pretty
quickly but it's just sad. Let's face it, they were just as New critical
as everyone else in the fifties. They all would assert the poetry was
not about them. It's about skimming the surface of the self. Using that
facility to shape the poem. My dirty secret has always been that it's of
course about me. But I have been educated to believe I'm no one so
there's a different self operating and I'm desperate to unburden my self
of my self so I'm coming from nowhere and returning. That's sort of
classic. You just cannot underestimate the massive difference in writing
out of female anonymity. It blows all the styles out of the water.

I think the form of the novel gives dignity to my shame. Sometimes I'm
just ashamed to block the sun. Performance, and I include readings in
that, makes the body be the container for the work. It is when you write
the words of course. The body always seems like the shame. The camera
must cut away to the trees, the animal is telling too much. The animal
doesn't want to die etc. If you've ever sat on a panel when some people
read and some people speak directly the difference in those two
deliveries in terms of the room's interaction with the speakers is
immeasurable. The dangerous loose canon is \emph{how} we conduct public
life---we want to stifle the stray remark and we want to house it. I'm
totally translating that impulse into a several hundred page thing, my
sorrowful body is now text. It's like a zoo.

One last thought. In the \textsc{Times} there's occasionally these
articles where the great artist, Brice Marden or whoever takes us into
the met to see the work that is interesting to them. I started reading
Robert Smithson a couple of years ago and his move out of the studio
helped me understand the poet's situation, having no such institution to
tour through, having no \textsc{Times} to conduct the tour in, and
possibly being invisible because of that. Smithson yearned to kick free
of that structure that has been happily plunked down ever more solidly.
It seems like a novel is the only way to prove I was here. Exactly how I
was here. It's a social poetic invention.

\hypertarget{daniel-nester-god-save-my-queen-the-freddie-fragments}{%
\chapter{Daniel Nester: God Save My Queen (The Freddie
Fragments)}\label{daniel-nester-god-save-my-queen-the-freddie-fragments}}

~

\hypertarget{working-notes-5}{%
\section*{Working Notes}\label{working-notes-5}}
\addcontentsline{toc}{section}{Working Notes}

\begin{quote}
"Mr.~Narrator, this is Bob Dylan to me."

--- D. Boon, The Minutemen, "History Lesson-Part II"
\end{quote}

\lettrine[lines=3, findent=0em, nindent=0.1em, lhang=0]{I}{n June 2000},
I had a crisis. In a club called "Fun," an dank, empty club deep in
Chinatown, I sat down for an interview for the VH1 show SoundAffects,
and talked about how Queen changed my life.

I left so much out in that interview that I tried to write more poems
about Queen, but hit a wall. Call it problems with verisimilitude,
whatever it was, poems weren't doing the trick.

So I started writing a confessional book, freeing myself from what I
think are very real aesthetic contraints on today's poem. The result is
a kind of prose poem, with narration leading to definitive moments where
I think my history and the object of my obsession's history-and I say
this with a straight face-intersect.

\begin{center}\includegraphics{images/"pandanus.png"}\end{center}

\hypertarget{god-save-my-queen-the-freddie-fragments}{%
\section*{God Save My Queen (The Freddie
Fragments)}\label{god-save-my-queen-the-freddie-fragments}}
\addcontentsline{toc}{section}{God Save My Queen (The Freddie
Fragments)}

The day I bought the \emph{Flash Gordon} soundtrack, the cashier, whom I
knew to be in a shitty suburban punk band, offered me my first record
store clerk snide comment.

"Nice bass line," he said, looking down at the banana yellow album on
the counter, punching up the price.

I walked away, meek and ashamed that my purchase had not passed muster.
He was, of course, referring to the single-note bass line in "Flash's
Theme," the single from on the soundtrack. But I was steamed. It was one
of those times you think of with the perfect comeback later that night,
and you lay there, beating yourself up for not coming up with it at the
time.

In 1977, when Queen were recording \emph{The News Of The World,} legend
has it that the Sex Pistols were recording in an adjacent studio. Sid
Vicious, wandering into the wrong room, as was his custom, came upon
Freddie Mercury, sitting at a piano. His customary four fingers of vodka
in hand.

``Ah, Freddie Mercury," Sid said, proud of his recognition.''Bringing
ballet to the masses then?"

Freddie looked up from the keys, and, not missing a beat, offered his
reply.

"Oh yes, Mister Ferocious. Well, we're doing our best, my dear."

\hypertarget{prophets-song}{%
\section*{\texorpdfstring{\emph{Prophet's
Song}}{Prophet's Song}}\label{prophets-song}}
\addcontentsline{toc}{section}{\emph{Prophet's Song}}

\begin{verse}
Chunky on the shag rug\\
I’m looking for my anthem\\
I’m looking for my headphones\\
I’m looking for the bare spot\\
On the rug to wallow side-\\
Stepped on the chair-stopped\\
Door, and as I blast my ears\\
Out with \emph{A Night at the Opera,}\\
I’m looking at pictures of you\\
My catholic princes, my mother\\
Father proxy, I’m in London\\
With you now, the koi pond, the flamingo,\\
The statues, the hymnal,\\
The Aretha Franklin song at the funeral,\\
Alone on a pew, watching the water,\\
Watching the bare spot on the rug\\
Filled with pictures on the floor,\\
I’m ignoring the knocks on the door\\
I’m ignoring the knocks on the door\\
Stepped and stripped on the chair-\\
Stopped door, I’m listening, listening.\\
\end{verse}

\begin{center}\includegraphics{images/"pandanus.png"}\end{center}

The Elektra butterfly, an indicia, the small caps Helvetica bold\\
font, the import Parlophones, their spare blank small writing. Those\\
first few moments playing the English 12-inch singles, fresh from the
thick plastic seal, realizing the beginning is too slow, that you\\
have to switch the RMP knob from 33 to 45. The silence between the\\
songs, readjusting in my chair, preparing for the first listen.\\
Rooting for them, noting who wrote which track, soaking up the\\
essence of each of them while listening to their song. And the order,
the plot that album took, the lead-off rocker, the last song on side\\
2's power ballad. The internal band politics I imagined in the\\
making of that order. The Greek lettering in a concentric dance on\\
the cheapo compilations.

\begin{center}\includegraphics{images/"pandanus.png"}\end{center}

When I think of the 1980s, when I think of the skinny ties, and the\\
big hair, teal drawstring pants, and Ronald Reagan, I think of Queen, of
course, and how, shifting to their big synthesizer singles, they\\
would graft a medley of past songs on the end of extended versions.\\
Or how their videos for these songs had montage footage of concerts\\
all over the world, all of which I found lame. I've come up with\\
three possible reasons:

\begin{enumerate}
\def\labelenumi{\arabic{enumi}.}
\item
  A retrospective impulse, not provided by the rock literary\\
  establishment.
\item
  New digital tinkering, combined with a desire to make\\
  connections within their own work, combined with the finest cocaine\\
  and vodka human beings can afford, a few way overdue divorces, and\\
  a general lack of interest in the current song at hand, pressure\\
  from EMI and Capitol to pad out singles with no non-album B-sides.
\item
  Because they can. They're Queen, the greatest fucking rock band\\
  there ever was or ever will be, and they can do whatever the hell\\
  they want.
\end{enumerate}

\begin{center}\includegraphics{images/"pandanus.png"}\end{center}

Sometimes I feel the need to summon the ghosts of the 70s crowds around
me, or maybe they come to me. These ghosts appear as people I meet by
chance, people who actually saw Queen in concert.

Adding insult to injury, these people usually have seen them, like,
four, five times. I'll corner them, say, at a barbeque, ask him---it's
always a guy---a battery of questions.

'What songs do you remember,' or, 'What was Freddie wearing" might be
lead-offs. I move on to set lists, how they sounded. Inevitably, these
people tire of my inquiries, and run back to their wives.

Once I bought a cherry copy of \emph{Live Killers} at a church yard sale
in the Williamsburg section of Brooklyn. The man who was selling them
also had a box set, \emph{The Complete Prestige Recordings of Miles
Davis.} Got it for ten dollars. The Queen album cost 50 cents.

My new copy of \emph{Live Killers,} the one with all four members
onstage, with arms raised receiving adulation from the crowd, had four
photos inserted inside one of the albums' slots. I remembered how I used
to stick girly photos inside an album I had designated for this purpose,
to retrieve them in private moments.

The dates on the back of the photos were 1979, the year the album came
out. In one picture, a guy is grabbing a blonde girl's bare tits from
behind, cupping them. In another, the same guy, his faint brown mustache
rides an ear-to-ear smile. His pants are completely down. The third and
fourth document the amount of coke the two had in their private basement
party. Both the boy and girl smiling, each taking a turn snorting three
and four-inch lines.

The way I find out about someone seeing Queen in concert is the
following. After my meek confession of my love for the band, after a
longer conversation about music, there's a set-up for the other person
to mention his own guilty generational pleasure. For instance, he liked
Styx at one time, he saw Journey once on the Escape tour. Then, my
reasoning goes, the path is made for me to say that my favorite band is
not some hip or critically acclaimed band, either; that it is, in fact,
Queen. It's like I'm outing myself, and then the guy--- always a
guy---tries to console me, tells me about the times he had seen Queen at
concert. Then, I get to ask the battery of questions. It's a safe
framing for the information that I mythopoetically need to hear. ``They
rocked the Garden or wherever," or''they blew my mind, man," while
off-hand remarks for some, is music to my ears.

During one easy-going outing/concert re-creation session, I brought up
these found photos I found in the yard sale copy of \emph{Live Killers.}
I was trying to spice up the conversation. I told him about the
pictures---the girl's tits, the guy's pants completely down, the smiling
pictures, the coke. We both laughed. I excused myself, offering to buy
him a drink at the bar.

When I came back to our seats, my recollect or muse had gone.

\begin{center}\includegraphics{images/"pandanus.png"}\end{center}

As if in my own benediction ceremony, I would lay out all of the Queen
albums, flush next to each other, in order of release, on my bedroom
floor. The 45s from each album would lay on top of them, in the lower
right-hand corner, also in order of release, from bottom to top. I would
then stand in front of this, drinking a wine cooler, as if I were Noah
in the \emph{Ten Commandments} movie, congratulatorily clasping my arms
behind my back, as if this was my ark, my own creation; that I had,
through my sheer accumulation and arrangement of these objects, some
part of creating them. Actually, back then, I am sure that I thought I
had created them, at least in the form of the configuration I was
looking at, and the Bartles \& James tasted sweet going down my throat,
and with my room clean and vacuumed, I would lay on my bed, jerking off.

\hypertarget{denial}{%
\section*{Denial}\label{denial}}
\addcontentsline{toc}{section}{Denial}

\begin{verse}
It was at the Queenvention\\
summer of ’91\\
when I met the old ladies\\
with the pasty photo albums\\
old 45s pressed binders\\
videotapes stuffing their purses.\\
We sat in hotel rooms\\
off the northeast extension\\
and drank Pennsylvania\\
wine, admired Freddie Mercury\\
and Brian in the lights\\
of their hometowns, smelled\\
Roger’s real tiger skin trousers.\\
What we didn’t know.\\
A hotel caterer grabs me\\
in the video room, she shakes me,\\
really, and asks *What in God’s name\\
is this?* I ask back: Haven’t you heard\\
crazy little thing called love? \emph{Oh, yeah,}\\
she says. \emph{They still around?}\\
I pull her aside\\
to make sure the Russians\\
don’t hear her. We look\\
at the new pictures, the ones\\
with the thin body draped\\
with thick sweatshirts.\\
\emph{Will they ever tour again?}\\
asks one. And I was always\\
of the consoling mind,\\
defining God by what he doesn’t do.\\
His impotence is His mercy,\\
and so forth. This, however,\\
didn’t wash with the ladies\\
of the Holiday Inn. *He’s just\\
getting skinny, is all.*\\
We danced that night,\\
haunted by the snapshots,\\
we hugged, huddled\\
round the boombox\\
for the message from London.\\
How they wished\\
they all could be there.\\
The proud ladies and I retire.\\
\end{verse}

\begin{center}\includegraphics{images/"pandanus.png"}\end{center}

Am I the biggest Queen fan? Certainly not. When I joined the Queen
International Fan Club in 1982, the newsletters were full of pictures of
fan club conventions, thousands of people taking weekend vacations to
watch videos, play trivia games, play tug-of-war, trade rare 45s. All of
this happened in exotic locations-Belgium, South America, Brighton.
Anywhere but America, it seemed. So I kept to myself with my collection
of 12 albums, 20 singles, and the few American rock magazines that
mentioned Queen at all. Then I would go out high school parties, listen
to Van Halen, new wave, Blondie 45s, Billy Joel, and Phil Collins, as if
that was the outside world, and my room was another.

One afternoon in 1984, when my parents were out, I called the Fan Club
in London, using the phone number from the album jacket of \emph{The
Works.} Ostensibly, I wanted to ask about my membership renewal, but I
also wanted to talk to someone else who actually liked my favorite band.
I would announce myself---"Hello from America!" ---ask the man who
answers about the weather, if a band member had been around that day.
\emph{Of course they hadn't,} I thought to myself, scoldingly. But of
course I had to ask: \emph{the man on the phone being so much closer to
them.} Getting down to business, I inquire about my poster and new fan
club newsletter, which I had yet to receive.

I say something about how I feel like I'm the only fan in America.

``Would be nice to live in England, where everyone's a fan," I say,
mopingly.''Well," the English-accented man says on the other end. "We
did have another person from America send something in," he says
off-handedly.

"Really!" I say. I am excited, panting. "Who?"

This could be someone I could correspond with, without running up my
parent's phone bill. Maybe, just maybe, it was a woman, killing two
birds with one stone: my love life and my cultural exile.

He looks up the person. ``Yes," he said, flipping through papers.''And
that person...that person...is...\emph{you}."

Both of us paused at this, the brutal, Spinal Tap pity of this, the
newest single, "Hammer to Fall," failing to even dent the charts in the
U.S.

I don't think I ever felt so alone.

\hypertarget{doug-rice-because-the-night-from-skin-prayers}{%
\chapter{\texorpdfstring{Doug Rice: Because the Night, from \emph{Skin
Prayers}}{Doug Rice: Because the Night, from Skin Prayers}}\label{doug-rice-because-the-night-from-skin-prayers}}

~

\emph{``Because the Night" was written out of a struggle to understand
my own uncertainties about gender and desire. For me, this battle has
always been sexual as well as religious. Traditional forms of
discourse-the memoir, realism, and so on-fail to engage in or express
this indeterminacy. I, thus, had to discover a language and syntax that
could express the confusion of the central character. I have never
believed in my own sex (that of a biological male) nor have I felt that
I am a woman''trapped" inside the body of a man. This queer space of my
own "becoming" demands different structures of representation.}

\begin{center}\includegraphics{images/"pandanus.png"}\end{center}

\lettrine[lines=3, findent=0em, nindent=0.1em, lhang=0]{A}{s a child} I
often awoke in the middle of the night startled by the noise of blood in
my body. Found myself lost inside burning dreams of walking across the
sands of the desert. Footprints dead in the wind. Fingers pressed hard
into my temples. I refused to believe in any of this. This, my body. The
stories told behind locked doors. Hungry to be yours. Bare legs streaked
by the cold air of winter in Pittsburgh. This mistaken tongue began,
without words, searching for the derangement of some other body.
Unfaithful words spoken aloud from my mouth. Words, refused by gods and
demons, thrown out against the dark into strange bodies reflected on the
bedroom wall. In the corner I saw God standing in fire. There, I saw
God. Silent. His awkward elbows and jaw. An open invitation to my own
mortality. In the beginning before the flood of language my blunt
fingers explored this body, mine, which I had been ordered to ignore.
Curious about the torn skin at the precise moment when I came to know
that God sees all and follows me. Into the attic. I climb wooden stairs.
God leaves me forsaken, father places his hand on my shoulder but I only
see a reflection in the mirror of some dear dead memory. Feel his cold
body close behind me. His gray hand trembles against the beauty of the
mirror. My belly warm and soft. When I try to move, the ropes tighten.
Into me. My father cries. Spilling blood and skin onto the hardwood
floor. I fall down on my back. The weight of God holds me to the earth.
Paralyzed in time. In some countries this is the sign for the awakening
of love. My thighs like the threat of open scissors. My lips tremble.
Speak, why does he not speak? So much dust.

I was taught to fear monsters under my bed, bogeymen in the walls, and
angry angels hanging upside down from the ceiling. ``They will," Daddy
told me,''bite off your toes and toss them into the Monongahela River."
Mother flicking and flicking the light switch. Three times. Me begging
her to turn off the lights. Be done with it. She stood in the shadows of
the doorway chewed away at her nails. Between her teeth. Pulled them off
her fingers. Spit them onto the floor. In the morning my feet bleed.
``The monsters will carry you off to hell,'' mother warns me. "In hell,
you can't ever be a girl." I keep my eyes open, wide-awake, deep through
these godless nights into mornings riddled by junk madness. In
stillness, not a breath, not a word of this to anyone. I wait never to
speak for my father to come home from work. Every minute, every night, I
wait for my father. He with his body uses his body. The one God had
given him at birth. In the name of his body, he interrupts my isolation.
He takes away my weariness. I never cry. Not once. I never shed a tear.
He likes that about me. I have heard that others have died because of
their tears.

Nine years old. 1966. The naked throat of a child. Cracked lips dry and
peeling. Born choking.

These guilty lips, mine, opening, stuck inside prayers of becoming a
girl. For my father. So that he need not use his fists so often. So that
his anger is more quiet. More gentle. Touch and let me touch you. I
offer him each desire of my body past before I became known as a boy.
She calls me her infant. But that is a dream, or it is years later and
we are in some parking lot in California. Her teeth are in my
collarbone, but this is my father. Now I lay me down to sleep. Each
night I lift my voice, whispering up at the plastic flesh of Jesus
nailed to a wooden cross. Forsaken body of childhood. Grace of the
Savior, my father. In his arms I want him to hurt me. To beat my
disobedient body into sudden desires and thrash my soul into some
unknown language. More. To want more. His knuckles in my blood. In the
image of. Stained hands and cold knees. I hope to be left alone on some
riverbank. Discarded. Abandoned beneath a bridge on the Southside given
over to the circle of homeless drunks sitting around their fires. Their
warm hands lighting my flesh. The burning prayer I can not say. His lips
on my lips her lips I remember. My back pressed into the bark of a tree.
Pray. This flesh, mine, my soul to keep, haunted day and night by
stunted freaks and biblical plagues. Locusts near my eyes scratched
their way into my dreams. The Lord. Never, Grandma warned me, again and
again, never look directly into the open mouth of God. If I die in the
eyes of God before I wake. Apologize. Say it. You must. Say this, the
cutting of your body, into words. My sore mouth close to the feet of
Jesus. A droplet of wounded blood, perfectly round. Forgive me, Son of
my Father. Pray without cutting into your skin the Lord. No.~I do not
fear the boys in the neighborhood. The ones spitting on street corners.
I only fear displeasing you. My soul this pain in here to take. Every
single night small men with dusty skin sneak through my bedroom window,
wake me from nightmares, from God's breath with their mad tales and
pointy fingers. Seven screaming horses. Their voices, pain escaping from
the shadows. Breathing their voices into my bones. I could not see
myself dressed in these clothes. Drunken nights on a boat. I never did
come to know their names, to understand how they had traveled from the
old country through the forest. Just one more lonely night. Pinned and
wriggling on the wall of a sawdust hotel. These men with their
uncontrollable betrayals mixing memory and desire spoke in tongues of a
world out past Babcock Boulevard. Carried rumors to me of a land that
was somehow not Pittsburgh. Me, silent and still, in my barren bed all
alone filled with waiting. Glow-in-the-dark familial lines of blood
being worn away by the passing of each minute. My mother's alarm clock
on the night stand breaks into my bones. One by one. This rib then that
one. Ticking my body into a nightly ruin. My all alone, nearly dead
eyes, squeezed shut. Desires of my body unable to forget those foreign
fingers pressed into my neck. Those tiny men plucking my eyelashes. Each
morning, my eyes inside a pain of not wanting to see. Cramped muscles.
Hot and hard. Tendons torn out of joint. All night long. And I wanted to
be entered. Dreamt of being entered the way any woman can be entered by
any man. With the passion of Jesus wandering through the desert, I tried
to push my eyes back into my soul. Push my eyes so deep down into the
inside of me that I would make myself into becoming blind. Wordless.
Invisible girl in this decayed bed. The bed that my grandmother escaped
from sometime in the late 70s. She had fled, screaming uncertain names,
into the forest. Words flung loose from between tight lips. White
knuckles. I stared up at the disappearing white of the ceiling, watched
as God turned against Satan, the teeth of Eve through the skin,
believing that if I looked hard enough beyond original sin then I could
lift my body right out of my bed through the ceiling straight into the
sky. My hands balled up into tiny fists. Frail fingers and prayers.
Shouts into the street at a disappointed God. Parasite girls in the far
off and long ago tugged at my lips. Sheets held to my chin. This throat,
my memory of speaking, raw and empty. Without a scream. His desire,
mine. Her fingers close to the secrets of lonely eyes. Bodies in motion
but never in flight. Perpetual motion. Once she had told me that bodies
at rest are compelled to remain at rest until another force acts upon
them. In them. Break me open. The other side her mouth. The making of
her will. Her forehead cracked the light. Wanting desire. Quiet.

Every night forgetful neighbors prowl around the shrubbery in our
backyard, try to steal glimpses of my body through the window. My mother
wants me to say in words why these neighbors circle our house each
night, fumbling around, kicking the dog, waiting their turn. She wants
me to explain my mouth to her. The meaning of life: to take it in your
mouth. Swallow. Close to dying. A death deeper than any darkness
visible. My mouth the slum of your wretched desires. You break my teeth
with your lies. Cut the corner of my mouth with a knife. You are that
big. I lift the blinds, stare outside at these neighbors. The ones with
all their eyes, trapped, looking in at me. Blinking through the
blindness, I search their bodies for signs. The window locked shut. My
mother seals the inside and the outside of my window with fresh caulking
every day and night. She checks the nails she has hammered into the
wooden frame of the window. Making certain that the window could never
be opened. No air. Just to breathe. A breath. Outside, leaves moving.
Concrete. The twigs of trees like the scratched bones of my soul. Soiled
underwear. The mud of the river of my father. I keep my knees locked
together and beg God to find my body unharmed. Forgive me my fingers. My
belly warm beneath. Rising to his touch. Skin ripped through by the
beginning, in the making of desires. I smell you on my body, Father. I
say my word against your suffering. Try to make noises. I say your name.
You refuse to hear my voice. Under your mouth. Say my name. You, then,
enter me like a clumsy foreigner seeking a new home in unfamiliar
streets. Your desire penetrates me here and there and here, again and
again, until my own name becomes incomprehensible to me. Every night my
father dreams. Sometimes his dreams make him speechless. Quiet mornings
listening to the radio report the weather, traffic, sports, news. My
mother moves from room to room turning on every light in the house. Even
the flashlights. The nightlight in the hall.

My father sits at the kitchen table sweating. The dead of winter and
still my father sweats. Nearly dies twice from heat exhaustion.

Day after day my sister stays alone in her bedroom. She is no longer one
of us. She has abandoned hope and become an invisible memory. Divided by
two. Tin soldiers and Nixon coming. In the summer of 1979 she burned all
her books. Said she would become a vegetarian mathematician with an
attitude. Would disprove every single idea that Einstein had ever
conceived. E=mc2, my sweet ass, she said every single night at dinner.
My father would slap her across the face with his hard bones and make
her go to her room for the night. Running down the hall, she would
scream that she, from this day forth, would only have sex with animals.
Only consent to wild, senseless fucking. Sudden and random.

At night, this father's body became another anonymous encounter with my
flesh.

The dried, flaking skin of my fingers. Rotten flesh. Have you ever
touched me in the winter? The miracle of her fingers inside my mouth.
But she looked away. Picked at the bark of an old tree. For a long time
I have been frightened of the uncleanliness of my Father. I am impure. I
say, ``I am impure." My body, I longing for one warm night. One night of
sleep. A way back inside where she forgets who I am. The dark room where
there are no lights. Nails, inherited from ancient women driven mad
against the moon and sons of fathers without speech, cut into my palms.
I live inside this house of silence where no one breathes. \emph{Not a
word of this to anyone.} My voice joined those of the ventriloquists in
the attic. But once there had been echoes. Noise from the streets of
history. On the other side of windows and mirrors. Persistent memories
locked inside some lonely cry to angels or deserted flesh against the
walls in some back alley. Left in that alley without a crumb trail.
Nothing to follow. No way back home. The policeman armed and ready
looking into that alley at me. Watching. Waiting for me to make my next
move. Hands and boots alert, waiting. Men in shirt sleeves leaning on
windowsills with all sorts of things to say about the Promised Land.
Sauerkraut and kielbassi. Rotten beer and cigars. My father gets me
drunk again. Old men in the shadows stare out at me. Look into my eyes
little boy. Come. I try to forget my name. Blisters. Cold sores. A mouth
alive with pain. I can barely unzip his pants. The scar on his thigh
terrifies me. He stands solid. Unmoving. Not even flinching. Smokes a
cigarette. Thinks of anyone but me. I keep thinking that he will fall
over on top of me. That he will die. Or that I will die. But he just
closes his eyes. Dreams. And neither of us seem able to die. My knees
hurt. Stone cold heart of gold. We go on living as if nothing is
happening. His pants at his ankles. Young boys turn their pockets inside
out and shrug their shoulders. Brilliant colors in the palms of their
hands all on fire. Offerings from nearly dead street gods. I buy a
stairway. Look left then right. Always look left first. Grandmother
logic from the old country. Knock on the side door to the Piano Bar.
These, my winter fingers, pick at the open diseases of my infant skin.
Cook up some damaged dreams. Mad bubbles boil on a narrow, tight spoon.
The needle takes on an eloquent life of its own. I tap a nearly
invisible line. A shot in the skies. Another lost angel. My eyes burn
white. Red scars. Infected memories. Rashes of my spirit. Blood and
blood. Skin of skin. Twisted hands. I cried beneath my desire. I will
not move from here until you have restored the inside of my throat. My
fingers I dip into the Holy Water. Receive communion down into me with
loving you. Small lips opened and waiting. Interrupted by kisses. Each
night more desperate than the night before.''Just one more time," I hear
him say. I cover my mouth with my hands. Try not to breathe so much.
Afraid I will awaken God. If only for some water. If only some water
beneath this red rock. My skin dried by the sun. Rocks. Rubbing. Toil.
He pulls my hands away.

My only begotten mother, oddball woman, this mother lived the life of a
stranger to my blood. A woman unknown to me. A deaf woman off in the
distance. Always forever standing in the doorway but never once entering
my room. Frozen eyes. Nothing better to do. Watching. Silenced lips
cursed my birth. Dry mouth. Moving her feet against the carpet. Static
electricity with no place to go. Her soul lost in the colors of some
forgotten desert. Misshapen words. Crippled mouth knowing only the lies
of alcohol burning her bed. Told me once-upon-a-time I would never be a
girl, a real girl, the kind men die for.

~

I creep down the hall. A hallway so dark that each time I walk down it,
I suffer the joy of amnesia. Forget my name. Connect nothing to nothing.
My sister, not me, trying to say I. This is my body to follow the will
of the Father. Once my father made me sleep in mud. Rain. God makes it
rain and the earth turns into mud and me unlocking the door I enter the
house. My body leaving marks on the kitchen floor. God's soil. Wet.
Footprints. Toes in the carpet. In my sleep, I open the door of my
parents' room. I always wanted to be a good girl. Steal my father with a
whisper. Take him away from my mother's body into the shadows with
promises of kisses. I drag him through the dark corridors and up onto my
bed. Pull off his clothes. And he begins the punishing of me for having
eyes too much like my mother's. My hands on his back. Fingers digging
into his skin. To die with his red blood under my fingernails. To
swallow him. His cum inside me. So my body will become a witness. I
throw my head back. Tossed scraggly blonde hair. Dust. Blind ancestors.
Holding tight his strong body. Desire in my hands. Gave me his love and
my hatred. I rub my lips. I rub these lips. My lips. I rub my lips. Wild
and free. Fear. Keep silent as before and before that. There are no
stories of my silence. I catch my breath. Hold still. Wait. Skin. My
eyes burn. Salt words, grandma called them. The desires of the body are
tears. I shed these tears onto the hard floor wanting you. In some other
life she told me there were rumors of maps. Told me that concrete was
not impossible. My eyes drifting. Out of me. Follow the white lines.
Across the mirror I breathe in the mad suffering mysteries of angels.
Drag the angry burn into my blood. Teresa and I call it bone surge. Away
from me. I look away from all that is happening to me. Searching for
roads, paths that lead into the woods. Dull morphine desires. Wanting.
Want. Something. I stare into my open hands, her blood, and try to
remember desire. She had bleed onto my fingers. Flushed forehead in dead
heat. Spit from my mouth, from someone, on her skin. Mixed with the red
of her cut. I put my lips to the center of the palm of her hand and kiss
her blood. My body feels dry. Or I feel my body is dry. Broken wrists.
The back of my head against the wall. The knees I spread. I give into
you. Naked hand on the thigh. The other holds the whip. You will listen
to my voice today.

I remember my father could not talk in my bed. Night after night, my
deaf and dumb father cracking open my forehead. Burning into my bones.
His elbows in my flesh. Fists everywhere. Purple morning skin. My
wrists, twisted behind my back, small in the traps of my father's hands.
Burning Indian. Blemishes of a father's hurried love. Promises made,
spoken in spit. Every inch of him on top of my skin and bones. Him. My
father weighed 183 pounds. This, my body will never forget. Sometimes he
weighed more. Sometimes less. His knees slightly above my knees. His
chin down on the top of my forehead. For years I practiced breathing
into his neck. His sweat on me. One day my father died. But that, his
death, comes years later. Or maybe it was yesterday. I make marks on the
wall. Counting something. Days. Nights. Both. Or desires.

\hypertarget{elizabeth-robinson-sleep-is-the-only-durable-narrative}{%
\chapter{Elizabeth Robinson: Sleep is the Only Durable
Narrative}\label{elizabeth-robinson-sleep-is-the-only-durable-narrative}}

\begin{verse}
A person is falling asleep.\\
There is some delay while the body shifts. Breathing\\
slows. He is\\
apparently asleep.\\
But then he moves again.\\
Legs above the blanket, legs below the blanket.\\
The changes in rhythm are interminable. But this\\
is the story.\\
\end{verse}

\begin{verse}
Sleeping itself. When he is\\
finally asleep, the body can resume\\
remaking itself or\\
call on a project of outright fabrication. He doesn’t\\
know his own story, except as he respires it.\\
Cells re-coat. Dreams are irrelevant. Connective\\
tissue lengthens, glossing itself.\\
\end{verse}

\begin{verse}
He is a wonderful entity, finally\\
voiceless and productive. Except were he\\
to fall apart in his sleep, a possibility\\
that has its own unity.\\
\end{verse}

\begin{verse}
This is mostly a string of overlapping\\
events gone unwitnessed. Unless someone\\
who cannot go to sleep is lying there too.\\
She goes through similar shifts and\\
sighs. What validity does\\
his sleep gain\\
\end{verse}

\begin{verse}
by virtue of her not sleeping— She decides to notice\\
his repose; she decides not to.\\
She considers that she is eavesdropping\\
somnolence, but it’s not so\\
\end{verse}

\begin{verse}
in a tale that cannot be permeated. Corpuscles\\
mending themselves. Insomniac,\\
she confesses that she is eavesdropping\\
her own unmaking, awake. This is not\\
an anti-story, but an alternative.\\
\end{verse}

\begin{verse}
White lights come in through the window,\\
begging a source. Now the frazzled body\\
rears up, claims that light as more than immanent, dying\\
to spy, at least, on some body. However atomic\\
or abrupt.\\
\end{verse}

\hypertarget{jono-schneider-from-the-parable}{%
\chapter{\texorpdfstring{Jono Schneider: From \emph{The
Parable}}{Jono Schneider: From The Parable}}\label{jono-schneider-from-the-parable}}

\hypertarget{working-notes-6}{%
\section*{Working Notes}\label{working-notes-6}}
\addcontentsline{toc}{section}{Working Notes}

\hypertarget{writing-the-tellings-wanting}{%
\subsection*{WRITING THE TELLING'S
WANTING}\label{writing-the-tellings-wanting}}
\addcontentsline{toc}{subsection}{WRITING THE TELLING'S WANTING}

\emph{Reading this excerpt from \textsc{The Parable}, what comes to mind
is the uncertainty of the "I" that is speaking, that wants to speak the
speech of narrative, the speech enacted by writing. Narrative is a
story, but to focus solely on the object is to diminish the subject that
tells even as the subject is writing the telling. And this telling is a
wanting to tell that is also the telling's failure. The telling fails to
achieve the aim of the wanting and diminishes the story being told, but
not the wanting. All the while the voice of the invisible speaker grows
louder, deafening the reader with the sound of words babbling over into
writing. French literary theory calls this babble} desire\emph{; in
\textsc{The Parable}, this desire is part of the opening, is what opens
up the story's opening to the reader, the hole in the fabric of the
story allowing writing to shine blankly onto the narrative of character.
``Of that I can be sure"---with this phrase the main character is the
only plot. The speech moves somewhere, introduces two additional
characters than that of the narrator, two colleagues, but they are the
reason the narrator has a story for the reader by way of the tone of his
words, which gather up a serene severity like the light in a Bergman
film. The narrator speaks as long as the speech s/he makes convinces the
reader that silence is the only possibility. The additional characters
only serve to distance the narrator from the rest of the world framed
within the speech being made, the world that this speech unfurls and
thus uses to gain the reader's confidence. It is this confidence that
the story elicits, but only at the reader's behest: "the meaning of this
is unclear,'' as in "still caught in an approach---and this approach is
a parable which I am writing." I would second the idea of narrative as
something that approaches a subject, that approaches language even, but
does not touch it, and this subject's being untouched is what
\textsc{The Parable} radiates outwards towards the reader. As in
Polanski's \textsc{The Tenant}, the narrator goes crazy because he's
alone in a room---the total devastation of narrative, what makes writing
jump into the still air.}

\begin{center}\includegraphics{images/"pandanus.png"}\end{center}

\lettrine[lines=3, findent=0em, nindent=0.1em, lhang=0]{H}{e had judged me}---of
that I can be sure. His back faced me, and his head was bent slightly
forward; I could not hear what he was saying. He did not stand up when I
entered, though the meaning of this is unclear, since it was she who had
invited me to join them. Perhaps he felt that I had interrupted them;
they were, after all, talking in her office, not his, meaning he had
approached her because he must have needed her confidence. I had broken
that trust with my appearance at the door; I shouldn't have come in at
all had the door not been open and had the sound of their voices not
reached me in the hall. I had other duties to attend to---this is
true---but I was unable to begin, and when I heard their voices, I
sought them out to forget these duties. A simple conversation with one's
colleagues often helps one to pass over the entire world in silence, and
I longed for the communion with others my work so rarely provides.

I was dressed horribly, and, by the way in which he looked at me, I knew
that he was glad that I was wearing the clothes I did: this set me apart
from them. He lowered his eyes when I sat down and kept them on the
ground for quite a long time---I had looked away, and I am not sure when
he finally raised his eyes to meet, to almost meet, mine. My clothes
bothered him, that was to be expected. I had tried other clothes on,
earlier, but I felt false in them, as if I were lowering my voice to
emphasize my authority, thus calling attention to my mere lack of it.
No, these were the clothes I must wear, and these were the clothes I
would be judged in; if these clothes added to the severity of his
judgment---well, I would readily accept that. I spoke too much---this is
also true---but how else could I have spoken? If I said too little, then
he would rush in and fill up the space between the three of us with his
silent judgment, and when I spoke, his look silenced me as well,
impressing upon me the fact that I could never say enough for him. His
silence, which he expressed calmly with his mouth, where he hung a slack
smile, further silenced me, so that I spoke more quickly, more urgently
than I felt or would have liked to feel. Yet I did not feel anxious; I
accepted his dislike of me because it was the easiest thing in the world
for me to return. I did not hate him---let me be absolutely clear on
this fact alone---but I could have hated him, I found it easier to hate
him than even my own clothes, which I admit were horrible: frayed at the
ends, ill fitted, and ballooning out from my body. Surely no one could
take me seriously in them.

He seemed concerned with me. He was not speaking to me---he would no
longer look in my direction now that my speech---my stuttering and
stunted soliloquy---had ended, although now it is easy for me to think
of what I should have said, what kind of person I wanted him to think I
was, and what kind of person I was when I was alone. I wanted him to
know how much I cared---not about him, to care for him would have been
going too far---but how much I cared about life itself, by which I meant
the people for whom I worked and to whom I spoke and would always be
speaking, the people who were the reason for this meeting---but I had
entered the room to listen to him speak about his own experiences in
dealing with them. Even though he was now speaking---and he was looking
directly at her, he had moved closer to her, their knees were almost
touching while I had somehow pushed my chair further into the
corner---even though he was now speaking only to her, and I couldn't
hear what he was saying, she was looking at me, she was smiling, she had
liked what I said. Yes, she liked me very much. She thought I was
intelligent, an intellectual. Yes, an intellectual. But I could not make
him turn around, I could not make him face me. I wanted him to tell me
if he had eaten.

"But I'm eating now; we're both eating now, don't you see?"

I saw no food on the table---but the table had just now appeared, or it
had just been cleared---and their mouths were moving words, they pushed
words through the air---perhaps these words were being tested. I could
hear these words now---perhaps the words were being tested---they were
faint against the table and the door, they were talking only of work, of
speaking, and how to speak soundly without raising one's voice. She was
not eating---she was looking at my pants; I had at least taken care of
them: I had ironed and pressed them, I had scrubbed until all the stains
had disappeared. But I was sweating now, and he had all the advantage
for it. I was smiling, I mimicked his smile and flashed it back at him,
she smiled, we were all smiling---but our faces no longer expressed
happiness. I felt slightly sick, but I was still hungry, and the fact
that I had not lost my appetite caused me to cry out---now I was feeling
slightly giddy.

"I'm now on my way to eat---" I stepped towards the door.

"But won't you join us?" She was holding out her plate of food to me.

I took it; he handed me a fork. I would eat, which would give me the
opportunity to listen; they were letting me in to listen, but only if I
kept my mouth still, which meant I had to fill it with food, and that I
would have to stop thinking. I was an intellectual, and I may have had
many thoughts which I needed to tell them, but for now I would have to
eat and be grateful for my work.

\begin{center}\includegraphics{images/"pandanus.png"}\end{center}

Gratitude was a state of being, a condition, towards which he aspired, I
could see it shining on his face and reflecting back at mine as he
strained to choose his words. His voice was deeper than mine, but he had
no advantage over me because of it, as she kept looking to me after he
finished each point to see if I approved of his ideas. Although I
thought that his intelligence was limited---and I still think this,
although my certainty is surely slipping---I sought to credit him; his
dislike for me registered itself in the air among us, and I struggled to
keep my breathing even and measured by grasping irregularly at the air
with my mouth. He spoke, and he would always speak, first.

``But you should eat more slowly! At the very least you should know
that!''

She had opened the windows. The air was slow in coming, but when it
finally reached me, I already felt that it had calmed me and that the
difficulty had passed. I had dropped my fork, and I saw my own face
reflecting back at me from the floor, mirrored by the prongs of the fork
that now lay at my feet. He was inspecting my shirt; I had left my shirt
untucked because I could work more easily without it stuffed into my
pants. He handed me a napkin and resumed, his talk easily inflated by
this gesture he had made only for me.

"One should always be present. Yes, at the very least, one should face
one's responsibilities!"

His hand had taken my plate away from me. How very odd this seems to me
now, I must say, for just as I had earlier seen no food, I had now grown
quite hungry, and I had no other desire than to lift the fork from the
plate to my mouth and back again. It seems to me that it was only when I
had begun to eat with emotion that he became interested in clearing my
plate, as if his manners masked a deeper, more mysterious motive. He had
turned his back to me again, now he was so close to her as to be
whispering. Her face had contorted at his unexpectedly close approach,
and I thought I noticed him trying to comfort her. My discomfort was not
caused by his hatred---this I can clearly understand, had I been him I
would have acted with less tact, that is certain---but from his attempts
at co-opting her into his hatred, which only furthered my desire to
destroy her. But I did not intend on lifting my hand against her; what I
needed was to get her away from here, I needed him on his own to make
him understand what I was saying, why I sat the way I did, why I liked
the food I did, etc. But as long as she was there, trying to speak to
him on my behalf, I would have no opportunity to teach him about me; as
long as she was there, he would continue, he did continue, to look at
her, to desire her approval, which, as all three of us were aware of,
she was lavishing upon me because she was an intellectual as well. He
might have also been an intellectual had he not spoken in such earnest,
had he not let his need to have someone else support him rise so clearly
to the surface of his speech. Yes, my speech had failed badly,
miserably, but I spoke with the detachment of a man who has yet to have
said, who is still to think, what it is most important for him to say.
His speech, you will say, was more honest, more in earnest than mine;
here I will not hesitate to agree with you. But I had not been trained
to be honest---I had been trained, I should say, in policy, and having
learned policy to the point where I could endlessly elaborate on it in
terms of form, content, and commentary, I was fearful whenever
conversation required that I venture beyond my training.

I wanted to continue eating, but it seems to me now that it was clearly
impossible to do so; as I said, his hand had taken my plate away from
me. But more than that, it was the way in which he was now speaking that
made eating impossible, for he had placed my plate--- he must have
filled it again after taking it from me, for I was so hungry that I had
cleaned the food off of it almost immediately--- he had placed my plate
on a small filing cabinet just to his right, where he had also laid his
glasses aside, and now he picked them up, although he didn't put them on
his face---he kept kneading them with his fingers while he launched into
a new speech about the endless dedication to one's work and how this
dedication must be final, the endpoint of one's goals in life. He spoke
about the future and how it could be filled with nothing but work itself
and those whom the work served. He was furious now, he was no longer
whispering, he was standing between the two of us, and the look on her
face told me it was my turn to act by replying to his emotions with the
logic, reason, and wit of an intellectual, but I was too preoccupied
with the food that he was keeping from me to think with the clarity that
his vicious speech demanded me to return. But it seems to me now that
this is an oversight on my part, that I could have easily stepped in
front of him and grabbed the food while he continued, and that he would
have thought nothing of it, because he was too involved in other
actions: he was mutilating his glasses, twisting the frames into shapes
that resembled the faces he was making at the time he was speaking to
emphasize this speech and the time that it took away from us. I very
well could have done this, which tells me that if this action was
possible, then I must have done it, for I have learned that one does
what it is possible for one to do in any situation and no more. So this
is what I must have done, since my clothes had already disarmed me --
although I was happy, perhaps I even felt ecstasy, in them. I reached
across him and took back my plate, but the food was cold, and I ate it
hastily and without pleasure, meaning that this moment, as did the
others, belonged to him, and he had invented it purely to vex me and
cause my immediate dismissal.

But I did not withdraw---if you read the transcripts, you will clearly
see this, there are three voices on that tape. It was she who kept
silent for most of the remaining time, smiling when I pointed out
options that he hadn't yet considered, and pulling her mouth taut when
he derailed the conversation with anger. The tapes tell the truth, that
is why I carry my recorder in my pocket. How was I to know that he was
taping all of this as well, and that he would hand his transcriptions
into the committee before I had even been given a chance to listen to
his tape? I was cautious, yes, and I had seized the opportunities that
were available to me, but I had not foreseen him. In fact, I knew next
to nothing about him before entering the room where they had been
quietly talking, I had gone in there to see her, but having heard two
voices, I was, as I said, quite anxious---too anxious, it is now
clear--- to be a part of a conversation which I was confident I could
add---and learn from, and yet I wanted to defend myself before any
rumors had begun about me. You will say that I am nervous, that it is my
nerves--- and I am quite the nervous type, I will not deny this---but I
have found that it is better to circumvent what is said about one than
to wait for the speech to arise and become writing, or policy, as they
say. They say this, I know, because now I can hear them talking in the
room adjacent to the empty cell in which I am writing this, and I have
heard them say this word---"policy"---no less than ten times in the
half-hour that I have been able to overhear them. Surely, you will say,
that ten times in half an hour is not that many times for a word to be
said--- once every three minutes---and this is true, but in relation to
all the other words that I can make out, this number is quite large,
since they have repeated not even one of those other words, other than
my name, which I hear said at the beginning of each sentence, and they
argue knowing I am listening to them speak, and they do not care.

\begin{center}\includegraphics{images/"pandanus.png"}\end{center}

He had handed me a toothpick---I am certain of this, else how would they
have found it in my pocket later, after everything had come to an end? I
remember thinking that I was surprised at how sharp it was--- it seemed
to me that someone had intentionally sharpened the toothpick to keep me
from putting it in my mouth---and perhaps this was what he had
done---instilling fear in me: fears of my tongue being cut, my speech
being disabled, and my eloquence being reduced to a whimper. So I put
the toothpick in my coat pocket, and I put my hand in my pocket from
time to time to make sure I had not lost or forgotten it---if I knew
anything, it was that I could keep this tool in perfect condition, I
would prove to him that it was possible that an intellectual could take
care of the body. She did not receive a toothpick, and while he handed
one to me, I saw the look on her face indicate a change in the room. Now
the window was wide open, the wind brought fresh air into the room, but
she had changed, too---his gesture of offering me something that was his
and his alone, something that preceded this meeting, had excluded her,
and from that moment on she was lost to us. I was very happy---I will
even go so far here as to say I was ecstatic, and the traces of this
feeling are still with me now, for it was no longer my responsibility to
get rid of her, a job I had feared from the outset because I had thought
that whatever I would do to her would not escape him, and I didn't want
him to know, or to guess at, anything else about me besides what I had
already told him. He may not have initially intended this effect, but he
surely could not have failed to notice that she was gone, she was
leaving, she would be leaving, and even if she didn't move from her
chair, we both knew that she would neither speak nor look at us for the
remainder of the afternoon. She only looked down at her hands,
alternating this with a look at the floor, successfully avoiding us, no
doubt, but only because she had to avoid us. So now we were entirely
free to say whatever we needed to.

And yet I soon missed her, for the talk was a tougher climb than ever,
as he asked me to repeat myself several times, and more than that, to
define certain words for him. Really, I was defining everything for him,
he couldn't understand me, and because I was an intellectual and he was
not, we had no language in common. And I was straining, I was at my very
limits---limits which it has, you must know, taken me years to build and
even longer to expand---and it seemed as if my words would no longer
hold me. I thought back to the quality of the speech between the two of
them before I had entered---there was something remarkably easy about
it, as if there were no need to test the value of language itself
because their talk was not rooted in ideas. No, they were merely two
people talking, and they still talk this way as far as I know; it was my
speech, the speech of an intellectual, that was the voice of another,
the voice that suggested others. But here you will point out that I have
said she was also an intellectual--- this is true---but she was only an
intellectual when she needed to be, when the situation demanded this of
her. This demand arose out of my presence, since any other speech
handcuffed me and angered me to the point of irrefutable testimony in
the form of logical arguments or "auditory editorializing" as I like to
call it. I have only this term to introduce to the world, and when I
told him of this discovery (my new way of working), he scowled in
disapproval and then launched into a prolonged battle cry that
exemplified my technique better than any instruction manual I could have
written for my employers. By hating me even further, he had miraculously
understood me more fully than those for whom I had created, and still
dictate, policy; I had done this to make him obsolete, and now I was in
need of him even more, while I became an intellectual at an even higher
level, and I spoke to him from an even more acute distance, than before.

I am now well aware that there are still many questions that I am in
need of answering, and yet I press on because I must ask still more
questions in order to be fully prepared for the final inquiry which may
take place at any time, including the past. I can feel your impatience
on the other side of the page, but I urge you to calm yourself, for the
obscurity of my documentation hardly approaches my own impatience and
the harshness which rises up in me from suppressing this patience each
time I try to move ahead too quickly. Perhaps you can think of me in the
following terms: if you are asking me a question, I am asking the same
question to a dead man who is still dying, a man who has never been ill,
as his illness has not yet reached him, it is still caught in an
approach---and this approach is a parable which I am writing.

~

\hypertarget{kathy-lou-schultz-chapters-5-7-and-8-from-some-vague-wife}{%
\chapter{\texorpdfstring{Kathy Lou Schultz: Chapters 5, 7, and 8 from
\emph{Some Vague
Wife}}{Kathy Lou Schultz: Chapters 5, 7, and 8 from Some Vague Wife}}\label{kathy-lou-schultz-chapters-5-7-and-8-from-some-vague-wife}}

~

\hypertarget{both-and-a-statement-for-narrativity}{%
\section*{"Both And": A Statement for
Narrativity}\label{both-and-a-statement-for-narrativity}}
\addcontentsline{toc}{section}{"Both And": A Statement for Narrativity}

\emph{Bob Glück gave me back the sentence, a line of elastic stretching
across the page. I had stood at the edge of the lyric and fallen in. I
had swallowed the lyric and choked on it. I had been lured into the
lyric's boudoir and slammed my elbows against the walls just trying to
turn around. I was out of metaphors. Though I still engage the lyric,
writing in prose allows me a kind of freedom to explore the nether
regions of sexuality and gender, to stroke the unnamable places, to
crack my head up against the jagged edge. My world is troubled by
"either ors": straight or gay, male or female, black or white, poetry or
prose. I'm lodged in the cracks between them, dynamiting my way out.}

\begin{center}\includegraphics{images/"pandanus.png"}\end{center}

\hypertarget{adjustable-perches}{%
\section*{5. Adjustable Perches}\label{adjustable-perches}}
\addcontentsline{toc}{section}{5. Adjustable Perches}

My boy-girl princess all groovy this evening, sharp with elbows and
shoulder blades. Desire across boundaries too sharp to cross. 24 hours
of comfortable heat.

Our skins slide across. Sweet-cheeked girl peck. A baby's unformed milky
bones.

A pact to reflect daily resonances. ``I miss you," means "Do you miss
me?''

A cupboard and a roof. Rough wooden shelves soaked with white paint. The
smell of oil flares our nostrils.

Your hand encloses mine, ratifying physical laws when you dismantle me
step by step. My weeping fresh and startling, you contract your
awe-struck mouth. Anticipatory like the smell of rain.

The temptation to bring you inside, where the light slants across my
desk, warming my pale skin. Afternoon lassitude of linguistic junk.
White boats across the lake.

Returning is a grand gesture emphasizing departure. When I finally lose
sight of you, the earth still suspends my body in its animation.

My freshly parted hair still damp from the morning, when you turned,
wanting me, all sleepy-eyed softness.

The road ahead craggy and distant multiples our longing, my breath in
hiccups, false stops until we find the wave that continues.

We proceed as if finding bread crumbs on a trail.

The containment I experience through your matter-of-fact nod proceeds
outbursts designed to unloose your fibrillar emotions. Dichotomies which
bore me.

What's the gift, the partial barter, which promotes our stasis---you
waiting for disintegration, and I for a stable distance. I can accept
your intimacy in your absence.

Words an elevated foreground, everywhere falling fictive, furious
scatter shot. I slip the pellet underneath my tongue, hidden in an
enclosure of human temperature.

In our shared understanding, there are few formulas to interpret this
new-found wound, wound round characters you describe as real. Goals
program the future requiring the abstraction of my body in front of you.

Time telescopes individual events, cheating clarity. Metaphysical
photos, faces sunk in relief. For now, kiss me in the way I remember and
let us look into the light of the morning.

~

\hypertarget{bait-and-switch}{%
\section*{7. Bait and Switch}\label{bait-and-switch}}
\addcontentsline{toc}{section}{7. Bait and Switch}

Or simply that it's me and not other other. Even if it is.

Dodie turns up the "D" in degenerate. The irony that fucking boys--- not
girls---made me into a first-class pervert. I laid claim to his dick
like it was my own---a rightful part of myself gone missing. This one
said he couldn't fuck me without getting emotionally involved. He could
clasp both of my wrists in one muscled hand. Prone, I lay on top of him,
my body floating on an island.

In a semi-trance state he discovers his desire to kill his mother. He
says, "Many women had paid for what you have done." In fairy tales, Mom
is always to blame. My left thumbnail refuses to adhere. If I were to
fall in love, I'd have to bring the white trash with me. I hide my
guilty desire to spend long periods of time alone.

Of course, when I wanted something, I always knew it. The very intensity
of my desire keeping it out of reach (or so I believed). She kept
herself from fear or disappointment by turning off her own desire
completely---while I was filled voraciously with need. Greed, while
unbecoming in a woman, is prized in men.

Another claimed that my very presence increased his need to be more
"manly." He opened doors for me, refused cute pet names, and fucked me
righteously.

My body demanded something to show for itself. Writing is thievery, as
in stealing time.

Anxiety is a sticky substance infused with fear. Dollar for dollar. Or,
for instance, poverty. My own collusion in bourgeois appearances
bleeding me dry. The need to be seen or recognized outweighing other
emotional vaunts. Every song he laid in my ear keeping my mouth and nose
above water.

But respiration takes place in lack of desire. Once your head hits the
floor, the body resumes intake and expiration. The old in and out.

Anger, not sex, is what makes a girl really bad.

~

\hypertarget{epilogue}{%
\section*{8. Epilogue}\label{epilogue}}
\addcontentsline{toc}{section}{8. Epilogue}

Her indestructible thighs. A vertical line going south. From belly
button to pubic bone. Brand-new scar.

She wears high heels like in dreams. Stunned. Apostate.

I could have said he only fell for tragic women, but I was falling for
tragic men. You know, their darkness, their pain. ``You light up my
life, you give me hope, to carry on,'' Debby Boone was singing to me.

I turned to my girlfriends. They had strong shoulders and clear
thoughts. Wedged in the back seat, her arm brushed against my breast and
rested there.

A man who saw the world through film, but experienced a lack of vision.
A man I encounter through dailiness, through habit. A man with the hands
of my father.

Distance spreads. Across your lips.

Your hips move me. Hand in glove.

Pink. Drawn into the skin. All summer. Or all winter. Scripted. Cut.

Rushing headlong in with the body and checking positions later.
Velocity.

Or when his dreadlocks were a lot less polite looking. That it was all
fair game. Down to the very last vowel. Avowal.

Two loulous wrapped rapt in bed sheets.

Baby's got Mingus Fingers. All around and through me. To thrill me.

Baby done gone and you know it know. Something to show for. Here and
before now. Baby done proved it.

All in a groove now. Here and before now. Oh, baby, you knew it.

Girl to girlness reveals a surety surfeit. Our beliefs clamoring for
real estate.

Churl. Countryman. Highlander.

Every possible apology or minimization. Privation. Dilation.

She dressed with her back to him, then left the motel room and walked
down to the water. Bare-legged and bare-armed in her cotton jersey
dress.

She had a sense that the dream had ended and she had slept on. Further
and farther. Up and away.

The first time he entered me he said, ``I think we've got a match
baby,'' and I sealed this sentence in my heart and swallowed the key.

Unmatched by, ``I think I'm falling in love with you, you
son-of-a-bitch,'' and later this made me smile, laughing at myself.

What she was staving off or holding out against. A grammar of
indeterminancy. A cold water walk-up.

She had watched them in supermarkets and knew the signs. All the indices
of the idle lonely.

Happiness came like an unmarked box without directions. She kept trying
to catch herself in it, afraid it would pass without her noticing.

Inscribing my language upon him. Brush strokes along the abdomen.
Purple-black nipples. Smooth, round ass.

Endless expanse of what if. Figure brought into fruition. A thorough
dressing down. Round or through me. A voice on the wire.

Or if I wrote "wife." Rife. Theocracy. Meritorious.

Mercy. For the woman's clandestine, zippered heart. Miles I have come
before you. Previously unknown.

Everyone got what he came for. A schedule to rival the master's.
Departures and arrivals. A vagabond. A conductor.

Miles I crossed the desert for water. Previously alone.

At the end of it all a ringing phone and she does not rise to answer it.
It twitters at the edges of her sleep.

Verboten. Far gone.

She wonders if this will be one of those evenings when something
unimaginable would occur. And the language used to describe it later.

What he was willing to do for her, and the sexiness of that.
Eroticization of her pleasure.

He covers the receiver with his palm and turns away.

I arranged everything and at the appointed hour the three of them
arrived. Before any words were spoken, I recognized him by the sound of
his breathing.

She drove her new BMW up to the GoodWill and was confused by the people
staring. I smiled and tried to act "natural."

I was the second man to fuck him. The first he met while he was still in
school. "Don't you go to UCLA?" It was a kind of code.

At the end of the day, his body a destination. Miles on a train I have
traveled postponing pleasure. Our voices hang in the air, I touch
myself.

Miles I have traveled to hear his breathing. Loudly overtaking him in
sleep, a comforting wheeze which lulls me.

Prone. His locks falling forward across his face. My dick gets hard just
looking at his ass.

There can be no question or assumption. The girlfriend in L.A. The fits
of rage or depression. The unshakable self-centeredness. All waiting to
be read like a dime store novel.

My narrative spins out like a car in the rain.

Our lives a careening that could have taken us anywhere. The sheer
random force of it underscoring her belief.

She weaves a web of talk. Contentment, stranger to this snare.

Used to desire, but uncomfortable with need. She covers her head with
the blanket to keep from shivering.

The whiteness of my arms practically glowing in the dimmed light, up
against his dark, supple skin.

When I fucked him, a high, clear-pitched sigh rang out from his mouth
each time he exhaled. "Go easy, baby," he said. And I withdrew my dick,
softly, gently exploring his ass until I gave it to him once again
harder. He tossed his head back sharply, his locks raining off his face.

Fists of wanting. Emboldened by grief. A world apart. Of frost.

No sooner than. No later. Had she walked to the door. Bereft.

Solidity in solitude. A place in the life cycle. Searching on hands and
knees. He fucked me so hard that I bled. Four fingers in my cunt, nearly
a fist.

"What I feel for you now is all I will ever feel for you." And ever. A
straightjacket.

Given to histrionics. Or expected.

Protected more than risked.

Stutter. Choke. Recede.

My arms a fortress.

Each flag in a field.

The harder my dick, the wetter his pussy. The deeper I fucked him, the
more open he became.

\hypertarget{aaron-shurin-narrativity}{%
\chapter{Aaron Shurin: Narrativity}\label{aaron-shurin-narrativity}}

\begin{verse}
©1990 Aaron Shurin\\
First Published by Sun and Moon Press\\
6148 Wilshire Boulevard, Gertrude Stein Plaza\\
Los Angeles, CA 90048\\
\end{verse}

\begin{verse}
*This essay was first delivered as a talk at Painted Bride,\\
\end{verse}

Philadelphia, June 1989.*

\begin{center}\includegraphics{images/"pandanus.png"}\end{center}

I'm interested in the utilization of both poetic and narrative tensions:
the flagrant surfaces of lyric, the sweet dream of storied events, the
terror of ellipsis, the audacity of dislocation, the irreversible
solidity of the past tense, the incarnate lure of pronouns, the
refractability of pronouns, the simultaneity of times, the weights and
balances of sentences. I'm interested in lyric's authenticity of
demonstration and narrative's drama of integration; lyric, whose
operation is display, and narrative, whose method is seduction. I
describe a set of binary terms across which I see writing passing an
exchange of values, and it becomes a multiple texture/text---writing in
just those created tensions between surface vocalic tangibility and
referential transparency; between theme and emptiness, measure and
interruption, the eternal present and past of memory/future of dream;
all present, all heightened, operational. Such conflated writing would
be worthy of Barthes's definition of the text: ``not a coexistence of
meanings but a passage, an overcrossing; thus it answers not to an
interpretation, even a liberal one, but to an explosion, a
dissemination.'' \footnote{Barthes, Roland, \textsc{Image - Music -
  Text}, Hill and Wang, 1977, p.~159.} One seeks to be out of order, to
shiver out of subjectivity, to shake off the mask of the material and to
shimmy in its arms, to finally retreat from logic and advance by radial
maneuvers, gathering meaning. "To break the sentence," says Rachel Blau
DuPlessis---and here the sentence carries its overtone of imprisonment
without parole---

rejects not grammar especially, but rhythm, pace, flow, expression: the
structuring of the female voice by the male voice, female tone and
manner by male expectations, female writing by male emphasis, female
writing by existing conventions of gender---in short, any way in which
dominant structures shape muted ones. \footnote{DuPlessis, Rachel Blau,
  \textsc{Writing beyond the Ending}, Indiana University Press, 1985,
  p.~32.}

One looks for alternate methods to proceed, to use and subvert the codes
at hand: stanza, line break, character, plot, point of view.

In ``The New Sentence," Ron Silliman suggests ways in which the prose
poem has used combined and measured sentences to interiorize poetic
structure, foregrounding language operations and surface values in a
writing mode---prose---whose usual form is the syllogism, building
structures of projection and depth.''The torquing which is normally
triggered by linebreaks," he points out, "the function of which is to
enhance ambiguity and polysemy, has now moved into the grammar of the
sentence." \footnote{Silliman, Ron, "The New Sentence," in
  \textsc{Talks, Hills} 6/7, 1980, p.~214.} The paragraph as a unit of
quantity and the sentence as a unit of measure, altered sentence
structure, controlled and limited integration: these devices begin to
conflate the values of poetry with those of prose. Other writers have
pursued not just prose but narrative prose, and foregrounded narrative
codes to awaken a reader's attention to process as well as result. In
his novel, Jack the Modernist, Robert Glück uses metaphorical and
metonymic litanies side by side, showing off the writing as writing as
he demonstrates that the devices are not mutually exclusive.

I grab his cock, unpromising, and he says in mock bewilderment, "What's
that?" As it hardens I answer for him, "It's my appendicitis, my
inchworm, my slug, my yardstick, my viola da gamba, my World Trade
Center, my banana, my statutory rape, my late string quartet, my garden
god, my minaret, my magnum opus, my datebook, my hornet, my Giacometti,
my West Side Story, my lance, my cannon, my nose-job, my hot dog, my
little sparrow, my worm on the sidewalk after a storm, my candle, my
Bic, my unicorn, my drawbridge, my white whale\ldots{} \footnote{Glück,
  Robert, \textsc{Jack The Modernist}, Gay Presses of New York, 1985,
  p.~27.}

and on for another sixty substitutions, ``my cyclops \ldots{} my Venus
of Willendorf \ldots{} my Dark Tower.'' Four pages later the elaborative
metonymic process of prose takes over from the comparative metaphorical
process of poetry.

My troubles were too numerous to consider all at once, their sheer
quantity defeated me. My mom would say, ``Write a list, get a handle on
your problems, deprive them of their active ingredient, time.'' So I
found a clean page in my yellow legal tablet \ldots{} Nuclear
catastrophe, destitution, famine, additives, melanomas, losing face,
U.S. involvement in El Salvador and Nicaragua, Puerto Rico, South Korea,
Chile, Lebanon and Argentina, war in the Middle East, genocide of
Guatemalan Indians and extermination of the native peoples of Brazil,
Philippines, Australia, answering the telephone \ldots{} toxic waste,
snipers, wrinkles, cult murderers, my car \ldots{} \footnote{ibid.,
  p.32.}

Though these are both descriptive processes, they are not transparent;
the reader is aware of being in a list, enjoys the ingenuity of
elaboration and substitution, is held to the surface of the writing at
the same time she is integrating the lists into the larger structures of
the story. Speaking of description, Alexander Gelley writes,

This kind of stillness in the narrative may be likened to islands of
repose for the reader, moments of collection. The hold that the level of
plot, speech, and action exercises on him is loosened. His attention may
wander, but it may also adjust to a changed mode of apprehension. I am
suggesting that the more circumstantial the description and the more
separate from the narrative in which it is embedded, the greater will be
the reader's part, and the more he will be forced to assume a stance for
which the narrative proper offers little support\ldots{}. When the
familiar codes of narrative are blocked or diverted, reading/writing
becomes problematic, and the subject of/in the narrative shifts from the
characters or the author to the reader\ldots{}. \footnote{Gelley,
  Alexander, \textsc{Narrative Crossings}, The Johns Hopkins University
  Press, 1987, p.~14.}

This problematization forces the reader to ask questions, to become
active in the role of reader, and Glück reinforces this tendency by
confronting the reader directly in his stories, ``You'll understand my
fear,'' he says, "because television has trained us to understand the
fear of a running man;" and, "I can only give this story, which is the
same as sitting with my back to you;" and, "Tell me, given the options,
where would your anger have taken you---where has it taken you?" By
confronting the reader, Glück not only breaks the window of his
narrative but creates and engages an audience, creates a social
registration for his writing by direct address, by luring the "real"
time of the reader into the "dream" time of his story. The foregrounding
of devices and codes does not neutralize them, they are too full of
historical determination, but it can ritualize them, or expose their
ritualization; reveal them not as necessities but constructions---open
to change.

Writing might use narrativity without succumbing to its hegemonic orders
of linear development, unity of time/tense---and apart from the
modernist reconstructing modes of memory and dream. A prose whose
paragraphic groupings themselves might be based on measure, whose higher
integrations might be thematic or associational rather than
developmental. ``How tenacious is our happiness!" says Kevin Killian in
Shy.''Unlike narrative, it invents and eludes itself from moment to
moment; it lacks conventions; its shape has no outline, its formal
properties those of the cloud---numinous, portentous, hungry\ldots{}."
\footnote{Killian, Kevin, \textsc{Shy}, The Crossing Press, 1989, p.
  103.} And then goes on to produce a narrative with properties of the
cloud, numinous and hungry, where characters search for themselves
alongside the writer as a character himself, where persons encounter
each other but never stoop so low as to engage in a plot.

The ceiling was gray and smooth as the beach that Gunther Fielder lived
by. Flat, and peaceful, the way that ``now" is without a past or future
to rock it up any. He could focus on the gray and try to hypnotize
himself, closer towards death.''Do it," he demanded.

"My name is Harry Van," he said. It sounded so false. He said it over
and over, didn't ring true somehow. Like somebody else who you couldn't
remember. Well try again, something new.

``I'm David Bowie," he said, experimenting. "I have come to earth a
space invader, hot tramp, I love you so.'' Oh that was so suffragette,
trying to "be" a star.

He'd start again. "Hi, my name is Mark the dead boy," he said with great
difficulty.

Yes.

``Are you Kevin Killian," he replied. "Can I help you?'' Just like the
Hot Line!

These voices came out of his mouth from nowhere, between heaven and
earth, this conversation developing like a photograph pulled from its
tray full of crystal chemicals. Emergent. \footnote{ibid., p.~161.}

These voices attack the proposition that characters or author must be
unified presences, and suggest that self itself may not be locatable
along such a monochromatic line.

He is telling you now a story about narrativity, he is telling her
story. She finds the story as she looks at each other: so many faces.
She is crossing gender from the start, she wants you to know she is
Elizabeth Taylor---and has the Halloween photos to prove it. He is a boy
playing Puck in a high school production of A Midsummer Night's Dream
wearing ballet slippers forever. ``I have died and am in a novel and was
a lyric poet, certainly, who attracted crowds to mountaintops.''
\footnote{Palmer, Michael, \textsc{Sun}, North Point Press, 1989, p. 83.}
I come after Robert Duncan but before Norma Cole. My name is
``Broiling-Days-In-A-Little-Patch-Of-Shade." "For better or worse,''
says Flaubert,

\begin{quote}
it is a delicious thing to write, to be no longer yourself but to move
in an entire universe of your own creating. Today, for instance, as man
and woman, both lover and mistress, I rode in a forest on an autumn
afternoon under the yellow leaves, and I was also the horses, the
leaves, the wind, the words my people uttered, even the red sun that
made them almost close their love-drowned eyes. \footnote{Flaubert,
  Gustave, \textsc{The Letters of Gustave Flaubert, 1830-1857}, The
  Belknap Press of Harvard University Press, 1979, p.~203.}
\end{quote}

``Voice," "Person," "Point of View"---always singular---propose a
unified filter through which events may be organized, and as filters
screen properties, screen out toxins and tannins and pieces too big to
fit neatly. But pluralities are possible.''I see only from one point of
view," says Lacan, "but in my existence I am looked at from all sides."
\textsuperscript{\protect\hyperlink{WC11}{11}}
\protect\hypertarget{text11}{}{}Pronouns are known as shifters because
they are by nature unstable linguistic units, referring not to people
but to moving circumstances of speech and audition, visibility and
perception. As such they are fictional opportunities; unlike names they
permit a character to be subject and object, to ride the Wheel of
Person, speak and be spoken of with equal weight, inhabit simultaneity.
Here is a poem from Alice Notley's sequence, "Congratulating Wedge":

\begin{verse}
No I wouldn’t know why anyone would\\
want to write like that. I should never\\
have had to do it. We were used to this\\
other thing we always know like when we’re\\
here. And you have this clear head and you’re\\
seeing things and there they are. You don’t\\
notice they’re spelled. That's how you\\
know you’re alive. I never saw you\\
looking like a dictionary definition and if I\\
did I wouldn’t tell nobody. People\\
aren’t like that. They say, Hey\\
asshole motherfucker turn that radio\\
off! But the sun’s playing on it! But\\
it ain’t real, you dumb package!\\
I recognize every package the way it\\
comes. Now I’m mixed up. But I\\
always wanted to be a package, person\\
thinks. Do they? Or, I gotta de-\\
fine this package, me. Or, God if only\\
I was a package but I’m not.\\
\end{verse}

\footnote{Notley, \textsc{Alice, Margaret \& Dusty}, Coffee House Press,
  1985, p.~68.}

What are people like and what method correctly presents/represents them;
from what angles are they constructed and who construes the angles into
voice? In her mind as ``I," out of her mind as "she," confronting or
confronted by''you"; conspiratorially social and partial as ``we"; part
of one another, occasionally indistinct, certainly indiscreet, we are
and we are not separate people.''My premise, in general and in writing,"
says Leslie Scalapino, "is that I do not think there is a man, or woman,
or society, social construction; though it is there. It is not there."
\footnote{Scalapino, Leslie, in correspondence, 1989.} I have been
marginalized as a poet, homosexual, counterculture protester, drug
taker, transvestite and Jew; I am as interested in boundaries for what
lies outside them as in. I would like to drop my "characters" onto the
sharpened point of a gemstone, so that the radial fractures would
illuminate a comprehensive pluralistic image.

Syntax is the plot of the sentence, a systematic ordering of person and
event, of who does what to whom and when and to what end. Encoded in its
structure are a variety of fixed agreements that always end in a point
(.) Who will speak for beside-the-point? Nouns and verbs must have
parallel numbers, pronouns and verbs parallel persons; tenses must agree
to produce time that resembles progression. Business conveniences that
make of stories little prisons of discrete power relations with
seemingly invisible walls. I am not talking about referentiality
vs.~non-referentiality; I'm talking about how narrative referentiality
might be better served. Gender is foregrounded and elementalized,
digressions are trivialized, passive constructions frowned upon. As Sara
Schulman points out, try to tell a lesbian story without names: she came
into a room, she looked at her, she looked at her, she said---and aside
from homophobia, what terrors would such unlocations unleash? Normative
pronoun usage subjects self and other to power/dominance models of unity
and authority, of he over she and it beneath them. For pure syntax there
is Charley Shively's reduction of the phallocentric rule: "the subject
fucks the object."\footnote{Shively, Charley, \textsc{The Advocate},
  \#342, May 13, 1982, p.~24.}

Here is Leslie Scalapino writing:

\begin{quote}
The young person living there, having an intense tortured as if tearing
in half pain in the middle, waking lying asleep, though this had only
occurred this one time. The day and night being free of the one person,
who hadn't had this tortured sharp pain as if to tear her in half except
this one time, the man lying waking staying gently with her during it
through the soft darkness and then ending in the warm balmy day with the
people around who go down the street.\footnote{Scalapino, Leslie,
  \textsc{The Return of Painting}, manuscript.}
\end{quote}

Passive participial constructions which don't inhabit time, genderless
and then confusing gender assignations, unlocated relative pronouns,
erratic time shifts without one simple present tense: an amalgam of
person and event that keeps elements suspended and active, ``an
explosion, a dissemination" of meaning. "His mouth are everywhere,'' I
wrote erotically in "Honor Roll," insisting that the plural verb was
truer to the polyvalence of desire.

And I have neither a coherent story to tell nor can I cop a coherent
attitude to give my voice a characteristic singularity. I was born in
sleep and raised in sleep and wake up to find myself sleepwalking. The
figures I know all have shadows; some figures are smaller than their
shadows. In the first photo I am a soft blasted thing, mouth open tongue
hanging, blotto. Six weeks premature, I was still ``in here'' out there.
The world was unformed, coalescent. His story is the story of an
intuited world, a story where digressions may be the point, where
ellipsis is an accurate representation of what there is.

This world in its order decomposes into air, simultaneously present and
absent. A writing, then, of enmeshed simultaneities, which gives
sufficient weight to its constituent presences so that they verge upon
each other. The material relations of the Unknown. ``The stuff of the
psyche,'' says Herakleitos, "is a smoke-like substance of finest
particles, that give rise to all other things\ldots{}. it is constantly
in motion: only movement can know movement." \footnote{Herakleitos,
  \textsc{Herakleitos and Diogenes}, trans. Davenport, Grey Fox Press,
  1976, p.~18.} His story pulls the reader down from the surface of
language not to rest but to ride back and forth between the manifest and
imaginary worlds, among selves. ``I wanted to write a story,'' he
begins, "to talk about the outside world and escape my projections, but
the outside world could not escape from my projections. I wanted to
write not 'my' story but 'theirs'; I wanted to write about evil." He
looks at his fingers to escape your accusations; a sunbeam deconstructs
him into motes. He is happy dissolved there, and wants to write from
such dissolutions, melting into the grain of his lover's nipples. He has
no lover; he has entered an argument about narrative and political
ruination. ``Tell me your story,'' he asks, and you do.

Here in this dialogue writing relies less on information, as Walter
Benjamin shows, than on the moral power of interpretation, ``to keep a
story free from explanation.'' It is left up to the reader to
``interpret things the way he understands them, and thus narrative
achieves an amplitude that information lacks.''\footnote{Benjamin
  Walter, \textsc{Illuminations}, Schocken Books, 1969, p.~89.} Here a
fabricated house open to the wind is both a shelter and a sharpener of
the wind's bite, a house of shadows and a moving shadow that resembles a
house. Narrativity, the action not the thing, a happening semblance that
is and is not a story, a gift given and taken away so that one must
finally stand fulfilled by transgression. Narrativity, a process of
integration not linear but aggregate, circular, partial---and so,
complete.

~

\hypertarget{brian-strang-the-kingdom}{%
\chapter{Brian Strang: the kingdom}\label{brian-strang-the-kingdom}}

~

\hypertarget{on-narrative}{%
\section*{On Narrative}\label{on-narrative}}
\addcontentsline{toc}{section}{On Narrative}

\hypertarget{am}{%
\subsection*{3:28am}\label{am}}
\addcontentsline{toc}{subsection}{3:28am}

\lettrine[lines=3, findent=0em, nindent=0.1em, lhang=0]{T}{he temporal element}
of narrative humanizes, introducing history, biography and, therefore,
mortality; as time is presenced, so is humanity and, so then, is the
passing of time. What are the implications for the conceptualization of
human history if, through temporality, pronouns become subjective and
indeterminate, if temporal elements in narrative are disjoined or occur
simultaneously? What are the implications if we are seduced by
narrative, by the promise of passing time (in an eroticized dance with
death), and are not led toward syllogism and deduction but discover,
instead, induction, investigation, participation?

Aaron Shurin, in ``Narrativity," proposes that "writing might use
narrative without succumbing to its hegemonic orders of linear
development, unity of time/tense-and apart from the modernist
reconstructing modes of memory and dream. A prose whose paragraphic
groupings themselves might be based on measure, whose higher
integrations might be thematic or associational rather than
developmental.''

When liberated from plodding development, can narrative, loaded with its
historical implications, move us away from the discrete series and veer
toward witness, philosophy, discovery?

\begin{center}\includegraphics{images/"pandanus.png"}\end{center}

\hypertarget{am-1}{%
\section*{2:48am}\label{am-1}}
\addcontentsline{toc}{section}{2:48am}

Turn your face toward the back seat. Walk up the stairs to an apartment
set in stone. Listen to the parking lot below. A man lies on the outside
of a fence, exhaling. This is the view of others' rooftops, making them
into little pieces of pie. Rows of fluorescent tubes. On the islands of
a nearly drained sea, turtles lie exposed and cracked roots tremble in a
particularly strong wind. The mail is delivered with surprising accuracy
down the storm drain. You're alarmed by the presence of gunmen in your
hotel room.

\hypertarget{am-2}{%
\section*{2:49am}\label{am-2}}
\addcontentsline{toc}{section}{2:49am}

A policeman on the street looks through layers of glass. A very dead
thing. Motion toward the approaching car. Now you would like to leave by
way of the great hall. Strange ways and the lack of knowledge of local
customs prevent you from negotiating stillness. Manage nothing. Reptile.

\hypertarget{am-3}{%
\section*{4:58am}\label{am-3}}
\addcontentsline{toc}{section}{4:58am}

Lying on the night flattened. Maps of territory barely existent. Pools
cover the floor. A single drip from the middle of the air. You think
this might be the kingfisher. Two figures in the middle of the
room-translucent. Thin warmth of blue light in a very small town. As
padding footsteps structure the quiet, a man looks from the top of a
house. They lie in a casual circle sleeping on the ocean floor.

\hypertarget{am-4}{%
\section*{4:59am}\label{am-4}}
\addcontentsline{toc}{section}{4:59am}

He pats his suit pocket, fingers the traveler's checks in his shoes.
Trucks begin to unload themselves all around. Drop off to sleep. Many
floors up the windows are missing. Entry is enclosed. Water continues to
drip.

\hypertarget{am-5}{%
\section*{3:29am}\label{am-5}}
\addcontentsline{toc}{section}{3:29am}

Walk below the surface and pick at the matrix of possibilities. Shifting
focus, a spray of bubbles through water, padded air on the ceiling. A
bent neck, S shaped, springs forward when needed, weeding out the sick
and the old.

\hypertarget{am-6}{%
\section*{5:13am}\label{am-6}}
\addcontentsline{toc}{section}{5:13am}

Synapse. The kingfisher hovers. Get together. A crush of grace. The
beats between are short. A real taste. A hand's distance from the floor.
Wings over balconies, through glass. On the inside of the skin, very
sweet and it becomes petrified.

\hypertarget{am-7}{%
\section*{2:16am}\label{am-7}}
\addcontentsline{toc}{section}{2:16am}

A car brings someone to this very corner. The greatest are able to brace
themselves, to support a moving supply. This too is style, you think.
Behind you in the car is a man. Turn around quickly to find the shell of
a turtle. You tremble at the lack of knowledge, at the blue light.
You're alarmed now by the lack of presence.

\hypertarget{pm}{%
\section*{4:22pm}\label{pm}}
\addcontentsline{toc}{section}{4:22pm}

Click of ribs. Drone, repeat. Footbridge. Highway. Landing zone agenda.
Corrective maneuvers. Unassailable.

\hypertarget{am-8}{%
\section*{6:12am}\label{am-8}}
\addcontentsline{toc}{section}{6:12am}

Dawn begins below the floor, sends razors, slits, leaves an overall
impression. You concentrate on the phenomenon. The television presents
you with a flat grey surface. Leave marks on the sandbar. Retract
yourself as the tide rises. Deactivated bleakness. Wiggle your fingers
as complexities weave into numb certainty. A fin is cutting through,
toward equilibrium. Paths intersect. A car door opens. Apartments.
Surprisingly nimble. Walk around inside. Retract.

\hypertarget{pm-1}{%
\section*{2:59pm}\label{pm-1}}
\addcontentsline{toc}{section}{2:59pm}

Somnambulence. Funicular shoes. The growth rate is a recursive
soliloquy. Some tetrahedral antipathy. Total shadow. Halitosis gloss.

\hypertarget{am-9}{%
\section*{5:55am}\label{am-9}}
\addcontentsline{toc}{section}{5:55am}

Unscathed, a memoir progressively installed. Move across the field of
vision. Turn ash piles. Bristle with teeth. Bone worms make their paths.
Spillover highway fanned out. Contrite motive. Recount sobs. Astonished
circuits in a coma. Damp faces. The nerve protectants.

\hypertarget{am-10}{%
\section*{2:12am}\label{am-10}}
\addcontentsline{toc}{section}{2:12am}

Open the door onto the waves. Drop beneath tension to the street. Find
your way along in the dark, fingering shells. Cling to a ledge, creating
shelter, always watching what hovers, what makes you burrow in wet soil.
The ceiling is glass. Peering from above. Under the gaze, crawl about
your business. A great neck dips, picks morsels from below. At any
moment. In the back seat of the car, in the security of a bed, on the
street, beneath the surface, on a small island, in the bones.

\hypertarget{lawrence-upton-tran-from-easy-kill}{%
\chapter{\texorpdfstring{Lawrence Upton: TRAN, from \emph{EASY
KILL}}{Lawrence Upton: TRAN, from EASY KILL}}\label{lawrence-upton-tran-from-easy-kill}}

~

\hypertarget{working-notes-7}{%
\section*{Working Notes}\label{working-notes-7}}
\addcontentsline{toc}{section}{Working Notes}

\emph{Prose and verse are rhythmic-semantic structures; both are
vehicles for poetry. The twisting space where the two meet is like a
path which is often hard to find, variously rocky, earthy and grassy.
Each kind makes particular demands and endows particular qualities.}

\textsc{\emph{EASY KILL / TRAN}} \emph{structures itself in paragraphs:
significant breaks come, not at the end of lines or sentences, but at
the end of paragraphs. Sub-structures are of variable length without
there being a strong base rhythm. It is laid out accordingly.}

\emph{Thus, the decision is technical, no more, no less.}

\begin{center}\includegraphics{images/"pandanus.png"}\end{center}

\lettrine[lines=3, findent=0em, nindent=0.1em, lhang=0]{T}{wenty seven men,}
slightly unformed, enter the building at the ground floor. They move
swiftly and smoothly, each minimizing use of space, as coins will as
they slide from production machines into circulation.

The group splits apart, rolling down corridors, rising up stairwells.
Reason is not engaged; locks and bolts hold no one back. They are
ruthless, even in the way that they sustain themselves.

When they have finished, twenty five leave as they entered. Two are
carried on stretchers, observed by a photographer who will be injured
later; the film crew never make it home.

He walks loose-limbed, does not drink and eats rarely. His skin seems
beetroot-stained; his hair greys in slippery afternoon light.

Is that your key? he asks, but his mind's not on the answer. Calm down,
she says, but that's another conversation. The coffee's beginning to
boil over.

Charts on the wall hang loose on their drawing pins. He turns to sit and
subsides into the chair, yet stays as bunched and clustered to himself
as a primrose. He begins to read, slowly, mouthing some of the words.

Gun emplacements fall apart. The cliff edges towards them. Roads widen;
building spreads; red ants extend across the bottom field; the climbing
rose gets rangy.

A woman makes a big noise, kissing her child; most of it's pre-contact.
The child responds with a startled pleasure yelp. A drill bit cuts
through steel. She rubs her nose against the child's nose. Some moisture
is transferred. I see it glisten as they pull back from each other. All
the colours of the rainbow, says the woman.

A glittering sphere swings with other reflected spheres below it,
doubling on the glittering surface of the wet floors, reflecting against
a wall. The spheres and their reflections merge.

Day's lights come on. The hall is well lit now. It's almost white.
Reflections make illusions of space, imaginary corridors between darker
reflections.

A bright empty day. Light from sea is brighter than light from rocks.
Both make it difficult to see. Steps come out of clouds. Something which
menaces hover\ldots{} more than two things, remaining one thing, a body
with paws and emerging tongues. An almost-silvery watery day.

A flock of birds, the letters of a shout, rising from a farm, among
plants growing there, and form themselves into a line which they fly
along. The sound makes several indistinct words, which form themselves
variously into a scarecrow and a windmill.

The scarecrow waves its arms in the wind and the windmill turns,
speaking of something flying and something explosive and something
breaking and useful rolling downhill, an old wrecked car in a quarry, a
navigation aid, a broken sign on the headland, a set of tools spilling
from a dropped tool box. A piece of soft paper blows through the
complex\ldots{}. The wind drops.

They are making this place their own, spreading their refuse! The walls
are large and solid; but the inscriptions are unreadable and they find
it hard to deduce purpose from the room structures - defensive walls and
none of the subsidiary signs of defensiveness, contained areas and no
utility to the containment. They think it is a temple. A break in the
walls as if from external damage or geological disturbance; and, within
that break, something that pulsed, unrecognizable, almost thoroughly
decayed, webs of fibres, doubling of sight to no purpose and with no
cause.

A flag of torn words, things said, a string of them tied in a narrative
that no one would or could believe, a verbal paling, a piece of song in
prison, a humming through boredom made bearable by concentration; a
flight of steps whose beginning and whose end has been abandoned,
excavated; carved into the steps, letters believed to be contemporaneous
with the steps, but not readable; a fence of broken words, torn from
library books and official buildings, on its side, after violent wind; a
series of shadows in the inner court; shells falling.

A roll of bunting falls across steps and cascades open. Near the bottom,
it spreads out and the figure of a man forms. He rises from a river and
disintegrates. The river becomes slogans and the slogans turn into the
figure of a man. This image breaks in several places.

He looks into shadow, out of bright day, and sees black; a smell and a
scuttling come out; the words of the tune he had been singing on a
screen of temporary blindness; smiles form between words; loved faces
and feared faces take on the words. Behind him, someone is calling
smiles into his camera. These words too shelter in the dark place;
lights in the head die\ldots{} The transformation ceases. It is still
too dark to see; the sight is fading, light from the bomb bay as it
opens; a city burning. It is black with no breath upon waters and no
thing flying over; not night because night had a beginning and will have
an end; not death because death flows out of living; not black, the
colour, because that has some amelioration, some grey suggestion or a
blueness so deep it does not seem blue; a silent black, loss of words,
sizeless whole; closing of eyes to pain before pain without light
filtered through eyelids, without memory of light; not any colour; but
absence of colour, and absence of a meaning to the word colour, and
absence of absence, because absence would imply a potential presence, or
the light which brings a knowledge of the colour, absence of wave and
wave length, absence of dimension, absence of genitive relationship,
absence of knowledge of relationship; and absence of knowledge, a single
tuple breaking as if dead wood if it were to exist, if there were a
place for it or a tree to exist, from a tree without existence in the
dark which has no existence. Look away.

Blow backs through the headlights speckle the windscreen; a voice on the
radio nags; the road disappears into a point; someone goes over the top
and dives, only his hand sticking from the tumble into white, and water
spills on the hard surfaces which would have seemed smooth and are now
seen, unevenly, as the fall buckles and splashes from its downward.
Caverns open in the hardness while rivulets flow into them. Dazzling
blank spaces break out and then break. Surfaces seem to merge under
speedy water. Voices bunch together and grow thorns we trip on.

Voices regulate themselves and march multi-legged. Avoid them. Film runs
in and out of sprockets, tearing and catching itself, projecting,
another film than the one expected and the one filmed. Film as film,
mutters the technician, sneering, as he joins broken sections and there
is little hope that anything will have changed - except this copy of the
film. He winds it backwards but the thing is tumbled in its reel.

A water wheel turns double imaged in an unprecedented flow of water.
Several stage lights stutter. A flint on a lighter sparks briefly and he
blows smoke into the eyes of one trying to help him. Someone has spilt
motor oil and it's everywhere.

Leaves and red metal. A sheet of red metal, a field of green
leaves\ldots{}

Words painted with yellow paint on sodden fencing wood, fallen
over\ldots{}

Roots of a tree which has been uprooted, hanging over the hole it had
itself made, displaced, and cut through straight eighteen inches up its
trunk\ldots{}

An old concrete platform towered over by flowering lupines\ldots{}

A bus goes by\ldots{}

A blue plastic crate\ldots{}

At the top of the hill, several stories of reflecting glass throw back
the hill and sky valley in unexpectedly-shaped segments, duplicating
disarraying simultaneously\ldots{}

A dovecote in slow process of falling down the side of a disused railway
embankment, collapsing as it does so. Paint flakes from it. Water
streaks it. Cobwebs fan out, catching and turning dim afternoon light.
Runnels of earth where the soil is washed away. Broken white enamel, a
matins bell, an upholstery armchair attracting flies in a village car
park. A flattened beer can rolls a little way across the Aldwych in
evening breeze. After a while it rolls not back but in a new direction
on a bus wind.

A man with a large beer gut and a cap over his eyes; a large black key;
an hysterical person between headlights; testicles and heads hanging
from butchers' hooks; a hole in perimeter fencing; warning signs in
empty corridors; a security guard exploding in gun fire; a woman
reaching for the moon; a bus turning in a circle, on and on; an
opportunity for metamorphosis; a clout hammer hard on a pair of hands;
cash; hard on; a body buried beneath a crumbling floor with only its
face showing; an unexpected fall; a shower of rain; a green olive doused
in olive oil; a stretch of silver sand!

A picture of a smiling woman, engineered to make her look younger than
she now is; a bar of chocolate; empty seats; a digital display showing a
binary code progression; a cave painting; a picture of a tomb; a smiling
robber; levels of pain; an army unit; false images, including misleading
images; row upon row upon row of a man calling, making excessive claims
in a high-pitched voice; blue water inside words you've used; a poseur
trying to organize a trip, invoking the spirit of Columbus; two men
rubbing their legs against each other's warily, fingering knives; a
period of blackness; there isn't much truth in anything you say.

Panic in the rowboats - there's so much steam hardly any of them can see
what they are doing - and they speak frantically, aggressive, but no one
gets hurt. They're trying to get aboard a large model boat.

The sand is green and the sea yellow, all out of focus. He floats,
somehow, in the cave. His feet are not in the sea, most certainly; but
he cannot see his feet.

No sound of sea. Jazz piano played fast. A spider swings on a web. A
snail on a twig, dewdrops hanging from surrounding branches. Forest
takes green from the beach. Which is sea? No sensation. A bit of forest
is in focus. A forest with a yellow sky.

Light white dust pouring through a down-angled window in a thick wall. A
vast cello bowed by someone even vaster. She cannot see across the room.
It's far too far. All she can see is the cello and a shadow of the
player. Towards her feet, the curvature of the Earth distorts her body's
shape.

Two suitcases, wide open, sprawled across the living room floor,
clothing items and other personal things hanging out of them; a vacuum
cleaner cable, the vacuum cleaner tube turned over itself; there's
hardly any place to sit; carpet predominantly red and brown; the people
hardly real, as if remembered from a long while back; murmur and mumble
of their voices, less clear the more one concentrates, less interesting
the more one hears.

A galleon upon the sea at evening, the sky black, running into gold and
red, and then blue, as she looks across the hemisphere, a single bird,
gliding, high up; distance collapses as the light goes, the whole fleet
becomes a shadow and the shadow merges with the atmosphere between,
called darkness; a light is on at the vanity table; a breeze lifts the
curtains; the wallpaper gleams in the indirect light, a cloying smell of
scent I can't quite place, the breath of someone.

A black and white photograph lying semi-crumpled on the floor picked out
in white torch light; bottom of a chest of drawers, a white strap
hanging down; a man hanging from a ceiling beginning to crack with the
weight.

A green yellow blue parrot painting itself intermittently and cleaning
its toes in a room got up in Turkish style; a video camera records
events; sometimes the parrot consults notes and sometimes it seems to
improvise; a potted fern sways as if in a breeze; the video camera
hardens; there is a smell of damp; muffled voices at a distance.

Freight yards stretch to the horizon; trains come here from all nodes on
the Earth's surface.

Bright light illuminates the public house opposite; but it is not the
light of day, rather a video screen of bright colours.

A trolley, high-sided, gloss black, building itself, two uprights
standing upright, spinning round the chassis and assembling in a street
of plane trees; some beggars, hands behind their backs, tied together,
heaped over a tarpaulin at the side of the bright street, watched by a
policeman who spits at them; the harsh wind blows into them; gloves with
the fingers torn out; the police sergeant is trying to talk sense but he
cannot make it; a man with a smart suit is holding up a framed mirror
around and about which he plays find the lady; his suit is getting
grubby; the wind up gramophone is slowing down; no one is paying any
attention to anything; a woman rocks a marionette in a pushchair; the
marionette is counting money; a child stamps in a puddle; in the puddle
one sees a reflection of the moon; from the moon the Earth might look
beautiful.

Everyone stands in the raked seating and the coffin is wheeled slowly
out of the hall to where, in the forest, someone has been demonstrating
what a railroad looks like, using twigs placed parallel; a head is
rolled down the tracks, bits of the throat still attached to it
trailing, coagulating blood running out, the terrified face aware; it
passes between my feet.

Almost the entire body, distorted towards the top or dwindling from head
to foot, a kind of homunculus, demonstrating a fish eye effect, follows
its head without assistance. She stays calm, the body, sliding down this
stupidly inadequate image; friction stops her; she begins to deteriorate
and, as her elements separate, rage takes her; she bangs her hands upon
the ground, she shakes and throws a silent tantrum, collapses into
leaves from the bottom of a forest; nearby a bonfire has been started.

A dark-skinned woman got up like a comic book sorceress brings down her
arm in a line towards people who have transgressed. They fall into
pieces, broken mannequins, and crumble into mouldy chocolate among
strands of tissue paper.

A scorpion spins on a sandy floor, chasing its own tail.

Blackness crystallizes out of fog moving forwards. They track ahead of
it, seventy, eighty miles per hour, shaded headlights, which drop as the
car transforms, layer below layer, concertinalike, to a row of tumbling
stacking chairs, in the very first instant of the tumble, eyes being
shot from their sockets, bulbs and lampholders from lampshades, guts
from torn bodies, presents dropped from Christmas stockings; and it
begins to lumber like an early attempt to demonstrate the reality of the
skeletons of giant animals; the road's been lost for ever; and the
voice, not just falsely jocular, but sounding false as it imitates a
falsely jocular voice, speaks: you turn upon your own centre like a
paper spill.

A male and a female walk together. He gestures to two male children
following obediently, throws a scarf counterclockwise around her neck
and moves forward striding high, features blotted out to protect the
guilty; the film merges with and is taken over by a second film quite
unrelated; the street is ended violently by cobbled nineteen nineties'
chic; a man walks a woman in a wheel chair up the centre of the street
which ends at a sort of knuckle joint, asserting thus a vortex of cobble
stones trying to force one away from further progress, cut into a
portrait of a monstrous wide-armed Christ above a city and thence into a
pasty white sun hanging in a mist; a silver screw dropping through a
Victorian municipal glass roof.

A foam cut out angel fights a cat; the cat's cut out, aggressive, back
bent to give it power; the angel, hides within her string bag, the one
expression she has appropriate because so ambiguous, a measure of
concern, a measure of dismay, a measure of remoteness.

A man sucking a sweet while he speaks apparently affably; an old man
scratching his head, dandruff rising.

%%%%%%

\end{document}
